%  Методы оптимизации — конспект лекций
%  Компиляция: pdflatex main.tex  (дважды, из-за ссылок и оглавления)

\documentclass[12pt,a4paper]{article}

%  Преамбула конспекта «Методы оптимизации»
%  Компилировать pdfLaTeX (кодировка utf8, шрифты T2A).

% Latin Modern не имеет T2A-начертаний, поэтому кириллица идёт через cm-super:
% семейство cmr в кодировке T2A.
\usepackage[T2A]{fontenc}
\usepackage[utf8]{inputenc}
\usepackage[english,russian]{babel}

\usepackage{amsmath,amssymb,amsthm,mathtools}
\usepackage{geometry}
\geometry{a4paper,margin=2.5cm}

\usepackage{tikz}
\usetikzlibrary{arrows.meta,patterns,decorations.pathmorphing,decorations.pathreplacing,shapes.geometric,calc,positioning}
\usepackage{pgfplots}
\usepgfplotslibrary{groupplots,fillbetween}
\pgfplotsset{compat=1.18}

\usepackage[most]{tcolorbox}
\tcbuselibrary{breakable}

\usepackage{enumitem}
\usepackage{booktabs}
\usepackage{hyperref}
\hypersetup{
  colorlinks=true,
  linkcolor=blue!55!black,
  citecolor=blue!55!black,
  urlcolor=blue!55!black,
  pdftitle={Методы оптимизации --- конспект}
}

%  Окружения для теорем
\theoremstyle{plain}
% Счётчики независимы: теоремы, леммы и утверждения нумеруются каждое своим рядом.
\newtheorem{theorem}{Теорема}
\newtheorem{proposition}{Утверждение}
\newtheorem{corollary}{Следствие}
\newtheorem*{corollary*}{Следствие}
\newtheorem{lemma}{Лемма}
\newtheorem{criterion}{Признак}
\newtheorem*{criterion*}{Признак}

\theoremstyle{definition}
\newtheorem{definition}{Определение}
\newtheorem{example}{Пример}
\newtheorem{exercise}{Упражнение}

\theoremstyle{remark}
\newtheorem*{remark}{Замечание}
\newtheorem*{reminder}{Напоминание}

\renewcommand{\qedsymbol}{$\blacksquare$}

%  Пометки конспекта
%    \begin{idea} ... \end{idea} — пояснения, которых не было на лекции
%    \begin{warn} ... \end{warn} — неразборчиво / описка / требует сверки
\newtcolorbox{idea}[1][]{
  breakable,
  enhanced,
  colback=blue!4!white,
  colframe=blue!45!black,
  boxrule=0.4pt,
  left=6pt,right=6pt,top=5pt,bottom=5pt,
  title={Пояснение},
  fonttitle=\bfseries\small,
  coltitle=white,
  attach boxed title to top left={xshift=6pt,yshift=-2pt},
  boxed title style={colback=blue!45!black,boxrule=0pt,arc=1pt},
  #1
}

\newtcolorbox{warn}[1][]{
  breakable,
  enhanced,
  colback=orange!6!white,
  colframe=orange!65!black,
  boxrule=0.4pt,
  left=6pt,right=6pt,top=5pt,bottom=5pt,
  title={Сверить с записью},
  fonttitle=\bfseries\small,
  coltitle=white,
  attach boxed title to top left={xshift=6pt,yshift=-2pt},
  boxed title style={colback=orange!65!black,boxrule=0pt,arc=1pt},
  #1
}

%  Пошаговая запись метода: \begin{algo}[title={Метод ...}] ... \end{algo}
\newtcolorbox{algo}[1][]{
  breakable,
  enhanced,
  colback=black!3!white,
  colframe=black!60,
  boxrule=0.4pt,
  left=6pt,right=6pt,top=5pt,bottom=5pt,
  title={Алгоритм},
  fonttitle=\bfseries\small,
  coltitle=white,
  attach boxed title to top left={xshift=6pt,yshift=-2pt},
  boxed title style={colback=black!60,boxrule=0pt,arc=1pt},
  #1
}

% Шаги внутри algo
\newlist{steps}{enumerate}{1}
\setlist[steps]{label=\textbf{Шаг \arabic*.}, leftmargin=*, labelsep=0.6em,
                topsep=4pt, itemsep=3pt}

% Отметка начала лекции
\newcommand{\lecture}[2]{%
  \bigskip
  \noindent\hfill\textit{\small #1\ifx\relax#2\relax\else, #2\fi}\\[-0.3em]
  \noindent\rule{\textwidth}{0.4pt}
  \medskip
}

%  Макросы
\newcommand{\R}{\mathbb{R}}
\newcommand{\NN}{\mathbb{N}}
\newcommand{\eps}{\varepsilon}
\newcommand{\T}{^{\mathsf{T}}}
\newcommand{\rank}{\operatorname{rank}}
\newcommand{\sgn}{\operatorname{sign}}
\newcommand{\norm}[1]{\left\lVert #1 \right\rVert}
\newcommand{\abs}[1]{\left\lvert #1 \right\rvert}
\DeclareMathOperator*{\argmin}{arg\,min}
\DeclareMathOperator*{\argmax}{arg\,max}

% Индексные множества канонической задачи.
%   N — индексы переменных, M — индексы ограничений,
%   N_0^+ / N_0^0 — положительные / нулевые компоненты допустимой точки,
%   N^+ — базисное множество (пополненное до |M|).
\newcommand{\Np}{N^{+}}
\newcommand{\Nz}{N^{0}}
\newcommand{\Nzp}{N_0^{+}}
\newcommand{\Nzz}{N_0^{0}}

% Сокращение «опорный вектор», как в тетради.
\newcommand{\ov}{о/в}

% Клетка транспортной таблицы: #1 — тариф c_ij (мелко, сверху),
% #2 — объём перевозки x_ij (обычным кеглем, снизу).
\newcommand{\tc}[2]{\shortstack[r]{\scriptsize $#1$\\ $#2$}}
% Пустая (незаполненная) клетка таблицы.
\newcommand{\tcx}[1]{\tc{#1}{\times}}

% Маркер «++ в кружке» из тетради: доказательство будет дописано позже.
\newcommand{\plusplus}{%
  \tikz[baseline=(pp.base)]{%
    \node[draw, ellipse, inner xsep=3pt, inner ysep=1pt, line width=0.4pt] (pp)
      {\footnotesize $+\,+$};}%
}

%  Быстрая сборка: \showfiguresfalse выключает все TikZ-рисунки.
\newif\ifshowfigures
\showfigurestrue

\newcommand{\inputfig}[1]{%
  \ifshowfigures
    \input{#1}%
  \else
    \begin{center}\fbox{\small рисунок выключен: \texttt{\detokenize{#1}}}\end{center}%
  \fi
}


\title{\textbf{Методы оптимизации}\\[0.3em]\large Конспект лекций}
\author{}
\date{}

\begin{document}

\maketitle
\thispagestyle{empty}

\begin{idea}[title={Обозначения в конспекте}]
Блок \textbf{«Пояснение»} — то, чего не было на лекции: интуиция, обоснование
пропущенных шагов, связи с другими темами.

\medskip
Маркер \plusplus{} вместо доказательства — так помечено в тетради; эти доказательства
дописываются позже.
\end{idea}

\begin{idea}[title={Индексная запись}]
Весь курс использует индексную запись векторов и матриц через множества индексов:
\begin{gather*}
  M = \{1, \dots, m\} \ \text{— индексы ограничений}, \\
  N = \{1, \dots, n\} \ \text{— индексы переменных} .
\end{gather*}
Тогда $x[N]$ — вектор-столбец длины $n$, $b[M]$ — вектор-столбец длины $m$,
$A[M,N]$ — матрица $m \times n$. Для подмножеств $N' \subset N$, $M' \subset M$
запись $x[N']$ означает подвектор из компонент с индексами из $N'$, а $A[M',N']$ —
подматрицу на пересечении соответствующих строк и столбцов. В частности,
$a_j[M]$ — $j$-й столбец матрицы $A$, а $A[M,\Np]$ — набор столбцов с индексами
из $\Np$. Размеры: $\abs{M} = m$, $\abs{N} = n$.
\end{idea}

\tableofcontents
\newpage

% Нумерация: section = глава, subsection = параграф (§)
\renewcommand{\thesubsection}{\S\arabic{subsection}}

\lecture{Лекция 1}{}

\section{Задачи линейного программирования}

\subsection{Угловые точки и многогранные множества}

\begin{definition}\label{def:extreme-point}
Пусть $S \subset \R^n$ — выпуклое множество. Точка $x \in S$ называется
\textbf{угловой} (крайней) точкой множества $S$, если её нельзя представить
в виде полусуммы двух \emph{различных} точек этого множества:
\[
  \nexists\ x_1, x_2 \in S,\ x_1 \ne x_2 : \qquad x = \tfrac{1}{2}\,(x_1 + x_2).
\]
\end{definition}

\begin{idea}
Эквивалентная формулировка: $x$ — угловая точка, если из
$x = \lambda x_1 + (1-\lambda) x_2$ с $x_1, x_2 \in S$ и $\lambda \in (0,1)$
следует $x_1 = x_2 = x$. То есть $x$ не лежит \emph{внутри} никакого отрезка,
целиком содержащегося в $S$. Форма с $\lambda = \tfrac12$ не слабее общей:
если $x$ делит отрезок в отношении $\lambda : (1-\lambda)$, то, укоротив длинную
сторону, получим отрезок множества с серединой в $x$.
\end{idea}

\inputfig{figures/fig-extreme-point}

\begin{definition}\label{def:polyhedral}
\textbf{Многогранное множество} — множество, образованное пересечением конечного
числа замкнутых полупространств и гиперплоскостей.
\end{definition}

\begin{idea}
Полупространство $\{x : \langle a, x\rangle \le \beta\}$ и гиперплоскость
$\{x : \langle a, x\rangle = \beta\}$ выпуклы и замкнуты, а пересечение выпуклых
замкнутых множеств выпукло и замкнуто. Значит, многогранное множество всегда
выпукло и замкнуто — определение~\ref{def:extreme-point} к нему применимо.
Ограниченным оно быть не обязано (например, $\{x \ge 0\}$).
\end{idea}

\subsection{Три формы задачи линейного программирования}

Задачи ЛП могут быть представлены в трёх формах.

\medskip
\noindent\textbf{(1) Каноническая форма:}
\[
  \begin{cases}
    c\T[N]\, x[N] \to \min, \\
    A[M,N]\, x[N] = b[M], \\
    x[N] \ge 0 .
  \end{cases}
\]

\noindent\textbf{(2) Симметрическая форма:}
\[
  \begin{cases}
    c\T[N]\, x[N] \to \min, \\
    A[M,N]\, x[N] \ge b[M], \\
    x[N] \ge 0 .
  \end{cases}
\]

\noindent\textbf{(3) Общая форма} — часть ограничений равенства, часть неравенства,
причём условие неотрицательности наложено не на все переменные:
\[
  \begin{cases}
    c\T[N]\, x[N] \to \min, \\
    A[M_1, N]\, x[N] = b[M_1], \\
    A[M_2, N]\, x[N] \ge b[M_2], \\
    x[N_1] \ge 0, \quad x[N \setminus N_1] \ \text{— свободные переменные},
  \end{cases}
\]
где $M = M_1 \sqcup M_2$.

Ограничения (без целевой функции) задают \textbf{допустимое множество} $S$ — это
многогранное множество в смысле определения~\ref{def:polyhedral}.

\begin{idea}
Все три формы равносильны — переход между ними делается стандартными приёмами:
\begin{itemize}[nosep,left=0pt]
  \item неравенство $\langle a, x\rangle \ge \beta$ превращается в равенство
        добавлением \emph{балансной} переменной:
        $\langle a, x\rangle - s = \beta$, $s \ge 0$;
  \item равенство разбивается на пару неравенств
        $\langle a, x\rangle \ge \beta$ и $-\langle a, x\rangle \ge -\beta$;
  \item свободная переменная заменяется разностью двух неотрицательных:
        $x_j = x_j^{+} - x_j^{-}$, $x_j^{\pm} \ge 0$;
  \item максимизация — это минимизация с $c \mapsto -c$.
\end{itemize}
Поэтому вся дальнейшая теория строится только для канонической формы.
\end{idea}

\subsection{Допустимая точка и разложение по столбцам}

Далее везде рассматривается \textbf{каноническая форма}. Примем два предположения:
\[
  \abs{M} \le \abs{N}, \qquad \rank A[M,N] = \abs{M} .
\]

\begin{idea}
Первое условие ($m \le n$) означает, что уравнений не больше, чем неизвестных, —
иначе система, как правило, переопределена. Второе (строки $A$ линейно независимы)
убирает зависимые уравнения: либо они следствия остальных и их можно вычеркнуть,
либо система несовместна. Вместе они гарантируют полный ранг по строкам, а значит
из $N$ всегда можно выбрать $\abs{M}$ линейно независимых столбцов.
\end{idea}

Пусть $x_0[N]$ — \textbf{допустимая точка}, то есть $A[M,N]\,x_0[N] = b[M]$ и
$x_0[i] \ge 0$ для всех $i \in N$. Введём обозначения для индексов положительных и
нулевых компонент:
\[
  \Nzp = \{\, i \in N \mid x_0[i] > 0 \,\}, \qquad
  \Nzz = \{\, i \in N \mid x_0[i] = 0 \,\} .
\]

Распишем систему ограничений по столбцам:
\[
  \left\{
  \begin{array}{ccccccc}
    a_{11} x_1 &+& a_{12} x_2 &+ \ \dots \ +& a_{1n} x_n &=& b_1, \\
    \vdots & & \vdots & & \vdots & & \vdots \\
    a_{m1} x_1 &+& a_{m2} x_2 &+ \ \dots \ +& a_{mn} x_n &=& b_m,
  \end{array}
  \right.
\]
где столбцы обозначены $a_1[M], a_2[M], \dots, a_n[M]$. В векторном виде это
\[
  \sum_{j \in N} x[j]\, a_j[M] = b[M] ,
\]
а для допустимой точки $x_0$ нулевые слагаемые выпадают, и остаётся
\begin{equation}\label{eq:column-decomp}
  \sum_{j \in \Nzp} x_0[j]\, a_j[M] = b[M] .
\end{equation}

\subsection{Опорные векторы}

\begin{definition}\label{def:support-vector}
Допустимая точка $x_0[N]$ называется \textbf{опорным вектором} (\ov), если столбцы
$\{a_j[M]\}_{j \in \Nzp}$, отвечающие положительным компонентам, линейно независимы.
\end{definition}

\begin{remark}
Из $\rank A[M,N] = \abs{M}$ следует $\abs{\Nzp} \le \abs{M}$: линейно независимых
столбцов в $\R^{m}$ не может быть больше $m$.
\end{remark}

\begin{definition}\label{def:degenerate}
Опорный вектор $x_0[N]$ называется \textbf{невырожденным}, если $\abs{\Nzp} = \abs{M}$,
и \textbf{вырожденным} иначе (то есть при $\abs{\Nzp} < \abs{M}$).
\end{definition}

Отсюда два случая.

\begin{enumerate}[label=(\arabic*),left=0pt]
  \item Если \ov{} невырожден, то $\det A[M, \Nzp] \ne 0$ для этого вектора:
        квадратная матрица $\abs{M} \times \abs{M}$ из линейно независимых столбцов
        невырождена. Здесь $\Np := \Nzp$.
  \item Если \ov{} вырожден, то набор $\Nzp$ можно \emph{пополнить} до множества $\Np$
        с $\abs{\Np} = \abs{M}$ так, что $\det A[M, \Np] \ne 0$.
\end{enumerate}

В терминах компонент самого вектора это даёт
\[
  (1) \implies x_0[i] > 0 \quad \forall i \in \Nzp = \Np ;
\]
\[
  (2) \implies
  \begin{cases}
    x_0[i] > 0, & i \in \Nzp, \\
    x_0[i] = 0, & i \in \Np \setminus \Nzp .
  \end{cases}
\]

\begin{idea}
Пополнение в случае (2) возможно именно из-за условия $\rank A[M,N] = \abs{M}$:
линейно независимый набор $\{a_j\}_{j \in \Nzp}$ можно достроить до базиса
пространства $\R^{m}$ столбцами той же матрицы $A$ (теорема о дополнении до базиса).

Разница между случаями — в том, сколько нулевых компонент попало в базис. У
невырожденного \ov{} базис состоит ровно из столбцов с положительными компонентами;
у вырожденного в базис добавляются «лишние» столбцы, которым отвечают нули. Поэтому
одному вырожденному \ov{} может отвечать \emph{несколько} разных базисов — это будет
видно в примере~\ref{ex:six-bases}.
\end{idea}

\begin{definition}\label{def:basis}
Набор столбцов $\{a_j[M]\}_{j \in \Np}$ называется \textbf{базисом опорного вектора}.
\end{definition}

\begin{lemma}[о числе опорных векторов]\label{lem:count}
Число опорных векторов конечно и не превосходит
\[
  \binom{n}{m} = \frac{n!}{m!\,(n-m)!} .
\]
\end{lemma}

\inputfig{figures/fig-matrix-selection}

\begin{proof}
Выделим из $A[M,N]$ всевозможные квадратные подматрицы вида $A[M,M']$,
составленные из $\abs{M} = m$ столбцов. Если $\det A[M,M'] = 0$, такая подматрица
исключается из рассмотрения; остаются только подматрицы с ненулевым определителем.
Всего наборов из $m$ столбцов, выбранных из $n$, ровно $\binom{n}{m}$, значит
невырожденных подматриц не больше $\binom{n}{m}$.

Для каждой оставшейся подматрицы решаем систему
\[
  A[M,M']\, x[M'] = b[M] ,
\]
полагая остальные компоненты нулями. Решение единственно, поскольку определитель
$\det A[M,M']$ отличен от нуля. Если полученный вектор удовлетворяет ограничению на знак,
$x[i] \ge 0$ для всех $i \in M'$, то он принимается к рассмотрению. Так получаются
все опорные векторы, и их не больше числа перебранных базисов.
\end{proof}

\begin{remark}
Одному взятому базису, составленному из столбцов $A$, соответствует \emph{ровно один}
опорный вектор. Обратное неверно: вырожденному \ov{} может отвечать несколько
базисов, поэтому оценка $\binom{n}{m}$ не точна.
\end{remark}

\begin{example}\label{ex:six-bases}
Рассмотрим задачу
\[
  \begin{cases}
    5x_1 + 5x_2 + 6x_3 + 4x_4 = 5, \\
    4x_1 + 5x_2 - 5x_3 - 4x_4 = 4, \\
    x_i \ge 0, \quad i = \overline{1,4} .
  \end{cases}
\]
Здесь $m = 2$, $n = 4$, столбцы матрицы
\[
  a_1 = \begin{pmatrix} 5 \\ 4 \end{pmatrix}, \quad
  a_2 = \begin{pmatrix} 5 \\ 5 \end{pmatrix}, \quad
  a_3 = \begin{pmatrix} 6 \\ -5 \end{pmatrix}, \quad
  a_4 = \begin{pmatrix} 4 \\ -4 \end{pmatrix} ,
\]
и по лемме~\ref{lem:count} перебираем не более $\binom{4}{2} = 6$ базисов.

\medskip
\noindent\textbf{1)} Базис $a_1, a_2$ (остальные переменные — нули):
\[
  \begin{cases} 5x_1 + 5x_2 = 5 \\ 4x_1 + 5x_2 = 4 \end{cases}
  \implies (x_1, x_2) = (1, 0),
  \qquad x_{\mathrm{I}}\T = (1, 0, 0, 0) \ \text{— \ov} .
\]

\noindent\textbf{2)} Базис $a_1, a_3$:
\[
  \begin{cases} 5x_1 + 6x_3 = 5 \\ 4x_1 - 5x_3 = 4 \end{cases}
  \implies (x_1, x_3) = (1, 0),
  \qquad x_{\mathrm{II}}\T = (1, 0, 0, 0) \ \text{— \ov} .
\]

\noindent\textbf{3)} Базис $a_1, a_4$:
\[
  \begin{cases} 5x_1 + 4x_4 = 5 \\ 4x_1 - 4x_4 = 4 \end{cases}
  \implies (x_1, x_4) = (1, 0),
  \qquad x_{\mathrm{III}}\T = (1, 0, 0, 0) \ \text{— \ov} .
\]

\noindent\textbf{4)} Базис $a_2, a_3$:
\[
  \begin{cases} 5x_2 + 6x_3 = 5 \\ 5x_2 - 5x_3 = 4 \end{cases}
  \implies (x_2, x_3) = \Big(\tfrac{49}{55},\ \tfrac{1}{11}\Big),
  \qquad x_{\mathrm{IV}}\T = \Big(0, \tfrac{49}{55}, \tfrac{1}{11}, 0\Big)
  \ \text{— \ov} .
\]

\noindent\textbf{5)} Базис $a_2, a_4$:
\[
  \begin{cases} 5x_2 + 4x_4 = 5 \\ 5x_2 - 4x_4 = 4 \end{cases}
  \implies (x_2, x_4) = \Big(\tfrac{9}{10},\ \tfrac{1}{8}\Big),
  \qquad x_{\mathrm{V}}\T = \Big(0, \tfrac{9}{10}, 0, \tfrac{1}{8}\Big)
  \ \text{— \ov} .
\]

\noindent\textbf{6)} Базис $a_3, a_4$:
\[
  \begin{cases} 6x_3 + 4x_4 = 5 \\ -5x_3 - 4x_4 = 4 \end{cases}
  \implies (x_3, x_4) = \Big(9,\ -\tfrac{49}{4}\Big),
  \qquad x_{\mathrm{VI}} \ \text{не существует, т.\,к. } -\tfrac{49}{4} < 0 .
\]
\end{example}

\begin{center}
\begin{tabular}{clll}
\toprule
№ & базис & решение & опорный вектор \\
\midrule
1 & $a_1, a_2$ & $(x_1, x_2) = (1, 0)$ & $(1,\,0,\,0,\,0)$ \\
2 & $a_1, a_3$ & $(x_1, x_3) = (1, 0)$ & $(1,\,0,\,0,\,0)$ \\
3 & $a_1, a_4$ & $(x_1, x_4) = (1, 0)$ & $(1,\,0,\,0,\,0)$ \\
4 & $a_2, a_3$ & $(x_2, x_3) = \big(\tfrac{49}{55}, \tfrac{1}{11}\big)$
  & $\big(0,\,\tfrac{49}{55},\,\tfrac{1}{11},\,0\big)$ \\
5 & $a_2, a_4$ & $(x_2, x_4) = \big(\tfrac{9}{10}, \tfrac{1}{8}\big)$
  & $\big(0,\,\tfrac{9}{10},\,0,\,\tfrac{1}{8}\big)$ \\
6 & $a_3, a_4$ & $(x_3, x_4) = \big(9, -\tfrac{49}{4}\big)$ & нет: $x_4 < 0$ \\
\bottomrule
\end{tabular}
\end{center}

\begin{idea}
Из шести базисов годными оказались пять, но \emph{различных} опорных векторов только
три:
\[
  (1,0,0,0), \qquad \Big(0, \tfrac{49}{55}, \tfrac{1}{11}, 0\Big),
  \qquad \Big(0, \tfrac{9}{10}, 0, \tfrac{1}{8}\Big) .
\]
Первый из них вырожден: $\abs{\Nzp} = 1 < 2 = \abs{M}$, и ему отвечают сразу три
базиса ($a_1,a_2$; $a_1,a_3$; $a_1,a_4$) — все они содержат обязательный столбец
$a_1$ и один «лишний», которому соответствует нулевая компонента. Два других \ov{}
невырождены, у каждого ровно один базис. Это и есть обещанная неточность оценки
$\binom{n}{m}$: пять годных базисов против трёх опорных векторов.
\end{idea}

\subsection{Связь опорных векторов и крайних точек}

\begin{theorem}[о совпадении опорного вектора и крайней точки]\label{th:support-extreme}
Пусть $x[N]$ — допустимая точка канонической задачи. Тогда
\[
  x[N] \ \text{— опорный вектор} \iff
  x[N] \ \text{— крайняя точка допустимого множества } S .
\]
\end{theorem}

\begin{idea}
В тетради в обоих пунктах написано $\Np$, но речь идёт именно о множестве
\emph{положительных} компонент: опорность по определению~\ref{def:support-vector}
проверяется для столбцов $\{a_j\}_{j \in \Nzp}$. Ниже везде пишем $\Nzp$ и $\Nzz$ —
это важно для корректности оценки $\theta > 0$ в первом пункте.
\end{idea}

\begin{proof}
\textbf{1) $(\Longleftarrow)$: крайняя точка $\implies$ опорный вектор.}
От противного: пусть $x[N]$ — \emph{не} опорный вектор. Тогда столбцы
$A[M,\Nzp]$ линейно зависимы, то есть
\[
  \exists\, u[\Nzp] \ne 0 : \qquad A[M, \Nzp]\, u[\Nzp] = 0 .
\]
Доопределим $u$ нулями вне $\Nzp$, то есть положим $u[i] = 0$ для всех
$i \in \Nzz$; тогда $A[M,N]\,u[N] = 0$.

Положим
\[
  x_1[N] = x[N] - \theta\, u[N], \qquad
  x_2[N] = x[N] + \theta\, u[N],
\]
где $\theta = \min\{\theta_1, \theta_2\}$ и
\[
  \theta_1 = \min_{\substack{i \in \Nzp \\ u[i] > 0}} \frac{x[i]}{u[i]} ,
  \qquad
  \theta_2 = \min_{\substack{i \in \Nzp \\ u[i] < 0}} \frac{-x[i]}{u[i]} .
\]
(Если один из наборов индексов пуст, соответствующий минимум считаем равным
$+\infty$.)

По подбору $\theta$ векторы $x_1$ и $x_2$ удовлетворяют ограничению на знак:
на $\Nzz$ обе поправки нулевые, а на $\Nzp$ сдвиг не выводит за ноль ни в одну,
ни в другую сторону. Кроме того $\theta > 0$, так как $x[i] > 0$ при $i \in \Nzp$,
и нулевых числителей там нет.

Проверим допустимость:
\begin{align*}
  A[M,N]\,x_1[N] &= A[M,N]\,x[N] - \theta\,\underbrace{A[M,N]\,u[N]}_{=\,0} , \\
  A[M,N]\,x_2[N] &= A[M,N]\,x[N] + \theta\,\underbrace{A[M,N]\,u[N]}_{=\,0} ,
\end{align*}
откуда
\[
  A[M,N]\,x_1[N] = A[M,N]\,x[N] = b[M] = A[M,N]\,x_2[N] .
\]
Значит $x_1, x_2 \in S$, причём $x_1 \ne x_2$ (так как $\theta > 0$ и $u \ne 0$), и
\[
  x[N] = \tfrac{1}{2}\big(x_1[N] + x_2[N]\big) .
\]
Получили противоречие с тем, что $x$ — крайняя точка.

\medskip
\textbf{2) $(\Longrightarrow)$: опорный вектор $\implies$ крайняя точка.}
От противного: пусть $x[N]$ — \emph{не} крайняя точка. Тогда
\[
  \exists\, x_1[N], x_2[N] \in S,\ x_1 \ne x_2 : \qquad
  x[N] = \tfrac{1}{2}\big(x_1[N] + x_2[N]\big) .
\]
Так как $x_1, x_2 \ge 0$, а на множестве $\Nzz$ полусумма равна нулю, обе точки
обращаются там в ноль:
\[
  x[\Nzz] = 0 \implies x_1[\Nzz] = 0, \quad x_2[\Nzz] = 0 .
\]
Поэтому ограничения для них сводятся к системам только по базисным индексам:
\[
  \begin{cases}
    A[M, \Nzp]\, x_1[\Nzp] = b[M], \\
    A[M, \Nzp]\, x_2[\Nzp] = b[M] .
  \end{cases}
\]
Вычитая, получаем
\[
  A[M, \Nzp]\big(x_1[\Nzp] - x_2[\Nzp]\big) = 0 ,
\]
причём $x_1[\Nzp] \ne x_2[\Nzp]$ (вне $\Nzp$ обе точки нулевые, так что различаться
они могут только там). Значит столбцы $A[M,\Nzp]$ линейно зависимы, то есть
$x[N]$ — не опорный вектор. Противоречие.
\end{proof}

\inputfig{figures/fig-theta-segment}

\subsection{Оптимальное решение достигается в опорном векторе}

\begin{theorem}\label{th:optimal-is-support}
Пусть $\bar{x}_*[N]$ — оптимальное решение канонической задачи \emph{с минимальным
числом положительных компонент}. Тогда $\bar{x}_*[N]$ — опорный вектор.
\end{theorem}

\begin{proof}
Рассмотрим разбиение индексов на $\Nzp$ и $\Nzz$ для точки $\bar{x}_*$. От
противного: пусть $\bar{x}_*$ — не опорный вектор. Тогда, как в
теореме~\ref{th:support-extreme},
\[
  \exists\, u[\Nzp] \ne 0 : \quad A[M,\Nzp]\, u[\Nzp] = 0 ,
\]
доопределённое нулями до $u[N]$, так что $A[M,N]\,u[N] = 0$.

Положим
\[
  x_\theta[N] = \bar{x}_*[N] - \theta\, u[N] .
\]
Эта точка допустима при всех $\theta \in [-\theta_0, \theta_0]$, где
\[
  \theta_0 = \min_{\substack{i \in \Nzp \\ u[i] \ne 0}}
             \frac{\bar{x}_*[i]}{\abs{u[i]}} > 0 :
\]
равенства $A x_\theta = b$ выполнены автоматически (поправка лежит в ядре $A$),
а ограничение на знак — по выбору $\theta_0$.

Значение целевой функции:
\[
  c\T[N]\, x_\theta[N] = c\T[N]\, \bar{x}_*[N] - \theta\, c\T[N]\, u[N] .
\]
Так как $x_\theta$ допустима при \emph{обоих} знаках $\theta$, знак $\theta$ можно
выбрать в зависимости от знака $c\T[N]\,u[N]$ так, чтобы
$\theta\, c\T[N]\,u[N] \ge 0$, то есть чтобы целевая функция не возрастала. Но по
определению оптимального решения не может быть
\[
  c\T[N]\, x_\theta[N] < c\T[N]\, \bar{x}_*[N] ,
\]
следовательно
\[
  c\T[N]\, u[N] = 0 \quad\implies\quad
  c\T[N]\, x_\theta[N] = c\T[N]\, \bar{x}_*[N]
  \quad \forall \theta \in [-\theta_0, \theta_0] ,
\]
то есть $x_\theta$ — тоже оптимальное решение.

Возьмём теперь $\theta = \theta_0$. Минимум в определении $\theta_0$ достигается на
некотором индексе $i_0 \in \Nzp$, и для него $x_{\theta_0}[i_0] = 0$. Остальные
компоненты остались неотрицательными, а компоненты из $\Nzz$ — нулевыми. Значит
$x_{\theta_0}$ — оптимальное решение, у которого положительных компонент
\emph{строго меньше}, чем у $\bar{x}_*$. Противоречие с минимальностью.
\end{proof}

\begin{corollary*}
Оптимальное решение можно найти, перебрав все крайние точки (= опорные векторы) и
посчитав в каждой значение целевой функции. По лемме~\ref{lem:count} перебор конечен.
\end{corollary*}

Но такой перебор может быть вычислительно трудным: число вариантов $\binom{n}{m}$
растёт экспоненциально.

\subsection{Симплекс-метод: идея и вспомогательная задача}

\begin{definition}[симплекс-метод]
\textbf{Симплекс-метод} — метод, основанный на \emph{направленном} переборе крайних
точек так, чтобы значение целевой функции в следующей точке было не хуже, чем в
предыдущей. Является итерационной схемой.
\end{definition}

\begin{idea}
Разница с полным перебором из следствия к теореме~\ref{th:optimal-is-support}: вместо
того чтобы обойти все $\binom{n}{m}$ базисов, на каждом шаге меняется ровно один
столбец базиса, причём так, чтобы целевая функция не возрастала. Монотонность
гарантирует, что ни один базис не повторится (при отсутствии вырожденности), а
конечность числа базисов — что алгоритм остановится. Подробный разбор шага —
в методичках.
\end{idea}

Чтобы запустить итерации, нужна стартовая крайняя точка. Её даёт вспомогательная
задача.

\subsubsection*{Вспомогательная задача симплекс-метода (метод искусственного базиса)}

Для исходной задачи
\[
  \begin{cases}
    c\T[N]\, x[N] \to \min, \\
    A[M,N]\, x[N] = b[M], \\
    x[N] \ge 0
  \end{cases}
\]
строим \textbf{вспомогательную задачу}, добавляя искусственные переменные $y[M]$:
\[
  \begin{cases}
    \displaystyle\sum_{i \in M} y[i] \to \min, \\[6pt]
    A[M,N]\, x[N] + E[M,M]\, y[M] = b[M], \\
    x[N] \ge 0, \quad y[M] \ge 0 ,
  \end{cases}
\]
где $E[M,M]$ — единичная матрица.

У вспомогательной задачи целевая функция \textbf{ограничена снизу} (она
неотрицательна), а множество допустимых точек \textbf{непусто}. Для неё сразу
известно начальное приближение:
\[
  x[N] = 0, \qquad y[M] = b[M] .
\]

К этому решению вспомогательной задачи применяем симплекс-метод. Если в результате
окажется, что $y[M] = 0$, то соответствующий $x[N]$ — опорный вектор исходной задачи,
с которого и стартуют основные итерации.

\begin{idea}
Начальное приближение $y[M] = b[M]$ требует $b[M] \ge 0$. Это не ограничение
общности: строку с $b_i < 0$ домножают на $-1$. Столбцы $E[M,M]$ линейно независимы,
поэтому точка $(x, y) = (0, b)$ — опорный вектор вспомогательной задачи, и запуск
корректен.

Дальше возможны два исхода:
\begin{itemize}[nosep,left=0pt]
  \item минимум вспомогательной задачи равен нулю, то есть $y[M] = 0$ — тогда
        найденный $x[N]$ допустим в исходной задаче и является её опорным вектором;
  \item минимум строго положителен — тогда допустимое множество исходной задачи
        \emph{пусто}, и исходная задача неразрешима.
\end{itemize}
Так вспомогательная задача одновременно даёт старт и проверяет совместность.
\end{idea}

\begin{idea}[title={Откуда этот материал}]
Дальнейшие параграфы этой главы на разобранных лекциях подробно не
записывались — вычислительная схема симплекс-метода была отнесена к
методичкам. Ниже она восстановлена полностью, в обозначениях курса, вместе с
доказательствами; это материал экзаменационного вопроса «Симплекс-метод.
Метод искусственного базиса».
\end{idea}

\subsection{Метод перебора крайних точек}
\label{subsec:enumeration}

Следствие из теоремы~\ref{th:optimal-is-support} даёт прямой, хотя и
неэффективный, способ решения.

\begin{algo}[title={Перебор крайних точек}]
\begin{steps}
  \item Выбрать набор базисных индексов $N_k \subset N$, $\abs{N_k} = m$.
  \item Построить квадратную матрицу $A[M, N_k]$.
  \item Проверить $\det A[M,N_k]$. Если $\det A[M,N_k] = 0$ — набор
        отбрасывается; иначе идти дальше.
  \item Решить систему $A[M,N_k]\,x[N_k] = b[M]$.
  \item Небазисные переменные положить нулю: $x[N \setminus N_k] = 0$.
  \item Проверить $x[N] \ge 0$. Если есть отрицательные компоненты — точка
        отбрасывается; иначе получена допустимая крайняя точка.
  \item Вычислить $F(x) = c\T[N]\,x[N]$.
  \item Повторить шаги 1--7 для всех наборов из $m$ столбцов.
  \item Среди найденных допустимых крайних точек выбрать точку с минимальным
        $F(x)$.
\end{steps}
\end{algo}

\begin{idea}
Метод конечен: наборов не больше $\binom{n}{m}$ (лемма~\ref{lem:count}).
Применим, когда задача записана в канонической форме, $\rank A = m$, $m < n$,
допустимое множество непусто и оптимум конечен. Если ни один допустимый
опорный вектор не найден — допустимое множество пусто. Неограниченность снизу
перебор сам по себе не обнаруживает: для этого нужен симплекс-метод.

Практическая негодность метода видна из оценки: при $n = 50$, $m = 25$ число
наборов $\binom{50}{25} \approx 1{,}26\cdot10^{14}$.
\end{idea}

\subsection{Симплекс-метод: вычислительная схема}
\label{subsec:simplex-scheme}

Симплекс-метод строит последовательность опорных векторов
$x^0[N], x^1[N], \dots$ так, чтобы значение целевой функции не возрастало:
\[
  c\T[N]\,x^{k+1}[N] \ \le\ c\T[N]\,x^{k}[N] .
\]

\subsubsection*{Обозначения}

Пусть на $k$-й итерации известен опорный вектор $x^k[N] \in S$ и его базис
$N_k \subset N$, $\abs{N_k} = m$, $\det A[M,N_k] \ne 0$. Обозначим
\[
  B[N_k, M] = A[M,N_k]^{-1} ,
  \qquad
  L_k = N \setminus N_k \quad \text{— небазисные индексы} ,
\]
при этом $x^k[L_k] = 0$, а базисные компоненты находятся из системы
$A[M,N_k]\,x^k[N_k] = b[M]$.

\subsubsection*{Вектор потенциалов и оценки}

Вектор $y^k[M]$ находится из системы
\[
  \big(y^k\big)\T[M]\, A[M,N_k] = c\T[N_k] ,
  \qquad\text{то есть}\qquad
  A\T[M,N_k]\, y^k[M] = c[N_k] ,
\]
откуда $\big(y^k\big)\T[M] = c\T[N_k]\,B[N_k,M]$. Для каждого небазисного
индекса $j \in L_k$ вычисляется \textbf{оценка}
\[
  d[j] = c[j] - \big(y^k\big)\T[M]\, a_j[M] ,
\]
где $a_j[M]$ — $j$-й столбец матрицы $A[M,N]$.

\begin{idea}
Это ровно та же конструкция, что и потенциалы транспортной задачи
(\S\ref{subsec:potentials-theory} и §7): $y^k$ — двойственные переменные,
$d[j]$ — «во сколько обойдётся» ввод переменной $x[j]$ в базис с учётом
вынужденного пересчёта базисных компонент. Отрицательная оценка означает, что
ввод выгоден.
\end{idea}

\begin{theorem}[критерий оптимальности]\label{thm:simplex-opt}
Если для текущего опорного вектора $x^k[N]$ выполняется $d[j] \ge 0$ для всех
$j \in L_k$, то $x^k[N]$ — оптимальное решение задачи ЛП.
\end{theorem}

\begin{proof}
Пусть $x[N] \in S$ — любое допустимое решение. Оба вектора удовлетворяют
системе, $A[M,N]x[N] = b[M] = A[M,N]x^k[N]$, поэтому
$A[M,N]\big(x[N] - x^k[N]\big) = 0$. Разбивая переменные на базисные и
небазисные,
\begin{gather*}
  A[M,N_k]\big(x[N_k] - x^k[N_k]\big) + A[M,L_k]\,x[L_k] = 0 , \\
  x[N_k] - x^k[N_k] = -B[N_k,M]\,A[M,L_k]\,x[L_k] .
\end{gather*}
Рассмотрим разность значений целевой функции:
\begin{align*}
  c\T[N]x[N] - c\T[N]x^k[N]
  &= c\T[N_k]\big(x[N_k] - x^k[N_k]\big) + c\T[L_k]\,x[L_k] \\
  &= -\,\underbrace{c\T[N_k]B[N_k,M]}_{=\ (y^k)\T[M]}A[M,L_k]\,x[L_k]
     + c\T[L_k]\,x[L_k] \\
  &= \sum_{j \in L_k}
     \underbrace{\Big(c[j] - (y^k)\T[M]\,a_j[M]\Big)}_{=\ d[j]}\, x[j] .
\end{align*}
При $d[j] \ge 0$ и $x[j] \ge 0$ сумма неотрицательна, значит
$c\T[N]x[N] \ge c\T[N]x^k[N]$: вектор $x^k[N]$ оптимален.
\end{proof}

\subsubsection*{Направление перехода}

Если существует индекс $j_k \in L_k$ с $d[j_k] < 0$, текущий опорный вектор не
оптимален, и переменная $x[j_k]$ \textbf{вводится в базис}. Направление
$u^k[N]$ строится так:
\[
  u^k[N_k] = B[N_k,M]\,a_{j_k}[M] ,
  \qquad
  u^k[j_k] = -1 ,
  \qquad
  u^k\big[L_k \setminus \{j_k\}\big] = 0 ,
\]
и рассматривается $x_\theta[N] = x^k[N] - \theta\,u^k[N]$, $\theta \ge 0$.

\begin{lemma}\label{lem:simplex-dir}
$A[M,N]\,u^k[N] = 0$, и потому $A[M,N]\,x_\theta[N] = b[M]$ при любом $\theta$.
\end{lemma}

\begin{proof}
По построению
\[
  A[M,N]u^k[N] = A[M,N_k]u^k[N_k] + a_{j_k}[M]\,u^k[j_k]
    + A\big[M, L_k \setminus \{j_k\}\big]\,u^k\big[L_k\setminus\{j_k\}\big] .
\]
Подставляя $u^k[N_k] = B[N_k,M]a_{j_k}[M]$, $u^k[j_k] = -1$ и нулевой остаток,
\[
  A[M,N]u^k[N] = \underbrace{A[M,N_k]B[N_k,M]}_{=\ E}\,a_{j_k}[M] - a_{j_k}[M] = 0 .
  \qedhere
\]
\end{proof}

\subsubsection*{Изменение целевой функции}

\[
  F(x_\theta) = c\T[N]x_\theta[N] = F(x^k) - \theta\,c\T[N]u^k[N] .
\]
Вычислим $c\T[N]u^k[N] = c\T[N_k]u^k[N_k] + c[j_k]u^k[j_k]
= c\T[N_k]B[N_k,M]a_{j_k}[M] - c[j_k]$. Так как
$c\T[N_k]B[N_k,M] = (y^k)\T[M]$,
\[
  c\T[N]u^k[N] = (y^k)\T[M]a_{j_k}[M] - c[j_k] = -\,d[j_k]
  \qquad\Longrightarrow\qquad
  F(x_\theta) = F(x^k) + \theta\, d[j_k] .
\]
Так как $d[j_k] < 0$, при $\theta > 0$ получаем $F(x_\theta) < F(x^k)$.

\subsubsection*{Выбор шага}

Для сохранения $x_\theta[N] \ge 0$ нужно $x^k[i] - \theta u^k[i] \ge 0$ при всех
$i \in N_k$.
\begin{itemize}[nosep,left=0pt]
  \item Если $u^k[i] \le 0$ для всех $i \in N_k$, то неотрицательность не
        нарушается ни при каком $\theta \ge 0$, а $F(x_\theta) \to -\infty$ при
        $\theta \to +\infty$ — \textbf{задача не ограничена снизу}.
  \item Иначе шаг выбирается по правилу минимального отношения
        \[
          \theta_k = \min_{\substack{i \in N_k \\ u^k[i] > 0}}
                     \frac{x^k[i]}{u^k[i]} .
        \]
\end{itemize}
Пусть минимум достигается на $i_k$. Тогда $x_{\theta_k}[i_k] = 0$: переменная
$x[i_k]$ \textbf{выходит из базиса}, $x[j_k]$ входит,
\[
  N_{k+1} = \big(N_k \setminus \{i_k\}\big) \cup \{j_k\} ,
  \qquad
  x^{k+1}[N] = x^k[N] - \theta_k u^k[N] .
\]

\begin{algo}[title={Симплекс-метод}]
Известен начальный опорный вектор $x^0[N]$ и его базис $N_0$.
\begin{steps}
  \item На $k$-й итерации известны $x^k[N]$, $N_k$, $L_k = N \setminus N_k$.
  \item Построить $B[N_k,M] = A[M,N_k]^{-1}$.
  \item Найти $y^k[M]$ из $\big(y^k\big)\T[M]A[M,N_k] = c\T[N_k]$.
  \item Для каждого $j \in L_k$ вычислить
        $d[j] = c[j] - \big(y^k\big)\T[M]a_j[M]$.
  \item Если $d[j] \ge 0$ для всех $j \in L_k$ — $x^k[N]$ оптимален, конец.
  \item Иначе выбрать $j_k \in L_k$ с $d[j_k] < 0$.
  \item Построить $u^k[N_k] = B[N_k,M]a_{j_k}[M]$, $u^k[j_k] = -1$,
        $u^k[L_k\setminus\{j_k\}] = 0$.
  \item Если $u^k[i] \le 0$ при всех $i \in N_k$ — задача не ограничена снизу,
        конец.
  \item Иначе найти
        $\theta_k = \min\limits_{i \in N_k,\, u^k[i]>0} x^k[i]/u^k[i]$;
        пусть минимум достигается на $i_k$.
  \item $x^{k+1}[N] = x^k[N] - \theta_k u^k[N]$.
  \item $N_{k+1} = \big(N_k \setminus \{i_k\}\big) \cup \{j_k\}$.
  \item $k := k+1$, перейти к шагу~2.
\end{steps}
\end{algo}

\subsubsection*{Вырожденность и зацикливание}

Опорный вектор вырожден, если $\abs{\Nzp} < m$ (определение~\ref{def:degenerate}).
Тогда некоторые базисные компоненты равны нулю, и при выборе шага может
получиться $\theta_k = 0$: опорный вектор не меняется, меняется лишь базис. В
вырожденных задачах возможно \textbf{зацикливание} — возврат к уже
встречавшемуся базису. Средство — \textbf{правило Бленда}: при нескольких
кандидатах (и при выборе вводимой, и при выборе выводимой переменной) брать
переменную с наименьшим индексом.

\begin{idea}
Правило Бленда гарантирует конечность: базисов конечное число
(лемма~\ref{lem:count}), а при правиле наименьшего индекса ни один базис не
повторяется. Платой является возможное замедление — обычно выгоднее брать
$j_k$ с наиболее отрицательной оценкой $d[j_k]$, и переключаться на правило
Бленда только при подозрении на зацикливание.
\end{idea}

\subsection{Метод искусственного базиса: обоснование}
\label{subsec:artificial-basis}

Вспомогательная задача построена в~\S4:
\[
  \begin{cases}
    \displaystyle\sum_{i \in M} y[i] \to \min , \\[6pt]
    A[M,N]\,x[N] + E[M,M]\,y[M] = b[M] , \\
    x[N] \ge 0 , \quad y[M] \ge 0 ,
  \end{cases}
  \qquad\text{старт: } x[N] = 0,\ y[M] = b[M] \ge 0 .
\]

\begin{theorem}\label{thm:artificial-basis}
Пусть вспомогательная задача решена и $\big(x^*[N], y^*[M]\big)$ — её
оптимальное решение. Тогда
\begin{enumerate}[label=\arabic*),left=0pt,nosep]
  \item если $\sum_{i \in M} y^*[i] > 0$, то исходная задача не имеет
        допустимых решений;
  \item если $\sum_{i \in M} y^*[i] = 0$, то $x^*[N]$ — допустимое решение
        исходной задачи.
\end{enumerate}
\end{theorem}

\begin{proof}
Для любой допустимой точки вспомогательной задачи $y[M] \ge 0$, поэтому
$\sum_{i\in M} y[i] \ge 0$: минимум не может быть отрицательным.

\medskip
\textbf{Случай 1.} Пусть $\sum_{i\in M} y^*[i] > 0$. От противного: пусть
исходная задача имеет допустимое решение $\hat{x}[N]$, то есть
$A[M,N]\hat{x}[N] = b[M]$, $\hat{x}[N] \ge 0$. Положим $\hat{y}[M] = 0$; тогда
\[
  A[M,N]\hat{x}[N] + E[M,M]\hat{y}[M] = A[M,N]\hat{x}[N] = b[M] ,
\]
то есть пара $\big(\hat{x}[N], \hat{y}[M]\big)$ допустима во вспомогательной
задаче, и значение её целевой функции равно нулю. Это противоречит
положительности минимума. Значит исходная задача несовместна.

\medskip
\textbf{Случай 2.} Пусть $\sum_{i\in M} y^*[i] = 0$. Так как $y^*[i] \ge 0$,
из равенства суммы нулю следует $y^*[M] = 0$. Подставляя в ограничения,
$A[M,N]x^*[N] = b[M]$, а из допустимости $x^*[N] \ge 0$. Значит
$x^*[N] \in S$.
\end{proof}

\begin{algo}[title={Метод искусственного базиса}]
Дана задача $c\T[N]x[N] \to \min$, $A[M,N]x[N] = b[M]$, $x[N] \ge 0$.
\begin{steps}
  \item Добиться $b[M] \ge 0$: строку с $b[i] < 0$ умножить на $-1$.
  \item Ввести искусственные переменные $y[M] \ge 0$.
  \item Составить вспомогательную задачу $\sum_{i \in M} y[i] \to \min$ при
        $A[M,N]x[N] + E[M,M]y[M] = b[M]$, $x[N] \ge 0$, $y[M] \ge 0$.
  \item Взять начальный искусственный базис $x[N] = 0$, $y[M] = b[M]$.
  \item Решить вспомогательную задачу симплекс-методом.
  \item Если $\sum_{i\in M} y^*[i] > 0$ — исходная задача несовместна, конец.
  \item Если $\sum_{i\in M} y^*[i] = 0$, то $y^*[M] = 0$, а $x^*[N]$ — допустимый
        вектор исходной задачи.
  \item Удалить искусственные переменные. Если какие-то из них остались в
        базисе с нулевым значением, заменить их столбцами исходной матрицы
        $A[M,N]$ так, чтобы базисные столбцы остались линейно независимыми.
  \item Полученный базис использовать как начальный для решения исходной задачи
        обычным симплекс-методом.
\end{steps}
\end{algo}

\begin{idea}
Шаг~8 — тонкое место. Искусственная переменная может остаться в базисе с
нулевым значением, если исходная система содержит линейно зависимую строку или
если допустимое множество вырождено. Замена столбца делается так: для
искусственного индекса $i$ ищут небазисный исходный столбец $a_j$ с ненулевым
элементом в $i$-й позиции преобразованной матрицы и вводят его в базис с
$\theta = 0$. Если такого столбца нет, соответствующее уравнение является
следствием остальных и его можно вычеркнуть.
\end{idea}

\subsubsection*{Численные примеры}

\begin{idea}[title={Откуда этот материал}]
На лекции симплекс-метод и метод искусственного базиса разбирались только в
общем виде. Ниже оба алгоритма прогнаны на числах — по тем же формулам, что в
\S\ref{subsec:simplex-scheme} и \S\ref{subsec:artificial-basis}.
\end{idea}

\begin{example}[Симплекс-метод]\label{ex:simplex-run}
\[
  \begin{cases}
    F(x) = -3x_1 - x_2 \to \min , \\
    x_1 + x_2 + x_3 = 4 , \\
    2x_1 + x_2 + x_4 = 5 , \\
    x_j \ge 0 , \quad j = \overline{1,4} .
  \end{cases}
\]
Здесь $m = 2$, $n = 4$,
\[
  A[M,N] = \begin{pmatrix} 1 & 1 & 1 & 0 \\ 2 & 1 & 0 & 1 \end{pmatrix} ,
  \qquad
  b[M] = \begin{pmatrix} 4 \\ 5 \end{pmatrix} ,
  \qquad
  c\T[N] = (-3, -1, 0, 0) .
\]
Переменные $x_3, x_4$ вошли в задачу как балансовые, поэтому начальный базис
виден сразу: $N_0 = \{3, 4\}$, $x^0[N] = (0, 0, 4, 5)\T$, $F(x^0) = 0$.

\medskip
\noindent\textbf{Итерация 0.} Так как $A[M,N_0] = E$, то и $B[N_0,M] = E$.
Потенциалы находятся из системы
\[
  A\T[M,N_0]\,y^0[M] = c[N_0] = (0,0)\T
  \qquad\Longrightarrow\qquad
  y^0[M] = (0,0)\T .
\]
Оценки небазисных индексов $L_0 = \{1,2\}$:
\[
  d[1] = c[1] - (y^0)\T a_1 = -3 - 0 = -3 , \qquad
  d[2] = c[2] - (y^0)\T a_2 = -1 - 0 = -1 .
\]
Обе отрицательны; берём наиболее отрицательную — $j_0 = 1$. Направление:
\[
  u^0[N_0] = B[N_0,M]\,a_1[M] = \begin{pmatrix} 1 \\ 2 \end{pmatrix} ,
  \qquad u^0[1] = -1 , \qquad u^0[2] = 0 ,
\]
то есть $u^0[3] = 1$, $u^0[4] = 2$. Шаг по правилу минимального отношения:
\[
  \theta_0 = \min\Big\{ \frac{x^0[3]}{u^0[3]},\ \frac{x^0[4]}{u^0[4]} \Big\}
           = \min\Big\{ \frac{4}{1},\ \frac{5}{2} \Big\} = \frac52 ,
  \qquad i_0 = 4 .
\]
Новый опорный вектор $x^1[N] = x^0[N] - \theta_0 u^0[N]$:
\[
  x^1[1] = 0 - \tfrac52(-1) = \tfrac52 , \quad
  x^1[2] = 0 , \quad
  x^1[3] = 4 - \tfrac52 = \tfrac32 , \quad
  x^1[4] = 5 - \tfrac52\cdot 2 = 0 ,
\]
\[
  x^1[N] = \big(\tfrac52,\ 0,\ \tfrac32,\ 0\big)\T ,
  \qquad N_1 = \{3, 1\} ,
  \qquad F(x^1) = -\tfrac{15}{2} .
\]
Контроль: $F(x^1) = F(x^0) + \theta_0 d[j_0] = 0 + \tfrac52 \cdot (-3)
= -\tfrac{15}{2}$ — совпадает.

\medskip
\noindent\textbf{Итерация 1.} Базисные столбцы $a_3 = (1,0)\T$, $a_1 = (1,2)\T$:
\[
  A[M,N_1] = \begin{pmatrix} 1 & 1 \\ 0 & 2 \end{pmatrix} ,
  \qquad
  B[N_1,M] = \begin{pmatrix} 1 & -\tfrac12 \\[2pt] 0 & \tfrac12 \end{pmatrix} .
\]
Потенциалы из $A\T[M,N_1]\,y^1[M] = c[N_1] = (0, -3)\T$:
\[
  \begin{cases}
    y^{(1)} = 0 , \\
    y^{(1)} + 2 y^{(2)} = -3 ,
  \end{cases}
  \qquad\Longrightarrow\qquad
  y^1[M] = \big(0,\ -\tfrac32\big)\T .
\]
Оценки по $L_1 = \{2,4\}$:
\[
  d[2] = -1 - \big(0\cdot 1 + (-\tfrac32)\cdot 1\big) = \tfrac12 \ge 0 ,
  \qquad
  d[4] = 0 - \big(0\cdot 0 + (-\tfrac32)\cdot 1\big) = \tfrac32 \ge 0 .
\]
Все оценки неотрицательны, значит по теореме~\ref{thm:simplex-opt}
\[
  x_*[N] = \big(\tfrac52,\ 0,\ \tfrac32,\ 0\big)\T ,
  \qquad F(x_*) = -\tfrac{15}{2} .
\]
\end{example}

\begin{idea}
Проверка геометрией: в исходных переменных задача — это
$3x_1 + x_2 \to \max$ при $x_1 + x_2 \le 4$, $2x_1 + x_2 \le 5$, $x \ge 0$.
Вершины допустимого многоугольника: $(0,0)$, $(0,4)$, $(1,3)$, $(2{,}5;\,0)$ со
значениями $0$, $4$, $6$, $7{,}5$. Максимум $7{,}5$ в точке $(2{,}5;\,0)$ — то
же, что дал метод за один шаг.
\end{idea}

\begin{example}[Метод искусственного базиса]\label{ex:artificial-run}
Требуется найти допустимое решение системы
\[
  \begin{cases}
    x_1 + x_2 = 3 , \\
    x_1 - x_2 = 1 , \\
    x_1, x_2 \ge 0 .
  \end{cases}
\]
Очевидного базиса нет. Правая часть уже неотрицательна, поэтому сразу вводим
искусственные переменные $y_1, y_2$ (индексы $3$ и $4$) и решаем вспомогательную
задачу
\[
  \begin{cases}
    W = y_1 + y_2 \to \min , \\
    x_1 + x_2 + y_1 = 3 , \\
    x_1 - x_2 + y_2 = 1 , \\
    x_1, x_2, y_1, y_2 \ge 0 ,
  \end{cases}
  \qquad
  A'[M,N] = \begin{pmatrix} 1 & 1 & 1 & 0 \\ 1 & -1 & 0 & 1 \end{pmatrix} ,
  \quad {c'}\T[N] = (0,0,1,1) .
\]
Старт: $N_0 = \{3,4\}$, $x^0[N] = (0,0,3,1)\T$, $W = 4$.

\medskip
\noindent\textbf{Итерация 0.} $A'[M,N_0] = E$, $B = E$, $y^0[M] = c'[N_0] = (1,1)\T$.
\[
  d[1] = 0 - (1\cdot 1 + 1\cdot 1) = -2 , \qquad
  d[2] = 0 - \big(1\cdot 1 + 1\cdot(-1)\big) = 0 .
\]
Выбираем $j_0 = 1$. Направление $u^0[N_0] = B\,a_1 = (1,1)\T$, то есть
$u^0[3] = u^0[4] = 1$, $u^0[1] = -1$, $u^0[2] = 0$. Шаг
\[
  \theta_0 = \min\Big\{ \frac31,\ \frac11 \Big\} = 1 , \qquad i_0 = 4
  \quad (\text{из базиса выходит } y_2) .
\]
\[
  x^1[N] = (1,\ 0,\ 2,\ 0)\T , \qquad N_1 = \{3,1\} , \qquad W = 2 .
\]

\medskip
\noindent\textbf{Итерация 1.} Базисные столбцы $a_3 = (1,0)\T$, $a_1 = (1,1)\T$:
\[
  A'[M,N_1] = \begin{pmatrix} 1 & 1 \\ 0 & 1 \end{pmatrix} ,
  \qquad
  B[N_1,M] = \begin{pmatrix} 1 & -1 \\ 0 & 1 \end{pmatrix} ,
  \qquad
  c'[N_1] = (1, 0)\T .
\]
Потенциалы: $y^{(1)} = 1$, $y^{(1)} + y^{(2)} = 0 \implies y^1[M] = (1,-1)\T$.
Оценки по $L_1 = \{2,4\}$:
\[
  d[2] = 0 - \big(1\cdot 1 + (-1)\cdot(-1)\big) = -2 < 0 ,
  \qquad
  d[4] = 1 - \big(1\cdot 0 + (-1)\cdot 1\big) = 2 .
\]
Выбираем $j_1 = 2$:
\[
  u^1[N_1] = B[N_1,M]\,a_2[M]
    = \begin{pmatrix} 1 & -1 \\ 0 & 1 \end{pmatrix}
      \begin{pmatrix} 1 \\ -1 \end{pmatrix}
    = \begin{pmatrix} 2 \\ -1 \end{pmatrix} ,
\]
то есть $u^1[3] = 2$, $u^1[1] = -1$, $u^1[2] = -1$, $u^1[4] = 0$. Положительна
только компонента при $i = 3$, поэтому
\[
  \theta_1 = \frac{x^1[3]}{u^1[3]} = \frac22 = 1 , \qquad i_1 = 3
  \quad (\text{из базиса выходит } y_1) .
\]
\[
  x^2[N] = (2,\ 1,\ 0,\ 0)\T , \qquad N_2 = \{1,2\} , \qquad W = 0 .
\]
Минимум вспомогательной задачи равен нулю, искусственные переменные вышли из
базиса, значит по теореме~\ref{thm:artificial-basis} вектор $(2,1)\T$ допустим в
исходной системе: $2 + 1 = 3$, $2 - 1 = 1$. Базис $N_2 = \{1,2\}$ и берётся
начальным для решения самой задачи обычным симплекс-методом.
\end{example}

\begin{idea}
Обе искусственные переменные вышли из базиса сами, поэтому шаг~8 алгоритма
(замена столбцов) не понадобился. Так бывает не всегда: если бы система
содержала зависимую строку, одна из $y_i$ осталась бы в базисе с нулевым
значением.
\end{idea}


\newpage
\lecture{Лекция 2}{}

\section{Транспортная задача}

\subsection{Постановка}

Пусть есть $m$ \textbf{пунктов хранения} (поставщики) $A_1, \dots, A_m$ и $n$
\textbf{пунктов назначения} (потребители) $B_1, \dots, B_n$. Заданы:
\begin{itemize}[nosep,left=0pt]
  \item $\{a_i\}_{i=1}^{m}$ — количество единиц груза в каждом пункте хранения
        (запасы);
  \item $\{b_j\}_{j=1}^{n}$ — количество единиц потребности в грузе в каждом
        пункте назначения;
  \item $\norm{c_{ij}}$ — матрица стоимости перевозок из $i$-го в $j$-й пункт
        (тарифы).
\end{itemize}

\medskip
\noindent\textbf{Задача:} составить план перевозок, удовлетворяющий потребности и
минимизирующий общую стоимость.

\medskip
Обозначим через $x_{ij}$ объём перевозок из $i$-го в $j$-й пункт. Тогда
\[
  \begin{cases}
    \displaystyle\sum_{i=1}^{m} \sum_{j=1}^{n} c_{ij}\, x_{ij} \to \min, \\[12pt]
    \displaystyle\sum_{j=1}^{n} x_{ij} \le a_i, \quad \forall i = \overline{1,m}, \\[12pt]
    \displaystyle\sum_{i=1}^{m} x_{ij} = b_j, \quad \forall j = \overline{1,n}, \\[10pt]
    x_{ij} \ge 0, \quad \forall i, j ,
  \end{cases}
\]
где знак «$\le$» в ограничении по запасам стоит у \textbf{открытой} задачи,
а «$=$» — у \textbf{закрытой}.

\inputfig{figures/fig-transport-table}

\medskip
\noindent\textbf{Условие корректности постановки задачи по исходным данным:}
\[
  \text{открытая:} \quad \sum_{i=1}^{m} a_i \ne \sum_{j=1}^{n} b_j ,
  \qquad\qquad
  \text{закрытая:} \quad \sum_{i=1}^{m} a_i = \sum_{j=1}^{n} b_j .
\]

\begin{idea}
Смысл условия: суммарный вывоз не может превысить суммарный запас, а суммарный
завоз обязан покрыть все потребности. Если $\sum a_i = \sum b_j$, то оба
ограничения автоматически обращаются в равенства — весь груз вывезен и все
потребности закрыты. Если $\sum a_i > \sum b_j$, то часть груза остаётся на
складах, и ограничение по запасам должно быть нестрогим. Если же
$\sum a_i < \sum b_j$, то равенства $\sum_i x_{ij} = b_j$ выполнить нельзя
вовсе — допустимое множество пусто, и задачу приходится переформулировать.
\end{idea}

\subsection{Сведение открытой задачи к закрытой}

Для открытой транспортной задачи производится преобразование её к закрытой.

\begin{enumerate}[label=\arabic*),left=0pt]
  \item Пусть $\displaystyle\sum_{i=1}^{m} a_i < \sum_{j=1}^{n} b_j$. Вводим
        \textbf{фиктивного поставщика} $A_{\text{фикт}}$ с запасом
        \[
          a_{\text{фикт}} = \sum_{j=1}^{n} b_j - \sum_{i=1}^{m} a_i ,
          \qquad x_{\text{фикт},\,j} \ge 0 \quad \forall j = \overline{1,n} .
        \]
        Недопоставки будут распределены оптимальным образом.
  \item Аналогично при $\displaystyle\sum_{i=1}^{m} a_i > \sum_{j=1}^{n} b_j$
        вводим \textbf{фиктивного потребителя} $B_{\text{фикт}}$ с потребностью
        \[
          b_{\text{фикт}} = \sum_{i=1}^{m} a_i - \sum_{j=1}^{n} b_j .
        \]
\end{enumerate}

\begin{idea}
Тарифы фиктивной строки (столбца) берут нулевыми: $c_{\text{фикт},\,j} = 0$ —
фиктивная перевозка ничего не стоит и не меняет значения целевой функции.
Поэтому величина $x_{\text{фикт},\,j}$ в оптимальном плане читается как
«сколько единиц груза потребитель $B_j$ недополучил», а
$x_{i,\,\text{фикт}}$ — как «сколько осталось лежать на складе $A_i$».
Размер задачи вырастает на одну строку (столбец), после чего работает обычный
алгоритм для закрытой задачи.
\end{idea}

\subsection{Методы получения первого приближения}

Все три метода разбираются на одном примере.

\begin{example}\label{ex:transport}
Тарифы $c_{ij}$, запасы $a_i$ и потребности $b_j$:
\begin{center}
\begin{tabular}{c|cccc|c}
\toprule
      & $B_1$ & $B_2$ & $B_3$ & $B_4$ & $a_i$ \\
\midrule
$A_1$ & 14 & 28 & 21 & 28 & 27 \\
$A_2$ & 10 & 17 & 15 & 24 & 20 \\
$A_3$ & 14 & 30 & 25 & 21 & 43 \\
\midrule
$b_j$ & 33 & 13 & 27 & 17 & 90 \\
\bottomrule
\end{tabular}
\end{center}
Задача закрытая: $\sum a_i = 27 + 20 + 43 = 90 = 33 + 13 + 27 + 17 = \sum b_j$.
\end{example}

\subsubsection*{1) Метод северо-западного угла}

Заполняем левую верхнюю («северо-западную») свободную клетку на максимально
возможный объём $\min\{a_i, b_j\}$, после чего вычёркиваем исчерпавшуюся строку
или столбец и повторяем.

\medskip
\noindent Шаги:
\[
  \begin{array}{llll}
    (1,1): & \min\{a_1, b_1\} = 27, &\qquad (2,3): & \min\{a_2, b_3\} = 1, \\
    (2,1): & \min\{a_2, b_1\} = 6,  &\qquad (3,3): & \min\{a_3, b_3\} = 26, \\
    (2,2): & \min\{a_2, b_2\} = 13, &\qquad (3,4): & \min\{a_3, b_4\} = 17 .
  \end{array}
\]

\begin{center}
\begin{tabular}{c|cccc|l}
\toprule
      & $B_1$ & $B_2$ & $B_3$ & $B_4$ & остаток $a_i$ \\
\midrule
$A_1$ & \tc{14}{27} & \tcx{28} & \tcx{21} & \tcx{28} & $27,\ 0$ \\
$A_2$ & \tc{10}{6}  & \tc{17}{13} & \tc{15}{1} & \tcx{24} & $20,\ 14,\ 1,\ 0$ \\
$A_3$ & \tcx{14}    & \tcx{30} & \tc{25}{26} & \tc{21}{17} & $43,\ 17,\ 0$ \\
\midrule
остаток $b_j$ & $33,6,0$ & $13,0$ & $27,26,0$ & $17,0$ & \\
\bottomrule
\end{tabular}
\end{center}

Стоимость полученного плана:
\[
  14 \cdot 27 + 10 \cdot 6 + 17 \cdot 13 + 15 \cdot 1 + 25 \cdot 26 + 21 \cdot 17
  = 1681 .
\]

\subsubsection*{2) Метод минимального элемента}

Заполняем ту свободную клетку, в которой тариф $c_{ij}$ минимален, на объём
$\min\{a_i, b_j\}$; дальше так же.

\medskip
\noindent Шаги:
\[
  \begin{array}{llll}
    (2,1): & \min\{20, 33\} = 20, &\qquad (3,4): & \min\{17, 43\} = 17, \\
    (1,1): & \min\{13, 27\} = 13, &\qquad (3,3): & \min\{26, 13\} = 13, \\
    (1,3): & \min\{14, 27\} = 14, &\qquad (3,2): & \min\{13, 13\} = 13 .
  \end{array}
\]

\begin{center}
\begin{tabular}{c|cccc|l}
\toprule
      & $B_1$ & $B_2$ & $B_3$ & $B_4$ & остаток $a_i$ \\
\midrule
$A_1$ & \tc{14}{13} & \tcx{28} & \tc{21}{14} & \tcx{28} & $27,\ 14,\ 0$ \\
$A_2$ & \tc{10}{20} & \tcx{17} & \tcx{15} & \tcx{24} & $20,\ 0$ \\
$A_3$ & \tcx{14}    & \tc{30}{13} & \tc{25}{13} & \tc{21}{17} & $43,\ 26,\ 13,\ 0$ \\
\midrule
остаток $b_j$ & $33,13,0$ & $13,0$ & $27,13,0$ & $17,0$ & \\
\bottomrule
\end{tabular}
\end{center}

Стоимость:
\[
  14 \cdot 13 + 21 \cdot 14 + 10 \cdot 20 + 30 \cdot 13 + 25 \cdot 13 + 21 \cdot 17
  = 1748 .
\]

\begin{idea}
Здесь метод минимального элемента дал план \emph{дороже}, чем северо-западный угол
($1748$ против $1681$). Это не ошибка: оба метода строят лишь \emph{начальное}
приближение, и гарантий оптимальности ни у одного нет. На шаге~2 тариф $14$
встречается дважды — в клетках $(1,1)$ и $(3,1)$; выбрана первая, и дальше
дешёвые клетки первой строки оказались «съедены», а остаток пришлось везти по
тарифам $30$ и $25$. Обычно минимальный элемент всё же ближе к оптимуму, но
решающий шаг — не выбор начального плана, а последующие итерации метода
потенциалов.
\end{idea}

\subsubsection*{3) Метод двойного предпочтения}

На лекции не разбирался (см. методички).

\begin{idea}
Кратко: в каждой строке отмечают клетку с минимальным тарифом, затем то же самое
делают в каждом столбце. Клетки, отмеченные \emph{дважды} (минимум и по строке,
и по столбцу), заполняют первыми, потом отмеченные один раз, потом остальные —
в каждом случае на объём $\min\{a_i, b_j\}$. Метод — компромисс между дешевизной
северо-западного угла и качеством минимального элемента.
\end{idea}

Далее рассмотрим метод поиска оптимального решения — \textbf{метод (разности)
потенциалов}.

\subsection{Метод потенциалов}

Для полученного начального приближения проверяется, что число заполненных клеток
равно $m + n - 1$.

\begin{idea}
Это условие невырожденности опорного плана. Ограничений в закрытой задаче
$m + n$, но они линейно зависимы — сумма «строчных» равенств совпадает с суммой
«столбцовых» (и то, и другое равно $\sum a_i = \sum b_j$). Значит независимых
ограничений $m + n - 1$, и опорный вектор имеет ровно столько положительных
компонент (см. определение~\ref{def:support-vector}). Если заполненных клеток
меньше, план вырожден, и в одну из пустых клеток вписывают «нулевую перевозку»,
чтобы восстановить базис.

В примере~\ref{ex:transport} $m + n - 1 = 3 + 4 - 1 = 6$ — ровно столько клеток
и заполнено обоими методами.
\end{idea}

\subsubsection*{Алгоритм}

\paragraph{1) Проверка оптимальности решения.}
Сопоставим каждому столбцу переменную $v_j$, а каждой строке — переменную $u_i$;
эти переменные называются \textbf{потенциалами}. Для \emph{заполненных} клеток
выписываем равенства $v_j - u_i = c_{ij}$.

Для плана, построенного методом северо-западного угла в примере~\ref{ex:transport}:
\[
  \begin{cases}
    v_1 - u_1 = 14, \quad v_1 - u_2 = 10, \\
    v_2 - u_2 = 17, \quad v_3 - u_2 = 15, \\
    v_3 - u_3 = 25, \quad v_4 - u_3 = 21 .
  \end{cases}
\]
Уравнений $m + n - 1 = 6$, неизвестных $m + n = 7$, поэтому одну из них задают
произвольно. Положим $v_1 = 0$, тогда
\[
  u_1 = -14, \quad u_2 = -10, \quad u_3 = -20,
  \qquad v_2 = 7, \quad v_3 = 5, \quad v_4 = 1 .
\]

\paragraph{2) Проверка по пустым (зачёркнутым) клеткам.}
\[
  v_j - u_i \le c_{ij} \quad \forall i, j
  \qquad\Longrightarrow\qquad \text{план оптимальный} .
\]

\paragraph{3) Построение цикла пересчёта.}
Выбирая клетку, у которой $\abs{v_j - u_i - c_{ij}}$ максимальна, быстрее придём
к оптимальному решению, но в общем случае берём любую пустую клетку и строим
ломаную (\textbf{цикл пересчёта}) по принципу:
\begin{itemize}[nosep,left=0pt]
  \item направление только параллельно рёбрам клеток;
  \item изменение направления происходит только в заполненных клетках;
  \item ломаная всегда возвращается в клетку старта.
\end{itemize}
Такую ломаную можно построить всегда и только одним способом.

Вершины ломаной нумеруем, чередуя знаки $\oplus$ и $\ominus$, начиная с $\oplus$
в стартовой клетке. Выбираем все клетки со знаком $\ominus$ и объёмы перевозок,
им соответствующие:
\[
  \theta = \min_{\ominus} x_{ij} .
\]

\inputfig{figures/fig-cycle}

В примере стартовой выбрана клетка $(1,4)$. Цикл:
\[
  (1,4)^{\oplus} \to (1,1)^{\ominus} \to (2,1)^{\oplus} \to (2,3)^{\ominus}
  \to (3,3)^{\oplus} \to (3,4)^{\ominus} \to (1,4) ,
\]
\[
  \theta = \min\{27,\ 1,\ 17\} = 1 .
\]

\paragraph{4) Пересчёт.}
К клеткам со знаком $\oplus$ прибавляем $\theta$, из клеток со знаком $\ominus$
вычитаем:
\[
  \begin{array}{lll}
    (1,4)^{\oplus}: & 0 + \theta = 1, &\qquad
    (2,3)^{\ominus}: \ 1 - \theta = 0, \\
    (1,1)^{\ominus}: & 27 - \theta = 26, &\qquad
    (3,3)^{\oplus}: \ 26 + \theta = 27, \\
    (2,1)^{\oplus}: & 6 + \theta = 7, &\qquad
    (3,4)^{\ominus}: \ 17 - \theta = 16 .
  \end{array}
\]
Клетка $(2,3)$ обнулилась и становится пустой, вместо неё заполнена $(1,4)$ —
число заполненных клеток снова $6$. Так получаем новый опорный план.

\paragraph{5) Переход на пункт 1).}

\begin{idea}[title={Что здесь на самом деле произошло}]
Проверим пункт~2 для всех пустых клеток, $u = (-14, -10, -20)$,
$v = (0, 7, 5, 1)$. Величина $\Delta_{ij} = v_j - u_i - c_{ij}$:
\begin{center}
\begin{tabular}{lrrrrrr}
\toprule
клетка & $(1,2)$ & $(1,3)$ & $(1,4)$ & $(2,4)$ & $(3,1)$ & $(3,2)$ \\
\midrule
$\Delta_{ij}$ & $-7$ & $-2$ & $-13$ & $-13$ & $+6$ & $-3$ \\
\bottomrule
\end{tabular}
\end{center}
Условие $\Delta_{ij} \le 0$ нарушено только в клетке $(3,1)$, значит план не
оптимален, и улучшает его именно она. Выбранная на лекции клетка $(1,4)$ имеет
$\Delta_{14} = -13 < 0$, поэтому пересчёт по ней стоимость не уменьшает, а
\emph{увеличивает} на $\abs{\Delta_{14}}\cdot\theta = 13 \cdot 1 = 13$:
было $1681$, стало $1694$. Цикл в лекции построен, чтобы показать механику
пересчёта, а не чтобы улучшить план.

Правильный шаг — стартовая клетка $(3,1)$. Её цикл:
\[
  (3,1)^{\oplus} \to (3,3)^{\ominus} \to (2,3)^{\oplus} \to (2,1)^{\ominus}
  \to (3,1) , \qquad \theta = \min\{26,\ 6\} = 6 ,
\]
после пересчёта $x_{31} = 6$, $x_{33} = 20$, $x_{23} = 7$, $x_{21} = 0$, и
стоимость падает на $\Delta_{31}\cdot\theta = 6 \cdot 6 = 36$: с $1681$ до $1645$.
Общее правило: пересчёт по клетке с $\Delta_{ij}$ меняет стоимость на
$-\Delta_{ij}\,\theta$, поэтому выбирать надо клетки с $\Delta_{ij} > 0$, а среди
них — с наибольшим $\Delta_{ij}$.
\end{idea}

\subsection{Явный вид ограничений и условие оптимальности}
\label{subsec:transport-duality-anchor}

Запишем в явном виде ограничения закрытой задачи:
\[
  \underbrace{
  \begin{cases}
    x_{11} + x_{12} + \dots + x_{1n} = a_1, \\
    x_{21} + x_{22} + \dots + x_{2n} = a_2, \\
    \quad \dots \\
    x_{m1} + x_{m2} + \dots + x_{mn} = a_m
  \end{cases}}_{m \ \text{«строчных» равенств}}
  \qquad
  \underbrace{
  \begin{cases}
    x_{11} + x_{21} + \dots + x_{m1} = b_1, \\
    x_{12} + x_{22} + \dots + x_{m2} = b_2, \\
    \quad \dots \\
    x_{1n} + x_{2n} + \dots + x_{mn} = b_n
  \end{cases}}_{n \ \text{«столбцовых» равенств}}
\]

Если переменные $x_{ij}$ вытянуть в один вектор $x[N]$ по строкам
($N$ — $mn$ индексов), то матрица ограничений имеет блочный вид
\[
  A[M,N] =
  \begin{pmatrix}
    \mathbf{1}\T & 0 & \cdots & 0 \\
    0 & \mathbf{1}\T & \cdots & 0 \\
    \vdots & \vdots & \ddots & \vdots \\
    0 & 0 & \cdots & \mathbf{1}\T \\[4pt]
    E & E & \cdots & E
  \end{pmatrix},
  \qquad
  \begin{aligned}
    &\mathbf{1}\T = (1\ 1\ \dots\ 1) \in \R^{n}, \\
    &E \ \text{— единичная } n \times n ,
  \end{aligned}
\]
размера $(m+n) \times mn$: первые $m$ строк — «строчные» равенства (по $n$ единиц
подряд), последние $n$ строк — «столбцовые» (единичные блоки). При этом
\[
  \sum_{i=1}^{m} a_i = \sum_{j=1}^{n} b_j = P .
\]

\medskip
\noindent\textbf{Условие оптимальности:}
\[
  c\T[N] - y\T[M]\, A[M,N] \ge 0 ,
  \qquad
  \big(c\T[N] - y\T[M]\, A[M,N]\big)\, x[N] = 0 .
\]

\begin{idea}[title={Откуда берутся $v_j - u_i \le c_{ij}$ и $v_j - u_i = c_{ij}$}]
Это и есть ответ на вопросительные знаки в тетради. Условие оптимальности выше —
стандартная пара «двойственная допустимость $+$ дополняющая нежёсткость»; надо
лишь расписать её покомпонентно.

\medskip
\textbf{Шаг 1: что такое $y[M]$.} Вектор $y[M]$ — двойственные переменные, по одной
на каждое ограничение. Ограничений $m + n$: первые $m$ отвечают строкам (пунктам
хранения), последние $n$ — столбцам (пунктам назначения). Обозначим их
\[
  y[M] = (\,\underbrace{y_1, \dots, y_m}_{\text{строки}},\
            \underbrace{y_{m+1}, \dots, y_{m+n}}_{\text{столбцы}}\,) .
\]

\medskip
\textbf{Шаг 2: как выглядит столбец $A$.} Переменная $x_{ij}$ входит ровно в два
ограничения: в $i$-е «строчное» и в $j$-е «столбцовое». Значит её столбец в
матрице $A$ имеет две единицы — в позициях $i$ и $m+j$ — и нули в остальных.
Поэтому компонента произведения $y\T A$, отвечающая $x_{ij}$, равна
\[
  \big(y\T A\big)_{ij} = y_i + y_{m+j} .
\]

\medskip
\textbf{Шаг 3: покомпонентная запись.} Первое условие
$c\T - y\T A \ge 0$ означает
\[
  c_{ij} - \big(y_i + y_{m+j}\big) \ge 0 \qquad \forall i, j .
\]
Положим
\[
  u_i = -\,y_i , \qquad v_j = y_{m+j} .
\]
Тогда $y_i + y_{m+j} = v_j - u_i$, и условие превращается ровно в
\[
  \boxed{\,v_j - u_i \le c_{ij} \quad \forall i, j\,}
\]
— то самое неравенство из пункта~2 алгоритма. Так потенциалы $u_i$ и $v_j$
оказываются просто двойственными переменными транспортной задачи, взятыми с
удобными знаками.

\medskip
\textbf{Шаг 4: вторая строка — дополняющая нежёсткость.} Произведение
$\big(c\T - y\T A\big) x$ — это сумма
\[
  \sum_{i,j} \underbrace{\big(c_{ij} - (v_j - u_i)\big)}_{\ge\, 0}\,
             \underbrace{x_{ij}}_{\ge\, 0} = 0 .
\]
Сумма неотрицательных слагаемых равна нулю только если каждое слагаемое нулевое.
Значит для всякой клетки с $x_{ij} > 0$ (то есть \emph{заполненной}) обязано
выполняться
\[
  \boxed{\,v_j - u_i = c_{ij}\,}
\]
— а это и есть система, из которой в пункте~1 алгоритма находят потенциалы.

\medskip
\textbf{Итог.} Две строки условия оптимальности распадаются ровно на два шага
метода потенциалов: равенства по заполненным клеткам дают $u_i$ и $v_j$,
неравенства по пустым проверяют оптимальность. Если хоть в одной пустой клетке
$v_j - u_i > c_{ij}$, двойственная допустимость нарушена, план не оптимален, и
эту клетку вводят в базис циклом пересчёта.

Система на потенциалы недоопределена на одну переменную (уравнений $m+n-1$,
неизвестных $m+n$) именно потому, что строки матрицы $A$ линейно зависимы:
сумма первых $m$ строк равна сумме последних $n$. Поэтому один потенциал
назначают произвольно — например $v_1 = 0$.
\end{idea}

\subsection{Разрешимость транспортной задачи}

\begin{theorem}[о транспортной задаче]\label{th:transport-solvable}
Любая закрытая транспортная задача имеет решение.
\end{theorem}

\begin{proof}
Допустимое множество непусто: положим для всех $i, j$
\[
  x_{ij} := \frac{a_i\, b_j}{P} .
\]
Подстановка этих значений показывает допустимость вектора:
\[
  \sum_{j=1}^{n} x_{ij} = \frac{a_i}{P} \sum_{j=1}^{n} b_j = \frac{a_i P}{P} = a_i ,
  \qquad
  \sum_{i=1}^{m} x_{ij} = \frac{b_j}{P} \sum_{i=1}^{m} a_i = b_j ,
\]
и очевидно $x_{ij} \ge 0$.

Целевая функция ограничена сверху и снизу. Действительно, существуют
\[
  \underline{c} = \min_{i,j} c_{ij} , \qquad \overline{c} = \max_{i,j} c_{ij} ,
\]
а суммарный объём перевозок у любого допустимого плана один и тот же:
$\sum_{i,j} x_{ij} = P$. Отсюда
\[
  \underline{c}\,P \ \le\ \sum_{i}\sum_{j} c_{ij} x_{ij} \ \le\ \overline{c}\,P .
\]
Непустое допустимое множество плюс ограниченная снизу целевая функция дают
существование оптимального решения.
\end{proof}

\begin{idea}
Строго говоря, ограниченности снизу мало — нужна ещё замкнутость допустимого
множества и достижимость минимума. Здесь это выполнено: допустимое множество
закрытой задачи — многогранное и \emph{ограниченное} ($0 \le x_{ij} \le \min\{a_i, b_j\}$),
то есть компакт, а линейная функция на компакте достигает минимума. Иначе это же
следует из общей теории: по теореме~\ref{th:optimal-is-support} оптимум достигается
в одном из конечного числа опорных векторов.
\end{idea}

Для приближения классической постановки транспортной задачи к реальным
рассматриваются усложнения транспортной задачи.

\subsection{Обоснование метода потенциалов}
\label{subsec:potentials-theory}

\begin{idea}[title={Откуда этот материал}]
Теорема об оптимальности плана и оценка улучшения на лекции не формулировались
— разбирался только сам алгоритм (§7). Ниже обе недостающие части
восстановлены; они закрывают экзаменационный вопрос о методе потенциалов.
\end{idea}

\begin{theorem}[об оптимальности плана]\label{thm:potentials-opt}
Пусть для допустимого базисного плана $X = (x_{ij})$ найдены потенциалы
$u_i$ и $v_j$ такие, что
\[
  v_j - u_i = c_{ij} \ \text{для заполненных клеток},
  \qquad
  v_j - u_i \le c_{ij} \ \text{для пустых} .
\]
Тогда план $X$ оптимален.
\end{theorem}

\begin{proof}
Рассмотрим произвольный допустимый план $Y = (y_{ij})$:
\[
  \sum_{j=1}^{n} y_{ij} = a_i , \qquad
  \sum_{i=1}^{m} y_{ij} = b_j , \qquad y_{ij} \ge 0 .
\]
Для всех клеток выполнено $v_j - u_i \le c_{ij}$, то есть
$c_{ij} \ge v_j - u_i$. Умножая на $y_{ij} \ge 0$ и суммируя,
\[
  \sum_{i=1}^{m}\sum_{j=1}^{n} c_{ij} y_{ij}
  \ \ge\ \sum_{i=1}^{m}\sum_{j=1}^{n} (v_j - u_i)\, y_{ij}
  = \sum_{j=1}^{n} v_j \underbrace{\sum_{i=1}^{m} y_{ij}}_{=\ b_j}
    - \sum_{i=1}^{m} u_i \underbrace{\sum_{j=1}^{n} y_{ij}}_{=\ a_i} ,
\]
откуда
\begin{equation}\label{eq:pot-bound}
  F(Y) \ \ge\ \sum_{j=1}^{n} v_j b_j - \sum_{i=1}^{m} u_i a_i .
\end{equation}
Теперь возьмём сам план $X$. Для заполненных клеток $v_j - u_i = c_{ij}$, а для
пустых $x_{ij} = 0$, поэтому в сумме можно заменить $c_{ij}$ на $v_j - u_i$
\emph{во всех} клетках:
\[
  F(X) = \sum_{i=1}^{m}\sum_{j=1}^{n} c_{ij} x_{ij}
       = \sum_{i=1}^{m}\sum_{j=1}^{n} (v_j - u_i)\, x_{ij}
       = \sum_{j=1}^{n} v_j b_j - \sum_{i=1}^{m} u_i a_i .
\]
Сравнивая с~\eqref{eq:pot-bound}, получаем $F(Y) \ge F(X)$ для любого
допустимого $Y$, то есть $X$ оптимален.
\end{proof}

\begin{idea}
Величина $\sum_j v_j b_j - \sum_i u_i a_i$ — значение двойственной задачи
(см.~\S\ref{subsec:transport-duality-anchor}); неравенство~\eqref{eq:pot-bound}
есть слабая двойственность, а совпадение на плане $X$ — сильная. Поэтому
проверка $\Delta_{ij} = v_j - u_i - c_{ij} \le 0$ и есть проверка двойственной
допустимости.
\end{idea}

\subsubsection*{Насколько улучшается план при пересчёте}

\begin{theorem}\label{thm:potentials-improve}
Пусть в базис вводится пустая клетка $(p,q)$, для которой
$\Delta_{pq} = v_q - u_p - c_{pq} > 0$, и построен цикл пересчёта с величиной
шага $\theta > 0$. Тогда изменение целевой функции равно
\[
  \Delta F = -\,\theta\, \Delta_{pq} \ <\ 0 .
\]
\end{theorem}

\begin{proof}
Обозначим через $C_+$ и $C_-$ множества клеток цикла со знаками $\oplus$ и
$\ominus$. К клеткам из $C_+$ прибавляется $\theta$, из $C_-$ вычитается,
поэтому
\[
  \Delta F = \theta \Big( \sum_{(i,j) \in C_+} c_{ij}
                        - \sum_{(i,j) \in C_-} c_{ij} \Big) .
\]
Все клетки цикла, кроме стартовой $(p,q)$, заполнены, значит для них
$c_{ij} = v_j - u_i$; для $(p,q)$ имеем
$c_{pq} = v_q - u_p - \Delta_{pq}$. Подставим и заметим, что цикл замкнут: в
нём каждая задействованная строка и каждый задействованный столбец встречаются
ровно дважды — один раз со знаком $\oplus$ и один раз со знаком $\ominus$,
поэтому все $u_i$ и $v_j$ взаимно сокращаются. Остаётся только добавка от
стартовой клетки:
\[
  \sum_{(i,j)\in C_+} c_{ij} - \sum_{(i,j)\in C_-} c_{ij} = -\,\Delta_{pq} ,
  \qquad\text{откуда}\qquad
  \Delta F = -\,\theta\,\Delta_{pq} .
\]
Так как $\theta > 0$ и $\Delta_{pq} > 0$, получаем $\Delta F < 0$.
\end{proof}

\begin{idea}
Эта формула и была проверена численно в примере §7: там выбиралась клетка с
$\Delta_{14} = -13$, и стоимость выросла ровно на
$-\theta\Delta_{14} = 13$; а правильная клетка $(3,1)$ с $\Delta_{31} = +6$ и
$\theta = 6$ дала падение на $36$. Отсюда же правило выбора: брать клетку с
\emph{наибольшим} $\Delta_{ij} > 0$ — при прочих равных это даёт наибольшее
уменьшение стоимости за шаг.
\end{idea}

\subsection{Усложнения транспортной задачи}
\label{subsec:transport-extensions}

Классическая постановка предполагает, что перевозится один однородный продукт,
все маршруты разрешены, а ограничения заданы только запасами и потребностями.
В реальных задачах появляются дополнительные условия.

\paragraph{1. Многопродуктовая модель.}
Перевозится несколько видов продукции. Для продукта $k$ вводятся переменные
$x_{ij}^{(k)}$; ограничения записываются отдельно по каждому виду продукции
либо совместно, с учётом общих ресурсов транспорта:
\[
  \sum_{j} x_{ij}^{(k)} = a_i^{(k)} , \qquad
  \sum_{i} x_{ij}^{(k)} = b_j^{(k)} , \qquad
  \sum_{k} x_{ij}^{(k)} \le d_{ij} ,
\]
где последнее неравенство — общая пропускная способность маршрута.

\paragraph{2. Запрет на поставку.}
Если из пункта $A_i$ нельзя поставлять груз в $B_j$, клетка $(i,j)$
запрещается. Её либо исключают из рассмотрения, либо задают очень большой
тариф $c_{ij} = M$ ($M$ — большое число). Тогда метод минимизации не
использует этот маршрут, пока есть другие допустимые.

\paragraph{3. Обязательная поставка.}
Если по маршруту $(i,j)$ перевозка должна быть, вводится условие
$x_{ij} > 0$; на практике его заменяют ограничением снизу
$x_{ij} \ge K_{ij}$, где $K_{ij} > 0$ — обязательный объём.

\paragraph{4. Поставка не менее заданного количества.}
Ограничение $x_{ij} \ge K_{ij}$ учитывается заранее: обязательную поставку
выделяют и уменьшают соответствующие запас и потребность,
\[
  a_i := a_i - K_{ij} , \qquad b_j := b_j - K_{ij} ,
\]
после чего задача снова становится обычной.

\paragraph{5. Поставка не более заданного количества.}
Ограничение $x_{ij} \le K_{ij}$ означает пропускную способность маршрута.

\paragraph{6. Сменно-суточный характер поставок.}
Если перевозки зависят от времени, вводятся дополнительные индексы:
$x_{ij}^{(t)}$ — перевозка из $A_i$ в $B_j$ в период $t$. Учитываются
ограничения по сменам, суткам, времени доставки и доступности транспорта.

\paragraph{7. Возврат тары и минимизация порожнего пробега.}
Если после доставки нужно вернуть тару или транспорт, добавляются обратные
перевозки $x_{ji}^{\text{back}}$; может ставиться задача минимизации порожнего
пробега, когда транспорт не должен возвращаться пустым.

\begin{idea}
Первые пять усложнений не выводят задачу из класса транспортных: запрет и
верхняя граница — это ограничения на отдельные клетки (их обрабатывают
«запрещённым» тарифом или переходом к транспортной задаче с пропускными
способностями), а нижняя граница снимается сдвигом переменной
$x_{ij}' = x_{ij} - K_{ij}$. Пункты~6 и~7 увеличивают размерность и, как
правило, уже требуют решения общей задачи ЛП симплекс-методом
(\S\ref{subsec:simplex-scheme}), а не табличного метода потенциалов.
\end{idea}


\newpage
\lecture{Лекция 3}{}

\subsection{Методы одномерной минимизации: постановка}

Рассматривается задача
\[
  f(x) \longrightarrow \min, \qquad x \in [a,b] ,
\]
где $[a,b]$ — \textbf{исходный интервал неопределённости}.

\medskip
\noindent\textbf{Задача.} Локализовать точку минимума $x_*$ в как можно более
узком \textbf{конечном интервале неопределённости} $[a_k, b_k]$.

\subsubsection*{Откуда берётся одномерная задача}

В многомерной минимизации $\varphi(x) \to \min$, $x \in \R^n$, работает
двухшаговая схема: из начального приближения $x_0$ строится
последовательность
\[
  x_{k+1} = x_k + \alpha_k p_k , \qquad \varphi(x_{k+1}) < \varphi(x_k) ,
\]
причём на каждой итерации
\begin{enumerate}[label=\arabic*), leftmargin=2em]
  \item находим направление $p_k$;
  \item находим шаг $\alpha_k$.
\end{enumerate}

\inputfig{figures/fig-descent-step}

Поиск $\alpha_k$ — подзадача одномерной минимизации: минимизируется функция
одного переменного $\alpha \mapsto \varphi(x_k + \alpha p_k)$. Особый интерес
представляют методы, \textbf{не требующие подсчёта производной}: именно их мы и
разбираем.

\subsubsection*{Унимодальность}

\begin{definition}
Функция $f(x)$ \textbf{унимодальна} на $[a,b]$, если существует единственная
точка глобального минимума $x_*$ на $[a,b]$, причём
\[
  f \ \text{строго убывает на} \ [a, x_*) ,
  \qquad
  f \ \text{строго возрастает на} \ (x_*, b] .
\]
\end{definition}

\begin{idea}
Унимодальность — это ровно то свойство, которое позволяет по двум вычислениям
функции выбросить часть отрезка. Непрерывность или гладкость здесь не нужны:
всё держится на монотонности по обе стороны от минимума.
\end{idea}

\begin{lemma}\label{lem:unimodal}
Пусть $f$ унимодальна на $[a,b]$ и $x_1, x_2 \in [a,b]$, $x_1 < x_2$. Тогда
\[
  1^\circ\quad f(x_1) \ge f(x_2) \ \Longrightarrow\ x_* \in [x_1, b] ,
\]
\[
  2^\circ\quad f(x_1) < f(x_2) \ \Longrightarrow\ x_* \in [a, x_2] .
\]
\end{lemma}

\inputfig{figures/fig-unimodal}

\begin{proof}
$1^\circ$. Пусть, напротив, $x_* < x_1$. Тогда обе точки $x_1 < x_2$ лежат в
промежутке $(x_*, b]$, где $f$ строго возрастает, откуда $f(x_1) < f(x_2)$ —
противоречие с условием $f(x_1) \ge f(x_2)$. Значит $x_* \ge x_1$.

\medskip
$2^\circ$. Пусть, напротив, $x_* > x_2$. Тогда обе точки лежат в промежутке
$[a, x_*)$, где $f$ строго убывает, откуда $f(x_1) > f(x_2)$ — противоречие с
условием $f(x_1) < f(x_2)$. Значит $x_* \le x_2$.
\end{proof}

\subsubsection*{Как сравнивать методы}

\begin{idea}[title={Критерий сравнения}]
Методы одномерной минимизации сравниваются по \textbf{числу обращений к
вычислению функции}, требуемому для достижения заданной точности.

\medskip
Логика такая: сама $f$ обычно считается дорого (в многомерной схеме каждое
обращение — это вычисление $\varphi$ в новой точке $\R^n$), а вся арифметика
метода по сравнению с этим бесплатна. Поэтому метод тем лучше, чем сильнее он
сокращает интервал неопределённости \emph{на одно вычисление функции}.
\end{idea}

\subsection{Методы, не требующие производной}
\label{subsec:direct-search-1d}

Все методы ниже устроены по одной схеме: вычислить $f$ в нескольких пробных
точках, применить лемму~\ref{lem:unimodal} и выбросить часть отрезка. Различие
только в том, \emph{где} ставить пробные точки, — и именно оно определяет цену
метода в числе обращений к $f$.

\begin{idea}[title={О форме записи алгоритмов}]
В тетради алгоритмы записаны схематично: для большинства методов выписано
только правило сокращения отрезка, без явных формул инициализации и критерия
остановки. Ниже каждый метод оформлен как полная пошаговая процедура; места,
где запись в тетради расходится с логикой леммы, отмечены отдельно.
\end{idea}

\subsubsection*{Метод равномерного поиска}

Отрезок разбивается равномерной сеткой, число разбиений $n$ задано заранее.

\begin{algo}[title={Метод равномерного поиска}]
\textbf{Вход:} $[a,b]$, число разбиений $n$, требуемая точность $\delta$.

\begin{steps}
  \item Положить $h = \dfrac{b-a}{n}$ и $x_i = a + h\,i$, $i = \overline{0,n}$.
  \item Вычислить $f(x_i)$ для всех $i$ и найти номер $j$, на котором
        достигается минимум по сетке:
        \[
          f(x_j) = \min_{i} f(x_i) .
        \]
  \item Тогда
        \[
          \begin{aligned}
            f(x_{j-1}) &> f(x_j) &&\Longrightarrow&& x_* \in [x_{j-1},\, b] , \\
            f(x_{j+1}) &\ge f(x_j) &&\Longrightarrow&& x_* \in [a,\, x_{j+1}] ,
          \end{aligned}
        \]
        и, пересекая, получаем
        \[
          x_* \in [x_{j-1},\, x_{j+1}] .
        \]
  \item Если $2h \le \delta$ — остановиться. Иначе положить
        $[a,b] := [x_{j-1},\, x_{j+1}]$ и вернуться к шагу~1.
\end{steps}
\end{algo}

\inputfig{figures/fig-uniform-grid}

\begin{remark}
Если минимум по сетке достигается на краю ($j = 0$ или $j = n$), надо считать
$x_{-1} := a$ и $x_{n+1} := b$. Для унимодальной $f$ такой номер $j$
единственный.
\end{remark}

\begin{idea}[title={Какое $n$ брать — оценка (на лекции осталось вопросом)}]
За один проход интервал сокращается в $n/2$ раз, а стоит это примерно $n-1$
новых вычислений (значения на концах уже известны с прошлого прохода). Значит
«коэффициент сокращения на одно вычисление» равен
\[
  q(n) = \left(\frac{2}{n}\right)^{1/(n-1)} ,
\]
и его надо минимизировать:
\[
  q(3) \approx 0{,}816, \quad
  \mathbf{q(4) \approx 0{,}794}, \quad
  q(5) \approx 0{,}795, \quad
  q(6) \approx 0{,}803 .
\]
Оптимум — $n = 4$. Для сравнения: золотое сечение даёт $0{,}618$ на одно
вычисление, то есть равномерный поиск проигрывает ему при любом $n$. Выкладка
моя, на лекции её не было.
\end{idea}

\subsubsection*{Метод локализации экстремума}

Отличается от равномерного поиска тем, что не нужно вычислять все значения на
сетке: идём слева направо и останавливаемся на первом «провале».

\begin{algo}[title={Метод локализации экстремума}]
\textbf{Вход:} $[a,b]$, число разбиений $n$, точность $\delta$.

\begin{steps}
  \item Положить $h = \dfrac{b-a}{n}$.
  \item Последовательно для $i = 1, 2, \dots, n-1$ проверять пару условий
        \[
          \begin{cases}
            f\big(a + (i-1)h\big) > f\big(a + ih\big), \\[2pt]
            f\big(a + ih\big) \le f\big(a + (i+1)h\big),
          \end{cases}
        \]
        и остановиться на первом номере $i$, для которого они выполнены.
  \item Тогда
        \[
          x_* \in \big[\, a + (i-1)h, \ a + (i+1)h \,\big] .
        \]
  \item Если $2h \le \delta$ — остановиться. Иначе повторить процесс для
        отрезка $\big[a + (i-1)h,\ a + (i+1)h\big]$.
\end{steps}
\end{algo}

\begin{idea}
Для унимодальной функции найденный здесь номер $i$ — тот же самый, что и номер
$j$ в равномерном поиске: первый «провал» и есть минимум по сетке. Выигрыш в
том, что в среднем значений вычисляется меньше, чем $n+1$: до минимума функция
убывает, и обход прекращается сразу после него.
\end{idea}

\subsubsection*{Метод золотого сечения}

Здесь пробные точки ставятся симметрично и так, чтобы \textbf{одна из них
переиспользовалась} на следующем шаге.

На $k$-м шаге на отрезке $[a_k, b_k]$ берутся
\[
  \lambda_k = a_k + \alpha\,(b_k - a_k), \qquad
  \mu_k = a_k + b_k - \lambda_k = b_k - \alpha\,(b_k - a_k),
  \qquad \alpha \in \left(0,\tfrac12\right) ,
\]
так что выполняется симметрия
\[
  b_k - \lambda_k = \mu_k - a_k .
\]

\inputfig{figures/fig-golden-section}

Правило сокращения (лемма~\ref{lem:unimodal}):
\[
  f(\lambda_k) \ge f(\mu_k) \ \Longrightarrow\ [\lambda_k,\, b_k] ,
  \qquad\qquad (1)
\]
\[
  f(\lambda_k) < f(\mu_k) \ \Longrightarrow\ [a_k,\, \mu_k] .
  \qquad\qquad (2)
\]

\paragraph{Подбор $\alpha$.} Необходимо сократить число обращений к $f$, то
есть добиться
\[
  \lambda_{k+1} = \mu_k \quad\text{в случае (1)},
  \qquad
  \mu_{k+1} = \lambda_k \quad\text{в случае (2)} .
\]
Обозначим $L_k = b_k - a_k$. В случае~(1) имеем $a_{k+1} = \lambda_k$,
$b_{k+1} = b_k$, значит $L_{k+1} = (1-\alpha)L_k$ и
\[
  \lambda_{k+1} = a_{k+1} + \alpha L_{k+1}
  = \lambda_k + \alpha(1-\alpha)L_k
  = a_k + \alpha L_k + \alpha(1-\alpha)L_k .
\]
Приравнивая это к $\mu_k = a_k + (1-\alpha)L_k$, получаем
\[
  \alpha + \alpha(1-\alpha) = 1 - \alpha
  \quad\Longleftrightarrow\quad
  2\alpha - \alpha^2 = 1 - \alpha
  \quad\Longleftrightarrow\quad
  \alpha^2 - 3\alpha + 1 = 0 ,
\]
откуда
\[
  \alpha_{1,2} = \frac{3 \pm \sqrt5}{2},
  \qquad
  \boxed{\ \alpha = \frac{3 - \sqrt5}{2} \approx 0{,}382\ }
\]
(корень со знаком «$+$» больше единицы и не подходит).

\begin{idea}
Случай~(2) даёт то же самое уравнение: там $\mu_{k+1} = a_k + (1-\alpha)^2 L_k$
и условие $\mu_{k+1} = \lambda_k$ означает $(1-\alpha)^2 = \alpha$, то есть
снова $\alpha^2 - 3\alpha + 1 = 0$. Так что одним подбором $\alpha$
обслуживаются оба исхода.

\medskip
Коэффициент сокращения равен
\[
  1 - \alpha = \frac{\sqrt5 - 1}{2} \approx 0{,}618 ,
\]
и это ровно золотое сечение — отсюда название. На каждой итерации, кроме
первой, тратится \textbf{одно} новое вычисление функции.
\end{idea}

\begin{algo}[title={Метод золотого сечения}]
\textbf{Вход:} $[a_1,b_1] = [a,b]$, точность $\delta$,
$\alpha = \tfrac{3-\sqrt5}{2}$.

\begin{steps}
  \item Положить
        $\lambda_1 = a_1 + \alpha(b_1-a_1)$,
        $\mu_1 = b_1 - \alpha(b_1-a_1)$;
        вычислить $f(\lambda_1)$ и $f(\mu_1)$.
  \item Если $b_k - a_k \le \delta$ — остановиться, принять
        $x_* \approx \tfrac{a_k+b_k}{2}$.
  \item Если $f(\lambda_k) \ge f(\mu_k)$, положить
        \[
          a_{k+1} = \lambda_k, \quad b_{k+1} = b_k, \quad
          \lambda_{k+1} = \mu_k \ \text{(значение известно)},
        \]
        и вычислить одно новое значение
        $f(\mu_{k+1})$, где $\mu_{k+1} = b_{k+1} - \alpha(b_{k+1}-a_{k+1})$.
  \item Иначе положить
        \[
          a_{k+1} = a_k, \quad b_{k+1} = \mu_k, \quad
          \mu_{k+1} = \lambda_k \ \text{(значение известно)},
        \]
        и вычислить одно новое значение
        $f(\lambda_{k+1})$, где $\lambda_{k+1} = a_{k+1} + \alpha(b_{k+1}-a_{k+1})$.
  \item Увеличить $k$ и вернуться к шагу~2.
\end{steps}
\end{algo}

\subsubsection*{Метод дихотомии}

Пробные точки $x_1 < x_2$ берутся вплотную к середине отрезка. Правило
сокращения то же:
\[
  f(x_1) \ge f(x_2) \ \Longrightarrow\ [x_1,\, b] ,
  \qquad
  f(x_1) < f(x_2) \ \Longrightarrow\ [a,\, x_2] .
\]

\paragraph{Почему именно середина.} Длина оставшегося отрезка равна либо
$b - x_1$, либо $x_2 - a$; заранее неизвестно, какой исход реализуется,
поэтому точки выбираются из условия наилучшего худшего случая:
\[
  \min_{x_1,\, x_2} \ \max\big\{\, b - x_1,\ x_2 - a \,\big\} .
\]
Первая величина убывает по $x_1$, вторая растёт по $x_2$, поэтому минимум
максимума достигается при
\[
  x_1 = x_2 = \frac{a+b}{2} .
\]

\inputfig{figures/fig-dichotomy-minmax}

Совпадающие точки бесполезны, поэтому их разносят на малую величину
$\eps > 0$:
\[
  x_1 := \frac{a+b}{2} - \eps , \qquad
  x_2 := \frac{a+b}{2} + \eps .
\]

\begin{algo}[title={Метод дихотомии}]
\textbf{Вход:} $[a_1,b_1] = [a,b]$, точность $\delta$, параметр разведения
$\eps > 0$.

\begin{steps}
  \item Если $b_k - a_k \le \delta$ — остановиться.
  \item Положить
        $x_1 = \dfrac{a_k+b_k}{2} - \eps$, \ \
        $x_2 = \dfrac{a_k+b_k}{2} + \eps$;
        вычислить $f(x_1)$ и $f(x_2)$.
  \item Если $f(x_1) \ge f(x_2)$, положить $[a_{k+1},b_{k+1}] = [x_1,\, b_k]$;
        иначе $[a_{k+1},b_{k+1}] = [a_k,\, x_2]$.
  \item Увеличить $k$ и вернуться к шагу~1.
\end{steps}
\end{algo}

\begin{remark}[выбор $\eps$]
$\eps$ берётся настолько малым, насколько позволяет различимость значений
$f\!\left(\tfrac{a+b}{2} - \eps\right)$ и
$f\!\left(\tfrac{a+b}{2} + \eps\right)$: если они сольются в пределах точности
вычислений, сравнение потеряет смысл. Практический ориентир из тетради — доли
процента от рассматриваемых величин.
\end{remark}

\begin{idea}
Длина отрезка за итерацию становится равной $\tfrac{L_k}{2} + \eps$, то есть
почти вдвое меньше, но ценой \textbf{двух} вычислений функции. На одно
вычисление это даёт $\sqrt{1/2} \approx 0{,}707$ — хуже, чем $0{,}618$ у
золотого сечения.
\end{idea}

\subsubsection*{Метод Фибоначчи}

Та же идея симметричных точек с переиспользованием, но коэффициенты берутся не
постоянные, а из последовательности Фибоначчи
\[
  F_0 = F_1 = 1, \qquad F_{i+1} = F_i + F_{i-1} ,
\]
причём общее число шагов $n$ фиксируется \emph{заранее}.

На $k$-м шаге на отрезке $[a_k, b_k]$ берутся
\[
  \lambda_k := a_k + \frac{F_{n-k-1}}{F_{n-k+1}}\,(b_k - a_k) ,
  \qquad
  \mu_k := a_k + \frac{F_{n-k}}{F_{n-k+1}}\,(b_k - a_k) ,
\]
и $\lambda_k < \mu_k$, поскольку $F_{n-k-1} < F_{n-k}$. Правило сокращения
прежнее:
\[
  f(\lambda_k) \ge f(\mu_k) \ \Longrightarrow\ [\lambda_k,\, b_k] ,
  \qquad
  f(\lambda_k) < f(\mu_k) \ \Longrightarrow\ [a_k,\, \mu_k] .
\]

\begin{idea}[title={Проверка переиспользования — ответ на «???» в тетради}]
В тетради записано «$\mu_{k+1} = \lambda_k$ \textbf{???}». Проверяется прямой
подстановкой. Пусть $L_k = b_k - a_k$.

\medskip
\textbf{Случай (1):} $a_{k+1} = \lambda_k$, $b_{k+1} = b_k$, и, пользуясь
$F_{n-k+1} = F_{n-k} + F_{n-k-1}$,
\[
  L_{k+1} = L_k\left(1 - \frac{F_{n-k-1}}{F_{n-k+1}}\right)
          = L_k\,\frac{F_{n-k}}{F_{n-k+1}} .
\]
Тогда
\[
  \lambda_{k+1} = \lambda_k + \frac{F_{n-k-2}}{F_{n-k}}\,L_{k+1}
                = \lambda_k + L_k\,\frac{F_{n-k-2}}{F_{n-k+1}}
                = \mu_k ,
\]
поскольку $\mu_k - \lambda_k = L_k\,\dfrac{F_{n-k} - F_{n-k-1}}{F_{n-k+1}}
= L_k\,\dfrac{F_{n-k-2}}{F_{n-k+1}}$.

\medskip
\textbf{Случай (2):} $b_{k+1} = \mu_k$, $L_{k+1} = L_k\,\dfrac{F_{n-k}}{F_{n-k+1}}$,
и точно так же $\mu_{k+1} = \lambda_k$. То есть переиспользуется ровно одна
точка, как и в золотом сечении.
\end{idea}

\begin{algo}[title={Метод Фибоначчи}]
\textbf{Вход:} $[a_1,b_1] = [a,b]$, число шагов $n$, числа $F_0,\dots,F_n$.

\begin{steps}
  \item Положить $k = 1$, вычислить $f(\lambda_1)$ и $f(\mu_1)$ по формулам
        выше.
  \item Если $f(\lambda_k) \ge f(\mu_k)$, положить
        $[a_{k+1},b_{k+1}] = [\lambda_k,\, b_k]$ и $\lambda_{k+1} = \mu_k$
        (значение известно); вычислить $f(\mu_{k+1})$.
  \item Иначе положить $[a_{k+1},b_{k+1}] = [a_k,\, \mu_k]$ и
        $\mu_{k+1} = \lambda_k$ (значение известно); вычислить
        $f(\lambda_{k+1})$.
  \item Повторять, пока $k < n-1$.
\end{steps}
\end{algo}

\begin{remark}
На последнем шаге $k = n-1$ обе точки сливаются:
\[
  \lambda_{n-1} = \mu_{n-1} = \frac{a_{n-1} + b_{n-1}}{2} ,
\]
и сравнивать становится нечего. Поэтому одну из точек сдвигают на малое $\eps$
— ровно как в методе дихотомии. В тетради этот момент записан как «далее
возврат на предыдущий шаг».
\end{remark}

\begin{idea}
Итоговая длина интервала неопределённости равна
\[
  \frac{b-a}{F_n} ,
\]
и это \textbf{наилучший возможный} результат при заданном числе вычислений
функции — в этом смысле метод Фибоначчи оптимален. Золотое сечение — его
предельный вариант: отношение $F_{n-1}/F_n$ стремится к $0{,}618$, поэтому при
больших $n$ оба метода работают практически одинаково, но золотому сечению не
нужно знать $n$ заранее.
\end{idea}

\subsubsection*{Метод бисекции}

На отрезке $[a_k, b_k]$ берутся три точки, делящие его на четыре равные части:
\[
  x_1 = a_k + \frac{L_k}{4}, \qquad
  x_2 = \frac{a_k + b_k}{2}, \qquad
  x_3 = a_k + \frac{3L_k}{4},
  \qquad L_k = b_k - a_k .
\]

\begin{algo}[title={Метод бисекции}]
\textbf{Вход:} $[a_1,b_1] = [a,b]$, точность $\delta$.

\begin{steps}
  \item Если $b_k - a_k \le \delta$ — остановиться.
  \item Вычислить $f(x_1)$, $f(x_2)$, $f(x_3)$ в точках выше (значение в
        середине $x_2$ обычно уже известно с прошлой итерации).
  \item Если $f(x_1) < f(x_2)$, то $x_* \in [a_k,\, x_2]$ — взять новым
        отрезком $[a_k, x_2]$.
  \item Иначе $x_* \in [x_1,\, b_k]$, и сравнить $f(x_2)$ с $f(x_3)$:
        \begin{itemize}[leftmargin=1.4em, topsep=2pt]
          \item если $f(x_2) > f(x_3)$, то $x_* \in [x_2,\, b_k]$;
          \item если $f(x_2) \le f(x_3)$, то $x_* \in [x_1,\, x_3]$.
        \end{itemize}
  \item Увеличить $k$ и вернуться к шагу~1.
\end{steps}
\end{algo}

\begin{idea}
В каждом из трёх исходов новый отрезок вдвое короче прежнего, а его серединой
оказывается одна из уже вычисленных точек ($x_1$, $x_3$ или $x_2$
соответственно). Поэтому на итерацию приходится два новых вычисления функции —
столько же, сколько в дихотомии, и с тем же коэффициентом $0{,}707$ на
вычисление.
\end{idea}

\subsubsection*{Сводка}

\begin{center}
\begin{tabular}{lcc}
\toprule
Метод & сокращение за итерацию & на одно вычисление $f$ \\
\midrule
Равномерный поиск ($n=4$) & $1/2$ за $3$ вычисления & $\approx 0{,}794$ \\
Дихотомия                 & $\approx 1/2$ за $2$    & $\approx 0{,}707$ \\
Бисекция                  & $1/2$ за $2$            & $\approx 0{,}707$ \\
Золотое сечение           & $0{,}618$ за $1$        & $\approx 0{,}618$ \\
Фибоначчи                 & $(b-a)/F_n$ за $n$      & оптимум \\
\bottomrule
\end{tabular}
\end{center}

\begin{idea}[title={О сводной таблице}]
Таблица собрана мной по разобранным оценкам — на лекции её не было. Числа в
последнем столбце — коэффициент, во сколько раз сокращается интервал
неопределённости в пересчёте на одно обращение к $f$: чем меньше, тем лучше.
\end{idea}

\subsubsection*{Не разобрано}

\begin{idea}[title={Не разобрано на лекции}]
В конце тетради перечислены ещё два метода, оба со знаками вопроса, то есть на
лекции они не разбирались:
\begin{itemize}[leftmargin=1.6em, topsep=2pt]
  \item \textbf{метод пробных точек};
  \item \textbf{метод аппроксимации параболой}.
\end{itemize}
\end{idea}

\begin{idea}
Для ориентира: метод аппроксимации параболой — это единственный из списка,
использующий не только сравнения, но и сами значения. По трём точкам строится
квадратичная интерполяция, и за новое приближение берётся вершина параболы.
На гладкой функции вблизи минимума он сходится заметно быстрее любого метода
деления отрезка, но гарантий, как у перечисленных выше, не даёт.
\end{idea}

\subsubsection*{Численный пример: все шесть методов на одной функции}

\begin{idea}[title={Откуда этот материал}]
На лекции ни один из методов не был прогнан на числах. Ниже все шесть
применены к одной и той же задаче — так видно, чем они отличаются по числу
обращений к $f$ и по скорости сокращения интервала.
\end{idea}

\begin{example}\label{ex:one-dim-run}
\[
  f(x) = (x-2)^2 \ \longrightarrow\ \min , \qquad x \in [0, 5] .
\]
Функция унимодальна на $[0,5]$: она строго убывает на $[0,2]$ и строго
возрастает на $[2,5]$, так что $x_* = 2$ — единственная точка минимума.
Производная нигде не используется, только сравнения значений.

\medskip
\noindent\textbf{1) Равномерный поиск.} Возьмём $n = 5$, тогда
$h = \frac{5-0}{5} = 1$ и узлы $x_i = 0,1,2,3,4,5$:
\[
  f = 4,\ 1,\ 0,\ 1,\ 4,\ 9 .
\]
Минимум по сетке на $j = 2$ ($x_2 = 2$), новый интервал
\[
  [x_1,\, x_3] = [1,\, 3] , \qquad \text{длина } 2 ,
\]
ценой $6$ вычислений функции.

\medskip
\noindent\textbf{2) Локализация экстремума.} Тот же шаг $h = 1$. Идём слева:
\[
  i = 1: \quad f(0) = 4 > f(1) = 1 , \quad\text{но}\quad f(1) = 1 \not\le f(2) = 0 ;
\]
\[
  i = 2: \quad f(1) = 1 > f(2) = 0 \quad\text{и}\quad f(2) = 0 \le f(3) = 1 .
\]
Условия выполнены, обход прекращается: $x_* \in [1, 3]$. Результат тот же, что
в равномерном поиске, но понадобилось $4$ вычисления вместо $6$.

\medskip
\noindent\textbf{3) Дихотомия} ($\eps = 0{,}1$).

\emph{Итерация 1.} $[a_1,b_1] = [0,5]$, середина $2{,}5$:
\[
  x_1 = 2{,}4 , \quad x_2 = 2{,}6 ,
  \qquad f(x_1) = 0{,}16 , \quad f(x_2) = 0{,}36 .
\]
$f(x_1) < f(x_2)$, значит $[a_2,b_2] = [a_1,\, x_2] = [0;\ 2{,}6]$.

\emph{Итерация 2.} Середина $1{,}3$:
\[
  x_1 = 1{,}2 , \quad x_2 = 1{,}4 ,
  \qquad f(x_1) = 0{,}64 , \quad f(x_2) = 0{,}36 .
\]
Теперь $f(x_1) \ge f(x_2)$, значит $[a_3,b_3] = [x_1,\, b_2] = [1{,}2;\ 2{,}6]$.

Длины: $5 \to 2{,}6 \to 1{,}4$ за $4$ вычисления.

\medskip
\noindent\textbf{4) Бисекция.}

\emph{Итерация 1.} $[0,5]$, $L = 5$:
\[
  x_1 = 1{,}25 , \quad x_2 = 2{,}5 , \quad x_3 = 3{,}75 ,
\]
\[
  f(x_1) = 0{,}5625 , \quad f(x_2) = 0{,}25 , \quad f(x_3) = 3{,}0625 .
\]
Здесь $f(x_1) > f(x_2)$, и далее $f(x_2) \le f(x_3)$, значит берётся средняя
часть: $[a_2,b_2] = [x_1,\, x_3] = [1{,}25;\ 3{,}75]$.

\emph{Итерация 2.} $L = 2{,}5$:
\[
  x_1 = 1{,}875 , \quad x_2 = 2{,}5 , \quad x_3 = 3{,}125 ,
\]
\[
  f(x_1) = 0{,}015625 , \quad f(x_2) = 0{,}25 , \quad f(x_3) = 1{,}265625 .
\]
Теперь $f(x_1) < f(x_2)$, значит $[a_3,b_3] = [a_2,\, x_2] = [1{,}25;\ 2{,}5]$.

Длины: $5 \to 2{,}5 \to 1{,}25$ — ровно вдвое за итерацию, как и обещано.

\medskip
\noindent\textbf{5) Золотое сечение}, $\alpha = \frac{3-\sqrt5}{2} \approx 0{,}381966$.

\emph{Итерация 1.} $[a_1,b_1] = [0,5]$:
\[
  \lambda_1 = 0 + \alpha\cdot 5 \approx 1{,}90983 ,
  \qquad
  \mu_1 = 5 - \alpha\cdot 5 \approx 3{,}09017 ,
\]
\[
  f(\lambda_1) \approx 0{,}00813 , \qquad f(\mu_1) \approx 1{,}18847 .
\]
$f(\lambda_1) < f(\mu_1)$, значит $[a_2,b_2] = [a_1,\, \mu_1] = [0;\ 3{,}09017]$,
причём $\mu_2 = \lambda_1 \approx 1{,}90983$ — значение уже известно.

\emph{Итерация 2.} Вычисляется \emph{одна} новая точка:
\[
  \lambda_2 = a_2 + \alpha(b_2 - a_2) \approx 1{,}18034 ,
  \qquad f(\lambda_2) \approx 0{,}67079 .
\]
Теперь $f(\lambda_2) \ge f(\mu_2)$, значит
$[a_3,b_3] = [\lambda_2,\, b_2] = [1{,}18034;\ 3{,}09017]$.

Длины: $5 \to 3{,}09017 \to 1{,}90983$ за $3$ вычисления функции — по одному на
итерацию после первой.

\medskip
\noindent\textbf{6) Метод Фибоначчи}, $n = 5$:
\[
  F_0 = 1,\ F_1 = 1,\ F_2 = 2,\ F_3 = 3,\ F_4 = 5,\ F_5 = 8 .
\]

\emph{Итерация 1} ($k=1$): коэффициенты
$\frac{F_{n-k-1}}{F_{n-k+1}} = \frac{F_3}{F_5} = \frac38$ и
$\frac{F_{n-k}}{F_{n-k+1}} = \frac{F_4}{F_5} = \frac58$, поэтому
\[
  \lambda_1 = 0 + \tfrac38\cdot 5 = 1{,}875 ,
  \qquad
  \mu_1 = 0 + \tfrac58\cdot 5 = 3{,}125 ,
\]
\[
  f(\lambda_1) = 0{,}015625 , \qquad f(\mu_1) = 1{,}265625 .
\]
$f(\lambda_1) < f(\mu_1)$, значит $[a_2,b_2] = [0;\ 3{,}125]$ и
$\mu_2 = \lambda_1 = 1{,}875$.

\emph{Итерация 2} ($k=2$): коэффициенты $\frac{F_2}{F_4} = \frac25$ и
$\frac{F_3}{F_4} = \frac35$:
\[
  \lambda_2 = \tfrac25\cdot 3{,}125 = 1{,}25 ,
  \qquad
  \mu_2 = \tfrac35\cdot 3{,}125 = 1{,}875 \quad\text{(совпало, как и должно)} ,
\]
\[
  f(\lambda_2) = 0{,}5625 , \qquad f(\mu_2) = 0{,}015625 .
\]
$f(\lambda_2) \ge f(\mu_2)$, значит $[a_3,b_3] = [1{,}25;\ 3{,}125]$.

Длины: $5 \to 3{,}125 \to 1{,}875$.
\end{example}

\begin{idea}
Совпадение $\mu_2 = \lambda_1 = 1{,}875$ в методе Фибоначчи — это и есть
переиспользование точки, доказанное выше; поэтому на второй итерации
вычисляется только $f(\lambda_2)$.

Сравнение по длине интервала на одно вычисление функции (все методы стартовали
с $[0,5]$):
\begin{itemize}[nosep,left=0pt]
  \item равномерный поиск: $2$ за $6$ вычислений;
  \item локализация: $2$ за $4$;
  \item дихотомия: $1{,}4$ за $4$;
  \item бисекция: $1{,}25$ за $4$ (с переиспользованием середины — за $2$ новых);
  \item золотое сечение: $1{,}90983$ за $3$;
  \item Фибоначчи: $1{,}875$ за $3$.
\end{itemize}
Фибоначчи обходит золотое сечение именно потому, что число шагов задано
заранее и коэффициенты подобраны под него; при $n \to \infty$ разница исчезает.
\end{idea}


\newpage
\lecture{Лекция 4}{}

\subsection{Многомерная минимизация: общие понятия}
\label{subsec:multidim-basics}

Задача многомерной минимизации:
\[
  \begin{cases}
    f(x) \longrightarrow \min, & x \in \R^n, \\
    x \in S .
  \end{cases}
\]

\subsubsection*{Два принципиальных подхода}

\begin{enumerate}[label=\arabic*., leftmargin=2em]
  \item \textbf{Аналитический}: решать уравнение $\nabla f(x) = 0$.
  \item \textbf{Итерационный}: от начального приближения $x_0$ строить
        последовательность $\{x_k\}$ с
        \[
          f(x_{k+1}) \le f(x_k) ,
        \]
        то есть \textbf{последовательность спуска}.
\end{enumerate}

\begin{idea}
Первый подход «более аналитичен», но уже для умеренно сложной $f$ система
$\nabla f(x) = 0$ решается не проще исходной задачи, да и даёт она лишь
стационарные точки, среди которых могут быть максимумы и сёдла. Поэтому весь
дальнейший курс — про второй подход.
\end{idea}

\subsubsection*{Характеристики алгоритмов}

\begin{definition}
Алгоритм \textbf{глобально сходится}, если последовательность, построенная по
этому алгоритму, сходится к оптимальному решению \textbf{независимо от выбора
начального приближения}.
\end{definition}

\begin{definition}
\textbf{Порядок алгоритма} — максимальный порядок производных, привлекаемых к
его реализации.
\end{definition}

Речь дальше пойдёт о методах $0$-го, $1$-го и $2$-го порядка.

\begin{definition}
\textbf{Скорость сходимости} алгоритма оценивается по величине порядка $p$:
\[
  \frac{\norm{x_{k+1} - x_*}}{\norm{x_k - x_*}^{\,p}} \le q ,
  \qquad q \in (0,1) ,
\]
причём требуется лишь, чтобы существовало $K$, начиная с которого неравенство
выполнено: $\forall k > K$.
\end{definition}

Принятые названия:
\[
  \begin{aligned}
    p &= 1 &&\text{— линейная скорость;} \\
    p &> 1 &&\text{— сверхлинейная скорость;} \\
    p &= 2 &&\text{— квадратичная скорость.}
  \end{aligned}
\]

\begin{idea}
Почему при $p = 1$ требуется $q < 1$, а при $p > 1$ это уже не нужно: при
$p=1$ неравенство означает $\norm{x_{k+1}-x_*} \le q\norm{x_k - x_*}$, то есть
геометрическую прогрессию, и сходимость обеспечивает именно $q<1$. При $p > 1$
ошибка возводится в степень, и достаточно лишь попасть в достаточно малую
окрестность $x_*$ — поэтому сверхлинейные и квадратичные методы сходятся, как
правило, только локально, зато очень быстро.
\end{idea}

\subsection{Градиентные методы первого порядка}

\subsubsection*{Откуда берётся антиградиент}

Пусть строится последовательность $\{x_k\}$ и пробная точка
\[
  x = x_k + \alpha p_k .
\]
Разложение по формуле Тейлора:
\[
  f(x) = f(x_k) + \alpha\, \nabla\T f(x_k)\, p_k
       + \tfrac12 \alpha^2\, p_k\T H(\bar x_k)\, p_k ,
  \qquad
  \bar x_k = x_k + \theta\,(x - x_k), \quad \theta \in [0,1] .
\]

При \emph{малом} $\alpha$ знак приращения определяется линейным членом,
поэтому
\[
  f(x) < f(x_k)
  \qquad\Longleftarrow\qquad
  \nabla\T f(x_k)\, p_k < 0 .
\]

\inputfig{figures/fig-descent-cone}

Для градиентных методов первого порядка берут
\[
  \boxed{\ p_k = -\nabla f(x_k)\ } \qquad \text{— антиградиент},
  \qquad \norm{\nabla f(x_k)} \ne 0 .
\]

\begin{idea}
Проверка условия убывания для этого выбора тривиальна:
$\nabla\T f(x_k)\,p_k = -\norm{\nabla f(x_k)}^2 < 0$, пока градиент не нулевой.
Более того, среди всех направлений фиксированной длины антиградиент даёт
наибольшее по модулю отрицательное значение $\nabla\T f(x_k)\,p_k$ — отсюда и
название «наискорейший спуск» для варианта с точным шагом.
\end{idea}

\subsubsection*{Три метода}

Направление у всех одно и то же, \textbf{разные методы отличаются способом
выбора шага}:

\begin{enumerate}[label=\arabic*), leftmargin=2em]
  \item метод с дроблением шага;
  \item метод наискорейшего спуска;
  \item метод с постоянным шагом.
\end{enumerate}

Первые два — \textbf{шаговые методы}, устроенные по общей схеме:
\[
  1^\circ \ \text{находится направление спуска};
  \qquad
  2^\circ \ \text{находится шаг} .
\]

\begin{exercise}
Разобрать самостоятельно метод с постоянным шагом ($\alpha_k \equiv \alpha$).
\end{exercise}

\begin{idea}
Подсказка к упражнению: оценка~\eqref{eq:grad-main} ниже при $\alpha_k \equiv
\alpha$ сразу даёт
$f(x_{k+1}) - f(x_k) \le -\alpha(1-\alpha R)\norm{\nabla f(x_k)}^2$, так что
метод работает при $0 < \alpha < 1/R$, а наилучшая гарантия получается при
$\alpha = 1/(2R)$. Недостаток очевиден: нужна константа Липшица $R$, которую
на практике не знают.
\end{idea}

\subsubsection*{Градиентный метод с дроблением шага}

Итерация:
\[
  x_0 \ \longrightarrow\ x_{k+1} = x_k - \alpha_k \nabla f(x_k) .
\]
Шаг $\alpha_k$ подбирается \emph{дроблением}: начиная с $\alpha^{(0)}$,
пробное значение умножается на $\lambda \in (0,1)$ до тех пор, пока не будет
выполнено условие достаточного убывания.

\begin{algo}[title={Градиентный метод с дроблением шага}]
\textbf{Вход:} $x_0$, начальный шаг $\alpha^{(0)}$, коэффициент дробления
$\lambda \in (0,1)$, параметр $\eps \in (0,1)$, точность $\delta$.

\begin{steps}
  \item Если $\norm{\nabla f(x_k)} < \delta$ — остановиться.
  \item Положить $\alpha := \alpha^{(0)}$.
  \item Вычислить $x = x_k - \alpha \nabla f(x_k)$ и проверить условие
        дробления
        \[
          f(x) - f(x_k) \ \le\ -\eps\,\alpha\, \norm{\nabla f(x_k)}^2 .
        \]
  \item Если условие не выполнено — положить $\alpha := \alpha\lambda$ и
        вернуться к шагу~3.
  \item Иначе принять
        $\alpha_k = \alpha = \alpha^{(0)}\lambda^{\,i_k}$ (где $i_k$ — число
        дроблений) и $x_{k+1} = x$.
  \item Увеличить $k$ и вернуться к шагу~1.
\end{steps}
\end{algo}

\inputfig{figures/fig-step-halving}

\begin{idea}
Зачем нужен множитель $\eps$: без него условие $f(x) < f(x_k)$ выполнялось бы
при сколь угодно малом $\alpha$, и метод мог бы «застрять», делая всё более
мелкие шаги с всё меньшим выигрышем. Условие с $-\eps\alpha\norm{\nabla f}^2$
требует убывания, \emph{пропорционального} сделанному шагу, — именно это
даст в теореме оценку $\alpha_k$ снизу.
\end{idea}

\subsubsection*{Градиентный метод наискорейшего спуска}

Здесь шаг выбирается точно — решением задачи одномерной минимизации
(\S9--\S10):
\[
  \big(f(x_{k+1}) - f(x_k)\big) \longrightarrow \min_{\alpha}
  \quad\Longleftrightarrow\quad
  f(x_{k+1}) \longrightarrow \min_{\alpha}
\]
\[
  \Longleftrightarrow\quad
  f\big(x_k - \alpha \nabla f(x_k)\big) \longrightarrow \min_{\alpha \ge 0} .
\]

\begin{algo}[title={Градиентный метод наискорейшего спуска}]
\textbf{Вход:} $x_0$, точность $\delta$.

\begin{steps}
  \item Если $\norm{\nabla f(x_k)} < \delta$ — остановиться.
  \item Положить $p_k = -\nabla f(x_k)$.
  \item Решить задачу одномерной минимизации
        \[
          \alpha_k = \arg\min_{\alpha \ge 0} f(x_k + \alpha p_k)
        \]
        любым методом из \S10.
  \item Положить $x_{k+1} = x_k + \alpha_k p_k$, увеличить $k$ и вернуться
        к шагу~1.
\end{steps}
\end{algo}

\inputfig{figures/fig-zigzag}

\begin{idea}
Переход «$f(x_{k+1}) - f(x_k) \to \min$» $\iff$ «$f(x_{k+1}) \to \min$»
законен потому, что $f(x_k)$ от $\alpha$ не зависит — это просто константа,
сдвигающая минимизируемую функцию.
\end{idea}

\subsubsection*{Основная оценка}

Обе теоремы о сходимости опираются на одно неравенство, поэтому выведем его
отдельно.

\begin{lemma}\label{lem:grad-main}
Пусть $f \in C^1(\R^n)$ и градиент липшицев:
\[
  \norm{\nabla f(x) - \nabla f(y)} \le R\,\norm{x-y},
  \qquad R = \mathrm{const} > 0 .
\]
Тогда для любого $\alpha > 0$
\begin{equation}\label{eq:grad-main}
  f\big(x_k - \alpha \nabla f(x_k)\big) - f(x_k)
  \ \le\ -\alpha\,(1 - \alpha R)\, \norm{\nabla f(x_k)}^2 .
\end{equation}
\end{lemma}

\begin{proof}
Положим $x = x_k - \alpha \nabla f(x_k)$, так что
$x - x_k = -\alpha\nabla f(x_k)$. По формуле конечных приращений
\[
  f(x) - f(x_k) = \nabla\T f(x_{\mathrm{cp}})\,(x - x_k),
  \qquad
  x_{\mathrm{cp}} = x_k + \theta\,(x - x_k), \quad \theta \in [0,1] .
\]
Прибавим и вычтем $\nabla\T f(x_k)(x-x_k)$:
\[
  f(x) - f(x_k)
  = \underbrace{\nabla\T f(x_k)\,(x-x_k)}_{=\ -\alpha\norm{\nabla f(x_k)}^2}
  + \big(\nabla f(x_{\mathrm{cp}}) - \nabla f(x_k)\big)\T (x - x_k) .
\]
Второе слагаемое оцениваем неравенством Коши—Буняковского и условием
Липшица:
\[
  \big(\nabla f(x_{\mathrm{cp}}) - \nabla f(x_k)\big)\T (x - x_k)
  \le \norm{\nabla f(x_{\mathrm{cp}}) - \nabla f(x_k)}\cdot\norm{x-x_k}
  \le R\,\norm{x_{\mathrm{cp}} - x_k}\cdot\norm{x - x_k} .
\]
Но $\norm{x_{\mathrm{cp}} - x_k} = \theta\norm{x-x_k} \le \norm{x-x_k}$,
поэтому оценка продолжается до
\[
  R\,\norm{x-x_k}^2 = R\,\alpha^2 \norm{\nabla f(x_k)}^2 .
\]
Складывая,
\[
  f(x) - f(x_k)
  \le -\alpha\norm{\nabla f(x_k)}^2 + R\alpha^2\norm{\nabla f(x_k)}^2
  = \norm{\nabla f(x_k)}^2\,\big(R\alpha^2 - \alpha\big) ,
\]
что и есть~\eqref{eq:grad-main}.
\end{proof}

\subsubsection*{Сходимость метода с дроблением шага}

\begin{theorem}\label{thm:grad-split}
Пусть
\begin{itemize}[leftmargin=1.6em, topsep=2pt]
  \item $f \in C^1(\R^n)$;
  \item $f$ ограничена снизу;
  \item градиент липшицев с константой $R > 0$.
\end{itemize}
Пусть $\{x_k\}$ строится градиентным методом с дроблением шага. Тогда
\[
  \norm{\nabla f(x_k)} \xrightarrow[k \to \infty]{} 0
\]
независимо от выбора начального приближения.
\end{theorem}

\begin{proof}
\textbf{Шаг 1. Дробление конечно.} Условие приёма шага
\begin{equation}\label{eq:armijo}
  f(x) - f(x_k) \le -\eps\,\alpha\,\norm{\nabla f(x_k)}^2
\end{equation}
заведомо выполнено, если правая часть~\eqref{eq:grad-main} не превосходит
правой части~\eqref{eq:armijo}:
\[
  -\alpha(1-\alpha R) \le -\eps\alpha
  \iff 1 - \alpha R \ge \eps
  \iff \alpha \le \frac{1-\eps}{R} .
\]
Значит после конечного числа дроблений пробный шаг попадёт в этот диапазон и
будет принят.

\medskip
\textbf{Шаг 2. Монотонность.} Беря $\alpha = \alpha_k$, $x = x_{k+1}$,
получаем из~\eqref{eq:armijo}
\[
  f(x_{k+1}) - f(x_k) \le -\eps\,\alpha_k \norm{\nabla f(x_k)}^2 < 0 ,
\]
то есть $f(x_k)$ строго убывает. Ограниченная снизу убывающая
последовательность сходится, поэтому
\begin{equation}\label{eq:grad-diff0}
  f(x_k) - f(x_{k+1}) \xrightarrow[k \to \infty]{} 0 .
\end{equation}

\medskip
\textbf{Шаг 3. Оценка градиента сверху.} Переписывая~\eqref{eq:armijo}:
\begin{equation}\label{eq:grad-upper}
  \norm{\nabla f(x_k)}^2 \ \le\ \frac{f(x_k) - f(x_{k+1})}{\eps\,\alpha_k} .
\end{equation}

\medskip
\textbf{Шаг 4. Оценка шага снизу.} Пусть $i_k \ge 1$, то есть хотя бы одно
дробление было. Тогда предыдущее пробное значение
$\widetilde\alpha = \alpha_k/\lambda$ было \emph{отвергнуто}:
\[
  f\big(x_k - \widetilde\alpha \nabla f(x_k)\big) - f(x_k)
  \ >\ -\eps\,\widetilde\alpha\,\norm{\nabla f(x_k)}^2 .
\]
С другой стороны, по лемме~\ref{lem:grad-main} левая часть не превосходит
$-\widetilde\alpha(1 - \widetilde\alpha R)\norm{\nabla f(x_k)}^2$. Сравнивая,
\[
  -\widetilde\alpha(1 - \widetilde\alpha R) > -\eps\,\widetilde\alpha
  \iff 1 - \widetilde\alpha R < \eps
  \iff \widetilde\alpha > \frac{1-\eps}{R} ,
\]
откуда
\[
  \alpha_k = \lambda\,\widetilde\alpha > \lambda\,\frac{1-\eps}{R} .
\]
Если же $i_k = 0$, то $\alpha_k = \alpha^{(0)}$. В обоих случаях
\begin{equation}\label{eq:alpha-lower}
  \alpha_k \ \ge\ \bar\alpha := \min\left\{\alpha^{(0)},\ \lambda\,\frac{1-\eps}{R}\right\} > 0 .
\end{equation}

\medskip
\textbf{Шаг 5. Завершение.} Подставляя~\eqref{eq:alpha-lower}
в~\eqref{eq:grad-upper}:
\[
  \norm{\nabla f(x_k)}^2
  \ \le\ \frac{f(x_k) - f(x_{k+1})}{\eps\,\bar\alpha}
  \xrightarrow[k \to \infty]{} 0
\]
в силу~\eqref{eq:grad-diff0}. Значит $\norm{\nabla f(x_k)} \to 0$. Начальное
приближение нигде не использовалось, то есть сходимость глобальная.
\end{proof}

\begin{idea}[title={Зачем нужны обе оценки шага}]
В тетради доказательство обрывается на словах «\emph{одновременно получили
верхнюю и нижнюю оценки}» со знаками вопроса; шаг~5 восстановлен.

\medskip
Смысл восклицания такой: оценка $\alpha \le \frac{1-\eps}{R}$
\emph{сверху} гарантирует, что дробление когда-нибудь закончится, а оценка
$\alpha_k > \lambda\frac{1-\eps}{R}$ \emph{снизу} гарантирует, что принятый
шаг не окажется слишком мелким. Вторая нужна именно для того, чтобы в
неравенстве~\eqref{eq:grad-upper} можно было избавиться от $\alpha_k$ в
знаменателе: без неё дробь могла бы не стремиться к нулю.

\medskip
Случай $i_k = 0$ в тетради не оговорён — там выкладка неявно считает, что
хотя бы одно дробление произошло. Это дополнено выше.
\end{idea}

\begin{idea}
Обратите внимание, что доказано: стремление к нулю \emph{градиента}, а не
сходимость $x_k$ к точке минимума. Без дополнительных предположений большего
и не получить — $f$ может вообще не иметь минимума (например, монотонно
убывать к своей нижней грани). Сходимость самих точек даёт следующий
параграф, где добавляется сильная выпуклость.
\end{idea}

\subsubsection*{Сходимость метода наискорейшего спуска}

\begin{theorem}\label{thm:steepest}
Пусть выполнены условия теоремы~\ref{thm:grad-split}, а $\alpha_k$ выбирается
по принципу наискорейшего спуска. Тогда
\[
  \norm{\nabla f(x_k)} \xrightarrow[k \to \infty]{} 0
\]
независимо от выбора начального приближения.
\end{theorem}

\begin{proof}
По лемме~\ref{lem:grad-main} для любого $\alpha > 0$
\[
  f\big(x_k - \alpha\nabla f(x_k)\big) - f(x_k)
  \le -\alpha(1-\alpha R)\norm{\nabla f(x_k)}^2 .
\]
Правая часть как функция $\alpha$ — парабола $R\alpha^2 - \alpha$, умноженная
на $\norm{\nabla f(x_k)}^2$; её минимум достигается при
\[
  \alpha = \frac{1}{2R},
  \qquad\text{и равен}\qquad
  -\frac{1}{4R}\,\norm{\nabla f(x_k)}^2 .
\]
Шаг $\alpha_k$ выбран из условия наименьшего значения $f$, поэтому для любого
$\alpha$
\[
  f(x_k + \alpha_k p_k) - f(x_k) \ \le\ f(x_k + \alpha p_k) - f(x_k),
  \qquad p_k = -\nabla f(x_k) .
\]
Берём $\alpha = \tfrac{1}{2R}$:
\[
  f(x_{k+1}) - f(x_k) \ \le\ -\frac{1}{4R}\,\norm{\nabla f(x_k)}^2 .
\]
Отсюда $f(x_k)$ строго убывает и ограничена снизу, значит
$f(x_k) - f(x_{k+1}) \to 0$, а
\begin{equation}\label{eq:steepest-bound}
  \norm{\nabla f(x_k)}^2 \ \le\ 4R\,\big(f(x_k) - f(x_{k+1})\big)
  \xrightarrow[k\to\infty]{} 0 .
\end{equation}
\end{proof}

\begin{idea}
Фокус доказательства в том, что \emph{знать} оптимальный шаг не требуется:
$\alpha_k$ не вычисляется явно, он лишь не хуже «пробного» $\alpha = 1/(2R)$.
Именно поэтому константа Липшица входит в оценку, но не в сам метод.
\end{idea}

\subsection{Скорость сходимости при сильной выпуклости}

Теоремы \ref{thm:grad-split} и \ref{thm:steepest} дают только
$\norm{\nabla f(x_k)} \to 0$. Чтобы получить сходимость самих точек и оценку
скорости, добавим двустороннюю оценку на гессиан.

\begin{theorem}\label{thm:grad-rate}
Пусть $f \in C^2(\R^n)$ и существуют константы $m, M$ такие, что
\[
  m\norm{z}^2 \ \le\ z\T H(y)\, z \ \le\ M\norm{z}^2 ,
  \qquad 0 < m \le M, \qquad \forall\, z, y \in \R^n .
\]
Пусть $\{x_k\}$ строится градиентным методом с дроблением шага. Тогда при
любом начальном приближении $x_0$
\[
  x_k \xrightarrow[k \to \infty]{} x_* , \qquad
  f(x_k) \xrightarrow[k \to \infty]{} f(x_*) ,
\]
где $x_*$ — точка минимума, причём
\[
  f(x_{k+1}) - f(x_*) \le q\,\big(f(x_k) - f(x_*)\big), \qquad q \in (0,1) ,
\]
\[
  f(x_k) - f(x_*) \le q^k\big(f(x_0) - f(x_*)\big) ,
  \qquad
  \norm{x_k - x_*} \le C\, q^{k/2}, \quad C < +\infty .
\]
\end{theorem}

\begin{idea}
Условие $m\norm{z}^2 \le z\T H z$ — это \textbf{сильная выпуклость}: гессиан
равномерно положительно определён. Верхняя оценка $M$ говорит, что кривизна не
бывает и слишком большой. Отношение $M/m$ — число обусловленности задачи;
именно оно, как видно из итоговой оценки, и определяет скорость.

\medskip
Попутно отсюда следует, что $M \le R$: липшицевость градиента с константой
$R$ как раз ограничивает гессиан сверху.
\end{idea}

\subsubsection*{Шаг 1: множество Лебега ограничено}

Пусть $\bar x \in \R^n$ и
\[
  Y = \{\, x \mid f(x) \le f(\bar x) \,\} .
\]
Множество $Y$ замкнуто (функция $f$ непрерывна). Покажем, что оно ограничено.

Разложение по Тейлору в точке $\bar x$:
\[
  f(x) = f(\bar x) + \nabla\T f(\bar x)(x - \bar x)
       + \tfrac12 (x-\bar x)\T H(\hat x)(x - \bar x) ,
  \qquad \hat x = \bar x + \theta(x - \bar x), \ \theta \in [0,1] .
\]
Для $x \in Y$ имеем $f(x) \le f(\bar x)$, и, пользуясь нижней оценкой на
гессиан,
\[
  f(\bar x) \ \ge\ f(x)
  \ \ge\ f(\bar x) + \nabla\T f(\bar x)(x-\bar x)
         + \frac{m}{2}\norm{x - \bar x}^2 ,
\]
откуда
\[
  \nabla\T f(\bar x)(x-\bar x) + \frac{m}{2}\norm{x-\bar x}^2 \le 0
  \qquad\Longrightarrow\qquad
  \frac{m}{2}\norm{x-\bar x}^2 \le \norm{\nabla f(\bar x)}\cdot\norm{x-\bar x} ,
\]
то есть
\[
  \norm{x - \bar x} \ \le\ \frac{2}{m}\,\norm{\nabla f(\bar x)} .
\]

\begin{idea}
Итак, $Y$ замкнуто и ограничено, то есть компактно. Значит минимум $f$
достигается, а из строгой выпуклости (следствия $m > 0$) он единственный —
это и есть $x_*$. Заодно вся последовательность $\{x_k\}$ лежит в
компакте $Y$ с $\bar x = x_0$, поскольку $f(x_k)$ убывает.
\end{idea}

\subsubsection*{Шаг 2: две оценки через $m$ и $M$}

Разложим $f$ в точке минимума, где $\nabla f(x_*) = 0$:
\[
  f(x) = f(x_*) + \underbrace{\nabla\T f(x_*)(x - x_*)}_{=\,0}
       + \tfrac12 (x - x_*)\T H(\bar{\bar x})(x - x_*) ,
  \qquad \bar{\bar x} = x_* + \theta(x - x_*) ,
\]
откуда $f(x) - f(x_*) = \tfrac12 (x-x_*)\T H(\bar{\bar x})(x-x_*)$ и
\begin{equation}\label{eq:rate-2}
  \frac{m}{2}\norm{x - x_*}^2 \ \le\ f(x) - f(x_*)
  \ \le\ \frac{M}{2}\norm{x - x_*}^2 ,
\end{equation}
в частности
\begin{equation}\label{eq:rate-3}
  \norm{x - x_*}^2 \ \ge\ \frac{2}{M}\big(f(x) - f(x_*)\big) .
\end{equation}

Теперь разложим, наоборот, $f(x_*)$ вокруг произвольной точки $x$:
\[
  f(x_*) = f(x) + \nabla\T f(x)(x_* - x)
         + \tfrac12 (x_* - x)\T H(\bar x)(x_* - x) ,
  \qquad \bar x = x + \theta(x_* - x) ,
\]
то есть
\begin{equation}\label{eq:rate-1}
  \begin{aligned}
    f(x) - f(x_*)
    &= \nabla\T f(x)(x - x_*) - \tfrac12 (x - x_*)\T H(\bar x)(x - x_*) \\[2pt]
    &\le \norm{\nabla f(x)}\cdot\norm{x - x_*} - \frac{m}{2}\norm{x-x_*}^2 .
  \end{aligned}
\end{equation}

\subsubsection*{Шаг 3: оценка расстояния через градиент}

Соединяя левое неравенство~\eqref{eq:rate-2} с~\eqref{eq:rate-1}:
\[
  \frac{m}{2}\norm{x-x_*}^2
  \ \le\ \norm{\nabla f(x)}\cdot\norm{x-x_*} - \frac{m}{2}\norm{x-x_*}^2 ,
\]
то есть $m\norm{x-x_*}^2 \le \norm{\nabla f(x)}\cdot\norm{x-x_*}$, откуда
\begin{equation}\label{eq:rate-4}
  \norm{x - x_*} \ \le\ \frac{\norm{\nabla f(x)}}{m} .
\end{equation}

\subsubsection*{Шаг 4: градиент снизу}

Подставим~\eqref{eq:rate-4} в первое слагаемое~\eqref{eq:rate-1},
а~\eqref{eq:rate-3} — во второе:
\[
  f(x) - f(x_*)
  \ \le\ \frac{\norm{\nabla f(x)}^2}{m}
       - \frac{m}{2}\cdot\frac{2}{M}\big(f(x) - f(x_*)\big)
  = \frac{\norm{\nabla f(x)}^2}{m}
       - \frac{m}{M}\big(f(x) - f(x_*)\big) ,
\]
то есть $\left(1 + \dfrac{m}{M}\right)\big(f(x)-f(x_*)\big)
\le \dfrac{\norm{\nabla f(x)}^2}{m}$, окончательно
\begin{equation}\label{eq:rate-5}
  \norm{\nabla f(x)}^2 \ \ge\ m\left(1 + \frac{m}{M}\right)
  \big(f(x) - f(x_*)\big) .
\end{equation}

\subsubsection*{Шаг 5: геометрическое убывание}

Для метода с дроблением шага выполняются
\[
  x_{k+1} = x_k - \alpha_k \nabla f(x_k) ,
  \qquad
  f(x_{k+1}) - f(x_k) \le -\eps\,\alpha_k \norm{\nabla f(x_k)}^2 .
\]
Беря в~\eqref{eq:rate-5} точку $x = x_k$:
\begin{equation}\label{eq:rate-6}
  \norm{\nabla f(x_k)}^2 \ \ge\ m\left(1 + \frac{m}{M}\right)
  \big(f(x_k) - f(x_*)\big) .
\end{equation}
Тогда
\[
  \begin{aligned}
    \big(f(x_{k+1}) - f(x_*)\big) - \big(f(x_k) - f(x_*)\big)
      &= f(x_{k+1}) - f(x_k)
       \ \le\ -\eps\,\alpha_k \norm{\nabla f(x_k)}^2 \\[4pt]
      &\overset{\eqref{eq:rate-6}}{\le}\
       -\eps\,\alpha_k\, m\left(1 + \frac{m}{M}\right)
       \big(f(x_k) - f(x_*)\big) ,
  \end{aligned}
\]
откуда
\[
  \boxed{\
  f(x_{k+1}) - f(x_*) \ \le\ q\,\big(f(x_k) - f(x_*)\big) ,
  \qquad
  q = 1 - \eps\,\alpha_k\, m\left(1 + \frac{m}{M}\right) . \ }
\]

\subsubsection*{Шаг 6: итог}

Итерируя полученное неравенство,
\[
  f(x_k) - f(x_*) \ \le\ q^k \big(f(x_0) - f(x_*)\big) ,
\]
а из левой части~\eqref{eq:rate-2} при $x = x_k$
\[
  \frac{m}{2}\norm{x_k - x_*}^2 \le f(x_k) - f(x_*) \le q^k\big(f(x_0)-f(x_*)\big)
  \qquad\Longrightarrow\qquad
  \norm{x_k - x_*} \le C\, q^{k/2} ,
\]
где
\[
  C = \sqrt{\frac{2\big(f(x_0) - f(x_*)\big)}{m}} < +\infty .
\]
В частности $x_k \to x_*$ и $f(x_k) \to f(x_*)$. \hfill$\blacksquare$

\begin{idea}[title={Почему $q$ действительно лежит в $(0,1)$}]
\textbf{Сверху.} По оценке~\eqref{eq:alpha-lower} из
теоремы~\ref{thm:grad-split} шаг отделён от нуля: $\alpha_k \ge \bar\alpha > 0$.
Значит
\[
  q \le 1 - \eps\,\bar\alpha\, m\left(1 + \frac{m}{M}\right) < 1 ,
\]
и этот равномерный по $k$ множитель и надо брать в оценке $q^k$.

\medskip
\textbf{Снизу.} Неотрицательность $q$ получается даром. Действительно,
$f(x_{k+1}) \ge f(x_*)$, поэтому
\[
  \begin{aligned}
    f(x_k) - f(x_*)
      &\ \ge\ f(x_k) - f(x_{k+1})
       \ \ge\ \eps\,\alpha_k \norm{\nabla f(x_k)}^2 \\[4pt]
      &\ \overset{\eqref{eq:rate-6}}{\ge}\
       \eps\,\alpha_k\, m\left(1+\frac{m}{M}\right)\big(f(x_k) - f(x_*)\big) ,
  \end{aligned}
\]
и, сокращая на положительное $f(x_k) - f(x_*)$, получаем
$\eps\,\alpha_k m\left(1+\frac{m}{M}\right) \le 1$, то есть $q \ge 0$.
\end{idea}

\begin{idea}[title={Что это за скорость}]
По значениям функции убывание геометрическое со знаменателем $q$, по
расстоянию до минимума — со знаменателем $\sqrt q$. И то и другое — это
\textbf{линейная} скорость в смысле определения из \S11 (показатель $p=1$).

\medskip
Чем хуже обусловленность, тем медленнее: при $M \gg m$ множитель
$1 + m/M \to 1$, а сам $m$ мал, так что $q$ подбирается к единице. Это
аналитическое отражение того самого зигзага на вытянутых линиях уровня
(рис.~\ref{fig:zigzag}).
\end{idea}


\newpage
\lecture{Лекция 5}{}

\subsection{Уточнённые оценки шага и скорости}

Теорема~\ref{thm:grad-rate} оставила коэффициент сжатия в виде
$q = 1 - \eps\,\alpha_k\, m\big(1 + \frac{m}{M}\big)$ — через сам шаг $\alpha_k$.
Доведём выкладку до оценки, в которой $\alpha_k$ уже не участвует.

\subsubsection*{Одношаговая оценка через гессиан}

Запишем приращение за итерацию по формуле Тейлора с остаточным членом в форме
Лагранжа:
\begin{align*}
  f(x_{k+1}) - f(x_k)
  &= f\big(x_k - \alpha_k \nabla f(x_k)\big) - f(x_k) \\
  &= f(x_k) - \alpha_k \nabla\T f(x_k)\,\nabla f(x_k)
     + \tfrac{1}{2}\alpha_k^2\, \nabla\T f(x_k)\, H(\bar{x}_k)\, \nabla f(x_k)
     - f(x_k) ,
\end{align*}
где
\[
  \bar{x}_k = x_k + \theta\,(x_{k+1} - x_k)
            = x_k - \theta\,\alpha_k \nabla f(x_k) , \qquad \theta \in (0,1) .
\]
Подставляя верхнюю границу гессиана $z\T H z \le M\norm{z}^2$, получаем
\begin{equation}\label{eq:one-step-M}
  f(x_{k+1}) - f(x_k)
  \ \le\ -\alpha_k \norm{\nabla f(x_k)}^2
          + \frac{M}{2}\alpha_k^2 \norm{\nabla f(x_k)}^2
  = \norm{\nabla f(x_k)}^2 \Big( \frac{M}{2}\alpha_k^2 - \alpha_k \Big) .
\end{equation}

\begin{idea}
Это та же лемма~\ref{lem:grad-main}, но в более сильной форме. Там квадратичный
член оценивался через липшицеву константу градиента и давал $R\alpha^2$; здесь
используется разложение второго порядка, и коэффициент получается вдвое меньше:
$\frac{M}{2}\alpha^2$. Для дважды дифференцируемой $f$ константа Липшица
градиента как раз равна наибольшему собственному числу гессиана, $R = M$,
так что~\eqref{eq:one-step-M} строго точнее~\eqref{eq:grad-main}.
Поэтому и пороги для шага дальше получаются вдвое мягче:
$\frac{2}{M}(1-\eps)$ вместо $\frac{1-\eps}{R}$.
\end{idea}

\subsubsection*{Границы для принятого шага}

Напомним условие приёма шага~\eqref{eq:armijo}:
\[
  f(x_{k+1}) - f(x_k) \ \le\ -\eps\, \alpha_k \norm{\nabla f(x_k)}^2 ,
  \qquad \eps \in (0,1) .
\]
Из~\eqref{eq:one-step-M} видно, что для его выполнения достаточно
\[
  -\alpha_k + \frac{M}{2}\alpha_k^2 \ \le\ -\eps\,\alpha_k
  \qquad\Longrightarrow\qquad
\]
\begin{equation}\label{eq:alpha-upper-M}
  \alpha_k \ \le\ \frac{2}{M}\,(1 - \eps) .
\end{equation}

\begin{remark}
Алгоритм дробления: $\alpha_{k+1} = \alpha_k \lambda$, где $\lambda \in (0,1)$ —
параметр дробления. Оценки при этом остаются справедливы.
\end{remark}

\begin{idea}[title={Почему дробление не портит оценки}]
Повтор рассуждения шага~4 теоремы~\ref{thm:grad-split}, но уже с константой $M$.
Шаг, \emph{не} прошедший проверку, обязан нарушать~\eqref{eq:alpha-upper-M}:
\[
  -\eps\,\alpha < -\alpha + \frac{\alpha^2}{2}M
  \qquad\Longleftrightarrow\qquad
  \alpha > \frac{2}{M}(1 - \eps) .
\]
Последним отвергнутым был шаг $\alpha_k / \lambda$, поэтому
\begin{equation}\label{eq:alpha-lower-M}
  \frac{\alpha_k}{\lambda} > \frac{2}{M}(1 - \eps)
  \qquad\Longrightarrow\qquad
  \alpha_k > \frac{2}{M}\,\lambda\,(1 - \eps) .
\end{equation}
Итого принятый шаг зажат между двумя положительными константами:
сверху~\eqref{eq:alpha-upper-M} гарантирует, что дробление конечно,
снизу~\eqref{eq:alpha-lower-M} — что убывание функции за итерацию отделено от
нуля. Это и значит «оценки справедливы».
\end{idea}

\subsubsection*{Коэффициент сжатия без $\alpha_k$}

Подставим нижнюю границу шага и оценку~\eqref{eq:rate-5} в $q$ из
теоремы~\ref{thm:grad-rate}:
\[
  q = 1 - \eps\,\alpha_k\, m\Big(1 + \frac{m}{M}\Big)
    \ \le\ 1 - \eps(1-\eps)\,\frac{2m}{M}\Big(1 + \frac{m}{M}\Big) .
\]
Выражение $\eps(1-\eps)$ максимально при $\eps = \tfrac12$, где равно $\tfrac14$,
поэтому наилучшая оценка —
\begin{equation}\label{eq:q-min}
  \eps_{\max} = \frac{1}{2} ,
  \qquad
  q_{\min} = 1 - \frac{1}{2}\,\frac{m}{M}\Big(1 + \frac{m}{M}\Big) ,
  \qquad 0 < \frac{m}{M} \le 1 .
\end{equation}

\begin{idea}[title={Почему $0 \le q_{\min} < 1$}]
Обозначим $t = m/M$. Из $0 < m \le M$ следует $t \in (0,1]$, и
\[
  q_{\min} = 1 - \tfrac{1}{2}\,t\,(1 + t) .
\]
Функция $t \mapsto \tfrac12 t(1+t)$ строго возрастает на $(0,1]$ от $0$ до $1$,
значит $q_{\min} \in [0, 1)$ — сходимость геометрическая. Крайние случаи:
\begin{itemize}[nosep,left=0pt]
  \item $t = 1$ (то есть $m = M$, линии уровня — сферы): $q_{\min} = 0$ —
        метод попадает в минимум за один шаг;
  \item $t \to 0$ (сильно вытянутые линии уровня): $q_{\min} \to 1$ —
        сходимость сколь угодно медленная.
\end{itemize}
Множитель $\lambda$ из~\eqref{eq:alpha-lower-M} здесь опущен; строго говоря
$q = 1 - \eps(1-\eps)\frac{2m\lambda}{M}\big(1 + \frac{m}{M}\big)$, и
записанная оценка отвечает предельному случаю $\lambda \to 1$ (дробление мелкое,
принятый шаг почти максимально допустимый).
\end{idea}

Дальше всё как в шаге~6 теоремы~\ref{thm:grad-rate}:
\[
  f(x_k) - f(x_*) \ \le\ q\,\big(f(x_{k-1}) - f(x_*)\big)
  \ \le\ \dots \ \le\ q^{k}\big(f(x_0) - f(x_*)\big) ,
\]
а вместе с левой частью~\eqref{eq:rate-2}
\[
  \frac{m}{2}\norm{x_k - x_*}^2 \ \le\ f(x_k) - f(x_*)
  \ \le\ q^{k}\big(f(x_0) - f(x_*)\big)
  \quad\Longrightarrow\quad
  \norm{x_k - x_*} \le C\, q^{k/2} ,
\]
\[
  C = \sqrt{\frac{2}{m}\big(f(x_0) - f(x_*)\big)} .
\]

\subsection{Скорость наискорейшего спуска}

\begin{theorem}\label{thm:steepest-rate}
Для метода наискорейшего спуска справедливы оценки
\[
  f(x_{k+1}) - f(x_*) \ \le\ q\,\big(f(x_k) - f(x_*)\big) ,
  \qquad
  \norm{x_k - x_*} \ \le\ C\, q^{k/2} , \quad C < +\infty ,
\]
\[
  q = 1 - \frac{m}{2M}\Big(1 + \frac{m}{M}\Big) .
\]
\end{theorem}

\begin{proof}
Пусть $\alpha_k$ — шаг, полученный по методу наискорейшего спуска, то есть
$\alpha_k = \argmin_{\alpha \ge 0} f\big(x_k - \alpha \nabla f(x_k)\big)$,
а $\alpha$ — любое. По определению минимума и оценке~\eqref{eq:one-step-M},
записанной для произвольного $\alpha$,
\begin{align*}
  f(x_{k+1}) - f(x_k)
  &= f\big(x_k - \alpha_k \nabla f(x_k)\big) - f(x_k)
   \ \le\ f\big(x_k - \alpha \nabla f(x_k)\big) - f(x_k) \\
  &\le \Big(-\alpha + \frac{M}{2}\alpha^2\Big) \norm{\nabla f(x_k)}^2 .
\end{align*}
Неравенство верно при любом $\alpha$, поэтому возьмём лучшее:
\[
  -\alpha + \frac{M}{2}\alpha^2 \ \to\ \min_{\alpha > 0}
  \qquad\Longrightarrow\qquad
  \alpha_* = \frac{1}{M} , \qquad
  -\alpha_* + \frac{M}{2}\alpha_*^2 = -\frac{1}{2M} .
\]
Отсюда
\[
  f(x_{k+1}) - f(x_k) \ \le\ -\,\frac{\norm{\nabla f(x_k)}^2}{2M} ,
\]
а с учётом~\eqref{eq:rate-5}
\[
  f(x_{k+1}) - f(x_k)
  \ \le\ -\,\frac{1}{2M}\Big(1 + \frac{m}{M}\Big) m \big(f(x_k) - f(x_*)\big) .
\]
Наконец,
\[
  f(x_{k+1}) - f(x_*) = \big(f(x_{k+1}) - f(x_k)\big) + \big(f(x_k) - f(x_*)\big)
  \le \Big(1 - \frac{m}{2M}\Big(1 + \frac{m}{M}\Big)\Big)
      \big(f(x_k) - f(x_*)\big) ,
\]
то есть $q = 1 - \frac{m}{2M}\big(1 + \frac{m}{M}\big)$. Оценка по аргументу
получается из~\eqref{eq:rate-2} так же, как выше:
\[
  \norm{x_k - x_*} \le q^{k/2}\, C_* ,
  \qquad
  C_* = \sqrt{\frac{2}{m}\big(f(x_0) - f(x_*)\big)} .
  \qedhere
\]
\end{proof}

\begin{idea}[title={Почему $q$ совпало с $q_{\min}$ дробления}]
Это не совпадение. В методе дробления лучшее значение достигалось при
$\eps = \tfrac12$, а допустимый шаг при таком $\eps$ —
$\alpha \le \frac{2}{M}\big(1 - \tfrac12\big) = \frac{1}{M}$. Ровно это же
значение $\alpha_* = 1/M$ минимизирует квадратичную мажоранту
$-\alpha + \frac{M}{2}\alpha^2$ в доказательстве выше. Обе выкладки
оптимизируют одну и ту же верхнюю оценку, поэтому и дают один и тот же $q$.

Содержательно: наискорейший спуск по построению не хуже любого фиксированного
шага, в том числе $\alpha = 1/M$; оценка для него поэтому наследуется от оценки
для шага $1/M$. Реальная скорость наискорейшего спуска может быть лучше, но
доказать это через ту же мажоранту нельзя.
\end{idea}

\subsubsection*{Ортогональность соседних градиентов}

\begin{idea}[title={Откуда этот материал}]
На лекции зигзагообразная траектория наискорейшего спуска была нарисована
(рис.~\ref{fig:zigzag}), но не обоснована. Утверждение ниже объясняет её
точно: соседние градиенты перпендикулярны.
\end{idea}

\begin{theorem}\label{thm:steepest-orthogonal}
Пусть в методе наискорейшего спуска шаг $\alpha_k$ найден точно и минимум по
$\alpha$ достигается во внутренней точке ($\alpha_k > 0$). Тогда
\[
  \nabla\T f(x_{k+1})\,\nabla f(x_k) = 0 .
\]
\end{theorem}

\begin{proof}
Обозначим $g_k = \nabla f(x_k)$ и введём функцию одной переменной
\[
  \varphi(\alpha) = f\big(x_k - \alpha g_k\big) .
\]
По определению наискорейшего спуска $\alpha_k = \argmin_{\alpha > 0}
\varphi(\alpha)$, и если минимум внутренний, то $\varphi'(\alpha_k) = 0$. По
правилу дифференцирования сложной функции
\[
  \varphi'(\alpha) = \nabla\T f\big(x_k - \alpha g_k\big)\,(-g_k) ,
\]
а так как $x_{k+1} = x_k - \alpha_k g_k$, подстановка $\alpha = \alpha_k$ даёт
\[
  0 = \varphi'(\alpha_k) = -\nabla\T f(x_{k+1})\, g_k
  \qquad\Longrightarrow\qquad
  \nabla\T f(x_{k+1})\,\nabla f(x_k) = 0 . \qedhere
\]
\end{proof}

\begin{idea}
Отсюда и берётся зигзаг: направление очередного шага $-\nabla f(x_{k+1})$
перпендикулярно предыдущему, то есть траектория всё время поворачивает ровно на
прямой угол. На вытянутых линиях уровня (большое $M/m$) такие повороты очень
мелко «пилят» вдоль оврага, что и объясняет медленную сходимость с
$q \to 1$ из теоремы~\ref{thm:steepest-rate}.

Содержательно условие $\varphi'(\alpha_k) = 0$ означает: двигаться вдоль
текущего направления больше нечего — вся информация о нём исчерпана. Но на
следующем шаге метод про него забывает, и часть продвижения приходится
отыгрывать заново. Именно это исправляют методы сопряжённых градиентов
(\S\ref{subsec:conjugate-gradients}): там соседние направления делают
ортогональными не в обычном смысле, а в метрике гессиана, и для квадратичной
функции зигзаг исчезает вовсе.
\end{idea}

Итак, оценки получены и для шага, и для сходимости, и для её скорости. Отсюда
естественный вопрос: нельзя ли взять \emph{постоянный} шаг и не решать на каждой
итерации одномерную задачу?

\subsection{Метод с постоянным шагом}

\begin{theorem}\label{thm:const-step}
Пусть $\alpha_k = \alpha = \mathrm{const}$, $0 < \alpha < \dfrac{2}{M}$. Тогда для
скорости сходимости справедлива оценка
\[
  \norm{x_{k+1} - x_*} \ \le\ q\, \norm{x_k - x_*} ,
  \qquad
  q = \max\big\{\, \abs{1 - \alpha m},\ \abs{1 - \alpha M} \,\big\} ,
\]
причём
\[
  q_{\min} = \frac{M - m}{M + m}
  \quad \text{при} \quad
  \alpha = \frac{2}{M + m} ,
\]
и, кроме того,
\[
  f(x_k) - f(x_*) \ \le\ C\, q^{2k} , \qquad C < +\infty .
\]
\end{theorem}

\begin{proof}
Так как $\nabla f(x_*) = 0$, по формуле конечных приращений
\[
  \nabla f(x_k) - \nabla f(x_*) = H(\bar{x}_k)\,(x_k - x_*) ,
  \qquad \bar{x}_k = x_k + \theta\,(x_* - x_k), \quad \theta \in (0,1) .
\]
Поэтому
\begin{align*}
  \norm{x_{k+1} - x_*}
  &= \norm{x_k - \alpha \nabla f(x_k) - x_*}
   = \norm{x_k - x_* - \alpha\big(\nabla f(x_k) - \nabla f(x_*)\big)} \\
  &= \norm{x_k - x_* - \alpha H(\bar{x}_k)(x_k - x_*)}
   = \norm{\big(E - \alpha H(\bar{x}_k)\big)(x_k - x_*)} \\
  &\le \norm{E - \alpha H(\bar{x}_k)} \cdot \norm{x_k - x_*} .
\end{align*}
Обозначив $q = \norm{E - \alpha H(\bar{x}_k)}$, получаем
$\norm{x_{k+1} - x_*} \le q \norm{x_k - x_*}$.

Матрица $H(\bar{x}_k)$ симметрична, и её собственные числа лежат в $[m, M]$.
Тогда $E - \alpha H(\bar{x}_k)$ симметрична с собственными числами
$1 - \alpha \mu$, $\mu \in [m, M]$, и её спектральная норма равна
\[
  \norm{E - \alpha H(\bar{x}_k)}
  = \max_{\mu \in [m,M]} \abs{1 - \alpha \mu}
  = \max\big\{\, \abs{1 - \alpha m},\ \abs{1 - \alpha M} \,\big\} ,
\]
поскольку модуль линейной функции достигает максимума на конце отрезка. Эта
оценка не зависит от $\bar{x}_k$, то есть одна и та же на всех итерациях.

Минимум по $\alpha$ достигается там, где две ветви сравниваются:
\[
  1 - \alpha m = \alpha M - 1
  \qquad\Longrightarrow\qquad
  \alpha = \frac{2}{m + M} ,
  \qquad
  q_{\min} = 1 - \frac{2m}{m+M} = \frac{M - m}{M + m} .
\]

Наконец, из правой части~\eqref{eq:rate-2}
\[
  f(x_k) - f(x_*) \ \le\ \frac{M}{2}\norm{x_k - x_*}^2
  \ \le\ \frac{M}{2}\, q^{2k} \norm{x_0 - x_*}^2 = C\, q^{2k} ,
  \qquad
  C = \frac{M}{2}\norm{x_0 - x_*}^2 .
  \qedhere
\]
\end{proof}

\inputfig{figures/fig-q-alpha}

\begin{idea}[title={Два места в доказательстве, которые стоит разобрать}]
\textbf{1. Переход $\norm{A\,y} \le \norm{A}\cdot\norm{y}$.} Это определение
операторной (подчинённой) нормы матрицы:
$\norm{A} = \sup_{y \ne 0} \norm{Ay} / \norm{y}$, откуда
$\norm{Ay} \le \norm{A}\norm{y}$ для всех $y$ автоматически. Для евклидовой
нормы вектора подчинённая норма матрицы — спектральная, а для симметричной
матрицы она равна максимуму модулей собственных чисел. Именно поэтому в
следующей строке можно перейти от нормы матрицы к собственным числам гессиана.

\medskip
\textbf{2. Условие $0 < \alpha < 2/M$.} Метод сходится тогда и только тогда,
когда $q < 1$, то есть когда \emph{обе} величины $\abs{1 - \alpha m}$ и
$\abs{1 - \alpha M}$ меньше единицы. Так как $m \le M$, ограничение даёт именно
большая из них: $\abs{1 - \alpha M} < 1 \iff 0 < \alpha < 2/M$. При
$\alpha \ge 2/M$ шаг «перелетает» минимум и метод расходится.

\medskip
\textbf{Что показывает график.} Ветвь $1 - \alpha m$ убывает и обнуляется в
$\alpha = 1/m$; ветвь $\abs{1 - \alpha M}$ обнуляется раньше, в $\alpha = 1/M$,
и дальше растёт. Максимум из двух сначала идёт по первой ветви, после точки
пересечения — по второй, поэтому минимум максимума достигается ровно в точке
пересечения $\alpha = 2/(m+M)$.
\end{idea}

\begin{idea}
Сравним три метода. Все оценки выражаются через $t = m/M \in (0,1]$:
\begin{gather*}
  \text{дробление и наискорейший спуск:} \quad q = 1 - \tfrac{t}{2}(1+t) , \\
  \text{постоянный шаг:} \quad q^2 = \Big(\frac{1-t}{1+t}\Big)^2 .
\end{gather*}
Обе величины стремятся к единице при $t \to 0$, то есть все три метода одинаково
страдают от плохой обусловленности. Разница в цене итерации: наискорейший спуск
требует решения одномерной задачи, дробление — нескольких пробных вычислений
$f$, постоянный шаг — ни того, ни другого, но требует знания $m$ и $M$, которые
на практике редко известны. Радикально ситуацию меняет только переход к методам
второго порядка.
\end{idea}

\subsection{Обобщение направления спуска}
\label{subsec:descent-direction}

\begin{remark}
Вектор $p$ является направлением спуска в точке $x_k$ тогда и только тогда,
когда
\[
  \nabla\T f(x_k)\, p < 0 ,
\]
и, в частности, годится всякий вектор вида
\[
  p_k = -F(x_k)\, \nabla f(x_k) ,
\]
где $F(x_k)$ — некоторая матрица. Это и есть идея методов второго порядка.
\end{remark}

\begin{idea}[title={Откуда берётся вид $p_k = -F \nabla f$}]
Подставим $p_k = -F \nabla f(x_k)$ в критерий спуска (рис.~\ref{fig:descent-cone}):
\[
  \nabla\T f(x_k)\, p_k = -\,\nabla\T f(x_k)\, F\, \nabla f(x_k) < 0
\]
для любой положительно определённой $F$ и любого $\nabla f(x_k) \ne 0$.
То есть \emph{любая} положительно определённая
матрица $F$ порождает корректное направление спуска, и выбор $F$ — это и есть
выбор метода:
\begin{itemize}[nosep,left=0pt]
  \item $F = E$ — обычный градиентный спуск;
  \item $F = H^{-1}(x_k)$ — метод Ньютона (\S\ref{subsec:second-order});
  \item $F \approx H^{-1}(x_k)$, построенная по накопленным градиентам, —
        квазиньютоновские методы.
\end{itemize}
Смысл множителя $F$ — исправить перекос линий уровня: антиградиент указывает
направление наибыстрейшего убывания в \emph{евклидовой} метрике, а $F$ меняет
метрику так, чтобы линии уровня стали ближе к сферам. Именно отношение $m/M$,
портившее все оценки выше, при удачном $F$ перестаёт играть роль.
\end{idea}

\subsection{Метод циклического покоординатного спуска}

Задача $f(x_1, x_2, \dots, x_n) \to \min$, начальное приближение
$\big(x_1^{(0)}, x_2^{(0)}, \dots, x_n^{(0)}\big)\T$.

Упрощаем задачу, минимизируя по одной координате:
\[
  f\big(x_1,\, x_2^{(0)},\, \dots,\, x_n^{(0)}\big) \to \min_{x_1} ,
  \qquad \exists\, \bar{x}_1 \ \text{— решение} .
\]
Далее либо
\[
  \left[
  \begin{aligned}
    &f\big(x_1^{(0)},\, x_2,\, \dots,\, x_n^{(0)}\big) \to \min_{x_2} , \\
    &f\big(\bar{x}_1,\, x_2,\, \dots,\, x_n^{(0)}\big) \to \min_{x_2} ,
  \end{aligned}
  \right.
\]
и так далее по всем координатам. Процедура зацикливается тогда, когда
\[
  \norm{x^{(k)} - x^{(k-1)}} < \delta .
\]

\begin{idea}
Две записанные альтернативы — это два разных режима:
\begin{itemize}[nosep,left=0pt]
  \item минимизация по $x_2$ при \emph{старом} значении $x_1^{(0)}$ — все
        одномерные задачи решаются независимо от одной и той же точки
        (параллельный вариант, аналог метода Якоби);
  \item минимизация при уже обновлённом $\bar{x}_1$ — найденные координаты сразу
        идут в дело (последовательный вариант, аналог метода Гаусса–Зейделя).
\end{itemize}
Второй вариант обычно сходится быстрее, и именно он подразумевается под
«циклическим» покоординатным спуском. Производные здесь не нужны вовсе — метод
нулевого порядка, а каждая одномерная подзадача решается любым методом
из~\S\ref{subsec:direct-search-1d}.
\end{idea}

\subsection{Метод Хука--Дживса}

Состоит из двух шагов.

\paragraph{Шаг 1. Исследующий поиск.}
Вычисляется значение $f$ в точке $x^{(0)}$; поочерёдно каждая переменная
изменяется на свой шаг из набора $\{h_1, \dots, h_n\}$.

Для первой координаты:
\[
  \begin{aligned}
    &f\big(x^{(0)} + h_1 e_1\big) < f\big(x^{(0)}\big)
     &&\Longrightarrow\quad x^{(0)} := x^{(0)} + h_1 e_1 ; \\[4pt]
    &f\big(x^{(0)} + h_1 e_1\big) \ge f\big(x^{(0)}\big)
     &&\Longrightarrow\quad
       \text{проверяем } f\big(x^{(0)} - h_1 e_1\big) < f\big(x^{(0)}\big) :
  \end{aligned}
\]
\[
  \begin{aligned}
    &\text{да} &&\Longrightarrow\quad x^{(0)} := x^{(0)} - h_1 e_1 ; \\
    &\text{нет} &&\Longrightarrow\quad \text{координата остаётся прежней.}
  \end{aligned}
\]
Затем то же самое для $x_2$, и так далее, пока не переберём все переменные.
Полученная точка $x^{(1)}$ — новая \textbf{базовая точка}.

\begin{idea}
Логика пробного шага по координате: сначала пробуем сдвиг «вперёд» $+h_i e_i$;
если значение не улучшилось — пробуем «назад» $-h_i e_i$; если и это не помогло —
координату не трогаем и переходим к следующей. За один исследующий поиск каждая
координата меняется не более одного раза, и значение $f$ только невозрастает.
Если после полного прохода базовая точка не сдвинулась вовсе, шаги $h_i$
уменьшают (например, вдвое) и повторяют поиск — это и есть критерий остановки.
\end{idea}

\paragraph{Шаг 2. Поиск по образцу.}
\[
  x^{(k+1)} = x^{(k)} + \big(x^{(k)} - x^{(k-1)}\big) ,
  \qquad\text{то есть}\qquad
  x^{(k+1)} = 2x^{(k)} - x^{(k-1)} .
\]
\[
  \begin{aligned}
    &f\big(x^{(k+1)}\big) < f\big(x^{(k)}\big)
     &&\Longrightarrow\quad \text{строим дальше, продолжаем поиск;} \\
    &f\big(x^{(k+1)}\big) > f\big(x^{(k)}\big)
     &&\Longrightarrow\quad
       \text{откатываемся, последняя удачная точка — } x^{(k)} .
  \end{aligned}
\]

\inputfig{figures/fig-hooke-jeeves}

\begin{idea}
Смысл второго шага: исследующий поиск за цикл выявил удачное направление
$x^{(k)} - x^{(k-1)}$ — значит есть смысл пройти по нему ещё раз, не тратя
вычислений на пробные шаги. Это дешёвая замена одномерной минимизации: вместо
подбора длины шага метод просто повторяет уже сделанное смещение, пока функция
убывает. Отсюда характерная траектория — короткие «ощупывающие» зигзаги,
перемежающиеся длинными прыжками по образцу.
\end{idea}

\subsubsection*{Численный пример}

\begin{idea}[title={Откуда этот материал}]
Прогон метода Хука--Дживса на числах; на лекции метод был описан только
схемой.
\end{idea}

\begin{example}\label{ex:hooke-jeeves}
\[
  f(x_1, x_2) = (x_1 - 1)^2 + (x_2 - 2)^2 \ \longrightarrow\ \min ,
  \qquad x \in \R^2 .
\]
Минимум очевиден — $x_* = (1, 2)\T$, но метод до него доберётся, не вычисляя ни
одной производной. Начальные данные:
\[
  x^{(0)} = \begin{pmatrix} 0 \\ 0 \end{pmatrix} , \qquad
  h_1 = h_2 = 1 , \qquad \beta = \tfrac12 , \qquad \eps = 0{,}1 ,
\]
\[
  f\big(x^{(0)}\big) = (0-1)^2 + (0-2)^2 = 5 .
\]

\medskip
\noindent\textbf{Итерация 1. Исследующий поиск из $x^{(0)} = (0,0)\T$.}

Координата $x_1$, пробный шаг вперёд:
\[
  (0,0)\T + 1\cdot e_1 = (1,0)\T , \qquad f(1,0) = 0 + 4 = 4 < 5 ,
\]
шаг принимается, текущая точка $(1,0)\T$.

Координата $x_2$, пробный шаг вперёд:
\[
  (1,0)\T + 1\cdot e_2 = (1,1)\T , \qquad f(1,1) = 0 + 1 = 1 < 4 ,
\]
шаг принимается. Координаты кончились:
\[
  x^{(1)} = (1,1)\T , \qquad f\big(x^{(1)}\big) = 1 < 5 ,
\]
исследующий поиск удачен.

\medskip
\noindent\textbf{Поиск по образцу.}
\[
  x_{\text{обр}} = 2x^{(1)} - x^{(0)} = 2\begin{pmatrix} 1 \\ 1 \end{pmatrix}
    - \begin{pmatrix} 0 \\ 0 \end{pmatrix}
    = \begin{pmatrix} 2 \\ 2 \end{pmatrix} ,
  \qquad
  f(2,2) = 1 + 0 = 1 .
\]
Убывания нет: $f(x_{\text{обр}}) = 1 \not< f\big(x^{(1)}\big) = 1$ — точка по
образцу отвергается, новой базовой остаётся $x^{(1)} = (1,1)\T$.

\medskip
\noindent\textbf{Итерация 2. Исследующий поиск из $x^{(1)} = (1,1)\T$.}

Координата $x_1$, вперёд:
\[
  (2,1)\T , \qquad f(2,1) = 1 + 1 = 2 \ge 1 \quad\text{— отвергается;}
\]
назад:
\[
  (0,1)\T , \qquad f(0,1) = 1 + 1 = 2 \ge 1 \quad\text{— тоже отвергается,}
\]
координата $x_1$ не меняется.

Координата $x_2$, вперёд:
\[
  (1,2)\T , \qquad f(1,2) = 0 + 0 = 0 < 1 \quad\text{— принимается.}
\]
\[
  x^{(2)} = (1,2)\T , \qquad f\big(x^{(2)}\big) = 0 < 1 .
\]

\medskip
\noindent\textbf{Поиск по образцу.}
\[
  x_{\text{обр}} = 2x^{(2)} - x^{(1)} = (1,3)\T , \qquad f(1,3) = 0 + 1 = 1 > 0 ,
\]
откат, базовая точка $x^{(2)} = (1,2)\T$.

\medskip
Дальше исследующий поиск из $(1,2)\T$ при $h_i = 1$ улучшения не даёт (все
четыре пробные точки дают $f = 1 > 0$), поэтому шаги дробятся:
$h_i := \beta h_i = \tfrac12$, затем $\tfrac14$, $\tfrac18$ и так далее. Каждый
раз поиск безуспешен, пока не выполнится $\max_i h_i < \eps = 0{,}1$, и метод
останавливается в точке
\[
  x_* = (1, 2)\T , \qquad f(x_*) = 0 ,
\]
найденной уже на второй итерации. Значения в базовых точках:
$5 \to 1 \to 0$.
\end{example}

\begin{idea}
Характерная деталь: оба поиска по образцу были отвергнуты. Так получается
потому, что функция симметрична и шаг сетки крупный — удачное направление
$x^{(k)} - x^{(k-1)}$ «перелетает» через минимум. На вытянутых оврагах картина
обратная: именно шаги по образцу дают основное продвижение, а исследующий
поиск только подправляет направление.

Останов здесь определяется не близостью к $x_*$, а измельчением сетки: метод
не знает, что уже стоит в минимуме, и тратит ещё несколько безуспешных
исследующих поисков на дробление $h_i$ до $\eps$.
\end{idea}

\subsection{Методы второго порядка}
\label{subsec:second-order}

Будем считать, что $f(x)$ строго выпукла и дважды непрерывно дифференцируема,
так что существует гессиан $H_f(x)$. Разложим $f$ в окрестности точки $y$:
\[
  f(x) = \underbrace{f(y) + \nabla\T f(y)\,(x - y)
         + \tfrac{1}{2}(x-y)\T H(y)\,(x-y)}_{\psi(x)} + \dots
\]
Отбрасывая члены выше второго порядка, получаем
\[
  f(x) - f(y) \ \approx\ \psi(x) - f(y) ,
\]
где $\psi(x)$ — \textbf{квадратичная аппроксимация} функции $f$.

Минимизируем приближение по приращению $x - y$:
\[
  \frac{\partial\big(\psi(x) - f(y)\big)}{\partial (x-y)}
  = \nabla\T f(y) + H(y)\,(x - y) = 0
  \qquad\Longrightarrow\qquad
  x - y = -H^{-1}(y)\,\nabla f(y) ,
\]
откуда итерационная схема
\begin{equation}\label{eq:newton}
  x_{k+1} = x_k - H^{-1}(x_k)\, \nabla f(x_k) .
\end{equation}

\subsubsection*{Идея методов второго порядка}

\[
  x_0 \ \longrightarrow\ x_k :
  \quad
  \begin{cases}
    H(x_k) \ \longrightarrow\ H^{-1}(x_k), \\
    \nabla f(x_k)
  \end{cases}
  \quad \Longrightarrow \quad
  x_{k+1} = x_k - H^{-1}(x_k)\,\nabla f(x_k) .
\]

Если оставить схему~\eqref{eq:newton} как есть — это \textbf{метод Ньютона}.
При усилении
\[
  x_{k+1} = x_k - \alpha_k H^{-1}(x_k)\, \nabla f(x_k)
\]
получаем различные методы (выбор $\alpha_k$ — как в градиентных методах:
постоянный шаг, дробление, наискорейший спуск).

Шаг можно выбирать в виде
\[
  p_k = -\alpha_k H^{-1}(x_k)\,\nabla f(x_k)
  \qquad\Longleftrightarrow\qquad
  H(x_k)\, p_k = -\alpha_k \nabla f(x_k) ,
\]
то есть на каждой итерации либо обращение матрицы, либо решение СЛАУ.

\begin{idea}
Форма $H(x_k)p_k = -\alpha_k \nabla f(x_k)$ практически всегда предпочтительнее:
обращение матрицы стоит порядка $n^3$ операций и численно неустойчиво, а решение
системы с симметричной положительно определённой матрицей делается разложением
Холецкого — тот же порядок, но меньше констант и лучше точность. При больших $n$
систему решают приближённо, итерационно, не формируя $H$ целиком.

Связь с \S\ref{subsec:descent-direction}: здесь $F(x_k) = H^{-1}(x_k)$, и
положительная определённость $F$ (а значит корректность направления спуска)
следует из строгой выпуклости $f$. Если $f$ не выпукла, гессиан может оказаться
неположительно определённым — тогда ньютоновское направление перестаёт быть
направлением спуска, и матрицу приходится регуляризовать.

Выигрыш методов второго порядка — в том, что оценки~\eqref{eq:q-min} и
теоремы~\ref{thm:const-step}, упирающиеся в отношение $m/M$, здесь не возникают:
квадратичную функцию схема~\eqref{eq:newton} решает за один шаг при любой
обусловленности, а вблизи минимума даёт квадратичную скорость сходимости
($p = 2$ в смысле определения из~\S\ref{subsec:multidim-basics}). Плата —
вычисление и хранение $H$ на каждой итерации.
\end{idea}


\newpage
\lecture{Лекция 6}{}

\subsection{Разновидности методов второго порядка}
\label{subsec:newton-variants}

Общая схема
\[
  x_{k+1} = x_k + \alpha_k p_k ,
  \qquad\text{где}\qquad
  H(x_k)\, p_k = -\nabla f(x_k)
\]
(ньютоновское направление). Методы различаются выбором шага $\alpha_k$.

\begin{enumerate}[label=\arabic*),left=0pt]
  \item $\alpha_k \equiv 1$ — \textbf{чистый метод Ньютона}.
  \item $\alpha_0 = 1$, $\lambda \in (0,1)$, $\eps \in \big(0, \tfrac12\big]$ —
        принцип дробления по условию
        \begin{equation}\label{eq:newton-armijo}
          f(x_k + \alpha p_k) - f(x_k)
          \ \le\ \alpha\,\eps\,\nabla\T f(x_k)\, p_k
        \end{equation}
        — \textbf{демпфированный метод Ньютона}.
  \item $\big(f(x_k + \alpha p_k) - f(x_k)\big) \to \min\limits_{\alpha > 0}
        \iff f(x_k + \alpha p_k) \to \min\limits_{\alpha > 0}$ —
        наискорейший спуск вдоль ньютоновского направления.
\end{enumerate}

\inputfig{figures/fig-newton-quadratic}

\begin{idea}[title={Почему единичный шаг}]
Ньютоновское направление $p_k$ ведёт ровно в вершину квадратичной
аппроксимации $\psi_k$, построенной в точке $x_k$:
\[
  x_k + 1 \cdot p_k = \argmin_x \psi_k(x) .
\]
То есть $\alpha_k = 1$ — это «пройти весь путь до минимума параболоида».
Вблизи $x_*$ аппроксимация почти точна, и единичный шаг даёт квадратичную
сходимость. Вдали от минимума $\psi_k$ может плохо описывать $f$, и
единичный шаг способен увеличить $f$ — отсюда вариант~2) с дроблением.

Именно это и показывает рисунок с вложенными параболами: на каждой итерации
строится своя квадратичная аппроксимация, её вершина становится следующей
точкой, и чем ближе к минимуму, тем теснее парабола прилегает к $f$.
\end{idea}

\begin{idea}[title={Почему $\eps \le \tfrac12$}]
Условие~\eqref{eq:newton-armijo} — обычное условие достаточного убывания, но
записанное через $\nabla\T f(x_k) p_k < 0$ (а не через $\norm{\nabla f}^2$,
как в градиентных методах): теперь направление не антиградиентное, и убывание
измеряется вдоль него.

Ограничение $\eps \le \tfrac12$ нужно, чтобы единичный шаг проходил проверку.
Посчитаем для квадратичной $f$. Так как $H p = -\nabla f$, имеем
$p\T H p = -\nabla\T f\, p$, поэтому
\[
  f(x_k + \alpha p) - f(x_k)
  = \alpha\,\nabla\T f\, p + \frac{\alpha^2}{2}\, p\T H p
  = \nabla\T f\, p \left(\alpha - \frac{\alpha^2}{2}\right) .
\]
При $\alpha = 1$ это $\tfrac12 \nabla\T f\, p$, а правая
часть~\eqref{eq:newton-armijo} равна $\eps\,\nabla\T f\, p$. Поскольку
$\nabla\T f\, p < 0$, условие
\[
  \tfrac12 \nabla\T f\, p \ \le\ \eps\,\nabla\T f\, p
  \qquad\Longleftrightarrow\qquad
  \eps \ \le\ \tfrac12 .
\]
То есть при $\eps > \tfrac12$ дробление отвергало бы ньютоновский шаг даже
там, где он идеален, и метод терял бы квадратичную скорость, скатываясь к
линейной.
\end{idea}

\subsubsection*{Модификации: экономия на гессиане}

\begin{idea}[title={Откуда этот материал}]
На лекции перечислены только три способа выбора шага. Ниже — две стандартные
модификации, снимающие главный недостаток метода: пересчёт и обращение $H(x_k)$
на каждой итерации. Третья модификация — регуляризация
$H_k + \lambda_k E$ — разбирается отдельно как метод Марквардта
(\S\ref{subsec:marquardt}).
\end{idea}

\paragraph{Метод с фиксированной матрицей.}
Гессиан вычисляется один раз, в начальной точке:
\[
  H_0 = \nabla^2 f(x_0) ,
  \qquad
  x_{k+1} = x_k - \alpha_k H_0^{-1}\, \nabla f(x_k) .
\]
Система $H_0 p = -\nabla f(x_k)$ решается с одним и тем же разложением
матрицы, поэтому каждая итерация стоит как у градиентного метода.

\paragraph{Метод с периодическим обновлением матрицы.}
Гессиан фиксируется на $t$ итераций подряд, затем пересчитывается:
\[
  x_{lt+i+1} = x_{lt+i} - \alpha_{lt+i}\, H^{-1}(x_{lt})\, \nabla f(x_{lt+i}) ,
  \qquad i = \overline{0,\,t-1} ,
\]
где $l$ — номер «блока» итераций. При $t = 1$ это обычный метод Ньютона, при
$t = \infty$ — метод с фиксированной матрицей.

\begin{idea}
Чем платим. Обе схемы — это метод вида $p_k = -F \nabla f(x_k)$ из
\S\ref{subsec:descent-direction} с матрицей $F$, которая \emph{не} равна
$H^{-1}(x_k)$. Направление остаётся направлением спуска, пока $F \succ 0$ (а
$H_0^{-1} \succ 0$ в силу сильной выпуклости), поэтому сходимость не теряется.
Теряется квадратичная скорость: оценка теоремы~\ref{thm:newton-rate} опиралась
на то, что квадратичная модель построена \emph{в текущей} точке. С устаревшей
матрицей скорость становится линейной, причём тем хуже, чем дальше $x_k$ ушла
от $x_{lt}$.

Разумный компромисс — умеренное $t$: стоимость итерации падает в разы, а
скорость остаётся сверхлинейной на участке, где кривизна меняется слабо. Та же
идея, доведённая до конца, приводит к квазиньютоновским методам
(\S\ref{subsec:quasi-newton-anchor}): там матрица не замораживается, а дёшево
подправляется по наблюдаемым парам «сдвиг аргумента — сдвиг градиента».
\end{idea}

\subsection{Сходимость метода второго порядка}

\begin{theorem}\label{thm:newton-rate}
Пусть $f(x)$ дважды непрерывно дифференцируема и удовлетворяет условиям:
\begin{enumerate}[label=\arabic*),left=0pt]
  \item $\exists\, m, M$ $(0 < m \le M)$:
        $\ m\norm{z}^2 \le z\T H(y)\, z \le M\norm{z}^2$ \ $\forall z, y$;
  \item $\exists\, R > 0$:
        $\ \norm{H(x) - H(y)} \le R\,\norm{x - y}$ \ $\forall x, y$.
\end{enumerate}
Пусть $\{x_k\}$ строится методом второго порядка с выбором шага по методу
дробления. Тогда
\[
  \{x_k\} \xrightarrow[k \to \infty]{} \argmin f(x)
\]
с \textbf{квадратичной} скоростью независимо от выбора начального приближения
$x_0$, причём
\[
  \norm{x_{k+1} - x_*} \ \le\ \frac{R}{m}\,\norm{x_k - x_*}^2 .
\]
\end{theorem}

\begin{idea}
Условие~1) — то же, что в теореме~\ref{thm:grad-rate}: оно даёт существование и
единственность $x_*$, а заодно обратимость $H$ (значит ньютоновское направление
определено и является направлением спуска). Условие~2) — новое: липшицевость
\emph{гессиана}. Именно она отвечает за квадратичную оценку: ошибка квадратичной
аппроксимации имеет порядок $R\norm{x - x_*}^2$, и ровно столько же остаётся от
расстояния до минимума после шага.

Глобальность сходимости обеспечивает дробление: вдали от $x_*$ шаг уменьшается
и $f$ монотонно убывает, а вблизи — по предыдущему пояснению единичный шаг
проходит проверку при $\eps \le \tfrac12$, и включается квадратичная скорость.
Это и есть главный выигрыш по сравнению с методами первого порядка, где
оценки~\eqref{eq:q-min} и теоремы~\ref{thm:const-step} давали лишь линейную
скорость с $q$, зависящим от обусловленности $m/M$.
\end{idea}

\subsection{Квазиньютоновские методы}
\label{subsec:quasi-newton-anchor}

Пусть последовательность $\{x_k\}$ строится по схеме
\[
  x_k = x_{k-1} + \alpha_k p_k , \qquad k \in \NN ,
\]
причём направление спуска берётся в виде
\[
  p_k = -A_k \nabla f(x_{k-1}) = A_k\, \omega_k , \qquad k \in \NN ,
  \qquad\text{где}\qquad
  \omega_k = -\nabla f(x_{k-1}) .
\]
Здесь $A_k$ — положительно определённая матрица порядка $n$, обновляемая на
$k$-й итерации (заменитель гессиана), причём
\[
  A_k \xrightarrow[k \to \infty]{} H^{-1}(x_*) ,
  \qquad
  A_{k+1} = A_k + \Delta A_k , \quad k \in \NN ,
\]
где $\Delta A_k$ — поправочная матрица порядка $n$. Начальное значение
$A_1 = E$.

\begin{idea}
Это ровно схема $p_k = -F(x_k)\nabla f(x_k)$ из \S\ref{subsec:descent-direction},
где $F = A_k$. Идея: гессиан не вычислять вовсе, а накапливать приближение к
$H^{-1}$ по наблюдаемым парам «сдвиг аргумента — сдвиг градиента». Поправка
$\Delta A_k$ подбирается так, чтобы выполнялось \textbf{секущее уравнение}
\[
  A_{k+1}\,\Delta\omega_k = -\Delta x_k ,
\]
то есть чтобы $A_{k+1}$ действовала на уже известном направлении так же, как
действовал бы настоящий $H^{-1}$ (поскольку
$H\,\Delta x \approx \Delta \nabla f = -\Delta\omega$). Положительная
определённость $A_k$ нужна, чтобы $p_k$ оставалось направлением спуска.

Старт с $A_1 = E$ означает, что первая итерация — обычный градиентный шаг:
информации о кривизне ещё нет.
\end{idea}

Остальные методы различаются по способу построения поправочной матрицы.

\paragraph{ДФП-метод (Дэвидона--Флетчера--Пауэлла).}
\[
  A_{k+1} = A_k
    - \frac{\Delta x_k\, (\Delta x_k)\T}{\Delta\omega_k\T\, \Delta x_k}
    - \frac{A_k \Delta\omega_k\, (\Delta\omega_k)\T A_k\T}
           {\Delta\omega_k\T A_k \Delta\omega_k} .
\]

\paragraph{БФШ-метод (Бройдена--Флетчера--Голдфарба--Шенно).}
\[
  A_{k+1} = A_k
    - \frac{\Delta x_k\, (\Delta x_k)\T}{\Delta\omega_k\T\, \Delta x_k}
    - \frac{A_k \Delta\omega_k\, (\Delta\omega_k)\T A_k\T}
           {\Delta\omega_k\T A_k \Delta\omega_k}
    + \big(\Delta\omega_k\T A_k \Delta\omega_k\big)\, r^k (r^k)\T ,
\]
\[
  r^k = \frac{A_k \Delta\omega_k}{\Delta\omega_k\T A_k \Delta\omega_k}
      - \frac{\Delta x_k}{\Delta x_k\T \Delta\omega_k} .
\]

Также: \textbf{метод Пауэлла}, \textbf{метод Мак-Кормика}.

\begin{idea}
БФШ — это ДФП плюс одно слагаемое ранга единица; оба обновления сохраняют
симметрию и положительную определённость $A_k$ (при
$\Delta\omega_k\T \Delta x_k < 0$, что обеспечивается точным или достаточно
аккуратным одномерным поиском). На практике БФШ устойчивее ДФП к неточной
одномерной минимизации, поэтому именно он стал стандартным.

Знаки в формулах объясняются тем, что здесь работают с $\omega = -\nabla f$,
а не с самим градиентом: при подстановке $\Delta\omega_k = -\Delta \nabla f_k$
обе формулы превращаются в привычные записи ДФП и БФШ.
\end{idea}

\subsubsection*{Общий алгоритм}

\begin{algo}[title={Квазиньютоновский метод}]
\textbf{Предварительный этап (1):}
\begin{itemize}[nosep,left=0pt]
  \item задать $\eps$ (параметр в условии останова $\norm{\omega_k} < \eps$);
  \item выбрать начальную точку $x_0$;
  \item $A_1 = E$;
  \item $I_0 = \{n,\, 2n,\, \dots\}$ — моменты обновления алгоритма;
  \item $k := 1$;
  \item $\omega_1 := -\nabla f(x_0)$.
\end{itemize}

\medskip
\textbf{Шаг (2):}
\begin{steps}
  \item На $k$-м шаге проверка останова $\norm{\omega_k} < \eps$:\\
        если да — $x_* \approx x_{k-1}$, $f(x_*) \approx f(x_{k-1})$, конец.
  \item Вычисление направления спуска $p_k = A_k \omega_k$ и шага
        \[
          \psi_k(\alpha) = f(x_{k-1} + \alpha p_k) \to \min_\alpha ,
          \qquad
          x_k = x_{k-1} + \alpha_k p_k ,
        \]
        затем $\omega_{k+1} = -\nabla f(x_k)$.\\
        Если $k \in I_0$ — положить $A_{k+1} = E$, $k := k+1$ и вернуться
        к~шагу~1; иначе перейти к~шагу~3.
  \item $\Delta x_k := x_k - x_{k-1}$, \quad
        $\Delta\omega_k := \omega_{k+1} - \omega_k$.
  \item Вычислить $A_{k+1}$ по формуле метода, $k := k+1$, вернуться к~шагу~1.
\end{steps}
\end{algo}

\begin{idea}
Зачем сброс $A_{k+1} = E$ каждые $n$ итераций (множество $I_0$). Во-первых,
для квадратичной функции квазиньютоновские формулы восстанавливают $H^{-1}$
точно за $n$ шагов — дальше новой информации о кривизне уже не появляется.
Во-вторых, для неквадратичной $f$ накопленные поправки постепенно перестают
отвечать текущей кривизне, а ошибки округления могут нарушить положительную
определённость $A_k$. Сброс на единичную матрицу гарантированно возвращает
корректное направление спуска ценой одного градиентного шага.
\end{idea}

\subsubsection*{Свойства формул обновления}

\begin{idea}[title={Откуда этот материал}]
На лекции формулы ДФП и БФГШ были выписаны без обоснования. Ниже доказано, что
обе удовлетворяют секущему условию и сохраняют положительную определённость —
то есть что направление $p_k = -A_k \nabla f(x_k)$ остаётся направлением
спуска.
\end{idea}

Перейдём к стандартным обозначениям: $s_k = x_{k+1} - x_k$,
$y_k = \nabla f(x_{k+1}) - \nabla f(x_k)$ (в обозначениях лекции
$s_k = \Delta x_k$, $y_k = -\Delta\omega_k$). \textbf{Секущее
условие} — это требование
\[
  A_{k+1}\, y_k = s_k ,
\]
а \textbf{условие кривизны} — $s_k\T y_k > 0$.

\begin{theorem}[ДФП: секущее условие]\label{thm:dfp-secant}
Пусть
\[
  A_{k+1} = A_k + \frac{s_k s_k\T}{s_k\T y_k}
            - \frac{A_k y_k y_k\T A_k}{y_k\T A_k y_k} .
\]
Тогда $A_{k+1} y_k = s_k$.
\end{theorem}

\begin{proof}
Умножим формулу справа на $y_k$:
\[
  A_{k+1} y_k = A_k y_k + \frac{s_k \,(s_k\T y_k)}{s_k\T y_k}
    - \frac{A_k y_k\,(y_k\T A_k y_k)}{y_k\T A_k y_k} .
\]
Величины $s_k\T y_k$ и $y_k\T A_k y_k$ — числа, поэтому второе слагаемое равно
$s_k$, а третье равно $A_k y_k$. Следовательно
$A_{k+1} y_k = A_k y_k + s_k - A_k y_k = s_k$.
\end{proof}

\begin{theorem}[ДФП: положительная определённость]\label{thm:dfp-pd}
Пусть $A_k \succ 0$ и $s_k\T y_k > 0$. Тогда $A_{k+1} \succ 0$.
\end{theorem}

\begin{proof}
Для произвольного $z \ne 0$
\[
  z\T A_{k+1} z = z\T A_k z
    - \frac{\big(z\T A_k y_k\big)^2}{y_k\T A_k y_k}
    + \frac{\big(z\T s_k\big)^2}{s_k\T y_k} .
\]
Так как $A_k \succ 0$, выражение
$\langle u, v\rangle_{A_k} = u\T A_k v$ — скалярное произведение, и по
неравенству Коши--Буняковского
\[
  \big(z\T A_k y_k\big)^2 \le \big(z\T A_k z\big)\big(y_k\T A_k y_k\big)
  \qquad\Longrightarrow\qquad
  z\T A_k z - \frac{\big(z\T A_k y_k\big)^2}{y_k\T A_k y_k} \ \ge\ 0 .
\]
Второе слагаемое неотрицательно, так как $s_k\T y_k > 0$. Значит
$z\T A_{k+1} z \ge 0$.

Покажем, что равенство невозможно при $z \ne 0$. Если $z\T A_{k+1}z = 0$, то
оба неотрицательных слагаемых нулевые. Из обращения в равенство неравенства
Коши следует, что $z$ пропорционален $y_k$: $z = c\,y_k$. Из второго
слагаемого $\big(z\T s_k\big)^2 = 0$, то есть $z\T s_k = 0$, откуда
$c\,y_k\T s_k = 0$. Но $y_k\T s_k = s_k\T y_k > 0$, значит $c = 0$ и $z = 0$ —
противоречие. Следовательно $z\T A_{k+1} z > 0$ при всех $z \ne 0$.
\end{proof}

\begin{theorem}[БФГШ: секущее условие]\label{thm:bfgs-secant}
Пусть
\[
  A_{k+1} = \big(E - \rho_k s_k y_k\T\big) A_k \big(E - \rho_k y_k s_k\T\big)
            + \rho_k s_k s_k\T ,
  \qquad \rho_k = \frac{1}{y_k\T s_k} .
\]
Тогда $A_{k+1} y_k = s_k$.
\end{theorem}

\begin{proof}
Рассмотрим правый множитель:
\[
  \big(E - \rho_k y_k s_k\T\big) y_k = y_k - \rho_k y_k \big(s_k\T y_k\big)
  = y_k - y_k = 0 ,
\]
так как $\rho_k \big(s_k\T y_k\big) = 1$. Значит первое слагаемое обращается в
нуль, и остаётся
\[
  \rho_k s_k s_k\T y_k = \frac{1}{y_k\T s_k}\, s_k \big(s_k\T y_k\big) = s_k .
  \qedhere
\]
\end{proof}

\begin{theorem}[БФГШ: положительная определённость]\label{thm:bfgs-pd}
Пусть $A_k \succ 0$ и $y_k\T s_k > 0$. Тогда $A_{k+1} \succ 0$.
\end{theorem}

\begin{proof}
Обозначим $w = \big(E - \rho_k y_k s_k\T\big) z$; тогда
$z\T\big(E - \rho_k s_k y_k\T\big) = w\T$, и для произвольного $z$
\[
  z\T A_{k+1} z = w\T A_k w + \rho_k \big(z\T s_k\big)^2 .
\]
Оба слагаемых неотрицательны ($A_k \succ 0$ и $\rho_k > 0$), значит
$z\T A_{k+1}z \ge 0$.

Пусть $z\T A_{k+1}z = 0$. Тогда $w\T A_k w = 0$ и $\big(z\T s_k\big)^2 = 0$.
Из $A_k \succ 0$ следует $w = 0$; из второго равенства $s_k\T z = 0$. Но
\[
  w = z - \rho_k y_k \big(s_k\T z\big) = z ,
\]
и, поскольку $w = 0$, получаем $z = 0$. Значит $z\T A_{k+1}z > 0$ при
$z \ne 0$.
\end{proof}

\begin{idea}
Условие кривизны $s_k\T y_k > 0$ выполняется автоматически для сильно выпуклой
$f$ и при точном одномерном поиске; в общем случае его обеспечивают условиями
Вульфа на шаг. Если же $s_k\T y_k \le 0$, обе формулы могут потерять
положительную определённость — тогда обновление пропускают или делают сброс
$A_{k+1} = E$ (шаг~9 алгоритма §\ref{subsec:quasi-newton-anchor}).

БФГШ устойчивее ДФП именно из-за своей формы: обновление записано как
«сопряжение» $W A_k W\T$ плюс слагаемое ранга один, и положительная
определённость в нём не может потеряться из-за накопления ошибок округления,
тогда как в ДФП она держится на разности двух слагаемых.
\end{idea}

\subsubsection*{Численный пример (ДФП)}

\begin{idea}[title={Откуда этот материал}]
Прогон ДФП на числах; на лекции формулы были выписаны без примера. Обозначения
стандартные: $s_k = x_{k+1} - x_k$, $y_k = \nabla f(x_{k+1}) - \nabla f(x_k)$,
нумерация точек с нуля.
\end{idea}

\begin{example}\label{ex:dfp-run}
\[
  f(x_1, x_2) = x_1^2 + 2x_2^2 \ \longrightarrow\ \min ,
  \qquad
  \nabla f(x) = \begin{pmatrix} 2x_1 \\ 4x_2 \end{pmatrix} ,
  \qquad
  x_0 = \begin{pmatrix} 2 \\ 1 \end{pmatrix} , \quad A_0 = E .
\]
Шаг выбирается точной одномерной минимизацией.

\medskip
\noindent\textbf{Итерация 0.}
\[
  g_0 = \nabla f(x_0) = \begin{pmatrix} 4 \\ 4 \end{pmatrix} ,
  \qquad
  p_0 = -A_0 g_0 = \begin{pmatrix} -4 \\ -4 \end{pmatrix}
\]
(первый шаг — обычный антиградиентный, потому что $A_0 = E$). Одномерная
задача:
\[
  \varphi(\alpha) = f(x_0 + \alpha p_0)
    = (2 - 4\alpha)^2 + 2(1 - 4\alpha)^2
    = 6 - 32\alpha + 48\alpha^2 ,
\]
\[
  \varphi'(\alpha) = -32 + 96\alpha = 0
  \quad\Longrightarrow\quad
  \alpha_0 = \tfrac13 ,
\]
\[
  x_1 = x_0 + \alpha_0 p_0 = \begin{pmatrix} 2/3 \\ -1/3 \end{pmatrix} ,
  \qquad
  f(x_1) = \tfrac49 + \tfrac29 = \tfrac23
  \qquad (6 \to \tfrac23) .
\]

\medskip
\noindent\textbf{Пересчёт матрицы.}
\[
  g_1 = \nabla f(x_1) = \begin{pmatrix} 4/3 \\ -4/3 \end{pmatrix} ,
  \qquad
  s_0 = x_1 - x_0 = \begin{pmatrix} -4/3 \\ -4/3 \end{pmatrix} ,
  \qquad
  y_0 = g_1 - g_0 = \begin{pmatrix} -8/3 \\ -16/3 \end{pmatrix} .
\]
Условие кривизны:
\[
  s_0\T y_0 = \tfrac{32}{9} + \tfrac{64}{9} = \tfrac{96}{9} = \tfrac{32}{3} > 0 .
\]
Поскольку $A_0 = E$, имеем $A_0 y_0 = y_0$ и
$y_0\T A_0 y_0 = \tfrac{64}{9} + \tfrac{256}{9} = \tfrac{320}{9}$. Слагаемые
формулы ДФП:
\[
  \frac{s_0 s_0\T}{s_0\T y_0}
    = \frac{16/9}{32/3}\begin{pmatrix} 1 & 1 \\ 1 & 1 \end{pmatrix}
    = \frac16 \begin{pmatrix} 1 & 1 \\ 1 & 1 \end{pmatrix} ,
\]
\[
  \frac{A_0 y_0 y_0\T A_0}{y_0\T A_0 y_0}
    = \frac{9}{320}\begin{pmatrix} 64/9 & 128/9 \\ 128/9 & 256/9 \end{pmatrix}
    = \begin{pmatrix} 1/5 & 2/5 \\ 2/5 & 4/5 \end{pmatrix} .
\]
Итого
\[
  A_1 = A_0 + \frac{s_0 s_0\T}{s_0\T y_0}
            - \frac{A_0 y_0 y_0\T A_0}{y_0\T A_0 y_0}
      = \begin{pmatrix} 29/30 & -7/30 \\[2pt] -7/30 & 11/30 \end{pmatrix} .
\]
Проверка секущего условия (теорема~\ref{thm:dfp-secant}):
\[
  A_1 y_0
   = \begin{pmatrix} 29/30 & -7/30 \\ -7/30 & 11/30 \end{pmatrix}
     \begin{pmatrix} -8/3 \\ -16/3 \end{pmatrix}
   = \begin{pmatrix} -4/3 \\ -4/3 \end{pmatrix} = s_0 .
\]

\medskip
\noindent\textbf{Итерация 1.}
\[
  A_1 g_1
   = \begin{pmatrix} 29/30 & -7/30 \\ -7/30 & 11/30 \end{pmatrix}
     \begin{pmatrix} 4/3 \\ -4/3 \end{pmatrix}
   = \begin{pmatrix} 8/5 \\ -4/5 \end{pmatrix} ,
  \qquad
  p_1 = -A_1 g_1 = \begin{pmatrix} -8/5 \\ 4/5 \end{pmatrix} .
\]
\[
  \varphi(\alpha) = f(x_1 + \alpha p_1)
    = \Big(\tfrac23 - \tfrac85\alpha\Big)^2
      + 2\Big(-\tfrac13 + \tfrac45\alpha\Big)^2
    = \tfrac23 - \tfrac{16}{5}\alpha + \tfrac{96}{25}\alpha^2 ,
\]
\[
  \varphi'(\alpha) = -\tfrac{16}{5} + \tfrac{192}{25}\alpha = 0
  \quad\Longrightarrow\quad
  \alpha_1 = \tfrac{5}{12} ,
\]
\[
  x_2 = x_1 + \alpha_1 p_1
      = \begin{pmatrix} 2/3 - 2/3 \\ -1/3 + 1/3 \end{pmatrix}
      = \begin{pmatrix} 0 \\ 0 \end{pmatrix} = x_* ,
  \qquad f(x_2) = 0 .
\]
Значения: $6 \to \tfrac23 \to 0$.
\end{example}

\begin{idea}
Точное попадание в минимум за две итерации — не случайность. Для
квадратичной функции $n$ переменных квазиньютоновский метод с \emph{точным}
одномерным поиском завершается не более чем за $n$ шагов: накопленные пары
$(s_k, y_k)$ полностью восстанавливают гессиан, и после $n$-й итерации
$A_n = H^{-1}$. Здесь $n = 2$, поэтому хватило двух шагов. Проверить это можно
прямо: $H = \begin{pmatrix} 2 & 0 \\ 0 & 4 \end{pmatrix}$,
$H^{-1} = \begin{pmatrix} 1/2 & 0 \\ 0 & 1/4 \end{pmatrix}$ — и именно к этой
матрице сошлась бы $A_2$, если бы её потребовалось строить.

Сравните с примером~\ref{ex:one-dim-run} по духу: там тоже использовались
только значения функции вдоль прямой, но здесь направление каждый раз
корректируется накопленной кривизной, а не берётся по антиградиенту.
\end{idea}

\subsection{Метод Марквардта}
\label{subsec:marquardt}

\textbf{Идея:} соединить характеристики методов первого и второго порядков.

\begin{idea}
В литературе этот метод чаще называют \textbf{методом
Левенберга--Марквардта}, а замену $H_k$ на $H_k + \lambda_k E$ —
\textbf{регуляризацией} гессиана. Это третья (наряду с двумя из
\S\ref{subsec:newton-variants}) стандартная модификация метода Ньютона, и
единственная из них, которая чинит не стоимость итерации, а саму матрицу: она
применима и тогда, когда $H_k$ не положительно определена.
\end{idea}

\begin{algo}[title={Метод Марквардта}]
\textbf{Начальный этап:} задать $\eps > 0$, $\lambda_0 = 10^4$, начальное
приближение $x_0$, $k = 0$.

\medskip
\textbf{Основной этап:}
\begin{steps}
  \item Вычислить $\nabla f(x_k)$.
  \item Проверить $\norm{\nabla f(x_k)} < \eps$: \ да — конец; нет — к~шагу~3.
  \item Вычислить $H_k$ и найти $p_k$ из системы
        \[
          \big(H_k + \lambda_k E\big)\, p_k = -\nabla f(x_k) .
        \]
  \item $x_{k+1} := x_k + p_k$.
  \item Проверить $f(x_{k+1}) < f(x_k)$: \ да — к~шагу~6; нет — к~шагу~7.
  \item $\lambda_{k+1} := \dfrac{\lambda_k}{2}$, \ $k := k+1$, к~шагу~1.
  \item $\lambda_{k+1} := 2\lambda_k$, вернуться к~шагу~3.
\end{steps}
\end{algo}

\begin{idea}
Смысл параметра $\lambda_k$ виден из крайних случаев системы на шаге~3:
\begin{itemize}[nosep,left=0pt]
  \item $\lambda_k \to 0$: система превращается в $H_k p_k = -\nabla f(x_k)$ —
        чистый метод Ньютона, быстрый вблизи минимума;
  \item $\lambda_k \to \infty$: слагаемое $\lambda_k E$ доминирует, и
        $p_k \approx -\frac{1}{\lambda_k}\nabla f(x_k)$ — градиентный шаг с
        малой длиной, надёжный вдали от минимума.
\end{itemize}
Старт с $\lambda_0 = 10^4$ — это «начать как градиентный метод». Удачный шаг
уменьшает $\lambda$ вдвое (сдвиг в сторону Ньютона), неудачный — удваивает
(откат к осторожному градиентному шагу). Попутно $\lambda_k E$ делает матрицу
положительно определённой даже там, где $H_k$ таковой не является, то есть
направление $p_k$ всегда остаётся направлением спуска.

В шаге~7 возврат именно к~шагу~3, а не к~4: после изменения $\lambda_k$ систему
надо решать заново, иначе $p_k$ останется прежним и точка не сдвинется.
\end{idea}

\subsection{Методы сопряжённых градиентов}
\label{subsec:conjugate-gradients}

\[
  x^{(k+1)} = x^{(k)} - \alpha^{(k)} p^{(k)} , \qquad k \in \NN_0 ,
\]
\[
  p^{(0)} = \nabla f\big(x^{(0)}\big) ,
  \qquad
  p^{(k)} = \nabla f\big(x^{(k)}\big) + \beta^{(k)} p^{(k-1)} ,
  \qquad k \in \NN ,
\]
\[
  \beta^{(k)}
  = \frac{\norm{\nabla f\big(x^{(k)}\big)}^2}
         {\norm{\nabla f\big(x^{(k-1)}\big)}^2}
  = \left( \sum_{i=1}^{n}
      \left(\frac{\partial f\big(x^{(k-1)}\big)}{\partial x_i}\right)^{\!2}
    \right)^{\!-1}
    \sum_{i=1}^{n}
      \left(\frac{\partial f\big(x^{(k)}\big)}{\partial x_i}\right)^{\!2} ,
  \qquad k \in \NN ,
\]
\[
  \beta^{(mn)} = 0 , \qquad m \in \NN_0 .
\]

\begin{idea}[title={Что означает $\beta^{(mn)} = 0$}]
Это \textbf{рестарт}: каждые $n$ итераций коэффициент обнуляют, и очередное
направление берётся чисто градиентным, $p^{(mn)} = \nabla f\big(x^{(mn)}\big)$.

Причина та же, что у сброса $A_k = E$ в квазиньютоновских методах. Для
квадратичной функции $n$ направлений $p^{(0)}, \dots, p^{(n-1)}$ взаимно
сопряжены относительно гессиана, и метод приходит в точный минимум не более
чем за $n$ шагов — дальше строить сопряжённые направления просто не из чего.
Для неквадратичной $f$ гессиан меняется от точки к точке, накопленная
сопряжённость теряется, и продолжение рекуррентной формулы даёт направления,
всё хуже связанные с текущей кривизной. Рестарт каждые $n$ шагов обрывает
накопление ошибки и гарантирует, что метод не застрянет: направление спуска
заведомо корректно.

Записанная формула — выбор $\beta$ по \textbf{Флетчеру--Ривсу}; вторая строка
с суммами — та же дробь, расписанная покоординатно.
\end{idea}

\begin{idea}
Место метода среди прочих: производные второго порядка не нужны, хранить
матрицу $n \times n$ (как в квазиньютоновских методах) тоже не нужно — только
предыдущее направление. Поэтому метод сопряжённых градиентов — основной выбор
при очень больших $n$, где $H$ или $A_k$ просто не помещаются в память.
\end{idea}


\newpage
\section{Методы условной оптимизации}

Теперь задача решается не на всём пространстве, а на допустимом множестве:
\[
  \varphi_0(x) \longrightarrow \min_{x \in S} ,
  \qquad S \subset \R^n .
\]

\begin{idea}[title={Общая картина: зачем вообще нужны отдельные методы}]
Все методы предыдущих глав строили последовательность спуска, свободно двигаясь
по $\R^n$. Ограничение $x \in S$ ломает эту схему: шаг вдоль антиградиента может
вывести из $S$, а минимум $\varphi_0$ на $S$, как правило, лежит на границе,
где $\nabla\varphi_0 \ne 0$ — то есть привычный критерий остановки
$\norm{\nabla \varphi_0} < \eps$ не сработает никогда.

Отсюда три принципиально разных подхода, и вся оставшаяся часть курса — про них:
\begin{enumerate}[label=\arabic*.,left=0pt,itemsep=3pt]
  \item \textbf{Методы возможных направлений} — остаться внутри $S$ и на каждом
        шаге искать направление, которое \emph{одновременно} не выводит из $S$ и
        уменьшает $\varphi_0$. Представитель — метод Зойтендейка
        (\S\ref{subsec:zoutendijk}).
  \item \textbf{Методы отсечения} — заменить $S$ множеством, описываемым более
        простыми условиями (обычно многогранником), решить задачу на нём и
        отсечь часть, где оптимума заведомо нет; повторять.
  \item \textbf{Методы штрафов} — вообще отказаться от ограничений, встроив их
        в целевую функцию, и решать последовательность задач \emph{без}
        ограничений уже известными методами.
\end{enumerate}
\end{idea}

\subsection{Методы возможных направлений}

\begin{definition}\label{def:feasible-direction}
Пусть $x \in S$, $\sigma \in \R^n$. Вектор $\sigma$ называется
\textbf{возможным направлением} в точке $x$, если
\[
  \exists\, \delta_0 > 0 : \qquad x + \delta\sigma \in S
  \quad \forall \delta \in [0, \delta_0] .
\]
\end{definition}

\inputfig{figures/fig-feasible-direction}

\begin{idea}
Во внутренней точке $S$ возможным является любое направление — отсюда и
совпадение с безусловной задачей вдали от границы. Всё содержание определения
проявляется на границе: там конус возможных направлений сужается, и метод
обязан выбирать направление спуска \emph{внутри} этого конуса. Пересечение двух
условий — «возможное» ($\sigma$ не выводит из $S$) и «убывающее»
($\nabla\T\varphi_0(x)\sigma < 0$) — и есть то, что ищет метод Зойтендейка.
\end{idea}

\subsection{Методы отсечения}

Аппроксимируем $S$ множеством, описываемым более простыми условиями. Построим
вспомогательную задачу, которая позволит отсечь от этого множества множество
точек, в котором оптимума заведомо нет; потом вторую, и так далее.

\inputfig{figures/fig-cutting-plane}

\begin{idea}
Схема работы: исходное $S$ заменяют объемлющим многогранником $P_0 \supset S$,
на котором задача решается просто (например, линейным программированием). Если
найденный оптимум $x^{(0)}$ попал в $S$ — он и есть решение. Если нет, через
него проводят \textbf{отсекающую плоскость}: гиперплоскость, которая отрезает
$x^{(0)}$, но не отрезает ни одной точки $S$. Получается $P_1 \subset P_0$,
по-прежнему содержащее $S$, и всё повторяется. Множества $P_k$ сжимаются к $S$
в окрестности оптимума.
\end{idea}

\subsection{Методы штрафов}

\paragraph{Направление 1 (внешний штраф).}
Строим последовательность спусков в области, но неявно: записываем
вспомогательную задачу, содержащую \emph{штраф за выход} за границу области, и
отказываемся от решения задачи с ограничениями. При решении вспомогательной
задачи значение функции цели в области совпадает со значением в области
исходной задачи — значит можно решать задачу без ограничений.

\paragraph{Направление 2 (барьерные функции, внутренний штраф).}
Строим вспомогательную задачу так, чтобы \emph{не выпускать} точку за границу,
и решаем задачу без ограничений.

\inputfig{figures/fig-penalty-barrier}

\begin{idea}
Формально обе конструкции выглядят так:
\[
  \text{внешний штраф:}\quad
  \Phi_r(x) = \varphi_0(x) + r \sum_{i=1}^{m} \big(\max\{0, \varphi_i(x)\}\big)^2 ,
  \qquad r \to +\infty ;
\]
\[
  \text{барьер:}\quad
  \Phi_r(x) = \varphi_0(x) - r \sum_{i=1}^{m} \frac{1}{\varphi_i(x)} ,
  \qquad r \to +0 .
\]
Разница принципиальная:
\begin{itemize}[nosep,left=0pt]
  \item \emph{внешний} штраф равен нулю внутри $S$ и растёт снаружи, поэтому
        внутри области $\Phi_r \equiv \varphi_0$ — это и есть «значение функции
        цели совпадает». Приближения идут \emph{снаружи} $S$ и допустимыми
        становятся лишь в пределе;
  \item \emph{барьер} обращается в $+\infty$ при подходе к границе изнутри, то
        есть точка физически не может выйти из $S$. Зато он искажает $\varphi_0$
        и внутри области, и это искажение снимается лишь при $r \to 0$.
\end{itemize}
Платят за удобство обусловленностью: при $r \to \infty$ (или $r \to 0$)
вспомогательная функция становится всё более овражной, и безусловные методы на
ней работают всё хуже. Поэтому решают не одну задачу, а последовательность с
постепенно меняющимся $r$, беря решение предыдущей как начальное приближение
следующей.
\end{idea}

\subsection{Метод Зойтендейка}
\label{subsec:zoutendijk}

\[
  \varphi_0(x) \longrightarrow \min_{x \in S} ,
  \qquad
  S = \big\{\, x \in \R^n \ \big|\ \varphi_i(x) \le 0,\ i = \overline{1,m},
      \quad Ax = b \,\big\} ,
\]
где $A = A[l, n]$, $b = b[l]$, $l < n$.

\medskip
\noindent\textbf{Основная идея:} в точке $x \in S$ искать возможное направление
спуска (определение~\ref{def:feasible-direction}).

\subsubsection*{Предположения}

\begin{enumerate}[label=\arabic*),left=0pt]
  \item $\varphi_i(x)$ — выпуклые, $i = \overline{0,m}$;
  \item $\varphi_i(x)$ — непрерывно дифференцируемые, $i = \overline{0,m}$;
  \item $\norm{\nabla \varphi_i(x)} \le K$ \ $(\exists K > 0)$,
        $i = \overline{0,m}$;
  \item $\norm{\nabla\varphi_i(x) - \nabla\varphi_i(y)} \le R\norm{x - y}$,
        $R > 0$, $\forall x, y \in \R^n$;
  \item $S$ — компакт;
  \item $\exists\, \bar{x} \in S$: $\varphi_i(\bar{x}) < 0$, $i = \overline{1,m}$,
        $A\bar{x} = b$ — \textbf{условие Слейтера}.
\end{enumerate}

\subsubsection*{Активные и почти активные ограничения}

Пусть $M = \{1, 2, \dots, m\}$ и для $x \in S$
\[
  I_0(x) = \{\, i \in M \mid \varphi_i(x) = 0 \,\} ,
  \qquad
  I_1(x) = \{\, i \in M \mid \varphi_i(x) < 0 \,\} .
\]
Для $\delta > 0$ введём
\[
  I_\delta(x) = \{\, i \in M \mid -\delta \le \varphi_i(x) \le 0 \,\}
  \qquad \text{— почти активные ограничения.}
\]

\begin{idea}
Зачем нужен «зазор» $\delta$. Если учитывать только активные ограничения
$I_0(x)$, направление может оказаться возможным лишь формально: ограничение,
до которого остаётся ничтожно мало, тут же станет активным, и допустимая длина
шага окажется почти нулевой — метод заклинит у границы. Включая в рассмотрение
все ограничения, до которых меньше $\delta$, мы гарантируем шаг конечной длины.
Цена — направление получается более «осторожным»; поэтому $\delta$ по ходу
работы уменьшают (шаг~2$^\circ$ алгоритма ниже).
\end{idea}

\subsubsection*{Вспомогательная задача поиска направления}

Пусть $x \in S$, $\delta > 0$, $\xi_i > 0$, $i = \overline{0,m}$.
Вспомогательная задача:
\[
  \eta \longrightarrow \min_{\mathbb{H}} ,
\]
\[
  \mathbb{H} = \Big\{\, \eta \ \Big|\
    \nabla\T\varphi_i(x)\, s \le \xi_i\, \eta,\ i \in I_\delta(x) \cup \{0\} ;
    \quad A s = 0 ; \quad \norm{s}_\infty \le 1 \,\Big\} ,
\]
где $\norm{s}_\infty \le 1$ означает
$-1 \le s^{(i)} \le 1$, $i = \overline{1,n}$.

При $s = 0$, $\eta = 0$ все условия выполнены — пара допустима, поэтому решение
$s_\delta$, $\eta_\delta$ вспомогательной задачи удовлетворяет
\[
  \eta_\delta \le 0 .
\]

\begin{idea}
Переменная $\eta$ — это «насколько уверенно» направление $s$ одновременно
уменьшает цель и уводит от всех почти активных ограничений: неравенства
$\nabla\T\varphi_i(x)s \le \xi_i\eta$ при $\eta < 0$ означают, что по направлению
$s$ убывает и $\varphi_0$, и каждое почти активное $\varphi_i$. Минимизация
$\eta$ — поиск наилучшего такого направления.

Условие $As = 0$ держит точку на плоскости $Ax = b$, а $\norm{s}_\infty \le 1$
нормирует задачу: без него при $\eta < 0$ можно было бы бесконечно растягивать
$s$ и $\eta \to -\infty$. Коэффициенты $\xi_i > 0$ — веса, позволяющие
по-разному взвешивать цель и ограничения. Вся задача линейна по $(\eta, s)$,
то есть решается симплекс-методом из~\S1.

Случай $\eta_\delta = 0$ означает, что возможного направления спуска нет — это
и будет признаком оптимальности.
\end{idea}

\subsubsection*{Лемма о длине шага}

\begin{lemma}\label{lem:zoutendijk-step}
Если $\eta_\delta(x) < 0$, то существует $\alpha > 0$ такое, что
\begin{equation}\label{eq:zout-goal}
  \begin{cases}
    \varphi_0\big(x + \alpha s_\delta(x)\big)
      \le \varphi_0(x) + \tfrac12 \alpha\, \xi_0\, \eta_\delta(x) , \\[4pt]
    \varphi_i\big(x + \alpha s_\delta(x)\big) \le 0 , & i \in I_\delta(x) , \\[4pt]
    \varphi_i\big(x + \alpha s_\delta(x)\big) \le 0 , & i \notin I_\delta(x) .
  \end{cases}
\end{equation}
То есть $s_\delta(x)$ — возможное направление спуска.
\end{lemma}

\begin{proof}
По формуле конечных приращений
\[
  \varphi_0\big(x + \alpha s_\delta(x)\big)
  = \varphi_0(x) + \alpha\,\nabla\T\varphi_0(\tilde{x})\, s_\delta(x) ,
  \qquad
  \tilde{x} = x + \theta\,\alpha\, s_\delta(x) , \quad \theta \in [0,1] .
\]
Прибавим и вычтем $\alpha \nabla\T\varphi_0(x) s_\delta(x)$:
\[
  \varphi_0\big(x + \alpha s_\delta(x)\big)
  = \varphi_0(x) + \alpha\,\nabla\T\varphi_0(x)\, s_\delta(x)
    + \alpha\big(\nabla\varphi_0(\tilde{x}) - \nabla\varphi_0(x)\big)\T s_\delta(x) .
\]
Последнее слагаемое оцениваем неравенством Коши--Буняковского и условием~4):
\[
  \alpha\,\norm{\nabla\varphi_0(\tilde{x}) - \nabla\varphi_0(x)}\cdot
  \norm{s_\delta(x)}
  \le \alpha R \underbrace{\norm{\tilde{x} - x}}_{=\ \theta\alpha\norm{s_\delta(x)}}
      \norm{s_\delta(x)}
  \le \alpha^2 R\, \norm{s_\delta(x)}^2 .
\]
Наконец, по построению множества $\mathbb{H}$ выполнено
$\nabla\T\varphi_0(x)\, s_\delta(x) \le \xi_0\, \eta_\delta(x)$, откуда
\begin{equation}\label{eq:zout-2}
  \varphi_0\big(x + \alpha s_\delta(x)\big)
  \ \le\ \varphi_0(x) + \alpha\,\xi_0\,\eta_\delta(x)
         + \alpha^2 R\,\norm{s_\delta(x)}^2 .
\end{equation}

Для $i \in I_\delta(x)$ точно та же выкладка даёт
\begin{equation}\label{eq:zout-3}
  \varphi_i\big(x + \alpha s_\delta(x)\big)
  \ \le\ \varphi_i(x) + \alpha\,\xi_i\,\eta_\delta(x)
         + \alpha^2 R\,\norm{s_\delta(x)}^2 .
\end{equation}

Для $i \notin I_\delta(x)$, то есть $\varphi_i(x) < -\delta$, используем
условие~3):
\begin{equation}\label{eq:zout-4}
  \begin{aligned}
    \varphi_i\big(x + \alpha s_\delta(x)\big)
    &= \varphi_i(x) + \alpha\,\nabla\T\varphi_i(\tilde{x})\, s_\delta(x) \\
    &\le \varphi_i(x)
       + \alpha\,\norm{\nabla\varphi_i(\tilde{x})}\cdot\norm{s_\delta(x)}
     \ <\ -\delta + \alpha K \norm{s_\delta(x)} .
  \end{aligned}
\end{equation}

Теперь подберём $\alpha > 0$ так, чтобы выполнялось~\eqref{eq:zout-goal}.
\begin{itemize}[left=0pt]
  \item Из~\eqref{eq:zout-2} достаточно
        $\alpha\xi_0\eta_\delta(x) + \alpha^2 R\norm{s_\delta(x)}^2
         \le \tfrac12 \alpha \xi_0 \eta_\delta(x)$, то есть
        \begin{equation}\label{eq:zout-6}
          \alpha \ \le\ -\frac{1}{2}\cdot
          \frac{\xi_0\,\eta_\delta(x)}{R\,\norm{s_\delta(x)}^2} .
        \end{equation}
  \item Из~\eqref{eq:zout-3} при $i \in I_\delta(x)$ имеем $\varphi_i(x) \le 0$,
        поэтому достаточно
        $\alpha\xi_i\eta_\delta(x) + \alpha^2 R\norm{s_\delta(x)}^2 \le 0$, то есть
        \begin{equation}\label{eq:zout-7}
          \alpha \ \le\ -\frac{\xi_i\,\eta_\delta(x)}{R\,\norm{s_\delta(x)}^2} .
        \end{equation}
  \item Из~\eqref{eq:zout-4} при $i \notin I_\delta(x)$ достаточно
        $-\delta + \alpha K \norm{s_\delta(x)} \le 0$, то есть
        \begin{equation}\label{eq:zout-8}
          \alpha \ \le\ \frac{\delta}{K\,\norm{s_\delta(x)}} .
        \end{equation}
\end{itemize}
Так как $\eta_\delta(x) < 0$, правые части
\eqref{eq:zout-6}--\eqref{eq:zout-8} положительны, и годится любое
$\alpha > 0$, не превосходящее их минимума.

Равенства сохраняются при любом $\alpha$:
\[
  A\big(x + \alpha s_\delta(x)\big)
  = Ax + \alpha \underbrace{A s_\delta(x)}_{=\,0} = b .
  \qedhere
\]
\end{proof}

\begin{idea}
Три оценки имеют разный смысл. Условие~\eqref{eq:zout-6} требует, чтобы шаг был
достаточно коротким и квадратичная поправка $\alpha^2 R\norm{s_\delta}^2$ не
съела линейное убывание — остаётся хотя бы его половина. Условие
\eqref{eq:zout-7} не даёт выйти за почти активные ограничения, а
\eqref{eq:zout-8} — за «далёкие»: до них запас ровно $\delta$, а скорость
изменения $\varphi_i$ не больше $K\norm{s_\delta}$.

Именно здесь работает весь набор предположений: липшицевость градиентов~4) даёт
квадратичную поправку, ограниченность градиентов~3) — оценку
\eqref{eq:zout-8}, а $\eta_\delta < 0$ — положительность всех трёх границ.
\end{idea}

\begin{theorem}\label{thm:zoutendijk}
Если $x \ne x_*$, $x \in S$, то $\eta_\delta(x) < 0$ для достаточно малых
$\delta > 0$.
\end{theorem}

\begin{idea}
Содержательно: пока точка не оптимальна, возможное направление спуска
существует — надо лишь взять достаточно узкую полосу почти активных
ограничений. Вместе с леммой~\ref{lem:zoutendijk-step} это и даёт сходимость:
на каждом шаге либо есть направление с $\eta_k < 0$ и целевая функция
убывает на конечную величину, либо $\delta$ уменьшается. Условие Слейтера~6)
нужно именно здесь — без внутренней точки конус возможных направлений может
выродиться.
\end{idea}

\subsubsection*{Алгоритм}

Пусть заданы $\xi_0$ и $\{\xi_i\}_{i=1}^{m}$: $\xi_i > 0$ для всех
$i = \overline{0,m}$; $x_0$ — начальное приближение, $\delta_0 > 0$,
$\lambda \in (0,1)$.

\medskip
\noindent\textbf{Общий шаг в точке $x_k$.}

\begin{steps}
  \item \textbf{Поиск возможного направления спуска:}
        $\eta \to \min_{\mathbb{H}}$, где
        \[
          \mathbb{H} = \Big\{\, \eta \ \Big|\
            \nabla\T\varphi_0(x_k)\, s \le \xi_0 \eta ; \quad
            \nabla\T\varphi_i(x_k)\, s \le \xi_i \eta,\ i \in I_{\delta_k}(x_k) ;
        \]
        \[
            A s = 0 ; \quad \norm{s}_\infty \le 1 \,\Big\} .
        \]
        Решение обозначаем $\eta_{\delta_k} \equiv \eta_k$,
        $s_{\delta_k} \equiv s_k$.
  \item \textbf{Выбор шага, построение $x_{k+1}$, $\delta_{k+1}$:}
        \begin{itemize}[left=0pt,nosep]
          \item[$1^\circ$] $\eta_k < -\delta_k$ — далеко от границы, двигаемся:
                \[
                  x_{k+1} = x_k + \alpha_k s_k , \qquad \delta_{k+1} = \delta_k ,
                \]
                где $\alpha_k$ выбирается по принципу дробления;
          \item[$2^\circ$] $0 > \eta_k \ge -\delta_k$ — направление слишком
                слабое, сужаем полосу:
                \[
                  x_{k+1} = x_k , \qquad \delta_{k+1} = \lambda\,\delta_k .
                \]
        \end{itemize}
        В обоих случаях $k := k+1$ и возврат к шагу~1.
\end{steps}

\medskip
\noindent\textbf{Условие останова:}
\[
  \delta_k < \delta_{0k} ,
  \qquad
  \delta_{0k} = -\max_{i \notin I_0(x_k)} \varphi_i(x_k) ,
\]
и $\eta_k = 0$. Тогда $x_k$ — решение.

\begin{idea}
Величина $\delta_{0k}$ — расстояние (по значению функций) до ближайшего
\emph{неактивного} ограничения: для $i \notin I_0(x_k)$ все $\varphi_i(x_k) < 0$,
и $-\max$ по ним даёт наименьший запас. Неравенство $\delta_k < \delta_{0k}$
означает, что полоса почти активных ограничений сузилась настолько, что
$I_{\delta_k}(x_k) = I_0(x_k)$ — в рассмотрении остались только по-настоящему
активные ограничения. Если при этом $\eta_k = 0$, то возможного направления
спуска нет даже в точном смысле, и $x_k$ оптимальна.

Логика всего алгоритма: пункт $1^\circ$ продвигает точку, пункт $2^\circ$
уточняет модель. Каждый проход через $2^\circ$ уменьшает $\delta$ в $\lambda$
раз, так что либо точка движется с ощутимым убыванием $\varphi_0$, либо
$\delta_k \to 0$ и срабатывает условие останова.
\end{idea}

\subsubsection*{Как найти начальное приближение}

Решается вспомогательная задача
\[
  \eta \longrightarrow \min_{\eta \in \mathbb{H}} ,
  \qquad
  \mathbb{H} = \big\{\, \eta \ \big|\ \varphi_i(x) - \eta \le 0,\
    i = \overline{1,m} \ \ \&\ \ Ax = b \,\big\} .
\]
По условию Слейтера существует $\bar{x}$: $\varphi_i(\bar{x}) < 0$ и
$A\bar{x} = b$, поэтому $\min \eta < 0$. Берём приближение $x$, удовлетворяющее
$Ax = b$.

\begin{idea}
Это та же идея, что и у вспомогательной задачи симплекс-метода (\S4): вместо
поиска допустимой точки «наугад» минимизируется максимальная невязка
$\eta = \max_i \varphi_i(x)$. Отрицательное значение минимума означает, что
найденная точка строго внутри всех неравенств — ровно то, что нужно для старта
метода, работающего с почти активными ограничениями.
\end{idea}

\subsubsection*{Численный пример}

\begin{idea}[title={Откуда этот материал}]
На лекции метод Зойтендейка разбирался только теоретически. Ниже две итерации
на числах: видно, как вспомогательная задача выбирает направление вдоль
границы, пока ограничение активно, и переходит к антиградиенту, когда точка
уходит внутрь.
\end{idea}

\begin{example}\label{ex:zoutendijk-run}
\[
  \begin{cases}
    \varphi_0(x) = (x_1 - 2)^2 + (x_2 - 1)^2 \to \min , \\
    x_1 \ge 0 , \quad x_2 \ge 0 ,
  \end{cases}
\]
то есть в стандартной форме
\[
  \varphi_1(x) = -x_1 \le 0 , \qquad \varphi_2(x) = -x_2 \le 0 .
\]
Ограничений-равенств нет, поэтому условие $As = 0$ отсутствует. Берём
$\xi_0 = \xi_1 = \xi_2 = 1$, $\delta_0 = 0{,}1$. Градиенты:
\[
  \nabla\varphi_0(x) = \begin{pmatrix} 2(x_1-2) \\ 2(x_2-1) \end{pmatrix} ,
  \qquad
  \nabla\varphi_1 = \begin{pmatrix} -1 \\ 0 \end{pmatrix} ,
  \qquad
  \nabla\varphi_2 = \begin{pmatrix} 0 \\ -1 \end{pmatrix} .
\]
Старт: $x_0 = (0, 2)\T \in S$, $\varphi_0(x_0) = 4 + 1 = 5$.

\medskip
\noindent\textbf{Итерация 0.} Значения ограничений:
$\varphi_1(x_0) = 0$, $\varphi_2(x_0) = -2$. Первое активно, второе даже не
почти активно ($-2 < -\delta_0$), значит
$I_{\delta_0}(x_0) = \{1\}$. Градиент цели
$\nabla\varphi_0(x_0) = (-4, 2)\T$.

Вспомогательная задача:
\[
  \eta \to \min
  \qquad\text{при}\qquad
  \begin{cases}
    -4 s^{(1)} + 2 s^{(2)} \le \eta , \\
    -s^{(1)} \le \eta , \\
    \abs{s^{(1)}} \le 1 , \quad \abs{s^{(2)}} \le 1 .
  \end{cases}
\]
Из второго ограничения $\eta \ge -s^{(1)} \ge -1$, то есть $\eta$ не может быть
меньше $-1$. Значение $\eta_0 = -1$ достижимо: при $s^{(1)} = 1$ второе условие
обращается в равенство, а первое принимает вид $2s^{(2)} \le 3$ и выполнено при
любом допустимом $s^{(2)}$. Оптимум вырожден — годится целое ребро; берём
вершину
\[
  s_0 = \begin{pmatrix} 1 \\ -1 \end{pmatrix} ,
  \qquad \eta_0 = -1 < 0 ,
\]
дающую наибольшее убывание цели:
$\nabla\T\varphi_0(x_0)\,s_0 = -4 - 2 = -6 < 0$.

Шаг. Точка $x(\alpha) = (\alpha,\ 2-\alpha)\T$ допустима при
$0 \le \alpha \le 2$, и
\[
  \varphi_0\big(x(\alpha)\big) = (\alpha - 2)^2 + (1 - \alpha)^2
    = 2\alpha^2 - 6\alpha + 5 ,
  \qquad
  \varphi_0' = 4\alpha - 6 = 0 \implies \alpha_0 = \tfrac32 \in [0,2] .
\]
\[
  x_1 = \begin{pmatrix} 3/2 \\ 1/2 \end{pmatrix} ,
  \qquad
  \varphi_0(x_1) = \tfrac14 + \tfrac14 = \tfrac12
  \qquad \big(5 \to \tfrac12\big) .
\]

\medskip
\noindent\textbf{Итерация 1.} Теперь $\varphi_1(x_1) = -\tfrac32$,
$\varphi_2(x_1) = -\tfrac12$ — обе строго меньше $-\delta_0$, значит
$I_{\delta_0}(x_1) = \varnothing$: точка ушла внутрь допустимого множества.
Градиент $\nabla\varphi_0(x_1) = (-1, -1)\T$, вспомогательная задача
вырождается в
\[
  \eta \to \min
  \qquad\text{при}\qquad
  -s^{(1)} - s^{(2)} \le \eta , \quad \abs{s^{(i)}} \le 1 ,
\]
откуда $s_1 = (1, 1)\T$, $\eta_1 = -2$. Это ровно антиградиент, нормированный
условием $\norm{s}_\infty \le 1$: ограничения на направление не влияют, и метод
работает как в безусловной задаче.

Шаг. $x(\alpha) = \big(\tfrac32 + \alpha,\ \tfrac12 + \alpha\big)\T$ допустима
при всех $\alpha \ge 0$, и
\[
  \varphi_0\big(x(\alpha)\big)
   = \Big(\alpha - \tfrac12\Big)^2 + \Big(\alpha - \tfrac12\Big)^2
   = 2\Big(\alpha - \tfrac12\Big)^2
  \quad\Longrightarrow\quad \alpha_1 = \tfrac12 ,
\]
\[
  x_2 = \begin{pmatrix} 2 \\ 1 \end{pmatrix} = x_* ,
  \qquad \varphi_0(x_2) = 0
  \qquad \big(\tfrac12 \to 0\big) .
\]
Значения цели: $5 \to \tfrac12 \to 0$.
\end{example}

\begin{idea}
Что здесь существенно.

\begin{itemize}[nosep,left=0pt]
  \item На нулевой итерации чистый антиградиент
        $-\nabla\varphi_0(x_0) = (4,-2)\T$ вывел бы точку за границу
        $x_2 \ge 0$ только при большом шаге, но главное другое: в вершине
        $x_1 = 0$ он смотрит \emph{наружу} по первой координате. Условие
        $\nabla\T\varphi_1 s \le \xi_1 \eta < 0$ заставляет брать
        $s^{(1)} > 0$ — направление отталкивается от активного ограничения.
  \item Вырожденность оптимума вспомогательной задачи — типичная ситуация:
        $\eta$ упирается в самое «тугое» ограничение, а остальные координаты
        направления остаются недоопределёнными. Симплекс-метод вернёт
        какую-то вершину; все они дают $\eta_0 = -1$, но убывание цели вдоль
        них разное.
  \item Шаг здесь брался точной минимизацией, а не по оценкам
        \eqref{eq:zout-6}--\eqref{eq:zout-8} из леммы~\ref{lem:zoutendijk-step}:
        те оценки нужны для доказательства сходимости, на практике же
        достаточно убедиться, что точка осталась допустимой и цель убыла.
\end{itemize}
\end{idea}


\newpage
\lecture{Лекция 7}{}

\subsection{Матрица проектирования}

Рассматривается задача
\[
  \begin{cases}
    \varphi(x) \to \min , & x \in \R^n , \\
    Cx = d , & C \ \text{размера } l \times n,\ d \in \R^{l} , \\
    F(x) \le g , & F \ \text{размера } m \times n,\ g \in \R^{m} ,
  \end{cases}
\]
где $\varphi$ может быть нелинейной. К такой задаче применяется \textbf{метод
Розена}. Всюду далее $\varphi(x)$ — выпуклая и непрерывно дифференцируемая.

\begin{definition}\label{def:projector}
Матрица $P$ размера $n \times n$ называется \textbf{матрицей проектирования}
(проектором), если
\[
  1^\circ\ P = P\T ,
  \qquad\qquad
  2^\circ\ P \cdot P = P .
\]
\end{definition}

\begin{example}\label{ex:projector}
Пусть $A$ — матрица размера $m \times n$ полного строчного ранга. Тогда
\[
  P = \underbrace{A\T}_{n \times m}
      \underbrace{\big(A A\T\big)^{-1}}_{m \times m}
      \underbrace{A}_{m \times n}
  \qquad (n \times n)
\]
— матрица проектирования.
\end{example}

\begin{proof}
$1^\circ$ Транспонируем, пользуясь $(XYZ)\T = Z\T Y\T X\T$:
\[
  \Big(A\T \big(A A\T\big)^{-1} A\Big)\T
  = A\T \Big(\big(A A\T\big)^{-1}\Big)\T \big(A\T\big)\T
  = A\T \Big(\big(A A\T\big)\T\Big)^{-1} A
  = A\T \big(A A\T\big)^{-1} A = P .
\]
Здесь использовано, что матрица $A A\T$ симметрична,
$\big(A A\T\big)\T = \big(A\T\big)\T A\T = A A\T$, а обратная к симметричной сама
симметрична: $\big(X^{-1}\big)\T = \big(X\T\big)^{-1}$.

$2^\circ$
\[
  P \cdot P
  = A\T \big(A A\T\big)^{-1}
    \underbrace{A \cdot A\T \big(A A\T\big)^{-1}}_{=\ E} A
  = A\T \big(A A\T\big)^{-1} A = P .
  \qedhere
\]
\end{proof}

\begin{idea}
Обратимость $A A\T$ — это и есть требование полного строчного ранга $A$: строки
матрицы ограничений должны быть линейно независимы. Для метода Розена это
означает, что среди активных ограничений нет дублирующих друг друга — иначе
матрицу нужно предварительно чистить.

Геометрически $P = A\T(AA\T)^{-1}A$ проектирует на линейную оболочку строк $A$,
а дополнительный проектор $Q = E - P$ (см. лемму ниже) — на \emph{ядро} $A$,
то есть на множество направлений, вдоль которых $Ax$ не меняется. Именно $Q$
и понадобится в методе.
\end{idea}

\begin{lemma}\label{lem:projector}
Пусть $P$ — матрица проектирования размера $n \times n$. Тогда
\begin{enumerate}[label=$\arabic*^\circ$,left=0pt]
  \item $Q \equiv E - P$ — матрица проектирования;
  \item $x\T P x \ge 0$ \ $\forall x$ (то есть $P$ положительно полуопределена);
  \item множества
        \[
          L \equiv \{\, y \mid y = Px \,\} ,
          \qquad
          L' \equiv \{\, z \mid z = Qx \,\}
        \]
        — ортогональные подпространства, причём
        \[
          \forall x \ \ \exists!\ \langle y, z\rangle :
          \quad x = y + z , \quad y \in L , \quad z \in L' .
        \]
\end{enumerate}
\end{lemma}

\begin{proof}
\textbf{1.1)} $Q\T = E\T - P\T = E - P = Q$.

\medskip
\textbf{1.2)}
$Q \cdot Q = (E-P)(E-P) = E - P - P + P\cdot P = E - P = Q$.

\medskip
\textbf{2)} Пользуясь обоими свойствами проектора,
\[
  x\T P x = x\T P \cdot P x = x\T P\T P x = \norm{P x}^2 \ \ge\ 0 .
\]

\medskip
\textbf{3.1) Ортогональность.} Что $L$ и $L'$ — подпространства, очевидно
(образы линейных отображений). Пусть $y \in L$, $z \in L'$, то есть
$y = Px_1$, $z = Qx_2$. Тогда
\[
  y\T z = (P x_1)\T (Q x_2) = x_1\T P\T Q\, x_2
        = x_1\T P (E - P)\, x_2
        = x_1\T \big(P - P\cdot P\big) x_2 = 0 .
\]

\medskip
\textbf{3.2) Существование и единственность разложения.} Существование:
\[
  x = Px + (E - P)x = Px + Qx ,
\]
и первое слагаемое лежит в $L$, второе — в $L'$. Единственность: пусть
\[
  x = y + z , \quad y \in L,\ z \in L' ,
  \qquad
  x = y_1 + z_1 , \quad y_1 \in L,\ z_1 \in L' .
\]
Вычитая, получаем $0 = (y - y_1) + (z - z_1)$, где $y - y_1 \in L$,
$z - z_1 \in L'$. Два ортогональных вектора в сумме дают нуль только если оба
нулевые, значит $y = y_1$, $z = z_1$.
\end{proof}

\inputfig{figures/fig-projection-decomposition}

\begin{idea}
Пункт~2$^\circ$ — именно то, что понадобится для доказательства убывания:
\[
  \nabla\T\varphi\, P\, \nabla\varphi = \norm{P\nabla\varphi}^2 \ \ge\ 0 ,
\]
причём строго больше нуля, пока проекция градиента ненулевая.

Пункт~3$^\circ$ — геометрический смысл всей конструкции: пространство
распадается в прямую ортогональную сумму $\R^n = L \oplus L'$, и проектор $P$
просто отбрасывает составляющую вдоль $L'$. В методе Розена $L'$ — это ядро
матрицы активных ограничений, то есть множество направлений, остающихся на
грани допустимого множества.
\end{idea}

\subsection{Метод Розена}
\label{subsec:rosen}

\begin{idea}[title={Обозначения}]
В тетради буквой $S$ обозначено и допустимое множество, и строящееся
направление. Ниже допустимое множество — $S$, а направление — $s$.
\end{idea}

\subsubsection*{Активные и пассивные ограничения}

Пусть $x$ — допустимая точка. Разобьём неравенства на активные и пассивные:
\[
  \begin{cases}
    Cx = d , \\
    F_1 x = g_1 & \text{— активные,} \\
    F_2 x < g_2 & \text{— пассивные,}
  \end{cases}
  \qquad
  F = \begin{pmatrix} F_1 \\ \hline F_2 \end{pmatrix} ,
  \qquad
  g = \begin{pmatrix} g_1 \\ \hline g_2 \end{pmatrix} .
\]
Объединим равенства и активные неравенства в одну систему:
\[
  A(x) = \begin{pmatrix} C \\ \hline F_1 \end{pmatrix} ,
  \qquad
  b = \begin{pmatrix} d \\ \hline g_1 \end{pmatrix} ,
  \qquad\text{так что}\qquad
  \begin{cases}
    A x = b , \\
    F_2 x < g_2
  \end{cases}
  \qquad \forall x \in S .
  \tag{1}\label{eq:rosen-1}
\]

Составим матрицу проектирования
\[
  P = E - A\T \big(A A\T\big)^{-1} A ,
  \qquad
  s = -P\, \nabla\varphi(x) .
\]

\begin{idea}
Это проектор $Q = E - \big(A\T(AA\T)^{-1}A\big)$ из леммы~\ref{lem:projector},
то есть проектирование на ядро $A$. Направление $s$ — \textbf{проекция
антиградиента} на подпространство, вдоль которого активные ограничения
сохраняются как равенства. Отсюда название: метод проекции градиента.
\end{idea}

\subsubsection*{Допустимость шага}

Рассмотрим точку $x + \alpha s$:
\[
  A(x + \alpha s) = \underbrace{Ax}_{=\,b} + \alpha A s
  = b + \alpha A\big(-P \nabla\varphi(x)\big) = b ,
  \tag{2}\label{eq:rosen-2}
\]
поскольку
\[
  A P = A\Big(E - A\T \big(A A\T\big)^{-1} A\Big)
      = A - \underbrace{A A\T \big(A A\T\big)^{-1}}_{=\ E} A
      = A - A = 0 .
\]
То есть активные ограничения не накладывают условия на выбор шага: при любом
$\alpha$ точка остаётся на той же грани.

Пассивные ограничения $F_2 x < g_2$ дают условие на шаг:
\[
  F_2(x + \alpha s) = F_2 x - \alpha\, F_2\, P \nabla\varphi(x) < g_2 .
  \tag{3}\label{eq:rosen-3}
\]

\begin{idea}
Разница принципиальная. Активные ограничения выполняются \emph{тождественно} по
$\alpha$ — их и «встроили» в направление через проектор. Пассивные выполнены
строго, значит нарушатся лишь при достаточно большом шаге; именно
условие~\eqref{eq:rosen-3} и ограничивает $\alpha$ сверху. Поэтому в алгоритме
шаг дробят до выполнения двух вещей: убывания $\varphi$ и
неравенства~\eqref{eq:rosen-3}.
\end{idea}

\subsubsection*{$s$ — направление спуска}

\paragraph{а) $s \ne 0$.}
\[
  \varphi(x + \alpha s)
  = \varphi(x) + \alpha\,\nabla\T\varphi(x)\, s + \dots
  = \varphi(x) - \alpha \underbrace{\nabla\T\varphi(x)\, P\, \nabla\varphi(x)}_{\ge\ 0}
    + \dots \ \le\ \varphi(x) ,
\]
где неотрицательность подчёркнутого — пункт~2$^\circ$
леммы~\ref{lem:projector}: $\nabla\T\varphi\, P\, \nabla\varphi
= \norm{P \nabla\varphi}^2 = \norm{s}^2 > 0$ при $s \ne 0$.

\paragraph{б) $s = 0$.}
Тогда
\[
  0 = -\Big(E - A\T \big(A A\T\big)^{-1} A\Big)\nabla\varphi(x)
    = -\Big(\nabla\varphi(x) - A\T \underbrace{\big(A A\T\big)^{-1} A \nabla\varphi(x)}_{=\ -\omega}\Big) ,
\]
то есть
\[
  \nabla\varphi(x) + A\T \omega = 0 ,
  \qquad
  \omega = -\big(A A\T\big)^{-1} A \nabla\varphi(x) .
  \tag{4}\label{eq:rosen-4}
\]
Размерность $\omega$ совпадает с числом строк $A$, то есть с числом равенств
плюс число активных ограничений в точке $x$. Расписывая по блокам
\[
  A = \begin{pmatrix} C \\ \hline F_1 \end{pmatrix} ,
  \qquad
  \omega = \begin{pmatrix} v \\ \hline u \end{pmatrix}
  \qquad\Longrightarrow\qquad
  \nabla\varphi(x) + C\T v + F_1\T u = 0 .
  \tag{5}\label{eq:rosen-5}
\]

\begin{idea}[title={Что такое \eqref{eq:rosen-5}: это условия ККТ}]
Равенство~\eqref{eq:rosen-5} — в точности \textbf{условие стационарности ККТ}
(Каруша--Куна--Таккера) для исходной задачи:
\[
  \nabla\varphi(x) + \sum_{\text{равенства}} v_i \nabla(c_i\T x)
   + \sum_{i \in I_0(x)} u_i \nabla(f_i\T x) = 0 ,
\]
где $v$ — множители при равенствах $Cx = d$ (знак произвольный), $u$ — множители
при активных неравенствах $F_1 x = g_1$ (обязаны быть неотрицательными).
Пассивные ограничения входят с нулевыми множителями — это дополняющая
нежёсткость.

\medskip
\textbf{Почему у $v$ знак свободный.} Равенство $\varphi_i(x) = 0$
эквивалентно паре неравенств
\[
  \begin{cases} \varphi_i(x) \le 0, \\ -\varphi_i(x) \le 0 , \end{cases}
\]
каждое со своим неотрицательным множителем; их разность и даёт один множитель
произвольного знака.

\medskip
\textbf{Про условие Слейтера.} В общей выпуклой задаче условия ККТ необходимы
лишь при каком-нибудь условии регулярности — например при условии Слейтера
(существует точка со строгими неравенствами). Без него ККТ могут не
выполняться: на лекции это иллюстрировалось примером
\[
  \varphi(x) \to \min , \qquad \sum_{i \in I} a_i\T x \le b^{(i)} .
\]
Здесь, однако, все ограничения \emph{линейны}, а линейность сама по себе —
условие регулярности. Поэтому для задачи метода Розена условия ККТ необходимы и
достаточны (в силу выпуклости $\varphi$) без каких-либо дополнительных
требований, и проверка сводится ровно к знаку $u$.
\end{idea}

Отсюда разбор случая $s = 0$:
\begin{enumerate}[label=\arabic*.,left=0pt]
  \item $u \ge 0$ \ $\Longrightarrow$ \ $x$ — оптимальная точка (выполнены ККТ);
  \item $\exists\, u^{(j)} < 0$ \ $\Longrightarrow$ \ перестроим направление $s$.
\end{enumerate}

\inputfig{figures/fig-rosen-step}

\subsubsection*{Перестроение направления при $u^{(j)} < 0$}

Пусть $u^{(j)} < 0$. Уберём из $F_1$ $j$-ю строку $f_j$ и обозначим
\[
  \bar{F}_1 = F_1 \ \text{без строки } f_j ,
  \qquad
  \bar{A} = \begin{pmatrix} C \\ \hline \bar{F}_1 \end{pmatrix} ,
  \qquad
  \bar{b} = \begin{pmatrix} d \\ \hline \bar{g}_1 \end{pmatrix} ,
\]
\[
  \bar{P} = E - \bar{A}\T \big(\bar{A}\,\bar{A}\T\big)^{-1} \bar{A} ,
  \qquad
  \bar{s} = -\bar{P}\, \nabla\varphi(x) .
\]

\begin{idea}
Смысл: множитель $u^{(j)} < 0$ означает, что $j$-е активное ограничение
«мешает» — если его отпустить, целевую функцию можно уменьшить. Поэтому строку
$f_j$ выбрасывают из набора активных, проектируют антиградиент на более широкое
подпространство и получают направление $\bar{s}$, уводящее \emph{внутрь} от
$j$-й грани.
\end{idea}

\paragraph{$\bar{s} \ne 0$.}
Если бы было $\bar{s} = 0$, то по той же выкладке, что
дала~\eqref{eq:rosen-4},
\[
  \nabla\varphi(x) + \bar{A}\T \bar\omega = 0 ,
  \qquad
  \bar\omega = -\big(\bar{A}\,\bar{A}\T\big)^{-1} \bar{A} \nabla\varphi(x) .
\]
С другой стороны, уберём из~\eqref{eq:rosen-4} вклад $j$-й строки: обозначив
через $\tilde\omega$ вектор $\omega$ без $j$-й компоненты,
\[
  \nabla\varphi(x) + \bar{A}\T \tilde\omega + u^{(j)} f_j\T = 0 .
  \tag{6}\label{eq:rosen-6}
\]
Вычитая, получаем
\[
  \bar{A}\T\big(\tilde\omega - \bar\omega\big) + u^{(j)} f_j\T = 0 ,
  \qquad u^{(j)} \ne 0 ,
\]
то есть строка $f_j$ линейно выражается через строки $\bar{A}$. Это
противоречит линейной независимости строк $A$ (без неё не существовало бы
$\big(A A\T\big)^{-1}$). Значит $\bar{s} \ne 0$.

\paragraph{Допустимость направления $\bar{s}$.}
Так как
\[
  \bar{A}\,\bar{P}
  = \bar{A}\Big(E - \bar{A}\T\big(\bar{A}\bar{A}\T\big)^{-1}\bar{A}\Big)
  = \bar{A} - \bar{A}\,\bar{A}\T\big(\bar{A}\bar{A}\T\big)^{-1}\bar{A} = 0 ,
\]
получаем
\[
  \bar{A}\,\bar{s}
  = \begin{pmatrix} C \\ \hline \bar{F}_1 \end{pmatrix} \bar{s}
  = -\bar{A}\,\bar{P}\,\nabla\varphi = 0
  \qquad\Longrightarrow\qquad
  \begin{cases}
    C \bar{s} = 0 , \\
    \bar{F}_1 \bar{s} = 0 .
  \end{cases}
  \tag{7--8}\label{eq:rosen-78}
\]

Проверим, что точка $x + \alpha\bar{s}$ остаётся допустимой:
\[
  \begin{aligned}
    &C(x + \alpha\bar{s}) = d , &&(*) \\
    &F_1(x + \alpha\bar{s}) \le g_1 , &&(**) \\
    &F_2(x + \alpha\bar{s}) < g_2 . &&(\!*\!*\!*)
  \end{aligned}
\]
Условие $(*)$ выполнено в силу~\eqref{eq:rosen-78}. Для $(**)$:
\[
  F_1(x + \alpha\bar{s})
  = \underbrace{F_1 x}_{=\ g_1} + \alpha\, F_1 \bar{s} \ \le\ g_1
  \qquad\Longleftrightarrow\qquad
  F_1 \bar{s} \le 0 .
\]
Строки $\bar{F}_1$ дают ноль по~\eqref{eq:rosen-78}, так что остаётся проверить
единственную выброшенную строку: нужно $f_j \bar{s} \le 0$. Условие
$(\!*\!*\!*)$ — как и раньше, ограничение на длину шага.

\paragraph{Доказательство $f_j \bar{s} \le 0$.}
Умножим~\eqref{eq:rosen-6} слева на $\bar{s}\T = -\nabla\T\varphi\,\bar{P}$:
\[
  \underbrace{\bar{s}\T \nabla\varphi(x)}_{=\ -\nabla\T\varphi\,\bar{P}\,\nabla\varphi}
  + \underbrace{\bar{s}\T \bar{A}\T}_{=\ (\bar{A}\bar{s})\T\ =\ 0} \tilde\omega
  + u^{(j)}\, \bar{s}\T f_j\T = 0 ,
\]
откуда
\[
  u^{(j)}\,\big(f_j \cdot \bar{s}\big)
  = \underbrace{\nabla\T\varphi\,\bar{P}\,\nabla\varphi}_{\ge\ 0} ,
  \qquad u^{(j)} < 0
  \qquad\Longrightarrow\qquad
  f_j \cdot \bar{s} \le 0 .
\]
Итак, $\bar{s}$ — допустимое направление спуска.

\subsubsection*{Алгоритм}

Пусть $x_0$ — начальное приближение. $k$-й шаг алгоритма:

\begin{steps}
  \item Составляем матрицу активных ограничений
        $A_k = \begin{pmatrix} C \\ \hline F_k \end{pmatrix}$, где $F_k$ —
        строки $F$, активные в точке $x_k$.
  \item Вычисляем
        \[
          P_k = E - A_k\T \big(A_k A_k\T\big)^{-1} A_k ,
          \qquad
          s_k = -P_k \nabla\varphi(x_k) .
        \]
        \begin{itemize}[left=0pt,nosep]
          \item[а)] $s_k \ne 0$: $\ x_{k+1} = x_k + \alpha_k s_k$, где
                $\alpha_k$ выбирается по принципу дробления,
                $\alpha_k = \alpha_0 \lambda^{\,i_k}$ ($\lambda$ — параметр
                дробления, $\alpha_0$ — начальное значение), до выполнения
                \emph{двух} условий:
                \[
                  \text{а.1)}\ \varphi(x_k + \alpha_k s_k) \le \varphi(x_k)
                  \quad\text{— уменьшение;}
                \]
                \[
                  \text{а.2)}\ F_2(x_k + \alpha_k s_k) \le g_2
                  \quad\text{— за пассивные ограничения не вышли.}
                \]
          \item[б)] $s_k = 0$: вычисляем $\omega_k = (v_k; u_k)$
                по~\eqref{eq:rosen-4}:
                \[
                  \text{б.1)}\ u_k \ge 0 \implies x_k \ \text{— оптимальная точка;}
                \]
                \[
                  \text{б.2)}\ \exists\, u^{(j)}_k < 0 \implies
                  \text{строим } \bar{F}_k,\ \bar{P}_k,\ \bar{s}_k
                  \ \text{и продолжаем.}
                \]
        \end{itemize}
\end{steps}

\begin{idea}
Структура метода: пока проекция антиградиента ненулевая — двигаемся вдоль
текущей грани, дробя шаг так, чтобы не выскочить за пассивные ограничения.
Когда проекция обнулилась, движение вдоль грани исчерпано, и вопрос лишь в том,
оптимальна ли точка: ответ даёт знак множителей $u$. Отрицательный множитель —
сигнал отпустить соответствующее ограничение, то есть перейти на грань большей
размерности. За конечное число таких переходов набор активных ограничений
стабилизируется, и метод либо движется, либо останавливается по $u \ge 0$.
\end{idea}


\newpage
\lecture{Лекция 8}{}

\subsection{Метод отсекающей гиперплоскости}
\label{subsec:cutting-plane}

\begin{idea}[title={Обозначения}]
В этом параграфе $\Omega$ — допустимое множество \emph{исходной} задачи,
а $S_k$ — аппроксимирующие его многогранники. В тетради оба обозначены одной
буквой, что мешает читать выкладки.
\end{idea}

Рассматривается задача
\[
  \begin{cases}
    \varphi_0(x) \to \min , & x \in \Omega , \\
    \Omega = \{\, x \in \R^n \mid \varphi(x) \le 0 \,\} ,
  \end{cases}
\]
где $\varphi$ — выпуклая дифференцируемая, а целевая функция линейна:
$\varphi_0(x) = c\T x$, $c \in \R^n$.

\subsubsection*{Приведение задачи к этому виду}

\begin{enumerate}[label=\arabic*),left=0pt]
  \item Если допустимое множество задано системой
        $\Omega = \{\, x \mid \varphi_i(x) \le 0,\ i = \overline{1,m} \,\}$,
        то достаточно взять
        \[
          \varphi(x) = \max_{i = \overline{1,m}} \varphi_i(x)
        \]
        — одно ограничение вместо $m$.
  \item Если $\varphi_0(x)$ не является линейной (но выпукла), вводим
        вспомогательную переменную $x^{(n+1)}$ и задачу
        \[
          \begin{cases}
            x^{(n+1)} \to \min , \\
            \varphi(x) \le 0 , \\
            \varphi_0(x) - x^{(n+1)} \le 0 .
          \end{cases}
        \]
  \item Объединяя 1) и 2) — сначала реализуется 2), потом 1):
        \[
          \begin{cases}
            x^{(n+1)} \to \min , \\
            \varphi_i\big(x, x^{(n+1)}\big) \le 0 , \quad i = \overline{1,m} , \\
            \varphi_{m+1}\big(x, x^{(n+1)}\big) = \varphi_0(x) - x^{(n+1)} ,
          \end{cases}
          \qquad
          \varphi(x) = \max_{i} \varphi_i\big(x, x^{(n+1)}\big) ,
        \]
        то есть
        \[
          \begin{cases}
            x^{(n+1)} \to \min , \\
            \varphi\big(x, x^{(n+1)}\big) \le 0 .
          \end{cases}
        \]
\end{enumerate}

\begin{idea}
Зачем эти преобразования. Весь метод построен на том, что на каждой итерации
решается \emph{линейная} задача на многограннике — иначе никакого выигрыша от
аппроксимации нет. Пункт~2) — стандартный приём «эпиграфной» переформулировки:
минимизировать выпуклую $\varphi_0$ — то же самое, что минимизировать новую
координату $x^{(n+1)}$ при условии $\varphi_0(x) \le x^{(n+1)}$. Пункт~1)
сворачивает систему неравенств в одно, пользуясь тем, что максимум конечного
числа выпуклых функций выпуклен. После обоих шагов задача имеет ровно тот вид,
для которого описан алгоритм.
\end{idea}

\subsubsection*{Алгоритм}

Считаем, что $\Omega \ne \varnothing$ и $\Omega$ — компакт.

\begin{algo}[title={Метод отсекающей гиперплоскости}]
\textbf{I. Начальный этап.} Построим объемлющий многогранник
\[
  S = \big\{\, x \mid a_i\T x - b^{(i)} \le 0,\ i = \overline{-1,0} \,\big\} ,
  \qquad S \ \text{— компакт}, \quad \Omega \subset S .
\]
Положим $S_0 = S$, $k = 0$, решим вспомогательную задачу
$\varphi_0 \to \min$, $x \in S_0$; пусть $x_0$ — её решение.

\medskip
\textbf{II. Основной этап ($k$-й шаг).} Пусть $S_k$ построено; решаем
$\varphi_0(x) \to \min$, $x \in S_k$, и получаем $x_k$.

Строим \textbf{субградиент} $a_{k+1}$ функции $\varphi$ в точке $x_k$:
\[
  \varphi(x) \ \ge\ \varphi(x_k) + a_{k+1}\T (x - x_k) \qquad \forall x .
\]
Обозначив $b^{(k+1)} = -\varphi(x_k) + a_{k+1}\T x_k$, получаем
\[
  \varphi(x) \ \ge\ a_{k+1}\T x - b^{(k+1)} ,
\]
и определяем
\[
  S_{k+1} = S_k \cap \big\{\, a_{k+1}\T x - b^{(k+1)} \le 0 \,\big\}
          = \big\{\, x \in S_k \ \big|\ a_{k+1}\T x - b^{(k+1)} \le 0 \,\big\} .
\]

\medskip
\textbf{III. Итерационный переход:} $k \to k+1$.

\medskip
\textbf{IV. Условие останова:} $\norm{x_{k+1} - x_k} < \eps$.
\end{algo}

\inputfig{figures/fig-cutting-subgradient}

\begin{idea}
Почему отсечение корректно, то есть не задевает $\Omega$. Для любой допустимой
точки $x \in \Omega$ выполнено $\varphi(x) \le 0$, а неравенство субградиента
даёт $a_{k+1}\T x - b^{(k+1)} \le \varphi(x) \le 0$ — новое ограничение
выполняется автоматически. Текущая же точка $x_k$ отсекается: для неё
$a_{k+1}\T x_k - b^{(k+1)} = \varphi(x_k) > 0$, если $x_k \notin \Omega$. То
есть каждая итерация выбрасывает заведомо недопустимую точку, не теряя ни одной
допустимой. Это и доказывается ниже в лемме~\ref{lem:cut-nested}.
\end{idea}

\subsubsection*{Особенности алгоритма}

\paragraph{I. Линейная подзадача.}
На каждом итерационном шаге решается
\[
  \begin{cases}
    c\T x \to \min , \\
    a_i\T x - b^{(i)} \le 0 , \quad i = \overline{-1, k} ,
  \end{cases}
  \tag{$*$}\label{eq:cut-lp}
\]
то есть от шага к шагу меняется \emph{только число ограничений} — значит
применим симплекс-метод (\S4).

\paragraph{II. Двойственная задача.}
Задача~\eqref{eq:cut-lp} равносильна
\[
  \begin{cases}
    c\T x \to \min , \\
    -a_i\T x \ge -b^{(i)} , \quad i = \overline{-1,k} ,
  \end{cases}
\]
двойственная к которой —
\[
  \begin{cases}
    \displaystyle -\sum_{i = -1}^{k} b^{(i)} y^{(i)} \ \to\ \max
      \qquad\text{то есть}\qquad -b\T y \to \max , \\[8pt]
    \displaystyle -\sum_{i = -1}^{k} y^{(i)} a_i = c ,
    \qquad y^{(i)} \ge 0 , \quad i = \overline{-1,k} .
  \end{cases}
\]

\begin{idea}
Здесь и прячется главный вычислительный выигрыш. Добавление одного ограничения
к прямой задаче — это добавление одной \emph{переменной} к двойственной.
Поэтому предыдущее двойственное решение, дополненное нулевой компонентой,
остаётся допустимым для новой двойственной задачи и служит готовым начальным
приближением: симплекс-метод продолжает работу с того места, где остановился,
а не стартует заново. Это и есть смысл фразы «в качестве начального приближения
даёт $y$ текущей итерации, дополненное нулевой компонентой».
\end{idea}

\paragraph{III. Построение субградиента.}
Для построения $S_{k+1}$ ищется субградиент $\varphi(x)$ в точке $x_k$:
\[
  \varphi(x) \ \ge\ \varphi(x_k) + a_{k+1}\T (x - x_k) .
\]
Возможны два случая.

\begin{enumerate}[label=\arabic*),left=0pt]
  \item Задача пришла в исходной постановке с \emph{дифференцируемой}
        $\varphi(x)$ — тогда градиент совпадает с субградиентом,
        $a_{k+1} = \nabla\varphi(x_k)$.
  \item Задача сведена к максимуму,
        $\varphi(x) = \max\limits_{i=\overline{1,m}} \varphi_i(x)$. Обозначим
        \[
          I(x_k) = \Big\{\, i \ \Big|\ \varphi_i(x_k)
            = \max_{i=\overline{1,m}} \varphi_i(x_k) \,\Big\}
        \]
        — набор индексов, на которых максимум достигается. По этим ограничениям
        строится субградиент (все $\varphi_i$ дифференцируемы):
        \begin{equation}\label{eq:subgrad-max}
          a_{k+1} = \sum_{i \in I(x_k)} \lambda_i \nabla\varphi_i(x_k) ,
          \qquad
          \lambda_i \ge 0 \ \ \forall i \in I(x_k) ,
          \qquad
          \sum_{i \in I(x_k)} \lambda_i = 1 .
        \end{equation}
\end{enumerate}

\begin{lemma}\label{lem:subgrad-max}
Пусть $\varphi_i$, $i = \overline{1,m}$, выпуклы и дифференцируемы,
$\varphi(x) = \max_i \varphi_i(x)$, и
$I(x) = \{\, i \mid \varphi_i(x) = \varphi(x) \,\}$. Тогда всякий вектор вида
\[
  a = \sum_{i \in I(x)} \lambda_i \nabla\varphi_i(x) ,
  \qquad \lambda_i \ge 0 , \qquad \sum_{i \in I(x)} \lambda_i = 1 ,
\]
является субградиентом $\varphi$ в точке $x$.
\end{lemma}

\begin{idea}[title={Доказательство (на лекции не приводилось)}]
\begin{proof}
Возьмём произвольную точку $z \in \R^n$ и любой индекс $j \in I(x)$. Так как
$\varphi$ — максимум, а $\varphi_j$ выпукла и дифференцируема,
\[
  \varphi(z) \ \ge\ \varphi_j(z)
  \ \ge\ \varphi_j(x) + \nabla\T\varphi_j(x)\,(z - x) .
\]
Но $j \in I(x)$ означает $\varphi_j(x) = \varphi(x)$, поэтому
\[
  \varphi(z) \ \ge\ \varphi(x) + \nabla\T\varphi_j(x)\,(z - x)
  \qquad \forall j \in I(x) .
\]
Умножим неравенство с номером $j$ на $\lambda_j \ge 0$ и просуммируем по
$j \in I(x)$, пользуясь $\sum_j \lambda_j = 1$:
\[
  \varphi(z) = \sum_{j \in I(x)} \lambda_j\, \varphi(z)
  \ \ge\ \sum_{j \in I(x)} \lambda_j
         \Big( \varphi(x) + \nabla\T\varphi_j(x)(z-x) \Big)
  = \varphi(x) + a\T (z - x) .
\]
Это и есть определение субградиента. \qedhere
\end{proof}

\medskip
Обратное утверждение (всякий субградиент максимума имеет вид~\eqref{eq:subgrad-max},
то есть $\partial\varphi(x) = \operatorname{conv}\{\nabla\varphi_i(x)\}_{i \in I(x)}$)
тоже верно, но его доказательство существенно длиннее и на лекции не нужно: для
алгоритма достаточно уметь \emph{предъявить} хотя бы один субградиент, а проще
всего взять $\lambda_{i_0} = 1$ для какого-нибудь одного $i_0 \in I(x_k)$, то
есть $a_{k+1} = \nabla\varphi_{i_0}(x_k)$.
\end{idea}

\subsubsection*{Сходимость}

\begin{lemma}\label{lem:cut-nested}
$\Omega \subset S_k$ для всех $k \ge 1$.
\end{lemma}

\begin{proof}
Возьмём произвольный элемент $x \in \Omega$ и покажем, что он лежит в $S_k$.
По определению $x \in \Omega \iff \varphi(x) \le 0$. С другой стороны, по
неравенству субградиента для каждого $j = \overline{1,k}$
\[
  \underbrace{\varphi(x)}_{\le\ 0}
  \ \ge\ \varphi(x_{j-1}) + a_j\T (x - x_{j-1})
  = a_j\T x - b^{(j)} ,
\]
откуда $a_j\T x - b^{(j)} \le 0$ для всех $j = \overline{1,k}$, то есть
$x \in S_k$.
\end{proof}

\begin{remark}
Целевая функция вдоль итераций \emph{не убывает}:
\[
  \varphi_0(x_0) \le \varphi_0(x_1) \le \dots \le \varphi_0(x_k) \le \dots ,
\]
поскольку $S_{k+1} \subset S_k$, а минимум по меньшему множеству не меньше. При
этом по лемме~\ref{lem:cut-nested} $\Omega \subset S_k$, значит
\[
  \varphi_0(x_*) \ \ge\ \dots \ \ge\ \varphi_0(x_k) \ \ge\ \dots ,
\]
где $x_*$ — оптимальное решение исходной задачи. То есть $\varphi_0(x_k)$ —
возрастающая последовательность \emph{нижних} оценок оптимума.
\end{remark}

\begin{theorem}\label{thm:cutting-plane}
Пусть $\varphi(x)$ непрерывна и выпукла, $\Omega$ — компакт, исходная
аппроксимация $S$ — тоже компакт, и субградиенты ограничены:
\[
  \norm{a(x)} \le R , \quad x \in S , \qquad \exists R > 0 .
\]
Тогда любая предельная точка последовательности $\{x_k\}$ является оптимальной,
причём $\varphi(x_k) \to 0$.
\end{theorem}

\begin{proof}
\textbf{1) Предельные точки существуют.} $S_0 = S$ и $S_{k+1} \subset S_k$,
значит вся последовательность $\{x_k\}$ лежит в компакте $S$, и у неё есть
сходящиеся подпоследовательности.

\medskip
\textbf{2) Если ограничение выполнено на каком-то шаге.} Пусть
$\exists k \in \NN$: $\varphi(x_k) \le 0$. Тогда $x_k \in \Omega$, и потому
$\varphi_0(x_k) \ge \varphi_0(x_*)$. С другой стороны, по
лемме~\ref{lem:cut-nested} $\Omega \subset S_k$, а $x_k$ минимизирует
$\varphi_0$ на $S_k$, поэтому
\[
  \varphi_0(x_k) = \min_{x \in S_k} \varphi_0(x)
  \ \le\ \min_{x \in \Omega} \varphi_0(x) = \varphi_0(x_*) .
\]
Следовательно $\varphi_0(x_k) = \varphi_0(x_*)$, то есть $x_k$ — оптимальное
решение, и алгоритм закончил работу.

\medskip
\textbf{3) Общий случай: $\varphi(x_k) > 0$ для всех $k$.} Покажем, что тогда
существует подпоследовательность с $\varphi(x_{k_j}) \to 0$.

От противного: пусть существуют $\beta > 0$ и бесконечное множество индексов
$J$ такие, что $\varphi(x_k) \ge \beta$ при $k \in J$. Последовательность
$\{x_k\}_{k \in J}$ лежит в компакте $S$, значит содержит фундаментальную
подпоследовательность: для любого $\delta > 0$ найдутся индексы $p < q$ из $J$ с
$\norm{x_p - x_q} \le \delta$.

Возьмём $p < q$. Так как $x_q \in S_q \subset S_{p+1}$, точка $x_q$
удовлетворяет $(p{+}1)$-му отсечению:
\[
  0 \ \ge\ a_{p+1}\T x_q - b^{(p+1)}
  = a_{p+1}\T x_q + \varphi(x_p) - a_{p+1}\T x_p
  = \varphi(x_p) + a_{p+1}\T (x_q - x_p) ,
\]
откуда
\[
  \varphi(x_p) \ \le\ a_{p+1}\T (x_p - x_q)
  \ \le\ \norm{a_{p+1}} \cdot \norm{x_p - x_q}
  \ \le\ R\,\norm{x_p - x_q} .
\]
Выберем теперь $p < q$ из $J$ так, чтобы $\norm{x_p - x_q} \le \beta/(2R)$.
Получаем $\varphi(x_p) \le \beta/2 < \beta$ — противоречие с $p \in J$.

Значит предположение неверно, и существует подпоследовательность $\{x_{k_j}\}$,
вдоль которой $\varphi(x_{k_j}) \to 0$.

\medskip
\textbf{4) Предельная точка оптимальна.} Переходя при необходимости к
подпоследовательности, считаем $x_{k_j} \to \bar{x}$. По непрерывности
$\varphi$
\[
  \lim_{j \to \infty} \varphi(x_{k_j}) = \varphi(\bar{x}) = 0
  \qquad\Longrightarrow\qquad \bar{x} \in \Omega
  \qquad\Longrightarrow\qquad \varphi_0(\bar{x}) \ge \varphi_0(x_*) .
\]
С другой стороны, как в пункте~2), $\varphi_0(x_{k_j}) \le \varphi_0(x_*)$ для
всех $j$, и по непрерывности $\varphi_0(\bar{x}) \le \varphi_0(x_*)$. Сравнивая
два неравенства,
\[
  \varphi_0(\bar{x}) = \varphi_0(x_*) ,
\]
то есть $\bar{x}$ — оптимальное решение исходной задачи.
\end{proof}

\begin{idea}
Две оценки на $\varphi_0(\bar{x})$ приходят из разных мест, и именно это стоит
держать в голове:
\begin{itemize}[nosep,left=0pt]
  \item $\varphi_0(\bar{x}) \ge \varphi_0(x_*)$ — потому что $\bar{x}$
        \emph{допустима} ($\varphi(\bar x) = 0$), а $x_*$ — минимум по всем
        допустимым точкам;
  \item $\varphi_0(\bar{x}) \le \varphi_0(x_*)$ — потому что каждая $x_{k_j}$
        минимизирует $\varphi_0$ по \emph{более широкому} множеству
        $S_{k_j} \supset \Omega$.
\end{itemize}
Метод зажимает оптимум с двух сторон: снизу — значениями на релаксациях,
сверху — допустимостью предельной точки.
\end{idea}

\subsubsection*{Численный пример}

\begin{idea}[title={Откуда этот материал}]
Прогон метода отсечений на числах; на лекции пример не разбирался. Задача
специально взята в том виде, к которому алгоритм применим без преобразований:
цель линейна, ограничение одно и выпукло.
\end{idea}

\begin{example}\label{ex:cutting-plane-run}
\[
  \begin{cases}
    \varphi_0(x) = x_1 \to \min , \\
    \varphi(x) = x_1^2 + x_2^2 - 1 \le 0 ,
  \end{cases}
  \qquad
  \Omega = \{\, x \mid \norm{x} \le 1 \,\} ,
  \qquad
  c = \begin{pmatrix} 1 \\ 0 \end{pmatrix} .
\]
Ответ очевиден: $x_* = (-1, 0)\T$, $\varphi_0(x_*) = -1$. Начальный объемлющий
многогранник — квадрат
\[
  S_0 = \{\, x \mid -2 \le x_1 \le 2,\ -2 \le x_2 \le 2 \,\} \supset \Omega .
\]
Градиент ограничения $\nabla\varphi(x) = (2x_1,\ 2x_2)\T$ — он же субградиент,
так как $\varphi$ дифференцируема.

\medskip
\noindent\textbf{Итерация 0.} Линейная подзадача $x_1 \to \min$, $x \in S_0$
даёт $x_1 = -2$ при любом $x_2 \in [-2,2]$: минимум достигается на всём левом
ребре. Выберем $x_0 = (-2, 0)\T$ (см. замечание после примера). Проверка
допустимости:
\[
  \varphi(x_0) = 4 + 0 - 1 = 3 > 0 ,
\]
точка недопустима, строим отсечение:
\[
  a_1 = \nabla\varphi(x_0) = \begin{pmatrix} -4 \\ 0 \end{pmatrix} ,
  \qquad
  b^{(1)} = -\varphi(x_0) + a_1\T x_0 = -3 + 8 = 5 ,
\]
\[
  a_1\T x - b^{(1)} \le 0
  \quad\Longleftrightarrow\quad
  -4x_1 - 5 \le 0
  \quad\Longleftrightarrow\quad
  x_1 \ge -\tfrac54 .
\]
\[
  S_1 = S_0 \cap \big\{\, x_1 \ge -\tfrac54 \,\big\} .
\]
Отсечение корректно: для любой точки круга $\abs{x_1} \le 1$, то есть
$x_1 \ge -1 > -\tfrac54$, — ни одна допустимая точка не потеряна
(лемма~\ref{lem:cut-nested}), а сама $x_0$ выброшена.

\medskip
\noindent\textbf{Итерация 1.} Решение подзадачи на $S_1$:
$x_1 = (-\tfrac54,\ 0)\T$,
\[
  \varphi(x_1) = \tfrac{25}{16} - 1 = \tfrac{9}{16} > 0 ,
\]
\[
  a_2 = \nabla\varphi(x_1) = \begin{pmatrix} -\tfrac52 \\ 0 \end{pmatrix} ,
  \qquad
  b^{(2)} = -\tfrac{9}{16} + \Big(-\tfrac52\Big)\Big(-\tfrac54\Big)
          = -\tfrac{9}{16} + \tfrac{50}{16} = \tfrac{41}{16} ,
\]
\[
  -\tfrac52 x_1 - \tfrac{41}{16} \le 0
  \quad\Longleftrightarrow\quad
  x_1 \ge -\tfrac{41}{40} = -1{,}025 ,
  \qquad
  S_2 = S_1 \cap \big\{\, x_1 \ge -\tfrac{41}{40} \,\big\} .
\]

\medskip
\noindent\textbf{Итерация 2.} $x_2 = (-\tfrac{41}{40},\ 0)\T$,
\[
  \varphi(x_2) = \tfrac{1681}{1600} - 1 = \tfrac{81}{1600} \approx 0{,}0506 .
\]
Невязка ограничения падает быстро:
\[
  3 \ \to\ \tfrac{9}{16} = 0{,}5625 \ \to\ \tfrac{81}{1600} \approx 0{,}0506 ,
\]
а приближения к оптимуму идут снизу:
\[
  -2 \ \to\ -1{,}25 \ \to\ -1{,}025 \ \to\ \dots \ \to\ -1 .
\]
Все $x_k$ недопустимы, поэтому $\varphi_0(x_k) \le \varphi_0(x_*)$ — внешняя
аппроксимация даёт \emph{нижние} оценки оптимума, что и использовалось в
доказательстве теоремы~\ref{thm:cutting-plane}.
\end{example}

\begin{remark}[о неединственности подзадачи]
Минимум $x_1$ на $S_k$ достигается на целом отрезке, и симплекс-метод вернул бы
какую-то вершину, например $(-2,-2)\T$. Мы фиксируем $x_2 = 0$: задача
симметрична относительно оси $x_1$, и такой выбор удерживает все итерации на
одной прямой, что делает счёт обозримым. На скорость это не влияет — отсечение
всё равно получается вертикальной прямой $x_1 = \text{const}$.
\end{remark}

\begin{idea}
Любопытная деталь: на прямой $x_2 = 0$ ограничение превращается в
$g(t) = t^2 - 1 \le 0$, а отсечение — в
\[
  g(t_k) + g'(t_k)(t - t_k) \le 0
  \quad\Longleftrightarrow\quad
  t \ \ge\ \frac{t_k^2 + 1}{2 t_k}
  \qquad (t_k < 0) ,
\]
то есть $t_{k+1} = \tfrac12\big(t_k + \tfrac{1}{t_k}\big)$ — в точности метод
Ньютона для $\sqrt{1}$. Отсюда и наблюдаемая скорость: $-2$, $-1{,}25$,
$-1{,}025$, $-1{,}0003$ — число верных знаков удваивается за шаг. В общем
случае такой скорости у метода отсечений нет: она здесь берётся из гладкости и
одномерности, а не из самого метода.
\end{idea}

\subsection{Метод условного градиента}
\label{subsec:cond-gradient}

Это тот же метод возможных направлений, но относящийся ещё и к методам
\emph{линейной аппроксимации}.

\[
  \begin{cases}
    \varphi(x) \to \min , & x \in S , \\
    S \ \text{— выпуклый компакт} ,
  \end{cases}
\]
$\varphi$ — выпуклая, непрерывно дифференцируемая, с липшицевым градиентом
(константа $R$).

\subsubsection*{Алгоритм}

Выбираются $\alpha_0 \in (0,1)$ (иногда $\alpha_0 = 1$), $\lambda \in (0,1)$,
$x_0 \in S$. Строится последовательность $\{x_k\} \subset S$.

Пусть выполнен $k$-й шаг и $x_k$ построена. Решается \textbf{вспомогательная
задача}
\[
  \nabla\T\varphi(x_k)\, x \ \longrightarrow\ \min_{x \in S} .
\]
Пусть $y_k$ — её решение. Выбираем
\[
  s_k = y_k - x_k \quad\text{— направление спуска},
  \qquad
  x_{k+1} = x_k + \alpha_k s_k = x_k + \alpha_k (y_k - x_k) ,
\]
и обозначаем
\[
  \eta_k = \nabla\T\varphi(x_k)\,(y_k - x_k) .
\]

\begin{idea}
Вспомогательная задача — это минимизация \emph{линейной} аппроксимации
$\varphi$ (постоянное слагаемое $\varphi(x_k) - \nabla\T\varphi(x_k)x_k$ на
точку минимума не влияет). Отсюда два практических замечания, помеченные на
лекции стрелками:
\begin{itemize}[nosep,left=0pt]
  \item если $S$ задано линейными ограничениями, вспомогательная задача — это
        задача ЛП, и решается симплекс-методом (\S4);
  \item в любом случае минимум линейной функции на выпуклом компакте
        достигается на границе (и, если $S$ — многогранник, в вершине).
        Поэтому $y_k$ — всегда «крайняя» точка, и шаг $\alpha_k$ обязателен:
        идти в $y_k$ целиком, как правило, нельзя.
\end{itemize}
\end{idea}

\inputfig{figures/fig-frank-wolfe}

Так как $y_k$ минимизирует $\nabla\T\varphi(x_k)x$ по $S$, а $x_k \in S$,
\[
  \nabla\T\varphi(x_k)\, y_k \ \le\ \nabla\T\varphi(x_k)\, x_k
  \qquad\Longrightarrow\qquad
  \eta_k \le 0 .
\]
Возможны два случая:
\begin{enumerate}[label=\arabic*),left=0pt]
  \item $\eta_k < 0$ \ $\Longrightarrow$ \
        $\eta_k = \nabla\T\varphi(x_k)\, s_k < 0$, то есть $s_k$ — направление
        спуска;
  \item $\eta_k = 0$ \ $\Longrightarrow$ \ движение прекращается (см. условие
        останова ниже).
\end{enumerate}
Относительно $x_{k+1}$ требуется: 2.1)~$x_{k+1}$ допустима; 2.2)~$\{x_k\}$ —
последовательность спуска. Обоим требованиям и подчинён выбор $\alpha_k$.

\subsubsection*{Допустимость}

Пусть $\eta_k < 0$. Тогда
\[
  x_{k+1} = x_k + \alpha_k (y_k - x_k) = (1 - \alpha_k) x_k + \alpha_k y_k ,
\]
и так как $x_k \in S$, $y_k \in S$, при $\alpha_k \in [0,1]$ точка $x_{k+1}$
лежит в $S$ — множество $S$ выпукло.

\subsubsection*{Выбор шага}

По формуле конечных приращений
\[
  \varphi(x_k + \alpha_k s_k)
  = \varphi(x_k) + \alpha_k \nabla\T\varphi(\bar{x}_k)\, s_k ,
  \qquad
  \bar{x}_k = x_k + \theta \alpha_k s_k , \quad \theta \in [0,1] .
\]
Прибавляя и вычитая $\alpha_k \nabla\T\varphi(x_k) s_k$ и пользуясь
липшицевостью градиента,
\begin{align*}
  \varphi(x_k + \alpha_k s_k)
  &= \varphi(x_k) + \alpha_k \underbrace{\nabla\T\varphi(x_k) s_k}_{=\ \eta_k}
     + \alpha_k \big(\nabla\varphi(\bar{x}_k) - \nabla\varphi(x_k)\big)\T s_k \\
  &\le \varphi(x_k) + \alpha_k \eta_k
     + \alpha_k \norm{\nabla\varphi(\bar{x}_k) - \nabla\varphi(x_k)}
       \cdot \norm{s_k} \\
  &\le \varphi(x_k) + \alpha_k \eta_k
     + \alpha_k R \underbrace{\norm{\bar{x}_k - x_k}}_{=\ \theta\alpha_k\norm{s_k}}
       \cdot \norm{s_k}
   \ \le\ \varphi(x_k) + \alpha_k \eta_k + \alpha_k^2 R \norm{s_k}^2 .
\end{align*}
Потребуем
\begin{equation}\label{eq:fw-req}
  \varphi(x_k) + \alpha_k \eta_k + \alpha_k^2 R \norm{s_k}^2
  \ \le\ \varphi(x_k) + \tfrac12 \alpha_k \eta_k ,
\end{equation}
то есть $\tfrac12 \alpha_k \eta_k + \alpha_k^2 R \norm{s_k}^2 \le 0$, откуда
\begin{equation}\label{eq:fw-alpha}
  \alpha_k \ \le\ -\frac{1}{2}\, \eta_k \cdot \frac{1}{R \norm{s_k}^2} .
\end{equation}

\begin{remark}
Пусть $x_0$ допустима, построена $\{x_k\} \subset S$, и правило выбора шага
привело к~\eqref{eq:fw-alpha}. Тогда
\begin{equation}\label{eq:fw-1}
  \varphi(x_{k+1}) \ \le\ \varphi(x_k) + \tfrac12 \alpha_k \eta_k .
\end{equation}
\end{remark}

\begin{idea}
При $\eta_k = 0$ правая часть~\eqref{eq:fw-alpha} равна нулю: шага не делается,
и это корректно — именно $\eta_k = 0$ служит признаком остановки.

При $\eta_k < 0$ имеем $s_k \ne 0$, и шаг $\alpha_k$ может быть выбран.
Берём его по принципу дробления, $\alpha_k = \alpha_0 \lambda^{i_k}$: это первое
значение, удовлетворяющее~\eqref{eq:fw-alpha}. Тогда предыдущее пробное
значение $\alpha = \alpha_k / \lambda$ условию~\eqref{eq:fw-alpha} уже
\emph{не} удовлетворяло (само неравенство~\eqref{eq:fw-req} от выбора $\alpha$
не зависит по форме), то есть
\[
  \alpha_0 \lambda^{\,i_k - 1} \ >\ -\frac{1}{2}\,\eta_k \cdot
  \frac{1}{R\norm{s_k}^2}
  \qquad\Longrightarrow\qquad
  \alpha_k \ >\ -\frac{1}{2}\,\eta_k \cdot \frac{\lambda}{R\norm{s_k}^2} .
\]
Наконец, $\norm{s_k} = \norm{y_k - x_k}$, где $y_k, x_k \in S$, а $S$ —
компакт, поэтому существует $C$: $\norm{s_k} < C$. Итого
\begin{equation}\label{eq:fw-2}
  \alpha_k \ >\ -\frac{\lambda}{2RC^2}\, \eta_k .
\end{equation}
\end{idea}

\subsubsection*{Основная лемма}

\begin{lemma}\label{lem:fw}
Пусть $\{x_k\}$ построена по методу условного градиента. Тогда
\[
  1^\circ\ \{x_k\} \subset S ;
  \qquad
  2^\circ\ \varphi(x_k) \ \text{монотонно убывает} ;
  \qquad
  3^\circ\ \eta_k \xrightarrow[k \to \infty]{} 0 .
\]
\end{lemma}

\begin{proof}
$1^\circ$ доказано выше (выпуклость $S$ и $\alpha_k \in [0,1]$).

\medskip
$2^\circ$--$3^\circ$. Из~\eqref{eq:fw-1} и~\eqref{eq:fw-2}, учитывая
$\eta_k < 0$ (умножение на отрицательное число меняет знак неравенства),
\begin{equation}\label{eq:fw-3}
  \varphi(x_{k+1}) - \varphi(x_k)
  = \varphi(x_k + \alpha_k s_k) - \varphi(x_k)
  \ \le\ \tfrac12 \alpha_k \eta_k
  \ <\ -\frac{\lambda}{4RC^2}\, \eta_k^2 \ \le\ 0 ,
\end{equation}
то есть $\varphi(x_k)$ убывает. Просуммируем~\eqref{eq:fw-3} по
$k = \overline{0, m-1}$; слагаемые телескопируются:
\[
  \varphi(x_m) - \varphi(x_0)
  \ <\ -\frac{\lambda}{4RC^2} \sum_{k=0}^{m-1} \eta_k^2 .
\]
Так как $\varphi(x_m) \ge \varphi(x_*)$,
\[
  \varphi(x_*) - \varphi(x_0) \ \le\ \varphi(x_m) - \varphi(x_0)
  \ <\ -\frac{\lambda}{4RC^2} \sum_{k=0}^{m-1} \eta_k^2 ,
\]
откуда при $m \to +\infty$
\[
  \sum_{k=0}^{\infty} \eta_k^2
  \ \le\ \frac{4RC^2}{\lambda}\big(\varphi(x_0) - \varphi(x_*)\big) < +\infty .
\]
Ряд из $\eta_k^2$ сходится, значит $\eta_k^2 \to 0$ и $\eta_k \to 0$.
\end{proof}

\subsubsection*{Условие останова}

\begin{enumerate}[label=\arabic*),left=0pt]
  \item за конечное число шагов получено $\eta_k = 0$;
  \item последовательность $\{\varphi(x_k)\}$ перестала заметно убывать.
\end{enumerate}

Разберём первый случай. Пусть $\eta_k = 0$, то есть
$\nabla\T\varphi(x_k)(y_k - x_k) = 0$, значит
$\nabla\T\varphi(x_k) y_k = \nabla\T\varphi(x_k) x_k$. Но $y_k$ по построению
минимизирует линейную функцию:
\[
  \nabla\T\varphi(x_k)\, y_k \ \le\ \nabla\T\varphi(x_k)\, x
  \qquad \forall x \in S ,
\]
следовательно $\nabla\T\varphi(x_k) x \ge \nabla\T\varphi(x_k) x_k$, то есть
\[
  \nabla\T\varphi(x_k)\,(x - x_k) \ \ge\ 0 \qquad \forall x \in S .
\]
Это условие экстремума выпуклой функции на выпуклом множестве — значит $x_k$
оптимальна.

\subsubsection*{Теорема о сходимости}

\begin{theorem}[о сходимости метода условного градиента]\label{thm:fw}
Пусть выполнены условия применимости метода и $\{x_k\}$ построена по методу
условного градиента. Тогда в любой предельной точке последовательности
выполняется условие минимума задачи $\varphi(x) \to \min$, $x \in S$.
\end{theorem}

\begin{proof}
Пусть $\bar{x}_*$ — предельная точка $\{x_k\}$, то есть существует
подпоследовательность $x_{k_j} \to \bar{x}_*$. По лемме~\ref{lem:fw}
\[
  \eta_{k_j} = \nabla\T\varphi(x_{k_j})\,(y_{k_j} - x_{k_j})
  \xrightarrow[j \to \infty]{} 0 .
\]
Последовательность $\{y_{k_j}\}$ лежит в компакте $S$, поэтому, переходя при
необходимости к подпоследовательности, считаем $y_{k_j} \to y_*$. По
непрерывности $\nabla\varphi$
\begin{equation}\label{eq:fw-4}
  \nabla\T\varphi(\bar{x}_*)\,(y_* - \bar{x}_*) = 0 .
\end{equation}

Далее — так же, как при разборе условия останова. Для всех $x \in S$
\[
  \nabla\T\varphi(x_{k_j})\, y_{k_j} \ \le\ \nabla\T\varphi(x_{k_j})\, x ,
\]
и при $k_j \to \infty$
\[
  \nabla\T\varphi(\bar{x}_*)\, y_* \ \le\ \nabla\T\varphi(\bar{x}_*)\, x .
\]
По~\eqref{eq:fw-4} левая часть равна
$\nabla\T\varphi(\bar{x}_*)\,\bar{x}_*$, поэтому
\[
  \nabla\T\varphi(\bar{x}_*)\, \bar{x}_* \ \le\ \nabla\T\varphi(\bar{x}_*)\, x
  \qquad\Longrightarrow\qquad
  \nabla\T\varphi(\bar{x}_*)\,(x - \bar{x}_*) \ \ge\ 0 \quad \forall x \in S .
\]
Это условие экстремума в точке $\bar{x}_*$; в силу выпуклости $\varphi$ оно
достаточно, значит $\bar{x}_*$ — точка минимума. Отсюда
\[
  \lim_{k \to \infty} \varphi(x_k) = \varphi(x_*) ,
  \qquad \varphi(x_*) = \min_{x \in S} \varphi(x) .
  \qedhere
\]
\end{proof}

\begin{remark}
Кроме того, справедлива оценка скорости
\[
  \varphi(x_k) - \varphi(x_*) \ \le\ \frac{D}{k} ,
  \qquad D > 0 \ \text{— const} .
\]
\end{remark}

\begin{idea}
Оценка $O(1/k)$ — классический результат для метода условного градиента
(Франка--Вульфа); на лекции она приведена без доказательства. Константа $D$
выражается через диаметр $S$ и константу Липшица градиента:
$D \sim R\,\operatorname{diam}^2 S$. Важно, что скорость здесь \emph{сублинейная}
— заметно хуже геометрической сходимости методов из главы о безусловной
минимизации. Платой за это является крайняя дешевизна итерации: вместо
проектирования на $S$ (как в методе Розена) решается одна линейная задача.
\end{idea}

\subsubsection*{Численный пример}

\begin{idea}[title={Откуда этот материал}]
Прогон метода условного градиента на числах; на лекции примера не было.
Допустимое множество взято многогранным, чтобы вспомогательная задача была
обычной задачей ЛП.
\end{idea}

\begin{example}\label{ex:frank-wolfe-run}
\[
  \varphi(x) = (x_1 - 0{,}3)^2 + (x_2 - 0{,}7)^2 \ \longrightarrow\ \min ,
  \qquad
  S = [0,1] \times [0,1] .
\]
Множество $S$ — выпуклый компакт, $\varphi$ выпукла и дифференцируема,
\[
  \nabla\varphi(x) = \begin{pmatrix} 2(x_1 - 0{,}3) \\ 2(x_2 - 0{,}7) \end{pmatrix} .
\]
Старт: $x_0 = (1,1)\T \in S$, $\varphi(x_0) = 0{,}49 + 0{,}09 = 0{,}58$.

\medskip
\noindent\textbf{Итерация 0.}
\[
  \nabla\varphi(x_0) = \begin{pmatrix} 1{,}4 \\ 0{,}6 \end{pmatrix} .
\]
Вспомогательная задача:
\[
  1{,}4\,x_1 + 0{,}6\,x_2 \ \longrightarrow\ \min_{x \in S} .
\]
Оба коэффициента положительны, поэтому минимум достигается в вершине с
наименьшими координатами:
\[
  y_0 = \begin{pmatrix} 0 \\ 0 \end{pmatrix} ,
  \qquad
  s_0 = y_0 - x_0 = \begin{pmatrix} -1 \\ -1 \end{pmatrix} ,
  \qquad
  \eta_0 = \nabla\T\varphi(x_0)\,s_0 = -1{,}4 - 0{,}6 = -2 < 0 .
\]
Шаг:
\[
  \varphi(x_0 + \alpha s_0)
   = (0{,}7 - \alpha)^2 + (0{,}3 - \alpha)^2
   = 0{,}58 - 2\alpha + 2\alpha^2 ,
\]
\[
  -2 + 4\alpha = 0 \quad\Longrightarrow\quad \alpha_0 = 0{,}5 \in [0,1] ,
\]
\[
  x_1 = (1-\alpha_0)x_0 + \alpha_0 y_0
      = \begin{pmatrix} 0{,}5 \\ 0{,}5 \end{pmatrix} ,
  \qquad
  \varphi(x_1) = 0{,}04 + 0{,}04 = 0{,}08 .
\]

\medskip
\noindent\textbf{Итерация 1.}
\[
  \nabla\varphi(x_1) = \begin{pmatrix} 0{,}4 \\ -0{,}4 \end{pmatrix} ,
  \qquad
  0{,}4\,x_1 - 0{,}4\,x_2 \ \longrightarrow\ \min_{x \in S} .
\]
Коэффициент при $x_1$ положителен — берём $x_1 = 0$; при $x_2$ отрицателен —
берём $x_2 = 1$:
\[
  y_1 = \begin{pmatrix} 0 \\ 1 \end{pmatrix} ,
  \qquad
  s_1 = \begin{pmatrix} -0{,}5 \\ 0{,}5 \end{pmatrix} ,
  \qquad
  \eta_1 = -0{,}2 - 0{,}2 = -0{,}4 < 0 .
\]
Шаг:
\[
  \varphi(x_1 + \alpha s_1)
   = (0{,}2 - 0{,}5\alpha)^2 + (-0{,}2 + 0{,}5\alpha)^2
   = 2\,(0{,}2 - 0{,}5\alpha)^2 ,
\]
\[
  \alpha_1 = 0{,}4 \in [0,1] ,
  \qquad
  x_2 = \begin{pmatrix} 0{,}5 - 0{,}2 \\ 0{,}5 + 0{,}2 \end{pmatrix}
      = \begin{pmatrix} 0{,}3 \\ 0{,}7 \end{pmatrix} = x_* ,
  \qquad \varphi(x_2) = 0 .
\]
Значения: $0{,}58 \to 0{,}08 \to 0$.
\end{example}

\begin{idea}
Три вещи, которые этот счёт иллюстрирует.

\begin{itemize}[nosep,left=0pt]
  \item Точки $y_k$ — всегда вершины квадрата: минимум линейной функции на
        многограннике достигается в вершине. Идти в $y_k$ целиком нельзя
        (при $\alpha = 1$ из $x_0$ мы попали бы в $(0,0)$ со значением
        $0{,}58$ — ровно столько же, сколько было), поэтому шаг обязателен.
  \item Допустимость получается даром: $x_{k+1} = (1-\alpha_k)x_k + \alpha_k y_k$
        — выпуклая комбинация двух точек $S$ при $\alpha_k \in [0,1]$. Никакого
        проектирования, в отличие от метода Розена.
  \item Попадание точно в $x_*$ за два шага — следствие того, что $x_*$ лежит
        \emph{внутри} $S$ и функция круговая. В типичной ситуации, когда
        оптимум на границе, метод подходит к нему сублинейно, с оценкой
        $O(1/k)$ из замечания выше: направления $s_k$ всё время «дёргаются»
        между вершинами.
\end{itemize}
\end{idea}


\newpage
\subsection{Метод Гриффиса и Стюарта}
\label{subsec:griffith-stewart}

\textbf{Идея:} построить линейные аппроксимации \emph{и} для функции цели,
\emph{и} для ограничений задачи, а затем ограничить длину шага, чтобы
линеаризация оставалась осмысленной.

Метод разбирается на примере.

\begin{example}\label{ex:griffith}
\[
  (\mathrm{I})
  \qquad
  \begin{cases}
    \varphi_0(x) = x_1^2 + x_2^2 - 16 x_1 - 10 x_2 \to \min , \\[2pt]
    g_1(x_1, x_2) = x_1^2 - 6 x_1 + 4 x_2 - 11 \le 0 , \\[2pt]
    g_2(x_1, x_2) = e^{\,x_1 - 3} - x_1 x_2 + 3 x_2 - 1 \le 0 , \\[2pt]
    x_1 \ge 0 , \quad x_2 \ge 0 .
  \end{cases}
\]
\end{example}

\subsubsection*{Начальный шаг, $k = 0$}

\[
  x^{(0)} = \begin{pmatrix} 4 \\ 3 \end{pmatrix} ,
  \qquad
  \varphi_0\big(x^{(0)}\big) = -69 .
\]
Линеаризуем все функции в точке $x^{(0)}$. Для целевой функции
\[
  \tilde\varphi_0(x) = -69 +
  \begin{pmatrix} 2x_1 - 16 \\ 2x_2 - 10 \end{pmatrix}\T_{\!\!\big|\,x^{(0)}}
  \begin{pmatrix} x_1 - 4 \\ x_2 - 3 \end{pmatrix}
  = -69 - 8(x_1 - 4) - 4(x_2 - 3)
  = -8x_1 - 4x_2 - 25 ,
\]
и аналогично для ограничений. Получается задача ЛП
\[
  (\mathrm{II})
  \qquad
  \begin{cases}
    \tilde\varphi_0(x) = -8x_1 - 4x_2 - 25 \to \min , \\[2pt]
    \tilde g_1(x) = 2x_1 + 4x_2 - 27 \le 0 , \\[2pt]
    \tilde g_2(x) = -0{,}28 x_1 - x_2 + 2{,}84 \le 0 , \\[2pt]
    x_1 \ge 0 , \quad x_2 \ge 0 .
  \end{cases}
\]
Её решение —
\[
  x^{(1)} = \begin{pmatrix} 13{,}5 \\ 0 \end{pmatrix} ,
  \qquad\text{тогда как}\qquad
  x^{(0)} = \begin{pmatrix} 4 \\ 3 \end{pmatrix} .
\]

Точка ушла слишком далеко: линеаризация в окрестности $x^{(0)}$ ничего не знает
про поведение $g_1, g_2$ при $x_1 \approx 13{,}5$. Поэтому накладывается
ограничение на шаг, предотвращающее выход из области:
\begin{equation}\label{eq:gs-trust}
  \big| x_j^{(1)} - x_j^{(0)} \big| \ \le\ \delta_j^{(0)} ,
  \qquad j = \overline{1,2} .
\end{equation}
Положив $\delta^{(0)} = (0{,}5;\ 0{,}5)\T$, получаем
\[
  x^{(1)} = \begin{pmatrix} 4 + 0{,}5 \\ 3 - 0{,}5 \end{pmatrix}
          = \begin{pmatrix} 4{,}5 \\ 2{,}5 \end{pmatrix} .
\]

\subsubsection*{Шаг $k = 1$}

Строим линейную аппроксимацию задачи~$(\mathrm{I})$ в окрестности $x^{(1)}$:
\[
  (\mathrm{III})
  \qquad
  \begin{cases}
    \tilde\varphi_0(x) = -7 x_1 - 5 x_2 - 26{,}5 \to \min , \\[2pt]
    3 x_1 + 4 x_2 - 31{,}40 \le 0 , \\[2pt]
    1{,}98 x_1 - 1{,}50 x_2 - 5{,}42 \le 0 , \\[2pt]
    x_1 \ge 0 , \quad x_2 \ge 0 .
  \end{cases}
\]
Решение ЛП —
\[
  x^{(2)} = \begin{pmatrix} 5{,}45 \\ 3{,}58 \end{pmatrix}
  \qquad\text{— не подходит, снова применяем \eqref{eq:gs-trust}:}
\]
\[
  \delta_j^{(1)} = \lambda\, \delta_j^{(0)} , \qquad \lambda \in (0,1) .
\]
При $\delta^{(1)} = (0{,}4;\ 0{,}4)\T$
\[
  x^{(2)} = \begin{pmatrix} 4{,}90 \\ 2{,}90 \end{pmatrix} ,
  \qquad
  \varphi_0\big(x^{(2)}\big) = -75{,}1 .
\]

\begin{idea}[title={О числах в примере}]
Свободный член линеаризованной цели ($-25$ на первом шаге, $-26{,}5$ на втором)
на решение задачи ЛП не влияет — минимум линейной функции не меняется от сдвига
на константу. Поэтому расхождения в этом слагаемом несущественны; важны только
коэффициенты при переменных, то есть компоненты градиента.

Коэффициенты в $\tilde g_1, \tilde g_2$ получаются прямой подстановкой:
\[
  \nabla g_1 = \begin{pmatrix} 2x_1 - 6 \\ 4 \end{pmatrix} ,
  \qquad
  \nabla g_2 = \begin{pmatrix} e^{\,x_1-3} - x_2 \\ -x_1 + 3 \end{pmatrix} ,
\]
например в $x^{(0)} = (4,3)$: $\nabla g_1 = (2;\,4)\T$,
$g_1(x^{(0)}) = -7$, откуда $\tilde g_1 = -7 + 2(x_1-4) + 4(x_2-3)
= 2x_1 + 4x_2 - 27$; и $\nabla g_2 = (e - 3;\, -1)\T \approx (-0{,}28;\,-1)\T$,
$g_2(x^{(0)}) \approx -1{,}28$, откуда
$\tilde g_2 \approx -0{,}28x_1 - x_2 + 2{,}84$. Числа в $(\mathrm{III})$
отличаются от пересчёта в третьем знаке — это округление, на ход метода оно не
влияет.
\end{idea}

\begin{idea}[title={Смысл ограничения \eqref{eq:gs-trust}}]
Линеаризация точна только в малой окрестности точки разложения, а минимум
линейной функции на многограннике всегда достигается в вершине — то есть
«как можно дальше». Без дополнительного ограничения метод на каждом шаге
улетал бы в далёкую вершину, где аппроксимация уже не имеет отношения к
исходной задаче (так и вышло: $x^{(1)} = (13{,}5;\,0)$ вместо разумного
смещения от $(4;\,3)$).

Неравенства~\eqref{eq:gs-trust} задают вокруг текущей точки куб со стороной
$2\delta_j^{(k)}$ — \emph{область доверия} к линейной модели. Если шаг оказался
неудачным, размеры куба уменьшают: $\delta^{(k+1)} = \lambda\delta^{(k)}$,
$\lambda \in (0,1)$ — то же самое дробление, что и при выборе шага в
градиентных методах, только дробится не длина шага вдоль фиксированного
направления, а размер всей области, в которой ищется следующая точка.
\end{idea}

\subsection{Метод квадратичной аппроксимации}
\label{subsec:quad-approx}

\begin{idea}
В тетради название метода пропущено. По содержанию это \textbf{метод
квадратичной аппроксимации} — прямой аналог метода Ньютона для задач с
ограничениями: там, где метод условного градиента (\S\ref{subsec:cond-gradient})
минимизирует линейную модель, здесь минимизируется квадратичная. На нём же
основано семейство методов последовательного квадратичного программирования.
\end{idea}

\[
  \begin{cases}
    \varphi(x) \to \min , & x \in S , \\
    \varphi(x) \ \text{— выпуклая, дважды дифференцируемая} , \\
    S \ \text{— выпуклый компакт} .
  \end{cases}
\]

\begin{algo}[title={Метод квадратичной аппроксимации}]
\textbf{I. Начальный этап.} Задать $x_0 \in S$ — начальное приближение,
$k = 0$; итерации вида $x_{k+1} = x_k + \alpha_k s_k$, $\alpha_k > 0$.

\medskip
\textbf{II. $k$-й шаг.} Рассмотрим квадратичную аппроксимацию в окрестности
$x_k$:
\begin{steps}
  \item Положим
        \[
          \psi_k(x) = \nabla\T\varphi(x_k)\,(x - x_k)
            + \tfrac{1}{2}\,(x - x_k)\T H(x_k)\,(x - x_k)
        \]
        и решим вспомогательную задачу
        \[
          \psi_k \to \min , \qquad x \in S .
        \]
        Пусть $y_k$ — её решение; так как $x_k \in S$ и $\psi_k(x_k) = 0$,
        \[
          \eta_k = \psi_k(y_k) \le \psi_k(x_k) = 0 .
        \]
        В качестве направления берём $s_k = y_k - x_k$.
  \item Шаг выбирается как в методе второго порядка для задач без ограничений:
        заданы $\alpha_0 = 1$, $\eps \in (0,1)$, $\lambda \in (0,1)$, и
        $\alpha_k = \alpha_0 \lambda^{\,i_k}$ — первое значение, удовлетворяющее
        условию дробления
        \begin{equation}\label{eq:qa-armijo}
          \varphi(x_k + \alpha s_k) - \varphi(x_k) \ \le\ \alpha\,\eps\,\eta_k
        \end{equation}
        (либо $\varphi(x_{k+1}) \to \min_{\alpha > 0}$ — наискорейший спуск
        вдоль $s_k$).
\end{steps}
\end{algo}

\subsubsection*{Вспомогательные факты}

\begin{lemma}\label{lem:qa-basic}
В условиях задачи:
\begin{enumerate}[label=$\arabic*^\circ$,left=0pt]
  \item $x_k + \alpha s_k \in S$ при всех $\alpha \in [0,1]$;
  \item $\nabla\T\varphi(x_k)\, s_k \le \eta_k \le 0$;
  \item если дополнительно $\varphi$ сильно выпукла с константой $m$, то
        \[
          \eta_k \ \le\ -\tfrac{m}{2}\,\norm{s_k}^2 .
        \]
\end{enumerate}
\end{lemma}

\begin{proof}
$1^\circ$ $x_k + \alpha s_k = (1-\alpha)x_k + \alpha y_k$ — выпуклая комбинация
точек $x_k, y_k \in S$.

$2^\circ$ По определению $\eta_k = \psi_k(y_k)
= \nabla\T\varphi(x_k)s_k + \tfrac12 s_k\T H(x_k) s_k$, а второе слагаемое
неотрицательно в силу выпуклости $\varphi$ ($H \succeq 0$). Неравенство
$\eta_k \le 0$ — из $\psi_k(x_k) = 0$.

$3^\circ$ Функция $\psi_k$ сильно выпукла с той же константой $m$, а $y_k$ —
её минимум на $S$, поэтому $\nabla\T\psi_k(y_k)(x_k - y_k) \ge 0$. Тогда
\[
  \underbrace{\psi_k(x_k)}_{=\ 0}
  \ \ge\ \psi_k(y_k) + \underbrace{\nabla\T\psi_k(y_k)(x_k - y_k)}_{\ge\ 0}
    + \tfrac{m}{2}\norm{x_k - y_k}^2
  \ \ge\ \eta_k + \tfrac{m}{2}\norm{s_k}^2 . \qedhere
\]
\end{proof}

\begin{lemma}\label{lem:qa-step}
Дробление~\eqref{eq:qa-armijo} конечно, и принятый шаг отделён от нуля:
существует $\bar\alpha > 0$ с $\alpha_k \ge \min\{1, \bar\alpha\,\abs{\eta_k}\}$.
\end{lemma}

\begin{proof}
Гессиан непрерывен на компакте $S$, значит ограничен: $\norm{H(x)} \le M$.
Тогда
\[
  \varphi(x_k + \alpha s_k) \ \le\ \varphi(x_k)
    + \alpha \nabla\T\varphi(x_k) s_k + \tfrac{M}{2}\alpha^2 \norm{s_k}^2
  \ \overset{\text{лемма }\ref{lem:qa-basic}}{\le}\
  \varphi(x_k) + \alpha\eta_k + \tfrac{M}{2}\alpha^2\norm{s_k}^2 .
\]
Для выполнения~\eqref{eq:qa-armijo} достаточно
\[
  \alpha\eta_k + \tfrac{M}{2}\alpha^2\norm{s_k}^2 \le \alpha\eps\eta_k
  \qquad\Longleftrightarrow\qquad
  \alpha \ \le\ -\frac{2(1-\eps)\,\eta_k}{M \norm{s_k}^2} ,
\]
и правая часть положительна при $\eta_k < 0$. Значит после конечного числа
дроблений условие выполнится. Если же принятый шаг не равен $\alpha_0$, то
предыдущее пробное значение $\alpha_k/\lambda$ было отвергнуто, откуда
\[
  \alpha_k \ >\ -\frac{2\lambda(1-\eps)\,\eta_k}{M \norm{s_k}^2}
  \ \ge\ \frac{2\lambda(1-\eps)}{M C^2}\, \abs{\eta_k} ,
\]
где $C = \operatorname{diam} S \ge \norm{s_k}$.
\end{proof}

\subsubsection*{Теоремы о сходимости}

\begin{theorem}\label{thm:qa-1}
Пусть $\varphi(x)$ удовлетворяет условиям постановки задачи и $\{x_k\}$
построена по описанной схеме. Тогда независимо от выбора начального
приближения последовательность $\varphi(x_k)$ монотонно убывает и
\[
  \lim_{k \to \infty} \varphi(x_k) = \varphi(x_*) = \min_{x \in S}\varphi(x) .
\]
\end{theorem}

\begin{idea}[title={Доказательство (на лекции не приводилось)}]
\begin{proof}
\textbf{Шаг 1: убывание и $\eta_k \to 0$.} Из~\eqref{eq:qa-armijo}
\[
  \varphi(x_k) - \varphi(x_{k+1}) \ \ge\ \eps\,\alpha_k\,\abs{\eta_k} \ \ge\ 0 ,
\]
то есть $\varphi(x_k)$ не возрастает; она ограничена снизу ($S$ — компакт,
$\varphi$ непрерывна), значит сходится, и $\varphi(x_k) - \varphi(x_{k+1}) \to 0$.
По лемме~\ref{lem:qa-step}
\[
  \varphi(x_k) - \varphi(x_{k+1})
  \ \ge\ \eps \min\Big\{ \abs{\eta_k},\ \tfrac{2\lambda(1-\eps)}{MC^2}\eta_k^2 \Big\} ,
\]
и правая часть стремится к нулю только если $\eta_k \to 0$.

\medskip
\textbf{Шаг 2: предельная точка оптимальна.} Пусть $x_{k_j} \to \bar{x}$
(существует по компактности $S$). По определению $y_k$ как минимума $\psi_k$
на $S$,
\[
  \eta_k \ \le\ \psi_k(x)
  = \nabla\T\varphi(x_k)(x - x_k) + \tfrac12 (x - x_k)\T H(x_k)(x - x_k)
  \qquad \forall x \in S .
\]
Переходя к пределу по $k_j$ и пользуясь $\eta_{k_j} \to 0$ и непрерывностью
$\nabla\varphi$, $H$:
\[
  0 \ \le\ \nabla\T\varphi(\bar{x})(x - \bar{x})
    + \tfrac12 (x - \bar{x})\T H(\bar{x})(x - \bar{x})
  \qquad \forall x \in S .
\]
Подставим $x = \bar{x} + t(z - \bar{x})$, где $z \in S$, $t \in (0,1]$ (такая
точка лежит в $S$ по выпуклости):
\[
  0 \ \le\ t\,\nabla\T\varphi(\bar{x})(z - \bar{x})
    + \tfrac{t^2}{2}(z - \bar{x})\T H(\bar{x})(z - \bar{x}) .
\]
Разделив на $t > 0$ и устремив $t \to 0$, получаем
\[
  \nabla\T\varphi(\bar{x})\,(z - \bar{x}) \ \ge\ 0 \qquad \forall z \in S ,
\]
а это — необходимое и (в силу выпуклости $\varphi$) достаточное условие
минимума. Значит $\bar{x}$ — точка минимума, и
$\varphi(\bar{x}) = \varphi(x_*)$.

\medskip
\textbf{Шаг 3.} Последовательность $\varphi(x_k)$ монотонна и имеет предел;
по шагу~2 вдоль подпоследовательности этот предел равен $\varphi(x_*)$,
следовательно $\varphi(x_k) \to \varphi(x_*)$.
\end{proof}
\end{idea}

\begin{theorem}\label{thm:qa-2}
Если в дополнение к условиям теоремы~\ref{thm:qa-1} функция $\varphi(x)$
\textbf{сильно выпукла}, то есть
\[
  \exists\, m, M > 0,\ m \le M : \qquad
  m \norm{y}^2 \le y\T H(x)\, y \le M \norm{y}^2 ,
\]
то $\{x_k\} \to x_*$ со \textbf{сверхлинейной} скоростью.
\end{theorem}

\begin{idea}[title={Доказательство (на лекции не приводилось)}]
\begin{proof}
\textbf{Шаг 1: $x_k \to x_*$.} Из сильной выпуклости
$\tfrac{m}{2}\norm{x_k - x_*}^2 \le \varphi(x_k) - \varphi(x_*) \to 0$ по
теореме~\ref{thm:qa-1}, значит $x_k \to x_*$ (точка $x_*$ при сильной
выпуклости единственна).

\medskip
\textbf{Шаг 2: единичный шаг в конце концов принимается.} По
лемме~\ref{lem:qa-basic}\,$3^\circ$
$\norm{s_k}^2 \le \tfrac{2}{m}\abs{\eta_k} \to 0$. По формуле Тейлора
\[
  \varphi(x_k + s_k) - \varphi(x_k)
  = \nabla\T\varphi(x_k)s_k + \tfrac12 s_k\T H(x_k) s_k + o\big(\norm{s_k}^2\big)
  = \eta_k + o\big(\norm{s_k}^2\big) ,
\]
так что условие~\eqref{eq:qa-armijo} при $\alpha = 1$ равносильно
\[
  (1 - \eps)\,\abs{\eta_k} \ \ge\ o\big(\norm{s_k}^2\big) .
\]
Но $\abs{\eta_k} \ge \tfrac{m}{2}\norm{s_k}^2$, то есть левая часть имеет
порядок $\norm{s_k}^2$, а правая — строго меньший. Значит начиная с некоторого
номера $\alpha_k = \alpha_0 = 1$, и
\[
  x_{k+1} = x_k + s_k = y_k .
\]

\medskip
\textbf{Шаг 3: оценка расстояния.} Обозначим
$e_k = x_k - x_*$, $d_k = y_k - x_*$. Точка $y_k$ минимизирует
$q_k(x) = \varphi(x_k) + \psi_k(x)$ на выпуклом $S$, поэтому
\[
  \nabla\T q_k(y_k)\,(x - y_k) \ge 0 \quad \forall x \in S ,
  \qquad
  \nabla q_k(y) = \nabla\varphi(x_k) + H(x_k)(y - x_k) ,
\]
а точка $x_*$ минимизирует $\varphi$ на $S$:
$\nabla\T\varphi(x_*)(x - x_*) \ge 0$ $\forall x \in S$. Подставим в первое
неравенство $x = x_*$, во второе $x = y_k$ и сложим:
\[
  \big[\nabla\varphi(x_k) - \nabla\varphi(x_*) + H(x_k)(y_k - x_k)\big]\T
  (x_* - y_k) \ \ge\ 0 .
\]
По формуле Ньютона--Лейбница
$\nabla\varphi(x_k) - \nabla\varphi(x_*) = G_k e_k$, где
$G_k = \int_0^1 H(x_* + t e_k)\,dt$. Так как $y_k - x_k = d_k - e_k$,
последнее неравенство переписывается как
\[
  \big[H(x_k)\, d_k + \big(G_k - H(x_k)\big) e_k\big]\T d_k \ \le\ 0 ,
\]
откуда
\[
  m \norm{d_k}^2 \ \le\ d_k\T H(x_k)\, d_k
  \ \le\ \norm{G_k - H(x_k)} \cdot \norm{e_k} \cdot \norm{d_k} ,
\]
то есть
\begin{equation}\label{eq:qa-rate}
  \norm{y_k - x_*} \ \le\ \frac{1}{m}\,\norm{G_k - H(x_k)}\cdot\norm{x_k - x_*} .
\end{equation}

\medskip
\textbf{Шаг 4.} Гессиан непрерывен, а $x_k \to x_*$, поэтому
\begin{align*}
  \norm{G_k - H(x_k)}
  &= \Big\| \int_0^1 \big[H(x_* + t e_k) - H(x_k)\big] dt \Big\| \\
  &\le \sup_{t \in [0,1]} \norm{H(x_* + t e_k) - H(x_k)}
   \ \to\ 0 , \quad k \to \infty .
\end{align*}
С учётом шага~2 ($x_{k+1} = y_k$) и~\eqref{eq:qa-rate}
\[
  \norm{x_{k+1} - x_*} = o\big(\norm{x_k - x_*}\big) ,
\]
то есть сходимость сверхлинейная.
\end{proof}
\end{idea}

\begin{theorem}\label{thm:qa-3}
Если дополнительно к условиям теоремы~\ref{thm:qa-2} гессиан $H(\cdot)$
удовлетворяет условию Липшица,
\[
  \norm{H(x) - H(y)} \le R\,\norm{x - y} \qquad \forall x, y ,
\]
то скорость сходимости \textbf{квадратичная}:
\[
  \norm{x_{k+1} - x_*} \ \le\ \frac{R}{2m}\,\norm{x_k - x_*}^2 .
\]
\end{theorem}

\begin{idea}[title={Доказательство (на лекции не приводилось)}]
\begin{proof}
Повторяем шаги~1--3 доказательства теоремы~\ref{thm:qa-2} вплоть
до оценки~\eqref{eq:qa-rate}; остаётся уточнить $\norm{G_k - H(x_k)}$.
Точка $x_* + t e_k$ отстоит от $x_k = x_* + e_k$ на $(1-t)\norm{e_k}$, поэтому
по условию Липшица
\[
  \norm{G_k - H(x_k)}
  \ \le\ \int_0^1 \norm{H(x_* + t e_k) - H(x_k)}\,dt
  \ \le\ \int_0^1 R\,(1 - t)\,\norm{e_k}\,dt
  = \frac{R}{2}\,\norm{e_k} .
\]
Подставляя в~\eqref{eq:qa-rate} и пользуясь $x_{k+1} = y_k$,
\[
  \norm{x_{k+1} - x_*} \ \le\ \frac{R}{2m}\,\norm{x_k - x_*}^2 . \qedhere
\]
\end{proof}
\end{idea}

\begin{idea}
Три теоремы образуют лестницу, типичную для методов второго порядка:
\begin{itemize}[nosep,left=0pt]
  \item выпуклость $+$ дробление шага дают \emph{глобальную} сходимость по
        значению функции, но без оценки скорости;
  \item сильная выпуклость добавляет единственность $x_*$, сходимость самих
        точек и \emph{сверхлинейную} скорость — за счёт того, что вблизи
        минимума единичный шаг перестаёт отвергаться;
  \item липшицевость гессиана превращает сверхлинейную оценку в
        \emph{квадратичную}, в точности как в теореме~\ref{thm:newton-rate}
        для безусловной задачи.
\end{itemize}
Содержательно вся разница с методом Ньютона в том, что минимум квадратичной
модели берётся не по всему $\R^n$, а по $S$; ключевое
неравенство~\eqref{eq:qa-rate} получается не из явной формулы
$y_k = x_k - H^{-1}\nabla\varphi$, а из вариационного неравенства — условия
минимума на выпуклом множестве.
\end{idea}

\subsection{Методы штрафных функций}
\label{subsec:penalty}

\subsubsection*{Метод внешних штрафных функций (МВШФ)}

\[
  \begin{cases}
    \varphi(x) \to \min , & x \in S , \\
    S = \big\{\, x \mid \psi_i(x) \le 0,\ i = \overline{1,m} ;
        \quad \chi_j(x) = 0,\ j = \overline{1,l} \,\big\} , \\
    x \in V \subset \R^n , & V \ \text{— область определения} .
  \end{cases}
\]

\paragraph{1) Одно неравенство.}
Пусть $S = \{\, x \mid \psi(x) \le 0 \,\}$. Положим
\[
  \alpha(x) =
  \begin{cases}
    0 , & \psi(x) \le 0 , \\
    \big(\psi(x)\big)^p , & \psi(x) > 0 ,
  \end{cases}
  \qquad p \ge 1 \ (\text{обычно } p = 2) ,
\]
и составим функцию
\[
  \varphi(x) + \mu\,\alpha(x) , \qquad \mu > 0 .
\]
Здесь $\alpha(x)$ — \textbf{штрафная функция}, $\mu$ — \textbf{параметр
штрафа}. Основное свойство:
\[
  \mu\,\alpha(x) = 0 \quad \forall x \in S ,
  \qquad
  \mu\,\alpha(x) \ne 0 \quad \forall x \notin S .
\]

Вспомогательная задача:
\[
  \big(\varphi(x) + \mu\,\alpha(x)\big) \to \min
  \qquad\Longleftrightarrow\qquad
  \Big(\varphi(x)
    + \mu\, \underbrace{\max\nolimits^{p}\{0,\ \psi(x)\}}_{\alpha(x)}\Big)
  \to \min .
\]
Решается последовательность вспомогательных задач при возрастающем $\mu$; пусть
эти решения — $\{x_\mu\}$.

\begin{idea}[title={Почему $p \ge 1$ и почему обычно $p = 2$}]
При $p < 1$ функция $\max^p\{0,\psi\}$ не выпукла, и вспомогательная задача
перестаёт быть выпуклой даже для выпуклых $\varphi$ и $\psi$ — методы из
предыдущих глав становятся неприменимы.

При $p = 1$ штраф $\max\{0,\psi(x)\}$ выпукл, но \emph{негладок} в точках
границы ($\psi(x) = 0$) — там у него излом, а именно туда и сходятся
приближения. Градиентные методы на такой функции работают плохо.

При $p = 2$ функция $\max^2\{0,\psi(x)\}$ уже непрерывно дифференцируема
(производная равна $2\max\{0,\psi\}\nabla\psi$ и обращается в нуль на границе),
выпукла при выпуклой $\psi$ — и при этом растёт достаточно быстро. Это
компромисс между гладкостью и «силой» штрафа, поэтому $p = 2$ — стандартный
выбор.
\end{idea}

\paragraph{2) Одно равенство.}
Пусть $S = \{\, x \mid \chi(x) = 0 \,\}$. Берём
\[
  \alpha(x) = \abs{\chi(x)}^{\,p}
\]
и рассматриваем ту же вспомогательную задачу
$\varphi(x) + \mu\,\alpha(x) \to \min$, $\mu > 0$.

\paragraph{3) Общий случай.}
\[
  \alpha(x) = \sum_{i=1}^{m} \xi_i \max\nolimits^{p}\{0,\ \psi_i(x)\}
            + \sum_{j=1}^{l} \eta_j \abs{\chi_j(x)}^{\,p} ,
  \qquad \xi_i > 0 ,\ \eta_j > 0
\]
— \textbf{штрафные коэффициенты}. Функция $\alpha(x)$ — штрафная; решается
вспомогательная задача
\[
  \varphi(x) + \mu\,\alpha(x) \ \to\ \min ,
\]
и получается решение $x_\mu$; далее $\mu$ увеличивают и повторяют.

\begin{idea}
Коэффициенты $\xi_i, \eta_j$ и общий множитель $\mu$ играют разные роли.
Коэффициенты $\xi_i, \eta_j$ фиксированы и выравнивают ограничения между собой
(они могут быть записаны в разных единицах и иметь разный масштаб), а $\mu$
растёт по ходу работы и постепенно «ужесточает» всю систему сразу. Именно
увеличение $\mu$ и означает «улучшаем свойства этих задач»: при малом $\mu$
вспомогательная задача хорошо обусловлена, но её решение далеко от $S$; при
большом $\mu$ решение почти допустимо, но функция становится овражной. Поэтому
решают не одну задачу, а серию, беря решение предыдущей как начальное
приближение следующей.
\end{idea}

\subsubsection*{Свойства вспомогательных задач}

Рассмотрим общую постановку. Построим вспомогательную задачу
\[
  \varphi(x) + \mu\,\alpha(x) \to \min , \qquad \mu > 0 , \quad x \in V ,
\]
пусть $x_\mu$ — её решение, и
\[
  P(\mu) = \varphi(x_\mu) + \mu\,\alpha(x_\mu) .
\]

\begin{lemma}\label{lem:penalty}
Пусть $\varphi$, $\psi_i$, $\chi_j$ непрерывны (то есть все функции, входящие в
постановку задачи), и для всякого $\mu > 0$ существует решение $x_\mu$
вспомогательной задачи. Тогда
\begin{enumerate}[label=\arabic*),left=0pt]
  \item $\displaystyle \min_{x \in S} \varphi(x) \ \ge\ \sup_{\mu > 0} P(\mu)$;
  \item $\alpha(x_\mu)$ — невозрастающая функция от $\mu$;
  \item $\varphi(x_\mu)$ — неубывающая функция от $\mu$;
  \item $P(\mu)$ — неубывающая функция от $\mu$;
  \item $\displaystyle \lim_{\mu \to \infty} P(\mu) = \sup_{\mu > 0} P(\mu)$.
\end{enumerate}
\end{lemma}

\lecture{Лекция 9}{}

\begin{proof}
\textbf{1)} Для любой допустимой точки $x \in S$ штраф равен нулю:
$\alpha(x) = 0$. Так как $x_\mu$ минимизирует $\varphi + \mu\alpha$ на $V
\supset S$,
\[
  P(\mu) = \min_{z \in V}\big(\varphi(z) + \mu\alpha(z)\big)
  \ \le\ \varphi(x) + \mu\underbrace{\alpha(x)}_{=\,0} = \varphi(x) .
\]
Беря минимум по $x \in S$ и затем супремум по $\mu$, получаем требуемое.

\medskip
\textbf{2)--4)} Пусть $0 < \mu_1 < \mu_2$. По определению $x_{\mu_1}, x_{\mu_2}$
как решений:
\[
  \varphi(x_{\mu_1}) + \mu_1 \alpha(x_{\mu_1})
  \ \le\ \varphi(x_{\mu_2}) + \mu_1 \alpha(x_{\mu_2}) ,
\]
\[
  \varphi(x_{\mu_2}) + \mu_2 \alpha(x_{\mu_2})
  \ \le\ \varphi(x_{\mu_1}) + \mu_2 \alpha(x_{\mu_1}) .
\]
Складывая, получаем
$(\mu_1 - \mu_2)\alpha(x_{\mu_1}) \le (\mu_1 - \mu_2)\alpha(x_{\mu_2})$, и так
как $\mu_1 - \mu_2 < 0$,
\[
  \alpha(x_{\mu_1}) \ \ge\ \alpha(x_{\mu_2})
  \qquad \text{— пункт 2)} .
\]
Подставляя это в первое неравенство,
\[
  \varphi(x_{\mu_1}) \ \le\ \varphi(x_{\mu_2})
    + \mu_1\big(\alpha(x_{\mu_2}) - \alpha(x_{\mu_1})\big)
  \ \le\ \varphi(x_{\mu_2})
  \qquad \text{— пункт 3)} .
\]
Наконец,
\begin{align*}
  P(\mu_2) &= \varphi(x_{\mu_2}) + \mu_2\alpha(x_{\mu_2})
   \ \ge\ \varphi(x_{\mu_2}) + \mu_1\alpha(x_{\mu_2}) \\
  &\ge \varphi(x_{\mu_1}) + \mu_1\alpha(x_{\mu_1}) = P(\mu_1)
   \qquad \text{— пункт 4)} ,
\end{align*}
где последнее неравенство — это первое неравенство оптимальности.

\medskip
\textbf{5)} По пункту~4) функция $P$ не убывает, по пункту~1) она ограничена
сверху величиной $\min_S \varphi$. Монотонная ограниченная функция имеет предел
при $\mu \to \infty$, и этот предел равен её супремуму.
\end{proof}

\begin{idea}
Смысл пунктов леммы: $P(\mu)$ — растущая снизу оценка оптимума
$\min_S\varphi$ (пункты~1 и~4), приближения становятся всё «допустимее»
(штраф $\alpha(x_\mu)$ убывает, пункт~2), но значение самой целевой функции
при этом растёт (пункт~3) — приближения подходят к $S$ \emph{снаружи}, со
стороны меньших значений $\varphi$. Пункт~5 говорит, что последовательность
нижних оценок исчерпывает всё, что из этой схемы можно выжать.
\end{idea}

\subsubsection*{Теорема о сходимости метода внешних штрафов}

\begin{theorem}\label{thm:penalty}
Пусть
\begin{enumerate}[label=\arabic*),left=0pt,nosep]
  \item исходная задача имеет решение $x_*$;
  \item выполнены условия леммы~\ref{lem:penalty};
  \item $\{x_\mu\} \subset M$, где $M$ — компакт.
\end{enumerate}
Тогда
\begin{enumerate}[label=\arabic*),left=0pt,nosep]
  \item $\displaystyle \min_{x \in S}\varphi(x) = \lim_{\mu \to +\infty} P(\mu)$;
  \item любая предельная точка последовательности $\{x_{\mu_k}\}$
        ($\mu_k \to +\infty$) является решением исходной задачи;
  \item $\mu\,\alpha(x_\mu) \xrightarrow[\mu \to +\infty]{} 0$.
\end{enumerate}
\end{theorem}

\begin{proof}
\textbf{Шаг 1 (ключевое утверждение).} Покажем, что
\begin{equation}\label{eq:penalty-key}
  \forall \eps > 0 \quad \exists \mu > 0 : \qquad \alpha(x_\mu) \le \eps .
\end{equation}

От противного: пусть существует $\eps > 0$ такое, что
$\alpha(x_\mu) > \eps$ \emph{при всех} $\mu > 0$. Возьмём $\mu_1 = 1$ и
обозначим $x_1 = x_{\mu_1}$. По пункту~1) леммы~\ref{lem:penalty} для любой
точки $y \in S$
\begin{equation}\label{eq:penalty-2}
  \varphi(y) \ \ge\ \min_{x \in S}\varphi(x) \ \ge\ \sup_{\mu > 0} P(\mu)
  \ \ge\ P(\mu_1) = \varphi(x_1) + \mu_1 \alpha(x_1) \ \ge\ \varphi(x_1) .
\end{equation}
Выберем
\begin{equation}\label{eq:penalty-3}
  \mu \ \ge\ 1 + \frac{1}{\eps}\big(\varphi(y) - \varphi(x_1)\big) \ \ge\ 1 ,
\end{equation}
что возможно, так как правая часть конечна и не меньше единицы
по~\eqref{eq:penalty-2}. Так как $\mu \ge 1 = \mu_1$, по пункту~3) леммы
\begin{equation}\label{eq:penalty-4}
  \varphi(x_\mu) \ \ge\ \varphi(x_1) .
\end{equation}
Тогда, пользуясь $\alpha(x_\mu) > \eps$, \eqref{eq:penalty-3}
и~\eqref{eq:penalty-4},
\begin{align*}
  \min_{x \in S}\varphi(x)
  \ \ge\ \sup_{\mu>0} P(\mu)
  \ \ge\ \underbrace{\varphi(x_\mu)}_{\ge\ \varphi(x_1)}
       + \underbrace{\mu\,\alpha(x_\mu)}_{>\ \mu\eps}
  &\ >\ \varphi(x_1)
       + \Big(1 + \tfrac{1}{\eps}\big(\varphi(y) - \varphi(x_1)\big)\Big)\eps \\
  &\ =\ \varphi(x_1) + \eps + \varphi(y) - \varphi(x_1)
   \ =\ \varphi(y) + \eps .
\end{align*}
Неравенство верно для \emph{любой} $y \in S$; подставляя $y = x_*$, получаем
$\min_S \varphi > \min_S \varphi + \eps$ — противоречие.
Значит~\eqref{eq:penalty-key} верно.

\medskip
\textbf{Шаг 2.} По пункту~2) леммы функция $\alpha(x_\mu)$ не возрастает, она
неотрицательна, и по~\eqref{eq:penalty-key} принимает сколь угодно малые
значения. Следовательно
\[
  \alpha(x_\mu) \xrightarrow[\mu \to +\infty]{} 0 .
\]

\medskip
\textbf{Шаг 3 (пункты 1 и 2 теоремы).} Пусть $\mu_k \to +\infty$. Так как
$\{x_\mu\} \subset M$ — компакт, существует сходящаяся подпоследовательность
$x_{\mu_k} \to \bar{x}$. По непрерывности $\alpha$ и шагу~2
\[
  \alpha(\bar{x}) = 0 \qquad\Longrightarrow\qquad \bar{x} \in S
  \qquad\Longrightarrow\qquad \varphi(\bar{x}) \ge \varphi(x_*) .
\]
С другой стороны,
\[
  P(\mu_k) = \varphi(x_{\mu_k}) + \mu_k \alpha(x_{\mu_k}) \ \ge\ \varphi(x_{\mu_k})
  \qquad\Longrightarrow\qquad
  \lim_{k \to \infty} P(\mu_k) \ \ge\ \varphi(\bar{x}) .
\]
Собирая всё вместе и пользуясь пунктом~5) леммы,
\[
  \underline{\min_{x \in S}\varphi(x)}
  \ \ge\ \sup_{\mu > 0} P(\mu)
  \ =\ \lim_{\mu \to +\infty} P(\mu)
  \ \ge\ \varphi(\bar{x})
  \ \ge\ \underline{\varphi(x_*)} .
\]
Крайние члены цепочки равны, значит все неравенства в ней — равенства. Отсюда
сразу пункт~1) теоремы и $\varphi(\bar{x}) = \varphi(x_*)$, то есть $\bar{x}$ —
решение исходной задачи (пункт~2).

\medskip
\textbf{Шаг 4 (пункт 3 теоремы).} Из определения $P$
\[
  \mu\,\alpha(x_\mu) = P(\mu) - \varphi(x_\mu) .
\]
При $\mu \to +\infty$ имеем $P(\mu) \to \min_S\varphi = \varphi(x_*)$ (шаг~3).
Функция $\varphi(x_\mu)$ не убывает (пункт~3 леммы) и ограничена сверху:
$\varphi(x_\mu) \le P(\mu) \le \min_S\varphi$, значит имеет предел; вдоль
подпоследовательности из шага~3 этот предел равен
$\varphi(\bar{x}) = \varphi(x_*)$. Следовательно обе величины в правой части
стремятся к $\varphi(x_*)$, и $\mu\,\alpha(x_\mu) \to 0$.
\end{proof}

\begin{corollary*}
Если $\alpha(x_\mu) = 0$ для некоторого $\mu > 0$, то $x_\mu$ — решение исходной
задачи.
\end{corollary*}

\begin{proof}
Из $\alpha(x_\mu) = 0$ следует $x_\mu \in S$. Тогда по пункту~1) леммы
\[
  \varphi(x_\mu) = \varphi(x_\mu) + \mu\,\alpha(x_\mu) = P(\mu)
  \ \le\ \sup_{\mu > 0} P(\mu) \ \le\ \min_{x \in S}\varphi(x) ,
\]
то есть $\min_S\varphi \ge \varphi(x_\mu)$. Но $x_\mu \in S$, поэтому
$\varphi(x_\mu) \ge \min_S\varphi$. Значит $\varphi(x_\mu) = \min_S\varphi$.
\end{proof}

\begin{idea}
Содержательно следствие говорит: если на какой-то конечной величине штрафа
решение вспомогательной задачи само оказалось допустимым — дальше увеличивать
$\mu$ незачем, задача уже решена точно. На практике это случается редко
(именно поэтому $\mu$ и приходится гнать к бесконечности), но служит признаком
останова «по счастливому случаю».
\end{idea}

\subsubsection*{Алгоритм метода внешних штрафов}

\begin{algo}[title={Метод внешних штрафных функций}]
\textbf{I. Начальный этап.} Задать $\mu_1 > 0$, $\eps > 0$, $\lambda > 1$ и
начальную точку $x_1$. Положить $k = 1$.

\medskip
\textbf{II. Основной этап.}
\begin{steps}
  \item Решить вспомогательную задачу с параметром $\mu_k$, стартуя из точки
        $x_k$:
        \[
          \varphi(x) + \mu_k\,\alpha(x) \ \to\ \min ,
        \]
        пусть $x_{\mu_k}$ — её решение.
  \item Проверка:
        \[
          \begin{aligned}
            &\mu_k\,\alpha(x_{\mu_k}) < \eps
             &&\Longrightarrow\quad \text{останов} ; \\
            &\mu_k\,\alpha(x_{\mu_k}) \ge \eps
             &&\Longrightarrow\quad x_{k+1} := x_{\mu_k} , \quad
               \mu_{k+1} = \mu_k \lambda , \quad k := k+1 .
          \end{aligned}
        \]
        Во втором случае — возврат к шагу~1.
\end{steps}
\end{algo}

\begin{idea}
Почему критерий останова именно $\mu_k\alpha(x_{\mu_k}) < \eps$. По пункту~3)
теоремы~\ref{thm:penalty} эта величина стремится к нулю, то есть критерий
заведомо когда-нибудь сработает. При этом
\[
  \mu\,\alpha(x_\mu) = P(\mu) - \varphi(x_\mu) \ \ge\ 0
\]
— это зазор между двумя оценками оптимума, которые даёт текущая итерация, а обе
они по теореме сходятся к $\min_S\varphi$. Малость зазора означает
одновременно и то, что точка почти допустима ($\alpha$ мал), и то, что оценки
сошлись. Старт следующей вспомогательной задачи из $x_{\mu_k}$ — обычный тёплый
старт: при большом $\mu$ функция $\varphi + \mu\alpha$ становится овражной, и
хорошее начальное приближение критично.
\end{idea}

\begin{example}[работа алгоритма]\label{ex:penalty}
\[
  \varphi(x) = (x-3)^2 \to \min , \qquad \psi(x) = x - 1 \le 0 ,
\]
то есть $S = \{x \le 1\}$, очевидное решение $x_* = 1$,
$\min_S\varphi = 4$. Берём квадратичный штраф
$\alpha(x) = \max^2\{0,\ x-1\}$ и вспомогательную задачу
\[
  F_\mu(x) = (x-3)^2 + \mu \max\nolimits^2\{0,\ x-1\} \ \to\ \min .
\]
Минимум лежит в области $x > 1$ (при $x \le 1$ функция $(x-3)^2$ убывает), там
\[
  F_\mu'(x) = 2(x-3) + 2\mu(x-1) = 0
  \qquad\Longrightarrow\qquad
  x_\mu = \frac{3 + \mu}{1 + \mu} .
\]
Прямой подстановкой:
\[
  \varphi(x_\mu) = \frac{4\mu^2}{(1+\mu)^2} ,
  \qquad
  \alpha(x_\mu) = \frac{4}{(1+\mu)^2} ,
  \qquad
  \mu\,\alpha(x_\mu) = \frac{4\mu}{(1+\mu)^2} ,
  \qquad
  P(\mu) = \frac{4\mu}{1 + \mu} .
\]
Видно всё, что обещают лемма и теорема: $x_\mu \downarrow 1 = x_*$,
$\alpha(x_\mu) \downarrow 0$, $\varphi(x_\mu) \uparrow 4$,
$P(\mu) \uparrow 4 = \min_S\varphi$, $\mu\alpha(x_\mu) \to 0$.

\medskip
Прогон при $\mu_1 = 1$, $\lambda = 10$, $\eps = 0{,}1$:
\begin{center}
\begin{tabular}{cccccc}
\toprule
$k$ & $\mu_k$ & $x_{\mu_k}$ & $\alpha(x_{\mu_k})$
    & $\mu_k \alpha(x_{\mu_k})$ & проверка \\
\midrule
1 & $1$   & $2$        & $1$              & $1$       & $\ge \eps$ \\
2 & $10$  & $1{,}1818$ & $0{,}0331$       & $0{,}331$ & $\ge \eps$ \\
3 & $100$ & $1{,}0198$ & $3{,}92\cdot10^{-4}$ & $0{,}0392$ & $< \eps$ — останов \\
\bottomrule
\end{tabular}
\end{center}
Ответ: $x \approx 1{,}02$, $\varphi \approx 3{,}92$ при точных
$x_* = 1$, $\varphi(x_*) = 4$.
\end{example}

\begin{idea}
Пример показывает и слабое место метода: точка $x_{\mu_k}$ \emph{никогда} не
попадает в $S$ — она лишь приближается к границе снаружи. Если допустимость
обязательна (например, вне $S$ функция $\varphi$ вообще не определена),
внешний штраф не годится, и нужен барьер — следующий параграф.
\end{idea}

\subsection{Метод барьерных функций}
\label{subsec:barrier}

Рассмотрим задачу
\[
  \begin{cases}
    \varphi(x) \to \min , & x \in S , \\
    S = \{\, x \mid \psi_i(x) \le 0,\ i = \overline{1,m} \,\} .
  \end{cases}
\]

\begin{remark}
Метод барьерных функций \textbf{неприменим}, если в задаче есть жёсткие
ограничения (равенства).
\end{remark}

\begin{idea}
Причина проста: барьер должен расти к $+\infty$ при подходе к границе
\emph{изнутри}, то есть метод всё время держится внутри допустимого множества.
У множества, заданного равенством $\chi(x) = 0$, внутренности нет вовсе —
«внутри» просто негде находиться. Поэтому равенства обрабатывают внешним
штрафом (\S\ref{subsec:penalty}), а барьером — только неравенства; на практике
их часто комбинируют.
\end{idea}

Пусть $\beta(x)$ — \textbf{барьерная функция}:
\[
  \beta(x) \xrightarrow[\ x \to \partial S\ ]{} +\infty , \qquad x \in S .
\]
На её основе строится новая функция цели
\[
  \varphi_0(x) = \varphi(x) + M\,\beta(x) , \qquad M > 0 , \quad x \in S ,
\]
и рассматривается вспомогательная задача $\varphi_0(x) \to \min$, $x \in S$.

\paragraph{1) Одно неравенство.}
Пусть $S = \{\, x \mid \psi(x) \le 0 \,\}$. Тогда
\[
  \beta(x) = -\frac{1}{\psi(x)} ,
  \qquad\text{вспомогательная задача}\qquad
  \begin{cases}
    \varphi(x) + M\beta(x) \to \min , \\
    M \to +0 , \\
    x \in \mathring{S} = \{\, x \mid \psi(x) < 0 \,\} .
  \end{cases}
\]
Такая форма неудобна для дифференцирования; удобнее
\[
  \beta(x) = -\ln\big(-\psi(x)\big) .
\]

\paragraph{2) Общая постановка.}
\[
  \beta(x) = -\sum_{i=1}^{m} \xi_i \cdot \frac{1}{\psi_i(x)} ,
  \qquad \xi_i > 0 ,
  \qquad\text{или}\qquad
  \beta(x) = -\sum_{i=1}^{m} \xi_i \ln\big(-\psi_i(x)\big) .
\]

\begin{idea}
Почему логарифмический барьер удобнее обратного. У $-1/\psi$ производная
$\psi'/\psi^2$ растёт при подходе к границе как $1/\psi^2$, у
$-\ln(-\psi)$ — только как $1/\psi$. Логарифмический барьер «мягче»: его
градиент и гессиан ведут себя заметно лучше, вспомогательная задача
оказывается лучше обусловленной, и именно он лежит в основе современных
методов внутренней точки.

Общий множитель $M$ играет ту же роль, что $\mu$ у внешнего штрафа, но в
обратную сторону: там $\mu \to +\infty$ ужесточал наказание за выход, здесь
$M \to +0$ ослабляет барьер, позволяя приближениям подходить к границе всё
ближе.
\end{idea}

\begin{remark}
В этих методах создаётся блок проверки выхода за границу области; в случае
необходимости он вносит коррективы при решении задачи.
\end{remark}

\subsubsection*{Свойства вспомогательных задач и сходимость}

\begin{idea}[title={Откуда этот материал}]
Для внешнего штрафа лемма о свойствах $P(\mu)$ и теорема о сходимости
разбирались на лекции (\S\ref{subsec:penalty}); для барьера зеркальные
утверждения были только названы. Ниже они выписаны и доказаны — выкладки
полностью симметричны, с заменой $\mu \to +\infty$ на $M \to +0$.
\end{idea}

Обозначим через $x_M$ решение вспомогательной задачи
$\varphi(x) + M\beta(x) \to \min$, $x \in \mathring{S}$, и
\[
  \Theta(M) = \varphi(x_M) + M\,\beta(x_M) .
\]
Считаем $\beta \ge 0$ на $\mathring{S}$ (для $\beta = -\sum \xi_i/\psi_i$ это
выполнено).

\begin{lemma}\label{lem:barrier}
\begin{enumerate}[label=\arabic*),left=0pt,nosep]
  \item $\displaystyle\min_{x \in S}\varphi(x) \ \le\ \Theta(M)$ при всех $M > 0$;
  \item $\beta(x_M)$ не убывает при $M \to 0$;
  \item $\varphi(x_M)$ не возрастает при $M \to 0$;
  \item $\Theta(M)$ — неубывающая функция $M$;
  \item $\displaystyle\inf_{M>0}\Theta(M) = \lim_{M \to 0}\Theta(M)$.
\end{enumerate}
\end{lemma}

\begin{proof}
Пусть $M_1 > M_2 > 0$, $x_i = x_{M_i}$.

\textbf{1)} $x_M \in \mathring{S} \subset S$ и $\beta \ge 0$, поэтому
$\Theta(M) \ge \varphi(x_M) \ge \min_S\varphi$.

\textbf{2)} Из минимальности $x_1$ и $x_2$:
\[
  \varphi(x_1) + M_1\beta(x_1) \le \varphi(x_2) + M_1\beta(x_2) ,
  \qquad
  \varphi(x_2) + M_2\beta(x_2) \le \varphi(x_1) + M_2\beta(x_1) .
\]
Складывая, $(M_1 - M_2)\big(\beta(x_1) - \beta(x_2)\big) \le 0$, то есть
$\beta(x_{M_1}) \le \beta(x_{M_2})$.

\textbf{3)} Из второго неравенства
$\varphi(x_2) - \varphi(x_1) \le M_2\big(\beta(x_1) - \beta(x_2)\big) \le 0$.

\textbf{4)}
$\Theta(M_2) \le \varphi(x_1) + M_2\beta(x_1) \le \varphi(x_1) + M_1\beta(x_1)
 = \Theta(M_1)$.

\textbf{5)} $\Theta$ монотонна и ограничена снизу по пункту~1), поэтому нижняя
грань достигается в пределе.
\end{proof}

\begin{theorem}\label{thm:barrier}
Пусть $\mathring{S} \ne \varnothing$, функции непрерывны, при каждом $M > 0$
решение $x_M$ существует, $\{x_M\}$ лежит в компакте и всякая точка $S$
приближается точками $\mathring{S}$. Тогда предельные точки $\{x_M\}$ —
решения исходной задачи, причём
\[
  M\,\beta(x_M) \xrightarrow[M \to 0]{} 0 ,
  \qquad
  \min_{S}\varphi = \lim_{M \to 0}\Theta(M) .
\]
\end{theorem}

\begin{proof}
Обозначим $\varphi_* = \min_S\varphi$. По условию плотности для любого
$\eps > 0$ найдётся $y \in \mathring{S}$ с $\varphi(y) < \varphi_* + \eps$.
Тогда
\[
  \Theta(M) \le \varphi(y) + M\beta(y) \xrightarrow[M\to0]{} \varphi(y)
  \quad\Longrightarrow\quad
  \lim_{M\to0}\Theta(M) \le \varphi_* + \eps \quad \forall \eps ,
\]
то есть $\lim\Theta \le \varphi_*$; обратное неравенство — пункт~1) леммы.
Значит $\lim_{M\to0}\Theta(M) = \varphi_*$, откуда
\[
  0 \le M\beta(x_M) = \Theta(M) - \varphi(x_M)
    \le \Theta(M) - \varphi_* \longrightarrow 0 .
\]
Для предельной точки $\bar{x}$: $\bar{x} \in S$ по замкнутости, и
$\varphi(\bar{x}) = \lim\varphi(x_M) \le \lim\Theta = \varphi_*$, то есть
$\bar{x}$ — решение.
\end{proof}

\begin{idea}
Условие «всякая точка $S$ приближается точками $\mathring{S}$» существенно:
без него оптимум может лежать на изолированном куске границы, к которому
изнутри не подойти, и барьерный метод до него не доберётся. Для выпуклого $S$
с непустой внутренностью оно выполняется автоматически.
\end{idea}

\begin{idea}
Формально барьер не выпускает точку из $\mathring{S}$ — но только если шаг
достаточно мал. Любой одномерный поиск вдоль направления спуска может «перелететь»
границу, а за ней $\ln(-\psi)$ просто не определён (или $-1/\psi$ меняет знак и
становится притягивающим). Поэтому в реализацию добавляют явную проверку
допустимости пробной точки и, если она нарушена, дробят шаг. У внешнего штрафа
такой блок не нужен: там функция определена всюду.
\end{idea}

\begin{example}\label{ex:barrier}
\[
  \varphi(x) = x - 1 \ \to\ \min , \qquad \psi(x) = -x + 2 \le 0 ,
\]
то есть $S = \{x \ge 2\}$, решение $x_* = 2$, $\varphi(x_*) = 1$.

Берём $\beta(x) = -1/\psi(x) = 1/(x-2)$ и вспомогательную задачу
\[
  \Big( (x-1) - \frac{M}{-x+2} \Big) \ \to\ \min
  \qquad\text{то есть}\qquad
  F_M(x) = (x - 1) + \frac{M}{x - 2} \ \to\ \min .
\]
Условие экстремума:
\[
  F_M'(x) = 1 - \frac{M}{(x-2)^2} = 0
  \qquad\Longrightarrow\qquad
  (x-2)^2 = M
  \qquad\Longrightarrow\qquad
  x = 2 \pm \sqrt{M} ,
\]
причём $x_1 = 2 - \sqrt{M}$ лежит вне области, а
\[
  x_2 = 2 + \sqrt{M} = x_M .
\]
Проверим все свойства:
\[
  \begin{array}{ll}
    1)\ x_M = 2 + \sqrt{M} \xrightarrow[M \to +0]{} 2 = x_* ;
    &\qquad
    2)\ \varphi(x_M) = \sqrt{M} + 1 \ \text{— возрастает по } M ; \\[4pt]
    3)\ \varphi(x_M) \xrightarrow[M \to +0]{} \varphi(x_*) = 1 ;
    &\qquad
    4)\ \beta(x_M) = \dfrac{1}{\sqrt{M}} \xrightarrow[M \to +0]{} +\infty ; \\[8pt]
    \multicolumn{2}{l}{
    5)\ M\,\beta(x_M) = \dfrac{M}{\sqrt{M}} = \sqrt{M}
        \xrightarrow[M \to +0]{} 0 . }
  \end{array}
\]
\end{example}

\begin{idea}[title={Внешний штраф против барьера}]
Сравним пример~\ref{ex:penalty} и пример~\ref{ex:barrier} — картины
зеркальные:
\begin{center}
\begin{tabular}{lll}
\toprule
 & внешний штраф & барьер \\
\midrule
параметр        & $\mu \to +\infty$       & $M \to +0$ \\
приближения     & вне $S$                 & строго внутри $S$ \\
$\varphi(x_\cdot)$ & растёт к $\min_S\varphi$ \emph{снизу}
                & убывает к $\min_S\varphi$ \emph{сверху} \\
штраф/барьер    & $\alpha(x_\mu) \downarrow 0$
                & $\beta(x_M) \uparrow +\infty$ \\
произведение    & $\mu\alpha(x_\mu) \to 0$ & $M\beta(x_M) \to 0$ \\
равенства       & обрабатываются          & неприменим \\
\bottomrule
\end{tabular}
\end{center}
Общее у них то, что произведение «параметр $\times$ штраф» стремится к нулю —
именно оно и служит критерием останова в обоих методах. Разница в том, с какой
стороны приближения подходят к оптимуму, и это решает выбор метода: если вне
$S$ модель теряет смысл (отрицательные концентрации, корень из отрицательного
числа), берут барьер; если допустимое множество задано и равенствами —
внешний штраф.
\end{idea}


\newpage
\subsection{Метод Ньютона решения задач с ограничениями}
\label{subsec:newton-constrained}

\begin{idea}[title={Откуда этот материал}]
Этот параграф и три следующих на разобранных лекциях не записывались, но входят
в список экзаменационных вопросов. Они восстановлены в обозначениях курса.
\end{idea}

\[
  \varphi(x) \to \min , \qquad x \in \R^n ,
  \qquad\text{при}\qquad
  h_i(x) = 0 , \quad i = \overline{1,m} ,
\]
или в векторной форме $h(x) = 0$, $h = (h_1, \dots, h_m)\T$. Если в задаче есть
неравенства $g_j(x) \le 0$, то в методе активного набора активные ограничения
на текущей итерации рассматриваются как равенства $g_j(x) = 0$.

\subsubsection*{Функция Лагранжа и условия ККТ}

\[
  L(x, \lambda) = \varphi(x) + \lambda\T h(x) ,
  \qquad
  \lambda = (\lambda_1, \dots, \lambda_m)\T \quad\text{— множители Лагранжа},
\]
\[
  \nabla_x L(x,\lambda) = \nabla\varphi(x) + J\T(x)\,\lambda ,
  \qquad
  J(x) = h'(x) =
  \begin{pmatrix}
    \nabla\T h_1(x) \\ \vdots \\ \nabla\T h_m(x)
  \end{pmatrix}
  \quad\text{— матрица Якоби.}
\]

\begin{theorem}[необходимые условия Лагранжа]\label{thm:lagrange-nec}
Пусть $x_*$ — локальное решение задачи $\varphi(x) \to \min$, $h(x) = 0$;
$\varphi$, $h_i$ дифференцируемы, а градиенты ограничений линейно независимы:
$\rank J(x_*) = m$. Тогда существует $\lambda_*$ такой, что
\[
  \nabla\varphi(x_*) + J\T(x_*)\,\lambda_* = 0 ,
  \qquad h(x_*) = 0 .
\]
\end{theorem}

\begin{proof}
Так как $x_*$ — локальное решение на множестве $h(x) = 0$, допустимые
направления $s$ в первом приближении удовлетворяют $J(x_*)s = 0$. Если бы для
какого-то такого $s$ выполнялось $\nabla\T\varphi(x_*)s < 0$, то движение из
$x_*$ вдоль $s$ уменьшало бы $\varphi$, не нарушая ограничений в первом
порядке. Заменив $s$ на $-s$ (вектор $-s$ тоже лежит в ядре), получаем, что
обязано быть
\[
  \nabla\T\varphi(x_*)\, s = 0 \qquad\forall s : J(x_*)s = 0 ,
\]
то есть $\nabla\varphi(x_*)$ ортогонален ядру $J(x_*)$. Ортогональное дополнение
к ядру матрицы совпадает с линейной оболочкой её строк, поэтому существует
$\lambda_*$ с $\nabla\varphi(x_*) = -J\T(x_*)\lambda_*$. Допустимость
$h(x_*) = 0$ выполняется по условию.
\end{proof}

Необходимые условия записываются как система нелинейных уравнений
\begin{equation}\label{eq:kkt-system}
  F(x,\lambda) = 0 ,
  \qquad
  F(x,\lambda) =
  \begin{pmatrix}
    \nabla_x L(x,\lambda) \\ h(x)
  \end{pmatrix}
  =
  \begin{pmatrix}
    \nabla\varphi(x) + J\T(x)\lambda \\ h(x)
  \end{pmatrix} .
\end{equation}
Метод Ньютона для задач с ограничениями — это обычный метод Ньютона,
применённый к системе~\eqref{eq:kkt-system}.

\subsubsection*{Ньютоновская система}

Поправка $(p_k, \Delta\lambda_k)$ находится из $F'(x_k,\lambda_k)
\big(p_k;\ \Delta\lambda_k\big) = -F(x_k, \lambda_k)$. Матрица производной
имеет блочный вид
\[
  F'(x_k,\lambda_k) =
  \begin{pmatrix} W_k & J_k\T \\ J_k & 0 \end{pmatrix} ,
  \qquad
  J_k = J(x_k) , \qquad W_k = \nabla^2_{xx} L(x_k, \lambda_k) ,
\]
то есть
\begin{equation}\label{eq:newton-kkt}
  \begin{pmatrix} W_k & J_k\T \\ J_k & 0 \end{pmatrix}
  \begin{pmatrix} p_k \\ \Delta\lambda_k \end{pmatrix}
  = -\begin{pmatrix} \nabla_x L(x_k,\lambda_k) \\ h(x_k) \end{pmatrix} ,
\end{equation}
после чего $x_{k+1} = x_k + \alpha_k p_k$,
$\lambda_{k+1} = \lambda_k + \alpha_k \Delta\lambda_k$. В классическом варианте
$\alpha_k = 1$.

\begin{theorem}\label{thm:newton-kkt-lin}
Система~\eqref{eq:newton-kkt} получается линеаризацией условий ККТ
$\nabla_x L(x,\lambda) = 0$, $h(x) = 0$.
\end{theorem}

\begin{proof}
Разложим первую группу условий в окрестности $(x_k, \lambda_k)$:
\[
  \nabla_x L(x_k + p, \lambda_k + \Delta\lambda)
  \approx \nabla_x L(x_k,\lambda_k) + W_k p + J_k\T \Delta\lambda .
\]
Требование обращения в нуль даёт
$W_k p + J_k\T\Delta\lambda = -\nabla_x L(x_k,\lambda_k)$ — первая строка
системы. Линеаризуя ограничения, $h(x_k + p) \approx h(x_k) + J_k p$, и
требуя $h(x_k+p) = 0$, получаем $J_k p = -h(x_k)$ — вторая строка.
\end{proof}

\subsubsection*{Достаточное условие второго порядка}

Для точки $x_*$, удовлетворяющей условиям Лагранжа, вводится касательное
подпространство
\[
  T = \{\, s \mid J(x_*)\,s = 0 \,\} .
\]

\begin{theorem}\label{thm:second-order-suff}
Если $\nabla_x L(x_*,\lambda_*) = 0$, $h(x_*) = 0$ и для всякого ненулевого
$s \in T$
\[
  s\T \nabla^2_{xx} L(x_*, \lambda_*)\, s > 0 ,
\]
то $x_*$ — строгий локальный минимум задачи с ограничениями.
\end{theorem}

\begin{proof}
Пусть $x$ — допустимая точка вблизи $x_*$, $d = x - x_*$. Обе точки допустимы,
значит в первом приближении $J(x_*)d = 0$, то есть $d \in T$. Разложим функцию
Лагранжа:
\[
  L(x, \lambda_*) = L(x_*,\lambda_*)
    + \underbrace{\nabla_x\T L(x_*,\lambda_*)\,d}_{=\,0}
    + \tfrac12\, d\T \nabla^2_{xx}L(x_*,\lambda_*)\, d + o\big(\norm{d}^2\big) .
\]
Для допустимых точек $h(x) = h(x_*) = 0$, поэтому $L(x,\lambda_*) = \varphi(x)$
и $L(x_*,\lambda_*) = \varphi(x_*)$. Следовательно
\[
  \varphi(x) - \varphi(x_*)
  = \tfrac12\, d\T \nabla^2_{xx}L(x_*,\lambda_*)\, d + o\big(\norm{d}^2\big) > 0
\]
при достаточно малом ненулевом $d$, то есть $\varphi(x) > \varphi(x_*)$.
\end{proof}

\begin{algo}[title={Метод Ньютона для задачи с равенствами}]
\begin{steps}
  \item Задать $x_0$, $\lambda_0$, $\eps > 0$.
  \item Вычислить $\nabla\varphi(x_k)$, $J_k = J(x_k)$,
        $W_k = \nabla^2_{xx}L(x_k,\lambda_k)$.
  \item Вычислить невязку ККТ
        $r_k = \big(\nabla_x L(x_k,\lambda_k);\ h(x_k)\big)$.
  \item Если $\norm{r_k} < \eps$ — остановиться.
  \item Решить систему~\eqref{eq:newton-kkt}.
  \item $x_{k+1} = x_k + \alpha_k p_k$,
        $\lambda_{k+1} = \lambda_k + \alpha_k \Delta\lambda_k$.
  \item $k := k+1$, перейти к шагу~2.
\end{steps}
\end{algo}

\begin{idea}[title={Глобализация и неравенства}]
При $\alpha_k = 1$ метод сходится квадратично, но лишь локально: полный шаг
может и увеличить $\varphi$, и увести от допустимого множества. В
глобализованном варианте берут $0 < \alpha_k \le 1$ и дробят шаг по мере
невязки
\[
  \Phi(x,\lambda) = \tfrac12 \norm{F(x,\lambda)}^2 ,
  \qquad
  \alpha := \beta\alpha, \ \ 0 < \beta < 1 ,
\]
пока не выполнится $\Phi(x_k + \alpha p_k, \lambda_k + \alpha\Delta\lambda_k)
< \Phi(x_k,\lambda_k)$.

Неравенства $g_i(x) \le 0$ обрабатываются активным набором
$I(x) = \{\, i \mid g_i(x) = 0 \,\}$: активные ограничения считаются
равенствами, для них строится $L(x,\lambda) = \varphi(x)
+ \sum_{i \in I(x)} \lambda_i g_i(x)$ и решается та же система. После решения
проверяются знаки множителей: должно быть $\lambda_i \ge 0$; ограничение с
$\lambda_i < 0$ удаляется из активного набора. Это в точности та же логика,
что в методе Розена (\S\ref{subsec:rosen}).
\end{idea}

\begin{example}
$\varphi(x) = (x_1-2)^2 + (x_2-1)^2 \to \min$ при $h(x) = x_1^2 + x_2 - 1 = 0$.
Тогда
\[
  \nabla_x L = \begin{pmatrix} 2(x_1-2) + 2\lambda x_1 \\ 2(x_2-1)+\lambda \end{pmatrix} ,
  \qquad
  J(x) = \begin{pmatrix} 2x_1 & 1 \end{pmatrix} ,
  \qquad
  W = \begin{pmatrix} 2+2\lambda & 0 \\ 0 & 2 \end{pmatrix} .
\]
Старт $x_0 = (0;1)\T$, $\lambda_0 = 0$ (точка допустима, $h(x_0) = 0$).
Система~\eqref{eq:newton-kkt} даёт $p_0 = (2;0)\T$, $\Delta\lambda_0 = 0$,
откуда $x_1 = (2;1)\T$: значение упало с $4$ до $0$, но
$h(x_1) = 4 \ne 0$ — точка стала недопустимой.

Вторая итерация: $J_1 = (4\ \ 1)$, система даёт
$\Delta\lambda_1 = \tfrac{8}{17}$, $p_1 = \big(-\tfrac{16}{17};\,-\tfrac{4}{17}\big)\T$,
откуда $x_2 = \big(\tfrac{18}{17};\,\tfrac{13}{17}\big)\T$,
$h(x_2) = \tfrac{256}{289} \approx 0{,}886$. Невязка ограничения упала с $4$ до
$0{,}886$.
\end{example}

\begin{idea}
Пример показывает главную особенность метода: он решает \emph{систему ККТ}, а
не минимизирует функцию, поэтому промежуточные точки допустимыми не являются и
значение $\varphi$ может вести себя немонотонно. Сходимость контролируется не
по $\varphi$, а по норме невязки $\norm{F(x,\lambda)}$ — отсюда и мера
$\Phi$ в глобализованном варианте.
\end{idea}

\subsection{Метод приведённого градиента Вульфа}
\label{subsec:reduced-gradient}

\[
  \varphi(x) \to \min
  \qquad\text{при}\qquad
  A x = b , \quad x \ge 0 ,
\]
где $A$ — матрица $m \times n$, $\rank A = m$, $\varphi$ дифференцируема.
Метод строит последовательность допустимых точек $x^0, x^1, \dots$ с
$Ax^k = b$, $x^k \ge 0$ и убывающим значением цели.

\subsubsection*{Разбиение на базисные и небазисные переменные}

Выберем $m$ линейно независимых столбцов $A$; они образуют базисную матрицу
$B$, остальные — матрицу $N$. Соответственно
\[
  x = \begin{pmatrix} x_B \\ x_N \end{pmatrix} ,
  \qquad
  Ax = b \iff B x_B + N x_N = b
  \qquad\Longrightarrow\qquad
  x_B = B^{-1}\big(b - N x_N\big) .
\]
Базисные переменные выражаются через небазисные — это и есть «приведение»
задачи.

\begin{lemma}\label{lem:wolfe-feasible}
Если $Ax = b$ и $Ad = 0$, то $A(x + \alpha d) = b$ при любом $\alpha$.
\end{lemma}

\begin{proof}
$A(x + \alpha d) = Ax + \alpha A d = b + \alpha\cdot 0 = b$.
\end{proof}

Разбивая направление так же, $d = (d_B;\ d_N)$, из $Ad = 0$ получаем
$B d_B + N d_N = 0$, то есть
\begin{equation}\label{eq:wolfe-dB}
  d_B = -B^{-1} N d_N .
\end{equation}
Значит достаточно выбрать небазисную часть $d_N$, базисная определяется
однозначно.

\subsubsection*{Приведённый градиент}

Пусть $g = \nabla\varphi(x) = (g_B;\ g_N)$. Вектор
\[
  r_N = g_N - N\T B^{-\mathsf{T}} g_B
\]
называется \textbf{приведённым градиентом}.

\begin{theorem}\label{thm:reduced-grad}
Для всякого направления $d$ с $Ad = 0$ выполняется
$g\T d = r_N\T d_N$.
\end{theorem}

\begin{proof}
По~\eqref{eq:wolfe-dB}
\[
  g\T d = g_B\T d_B + g_N\T d_N
        = g_B\T\big(-B^{-1}N d_N\big) + g_N\T d_N
        = \big(g_N\T - g_B\T B^{-1}N\big) d_N .
\]
Транспонируя скобку, $g_N\T - g_B\T B^{-1}N
= \big(g_N - N\T B^{-\mathsf{T}} g_B\big)\T = r_N\T$, откуда
$g\T d = r_N\T d_N$.
\end{proof}

\begin{idea}
Смысл: после исключения базисных переменных задача зависит только от $x_N$, и
$r_N$ — градиент этой «приведённой» функции. Первое слагаемое $g_N$ — прямое
влияние небазисных переменных, второе $N\T B^{-\mathsf T}g_B$ — косвенное,
через вынужденный пересчёт $x_B$. Если ограничений-равенств нет ($m = 0$), то
$r_N = g$, и метод вырождается в обычный градиентный спуск с учётом $x \ge 0$.
\end{idea}

\subsubsection*{Направление Вульфа}

Если небазисная переменная не на границе, можно двигаться против приведённого
градиента; если она на границе $x_j = 0$, отрицательная компонента $d_j$
недопустима. Отсюда правило
\begin{equation}\label{eq:wolfe-dir}
  d_j =
  \begin{cases}
    -r_j , & x_j > 0 , \\
    -r_j , & x_j = 0,\ r_j < 0 , \\
    0 , & x_j = 0,\ r_j \ge 0 ,
  \end{cases}
  \qquad j \ \text{— небазисный},
\end{equation}
а базисная часть берётся из~\eqref{eq:wolfe-dB}.

\begin{theorem}\label{thm:wolfe-descent}
Если направление построено по правилу~\eqref{eq:wolfe-dir} и $d_N \ne 0$, то
$g\T d < 0$, то есть $d$ — направление спуска.
\end{theorem}

\begin{proof}
По теореме~\ref{thm:reduced-grad} $g\T d = r_N\T d_N = \sum_j r_j d_j$.
Разберём слагаемые:
\begin{itemize}[nosep,left=0pt]
  \item $x_j > 0$: $d_j = -r_j$, значит $r_j d_j = -r_j^2 \le 0$;
  \item $x_j = 0$, $r_j < 0$: $d_j = -r_j > 0$, снова $r_j d_j = -r_j^2 < 0$;
  \item $x_j = 0$, $r_j \ge 0$: $d_j = 0$, слагаемое нулевое.
\end{itemize}
Все слагаемые неположительны, а при $d_N \ne 0$ хотя бы одно строго
отрицательно. Значит $g\T d < 0$.
\end{proof}

\subsubsection*{Максимальный допустимый шаг}

Новая точка $x(\alpha) = x + \alpha d$; равенства сохраняются по
лемме~\ref{lem:wolfe-feasible}, остаётся неотрицательность. Для компонент с
$d_j < 0$ нужно $x_j + \alpha d_j \ge 0$, то есть $\alpha \le x_j/(-d_j)$,
откуда
\[
  \alpha_{\max} = \min_{j :\, d_j < 0} \frac{x_j}{-d_j} ,
\]
а если $d_j \ge 0$ при всех $j$, то $\alpha_{\max} = +\infty$.

\begin{theorem}\label{thm:wolfe-step}
Если $x \ge 0$, $Ad = 0$ и $0 \le \alpha \le \alpha_{\max}$, то
$x + \alpha d \ge 0$ и $A(x + \alpha d) = b$.
\end{theorem}

\begin{proof}
Равенства — лемма~\ref{lem:wolfe-feasible}. Для $d_j \ge 0$ очевидно
$x_j + \alpha d_j \ge x_j \ge 0$. Для $d_j < 0$ по определению $\alpha_{\max}$
имеем $\alpha \le x_j/(-d_j)$; умножая на $-d_j > 0$, получаем
$-\alpha d_j \le x_j$, то есть $x_j + \alpha d_j \ge 0$.
\end{proof}

\subsubsection*{Условие оптимальности}

\begin{theorem}\label{thm:wolfe-opt}
Пусть $\varphi$ выпукла и дифференцируема, $x$ допустима и для всех
небазисных индексов
\[
  x_j > 0 \implies r_j = 0 ,
  \qquad
  x_j = 0 \implies r_j \ge 0 .
\]
Тогда $x$ — глобальный минимум задачи.
\end{theorem}

\begin{proof}
Возьмём любую допустимую $y$. Так как $Ax = b$ и $Ay = b$, вектор
$d = y - x$ удовлетворяет $Ad = 0$, и по теореме~\ref{thm:reduced-grad}
$g\T(y - x) = r_N\T d_N$. Разберём слагаемые: при $x_j > 0$ имеем $r_j = 0$ и
вклад нулевой; при $x_j = 0$ из допустимости $y_j \ge 0$ следует
$d_j = y_j - x_j = y_j \ge 0$, а $r_j \ge 0$, значит $r_j d_j \ge 0$.
Следовательно $g\T(y-x) \ge 0$. По неравенству первого порядка для выпуклой
функции
\[
  \varphi(y) \ \ge\ \varphi(x) + g\T(y - x) \ \ge\ \varphi(x) .
\]
Это верно для любой допустимой $y$, значит $x$ — глобальный минимум.
\end{proof}

\begin{idea}
Записанное условие — это в точности условия ККТ задачи
$\varphi(x) \to \min$, $Ax = b$, $x \ge 0$: приведённый градиент играет роль
оценок $d[j]$ симплекс-метода (\S\ref{subsec:simplex-scheme}), а условие
«$x_j > 0 \Rightarrow r_j = 0$» — дополняющей нежёсткости. При линейной
$\varphi$ метод Вульфа превращается в симплекс-метод.
\end{idea}

\begin{algo}[title={Метод приведённого градиента Вульфа}]
\begin{steps}
  \item Задать допустимую $x^0 \ge 0$, $Ax^0 = b$; выбрать базис $B$; задать
        $\eps > 0$; $k = 0$.
  \item Разбить переменные $x^k = (x_B^k;\ x_N^k)$ и вычислить
        $g^k = \nabla\varphi(x^k) = (g_B^k;\ g_N^k)$.
  \item Вычислить приведённый градиент
        $r_N^k = g_N^k - N\T B^{-\mathsf T} g_B^k$.
  \item Построить $d_N^k$ по правилу~\eqref{eq:wolfe-dir}.
  \item Если $d_N^k = 0$ — допустимого направления спуска нет; в выпуклом
        случае $x^k$ оптимальна, конец.
  \item Найти $d_B^k = -B^{-1}N d_N^k$ и полное направление $d^k$.
  \item Вычислить $\alpha_{\max} = \min_{j:\,d_j^k<0} x_j^k / (-d_j^k)$.
  \item Найти $\alpha_k = \argmin_{0 \le \alpha \le \alpha_{\max}}
        \varphi(x^k + \alpha d^k)$.
  \item $x^{k+1} = x^k + \alpha_k d^k$. Если какая-то переменная обратилась в
        нуль, при необходимости сменить базис.
  \item $k := k+1$, перейти к шагу~2.
\end{steps}
\end{algo}

\begin{example}
$\varphi(x) = (x_1-0{,}2)^2 + (x_2-0{,}3)^2 + (x_3-0{,}5)^2 \to \min$ при
$x_1+x_2+x_3 = 1$, $x \ge 0$. Здесь $A = (1\ \ 1\ \ 1)$, $b = 1$. Возьмём
$x^0 = (0{,}6;\,0{,}1;\,0{,}3)\T$ и базисную переменную $x_3$, тогда $B = (1)$,
$N = (1\ \ 1)$ и $d_3 = -(d_1 + d_2)$.

Градиент $g^0 = (0{,}8;\,-0{,}4;\,-0{,}4)\T$, приведённый градиент
$r_N^0 = (0{,}8;\,-0{,}4)\T - (1;\,1)\T(-0{,}4) = (1{,}2;\,0)\T$. Обе небазисные
переменные положительны, значит $d_N^0 = -r_N^0 = (-1{,}2;\,0)\T$,
$d_3^0 = 1{,}2$. Проверка: $Ad^0 = 0$, $(g^0)\T d^0 = -1{,}44 < 0$.
Ограничение даёт $\alpha_{\max} = 0{,}6/1{,}2 = 0{,}5$; одномерная минимизация
$\varphi = 0{,}24 - 1{,}44\alpha + 2{,}88\alpha^2$ даёт $\alpha_0 = 0{,}25$,
откуда $x^1 = (0{,}3;\,0{,}1;\,0{,}6)\T$ и $\varphi$ падает с $0{,}24$ до
$0{,}06$.

Вторая итерация: $g^1 = (0{,}2;\,-0{,}4;\,0{,}2)\T$,
$r_N^1 = (0;\,-0{,}6)\T$, $d^1 = (0;\,0{,}6;\,-0{,}6)\T$,
$\alpha_{\max} = 1$, $\alpha_1 = 0{,}25$, откуда
$x^2 = (0{,}3;\,0{,}25;\,0{,}45)\T$ и $\varphi = 0{,}015$. Итого
$0{,}24 \to 0{,}06 \to 0{,}015$.
\end{example}

\subsection{Выпуклый симплексный метод Зангвилла}
\label{subsec:zangwill}

\[
  \varphi(x) \to \min
  \qquad\text{при}\qquad
  Ax = b , \quad x \ge 0 ,
\]
где $\varphi$ — дифференцируемая (желательно выпуклая) функция, $A$ — матрица
$m \times n$. Допустимое множество $S = \{\, x \mid Ax = b,\ x \ge 0 \,\}$.
Метод Зангвилла — обобщение симплекс-метода ЛП на случай нелинейной выпуклой
цели.

\subsubsection*{Линейная подзадача}

На $k$-й итерации $\varphi$ заменяется линейной аппроксимацией в точке $x^k$:
\[
  \varphi(x) \approx \varphi(x^k) + \nabla\T\varphi(x^k)\,(x - x^k) .
\]
Слагаемые $\varphi(x^k)$ и $-\nabla\T\varphi(x^k)x^k$ от $x$ не зависят,
поэтому минимизация аппроксимации сводится к задаче
\[
  g_k\T x \to \min_{x \in S} ,
  \qquad g_k = \nabla\varphi(x^k) ,
\]
то есть к задаче \emph{линейного программирования} — отсюда и название
«симплексный». Её решение обозначим $y^k$, направление движения
\[
  d^k = y^k - x^k .
\]

\begin{theorem}\label{thm:zangwill-descent}
Если $g_k\T\big(y^k - x^k\big) < 0$, то $d^k = y^k - x^k$ — направление спуска.
\end{theorem}

\begin{proof}
$\nabla\T\varphi(x^k)d^k = g_k\T\big(y^k - x^k\big) < 0$ — это и есть
определение направления спуска.
\end{proof}

\begin{theorem}\label{thm:zangwill-opt}
Пусть $\varphi$ выпукла и дифференцируема, $S$ выпукло, $y^k$ решает
$g_k\T x \to \min_{x\in S}$. Если $g_k\T\big(y^k - x^k\big) = 0$, то $x^k$ —
решение задачи $\varphi(x) \to \min_{x \in S}$.
\end{theorem}

\begin{proof}
Так как $y^k$ минимизирует линейную функцию на $S$, для любого $x \in S$
выполнено $g_k\T y^k \le g_k\T x$. Вычитая $g_k\T x^k$,
\[
  \underbrace{g_k\T\big(y^k - x^k\big)}_{=\ 0} \ \le\ g_k\T\big(x - x^k\big)
  \qquad\Longrightarrow\qquad
  \nabla\T\varphi(x^k)\big(x - x^k\big) \ \ge\ 0 \quad \forall x \in S .
\]
Это условие экстремума выпуклой функции на выпуклом множестве; в силу
выпуклости $\varphi$ оно достаточно, значит $x^k$ — глобальный минимум.
\end{proof}

\begin{theorem}\label{thm:zangwill-feasible}
Если $x^k \in S$, $y^k \in S$ и $S$ выпукло, то
$x^{k+1} = x^k + \alpha_k\big(y^k - x^k\big) \in S$ при любом
$\alpha_k \in [0,1]$.
\end{theorem}

\begin{proof}
$x^{k+1} = (1-\alpha_k)x^k + \alpha_k y^k$ — выпуклая комбинация двух точек
$S$; по выпуклости она лежит в $S$.
\end{proof}

Шаг выбирается одномерной минимизацией
\[
  \alpha_k = \argmin_{0 \le \alpha \le 1} \varphi\big(x^k + \alpha d^k\big) .
\]

\begin{algo}[title={Выпуклый симплексный метод Зангвилла}]
\begin{steps}
  \item Задать допустимую $x^0 \in S$, точность $\eps > 0$, $k = 0$.
  \item Вычислить $g_k = \nabla\varphi(x^k)$.
  \item Решить линейную подзадачу $y^k = \argmin_{x\in S} g_k\T x$ (симплекс-метод).
  \item Построить $d^k = y^k - x^k$.
  \item Если $g_k\T d^k \ge -\eps$ — остановиться; в выпуклом случае $x^k$
        принимается за приближённое решение.
  \item Выбрать $\alpha_k = \argmin_{0\le\alpha\le1}\varphi(x^k + \alpha d^k)$.
  \item $x^{k+1} = x^k + \alpha_k d^k$, $k := k+1$, перейти к шагу~2.
\end{steps}
\end{algo}

\subsubsection*{Классический вариант: смена одной небазисной переменной}

\begin{idea}[title={Откуда этот материал}]
Выше изложен вариант со вспомогательной задачей ЛП на всём $S$. В
экзаменационной формулировке под методом Зангвилла понимается \emph{классический}
выпуклый симплексный метод, в котором задача ЛП не решается вовсе: по правилу
$u$/$v$ выбирается ровно одна небазисная переменная, и шаг делается вдоль
одного ребра — ровно как в симплекс-методе ЛП. Он и разобран ниже.
\end{idea}

Пусть $x = (x_B, x_N)$, $A = (B,N)$, $B$ невырождена, $x_B > 0$. Как и в методе
Вульфа (\S\ref{subsec:reduced-gradient}), вводится приведённый градиент
\[
  r_j = \frac{\partial\varphi}{\partial x_j}
      - \nabla_B\T\varphi\,B^{-1}A_j , \qquad j \in N ,
\]
и \textbf{критерий оптимальности} (условия Куна--Таккера):
\[
  r_j \ge 0 \ \text{при } x_j = 0 ,
  \qquad
  r_j = 0 \ \text{при } x_j > 0 .
\]

\paragraph{Выбор переменной.}
\[
  u = \min_{j \in N} r_j = r_s ,
  \qquad
  v = \max_{j \in N} x_j r_j = x_t r_t .
\]
Останов при $u \ge 0$ и $v = 0$ — это в точности критерий выше. Иначе:
\begin{gather*}
  \abs{u} \ge v \ \Longrightarrow\ \text{увеличиваем } x_s \
    (\sigma = +1,\ q = s) ; \\
  \abs{u} < v \ \Longrightarrow\ \text{уменьшаем } x_t \
    (\sigma = -1,\ q = t) .
\end{gather*}

\paragraph{Направление.}
$d_q = \sigma$, остальные небазисные нули, $d_B = -\sigma B^{-1}A_q$. Тогда
$Ad = \sigma A_q - \sigma A_q = 0$, то есть равенства сохраняются, и
$\nabla\T\varphi\,d = r_q d_q$:
\begin{itemize}[nosep,left=0pt]
  \item $\sigma = +1$: $r_q = u < 0$, значит $\nabla\T\varphi\,d < 0$;
        увеличение $x_q \ge 0$ допустимо всегда;
  \item $\sigma = -1$: $r_q > 0$ и $x_q > 0$, значит
        $\nabla\T\varphi\,d = -r_q < 0$; уменьшать $x_q$ можно до нуля.
\end{itemize}
Базисные компоненты строго положительны, поэтому при малом шаге точка остаётся
допустимой, а максимальный шаг —
$\alpha_{\max} = \min\{-x_i/d_i : d_i < 0\}$ (при $\sigma = -1$ учитывается и
$x_q$).

\begin{theorem}\label{thm:zangwill-kkt}
Предельные точки последовательности $\{x^k\}$ — точки Куна--Таккера; при
выпуклой $\varphi$ они являются решениями задачи.
\end{theorem}

\begin{idea}[title={Доказательство (схема)}]
$\varphi(x^k)$ убывает и ограничена снизу, значит сходится. Базисов конечное
число, поэтому можно выделить подпоследовательность с фиксированным базисом;
при фиксированном базисе направление $d(x)$ непрерывно по $x$. Если бы для
предельной точки $\bar{x}$ было $d(\bar{x}) \ne 0$, нашёлся бы шаг
$\bar\alpha > 0$ с $\varphi(\bar{x} + \bar\alpha d(\bar{x})) < \varphi(\bar{x})$,
что противоречит сходимости $\varphi(x^k)$. Значит $d(\bar{x}) = 0$, то есть
$u \ge 0$ и $v = 0$.
\end{idea}

\begin{idea}[title={Чем это отличается от метода условного градиента}]
Вариант со вспомогательной задачей ЛП — это тот же метод Франка--Вульфа
(\S\ref{subsec:cond-gradient}), и теоремы~\ref{thm:zangwill-descent}--\ref{thm:zangwill-feasible}
дословно повторяют результаты того параграфа. Разница в акценте:
\begin{itemize}[nosep,left=0pt]
  \item в \S\ref{subsec:cond-gradient} множество $S$ — произвольный выпуклый
        компакт, и подзадача решается «как получится»;
  \item у Зангвилла $S$ задано как $Ax = b$, $x \ge 0$, подзадача — честная
        задача ЛП в канонической форме, и решается она симплекс-методом с
        тёплым стартом от предыдущего базиса.
\end{itemize}
Поэтому метод Зангвилла — это практическая реализация условного градиента для
линейных ограничений, где вся техника §1--§4 работает внутри каждой итерации.
Скорость та же сублинейная, $O(1/k)$.
\end{idea}

\begin{example}
$\varphi(x) = (x_1-0{,}2)^2 + (x_2-0{,}3)^2 + (x_3-0{,}5)^2 \to \min$ при
$x_1+x_2+x_3 = 1$, $x \ge 0$; старт $x^0 = (1;\,0;\,0)\T$,
$\varphi(x^0) = 0{,}98$.

Итерация 1: $g_0 = (1{,}6;\,-0{,}6;\,-1)\T$. Линейная подзадача
$\min(1{,}6x_1 - 0{,}6x_2 - x_3)$ на симплексе достигает минимума в вершине с
наименьшим коэффициентом, то есть $y^0 = (0;\,0;\,1)\T$. Направление
$d^0 = (-1;\,0;\,1)\T$, $g_0\T d^0 = -2{,}6 < 0$. Одномерная минимизация
$\varphi = 0{,}98 - 2{,}6\alpha + 2\alpha^2$ даёт $\alpha_0 = 0{,}65$, откуда
$x^1 = (0{,}35;\,0;\,0{,}65)\T$, $\varphi = 0{,}135$.

Итерация 2: $g_1 = (0{,}3;\,-0{,}6;\,0{,}3)\T$, минимальный коэффициент при
$x_2$, значит $y^1 = (0;\,1;\,0)\T$, $d^1 = (-0{,}35;\,1;\,-0{,}65)\T$,
$g_1\T d^1 = -0{,}9 < 0$. Минимизация
$\varphi = 0{,}135 - 0{,}9\alpha + 1{,}545\alpha^2$ даёт
$\alpha_1 = 0{,}9/3{,}09 \approx 0{,}2913$, откуда
$x^2 \approx (0{,}248;\,0{,}291;\,0{,}461)\T$, $\varphi \approx 0{,}0039$.
Итого $0{,}98 \to 0{,}135 \to 0{,}0039$.
\end{example}

\begin{idea}
Видно характерное для метода поведение: решение линейной подзадачи всегда
попадает в \emph{вершину} симплекса, и без одномерного поиска точка прыгала бы
между вершинами. Именно шаг $\alpha_k \in [0,1]$ позволяет приближаться к
внутреннему оптимуму $x^* = (0{,}2;\,0{,}3;\,0{,}5)\T$ как к выпуклой
комбинации вершин.
\end{idea}


\newpage
\section{Квадратичное программирование в методе опорных векторов}
\label{sec:svm}

\subsection{Постановка задачи классификации}

Дана обучающая выборка $(x_i, y_i)$, $i = \overline{1,l}$, где
$x_i \in \R^n$, $y_i \in \{-1, 1\}$. Нужно построить разделяющую гиперплоскость
\[
  w\T x + b = 0 ,
\]
где $w$ — вектор нормали, $b$ — свободный член; новая точка классифицируется
правилом $\sgn\big(w\T x + b\big)$.

Гиперплоскость разделяет выборку, если $y_i\big(w\T x_i + b\big) \ge 1$ при
всех $i$ (при $y_i = 1$ это $w\T x_i + b \ge 1$, при $y_i = -1$ —
$w\T x_i + b \le -1$).

\subsubsection*{Геометрический зазор}

Расстояние от $x_i$ до гиперплоскости равно
$\rho_i = \abs{w\T x_i + b} / \norm{w}$. При выполнении
$y_i(w\T x_i + b) \ge 1$ минимальное расстояние равно $\rho = 1/\norm{w}$, а
ширина полосы между опорными гиперплоскостями $w\T x + b = \pm 1$ равна
$2/\norm{w}$. Поэтому максимизация зазора равносильна минимизации
$\tfrac12\norm{w}^2$.

\subsubsection*{Прямая задача}

Для линейно разделимой выборки
\[
  \tfrac12 \norm{w}^2 \to \min_{w,b}
  \qquad\text{при}\qquad
  y_i\big(w\T x_i + b\big) \ge 1 , \quad i = \overline{1,l} .
\]
Если выборка не разделима линейно, вводят переменные ошибок $\xi_i \ge 0$ и
\textbf{мягкий зазор}:
\begin{equation}\label{eq:svm-primal}
  \tfrac12 \norm{w}^2 + C \sum_{i=1}^{l}\xi_i \ \to\ \min_{w,b,\xi} ,
  \qquad
  y_i\big(w\T x_i + b\big) \ge 1 - \xi_i , \quad \xi_i \ge 0 ,
\end{equation}
где $C > 0$ — параметр штрафа за ошибки.

\begin{idea}
Это задача \textbf{квадратичного программирования}: целевая функция
квадратична, ограничения линейны. Чем больше $C$, тем сильнее наказываются
нарушения зазора — при $C \to \infty$ получается жёсткий зазор, при малом $C$
модель допускает много ошибок ради широкой полосы. Именно $C$ управляет
балансом между переобучением и недообучением.
\end{idea}

\subsection{Двойственная задача}

Запишем ограничения в виде $1 - \xi_i - y_i(w\T x_i + b) \le 0$,
$-\xi_i \le 0$, введём множители $\alpha_i \ge 0$ и $\mu_i \ge 0$:
\[
  L(w,b,\xi,\alpha,\mu) = \tfrac12\norm{w}^2 + C\sum_{i=1}^{l}\xi_i
    + \sum_{i=1}^{l}\alpha_i\Big(1 - \xi_i - y_i\big(w\T x_i + b\big)\Big)
    - \sum_{i=1}^{l}\mu_i \xi_i .
\]

\begin{theorem}\label{thm:svm-w}
В седловой точке функции Лагранжа
\[
  w = \sum_{i=1}^{l}\alpha_i y_i x_i .
\]
\end{theorem}

\begin{proof}
$\nabla_w L = w - \sum_i \alpha_i y_i x_i$; условие стационарности
$\nabla_w L = 0$ даёт требуемое.
\end{proof}

Остальные условия стационарности:
\[
  \frac{\partial L}{\partial b} = -\sum_{i=1}^{l}\alpha_i y_i = 0
  \ \Longrightarrow\ \sum_{i=1}^{l}\alpha_i y_i = 0 ,
  \qquad
  \frac{\partial L}{\partial \xi_i} = C - \alpha_i - \mu_i = 0
  \ \Longrightarrow\ 0 \le \alpha_i \le C
\]
(последнее — из $\mu_i \ge 0$ и $\alpha_i \ge 0$). Подставляя
$w = \sum_i \alpha_i y_i x_i$ в $L$, после сокращений получаем
\textbf{двойственную задачу}
\begin{equation}\label{eq:svm-dual}
  W(\alpha) = \sum_{i=1}^{l}\alpha_i
    - \frac12 \sum_{i=1}^{l}\sum_{j=1}^{l}
      \alpha_i\alpha_j y_i y_j\, x_i\T x_j \ \to\ \max_\alpha ,
  \qquad
  \sum_{i=1}^{l}\alpha_i y_i = 0 , \quad 0 \le \alpha_i \le C .
\end{equation}

\subsubsection*{Ядровая форма}

Заменяя скалярное произведение ядром
$K(x_i,x_j) = \varphi\T(x_i)\varphi(x_j)$, получаем
\[
  \sum_{i}\alpha_i - \frac12\sum_{i}\sum_j \alpha_i\alpha_j y_iy_j K(x_i,x_j)
  \to \max_\alpha ,
  \qquad
  \sum_i \alpha_i y_i = 0 , \quad 0 \le \alpha_i \le C ,
\]
а решающая функция —
$F(x) = \sum_{i}\alpha_i y_i K(x_i, x) + b$, классификация $\sgn F(x)$.

\begin{theorem}\label{thm:svm-convex}
Двойственная задача SVM — выпуклая задача квадратичного программирования.
\end{theorem}

\begin{proof}
Запишем $W(\alpha) = \sum_i \alpha_i - \tfrac12 \alpha\T Q \alpha$, где
$Q_{ij} = y_i y_j K(x_i,x_j)$. Для произвольного $z$, обозначив
$u_i = z_i y_i$,
\[
  z\T Q z = \sum_{i}\sum_j z_i z_j y_i y_j K(x_i,x_j)
          = \sum_{i}\sum_j u_i u_j K(x_i, x_j) \ \ge\ 0 ,
\]
так как ядровая матрица Грама положительно полуопределена. Значит
$Q \succeq 0$, функция $-\tfrac12\alpha\T Q\alpha$ вогнута, линейная часть
$\sum_i\alpha_i$ и вогнута, и выпукла, поэтому $W$ вогнута. Максимизация
вогнутой функции при линейных ограничениях — выпуклая задача.
\end{proof}

\subsubsection*{Условия ККТ и опорные векторы}

\[
  y_i\big(w\T x_i + b\big) \ge 1 - \xi_i ,\quad \xi_i \ge 0 ;
  \qquad 0 \le \alpha_i \le C ;
  \qquad \sum_i \alpha_i y_i = 0 ;
\]
\[
  \alpha_i\Big(y_i\big(w\T x_i + b\big) - 1 + \xi_i\Big) = 0 ,
  \qquad
  \big(C - \alpha_i\big)\xi_i = 0 .
\]
Отсюда:
\begin{itemize}[nosep,left=0pt]
  \item $\alpha_i = 0$ — объект не влияет на гиперплоскость;
  \item $0 < \alpha_i < C$ — объект лежит точно на границе зазора,
        $y_i(w\T x_i + b) = 1$;
  \item $\alpha_i = C$ — объект нарушает зазор или классифицирован ошибочно.
\end{itemize}
Объекты с $\alpha_i > 0$ называются \textbf{опорными векторами}.

\subsection{Алгоритмы решения}

\paragraph{Метод активного набора.}
Предполагается, что часть ограничений активна ($\alpha_i = 0$ либо
$\alpha_i = C$); на каждой итерации решается задача при фиксированном активном
наборе, проверяются условия ККТ, набор корректируется. Применим, когда матрица
квадратичной формы положительно полуопределена, ограничения линейны и можно
решать соответствующие СЛАУ. Минус: при большом $l$ активный набор часто
меняется, а матрицу ядра приходится хранить целиком.

\paragraph{Метод внутренней точки.}
Ограничения $0 < \alpha_i < C$ заменяются логарифмическим барьером
(\S\ref{subsec:barrier})
\[
  -\mu\sum_{i=1}^{l}\ln\alpha_i - \mu\sum_{i=1}^{l}\ln(C - \alpha_i) ,
\]
и решается последовательность гладких задач при $\mu \to +0$. Минус: большие
системы линейных уравнений на каждом шаге и необходимость строго внутренней
стартовой точки.

\paragraph{Методы декомпозиции.}
Оптимизация ведётся не по всем $\alpha_1,\dots,\alpha_l$ сразу, а по небольшому
рабочему множеству $B_k \subset \{1,\dots,l\}$; остальные переменные
фиксируются. Это и позволяет работать при очень больших $l$.

\subsection{Алгоритм SMO}
\label{subsec:smo}

SMO — частный случай декомпозиции, где на каждой итерации оптимизируются
\emph{ровно две} переменные $\alpha_i, \alpha_j$.

\begin{idea}
Почему именно две. Из-за ограничения $\sum_s \alpha_s y_s = 0$ изменение одной
переменной почти всегда его нарушает; двух переменных достаточно, чтобы
скомпенсировать изменение. Двумерная подзадача при этом решается аналитически,
без всякого численного метода, — отсюда крайняя дешевизна итерации.
\end{idea}

\begin{theorem}\label{thm:smo-eq}
Если новые значения выбираются так, что
$y_i\alpha_i^{\text{нов}} + y_j\alpha_j^{\text{нов}}
 = y_i\alpha_i^{\text{стар}} + y_j\alpha_j^{\text{стар}}$,
то ограничение $\sum_s \alpha_s y_s = 0$ сохраняется.
\end{theorem}

\begin{proof}
Меняются только компоненты $i$ и $j$, поэтому
\[
  \sum_s \alpha_s^{\text{нов}} y_s - \sum_s \alpha_s^{\text{стар}} y_s
  = y_i\big(\alpha_i^{\text{нов}} - \alpha_i^{\text{стар}}\big)
  + y_j\big(\alpha_j^{\text{нов}} - \alpha_j^{\text{стар}}\big) = 0
\]
по условию. Значит новая сумма равна старой, то есть нулю.
\end{proof}

\subsubsection*{Формулы SMO}

Обозначим $F(x) = \sum_s \alpha_s y_s K(x_s,x) + b$ и ошибку на объекте
$E_i = F(x_i) - y_i$. Для выбранной пары $i, j$:
\[
  \eta = K(x_i,x_i) + K(x_j,x_j) - 2K(x_i,x_j) .
\]
При $\eta > 0$ новое значение до отсечения
\[
  \alpha_j^{\text{нов}} = \alpha_j^{\text{стар}}
    + \frac{y_j\big(E_i - E_j\big)}{\eta} ,
\]
затем оно отсекается отрезком $L \le \alpha_j^{\text{нов}} \le H$:
\[
  y_i \ne y_j : \quad
  L = \max\big(0,\ \alpha_j^{\text{стар}} - \alpha_i^{\text{стар}}\big) ,
  \quad
  H = \min\big(C,\ C + \alpha_j^{\text{стар}} - \alpha_i^{\text{стар}}\big) ;
\]
\[
  y_i = y_j : \quad
  L = \max\big(0,\ \alpha_i^{\text{стар}} + \alpha_j^{\text{стар}} - C\big) ,
  \quad
  H = \min\big(C,\ \alpha_i^{\text{стар}} + \alpha_j^{\text{стар}}\big) .
\]
После этого
\[
  \alpha_i^{\text{нов}} = \alpha_i^{\text{стар}}
    + y_i y_j\big(\alpha_j^{\text{стар}} - \alpha_j^{\text{нов}}\big) .
\]

\subsubsection*{Обновление свободного члена}

\[
  b_1 = b^{\text{стар}} - E_i
    - y_i\Delta\alpha_i K(x_i,x_i) - y_j\Delta\alpha_j K(x_i,x_j) ,
\]
\[
  b_2 = b^{\text{стар}} - E_j
    - y_i\Delta\alpha_i K(x_i,x_j) - y_j\Delta\alpha_j K(x_j,x_j) .
\]
Если $0 < \alpha_i^{\text{нов}} < C$, берут $b^{\text{нов}} = b_1$; если
$0 < \alpha_j^{\text{нов}} < C$ — берут $b_2$; если обе переменные на границах
— $b^{\text{нов}} = (b_1 + b_2)/2$.

\begin{example}
Линейный SVM, $C = 1$, ядро $K(x_i,x_j) = x_i x_j$, выборка в $\R$:
\[
  x_1 = -1,\ y_1 = -1 ; \quad x_2 = 0,\ y_2 = -1 ; \quad
  x_3 = 1,\ y_3 = 1 ; \quad x_4 = 2,\ y_4 = 1 .
\]
Старт $\alpha^0 = 0$, $b^0 = 0$, тогда $F \equiv 0$ и
$E = (1;\,1;\,-1;\,-1)$.

\textbf{Итерация 1}, пара $i=1$, $j=3$: $K_{11} = K_{33} = 1$,
$K_{13} = -1$, $\eta = 4$. Так как $y_1 \ne y_3$, то $L = 0$, $H = 1$.
Получаем $\alpha_3^{\text{нов}} = 0 + \tfrac{1\cdot(1-(-1))}{4} = \tfrac12$
(в отрезке), $\alpha_1^{\text{нов}} = \tfrac12$. Проверка:
$\sum\alpha_i y_i = 0$. Обе формулы дают $b_1 = b_2 = 0$, значит $b^1 = 0$,
$w = \tfrac12 + \tfrac12 = 1$, $F(x) = x$.

\textbf{Итерация 2}, пара $i=2$, $j=3$: $E_2 = 1$, $E_3 = 0$,
$K_{22} = 0$, $K_{33} = 1$, $K_{23} = 0$, $\eta = 1$. При $y_2 \ne y_3$
получаем $L = \tfrac12$, $H = 1$, и
$\alpha_3^{\text{нов}} = \tfrac12 + 1 = \tfrac32 \to$ после отсечения
$\alpha_3^{\text{нов}} = 1$; тогда $\alpha_2^{\text{нов}} = \tfrac12$.
Итог $\alpha^2 = \big(\tfrac12;\,\tfrac12;\,1;\,0\big)$,
$\sum\alpha_iy_i = 0$, $b^2 = -1$, $w = \tfrac32$,
$F(x) = \tfrac32 x - 1$.
\end{example}

\begin{idea}
Ограничения алгоритмов: матрица ядра имеет размер $l \times l$ и при больших
$l$ не помещается в память; активные методы и методы внутренней точки для
таких выборок слишком дороги; SMO делает очень дешёвые итерации, но их может
потребоваться много, и качество сильно зависит от эвристики выбора пары
$(i,j)$ — обычно берут пару, наиболее нарушающую условия ККТ.
\end{idea}


\end{document}
